\documentclass[11pt,reqno]{amsart}
\usepackage[T1]{fontenc}
\usepackage[utf8]{inputenc}
\usepackage{lmodern}
\usepackage{amsmath,amssymb,mathtools,mathrsfs}
\usepackage{array,booktabs,longtable,enumitem,microtype,needspace}
\usepackage[margin=1.05in]{geometry}
\usepackage{xurl}
\usepackage[colorlinks=true,linkcolor=blue,citecolor=blue,urlcolor=blue]{hyperref}
\numberwithin{equation}{section}
\newtheorem{theorem}{Theorem}[section]
\newtheorem{proposition}[theorem]{Proposition}
\newtheorem{lemma}[theorem]{Lemma}
\newtheorem{corollary}[theorem]{Corollary}
\theoremstyle{definition}
\newtheorem{definition}[theorem]{Definition}
\theoremstyle{remark}
\newtheorem{remark}[theorem]{Remark}
\newtheorem{example}[theorem]{Example}
\DeclareMathOperator{\conv}{conv}
\DeclareMathOperator{\tr}{tr}
\DeclareMathOperator{\diag}{diag}
\renewcommand{\Re}{\operatorname{Re}}
\renewcommand{\Im}{\operatorname{Im}}
\DeclareMathOperator{\Log}{Log}
\newcommand{\C}{\mathbb C}
\newcommand{\D}{\mathbb D}
\newcommand{\E}{\mathbb E}
\DeclareRobustCommand{\cert}[1]{\nolinkurl{#1}}
\allowdisplaybreaks[1]
\emergencystretch=2em
\setlist[enumerate,1]{label=\textup{(\roman*)},leftmargin=*}
\setcounter{tocdepth}{1}
\title[Schmeisser's conjecture in degree at least 6251]{Schmeisser's conjecture for polynomials of degree at least 6251}
\author{Qiyue Tang}
\address{Zhili College, Tsinghua University, Beijing 100084, China}
\email{tangqy24@mails.tsinghua.edu.cn}
\date{November 11, 2026}
\subjclass[2020]{Primary 30C15; Secondary 30C10, 26D10, 65G30}
\keywords{Schmeisser's conjecture, critical points, zeros of polynomials, convex hull, interval arithmetic}
\hypersetup{pdftitle={Schmeisser's conjecture for polynomials of degree at least 6251},pdfauthor={Qiyue Tang}}
\begin{document}
\begin{abstract}
We prove Schmeisser's conjecture for polynomials of degree at least $6251$.
\end{abstract}
\maketitle
\tableofcontents
\section{Introduction}\label{sec:introduction}

Let $p$ be a complex polynomial of degree $n\ge2$, and let $Z(p)$
denote its zeros, counted with algebraic multiplicity. The Gauss--Lucas
theorem places the critical points $Z(p')$ in the convex hull
$\conv Z(p)$. Schmeisser's conjecture asserts that, when
$Z(p)\subseteq\overline\D=\{z\in\C:|z|\le1\}$, the closed unit disks
centered at the critical points cover this entire convex hull.
We prove the following degree range.

\begin{theorem}\label{thm:main}
Let $p$ be a complex polynomial of degree $n\ge6251$, all of whose
zeros lie in $\overline\D$. Then
\[
 \conv Z(p)\subseteq
 \bigcup_{\xi\in Z(p')}\{z\in\C:|z-\xi|\le1\}.
\]
\end{theorem}

The radius is sharp in every degree: for $p(z)=z^n-1$, the only
critical point is zero, while every zero has modulus one.

\subsection{The conjecture and its early development}

Sendov's conjecture concerns the same distance bound at each zero of
$p$. Sendov proposed it in 1958; Hayman's 1967 collection included it
as Problem~4.5 under Ilieff's name, through whom the problem had been
communicated \cite{Hayman1967,Schmeisser1971,Schmeisser2003}.
Rubinstein \cite{Rubinstein1968} proved the assertion at zeros on the
unit circle and for polynomials of degrees three and four.
Meir and Sharma \cite{MeirSharma1969} subsequently treated degree five.

A stronger localization already appeared in Schmeisser's early work.
For degrees at most four, he proved that for every zero $a$ there is
a critical point $\xi$ satisfying
\begin{equation}\label{in:early-localization}
             |\xi-a/2|\le1-|a|/2
\end{equation}
\cite{Schmeisser1969}. Gacs extended this localization to degree five
\cite{Gacs1971}. It implies the convex-hull covering directly.
Indeed, if $w=\sum_i t_i z_i$, with $t_i\ge0$ and $\sum_i t_i=1$,
then
\[
 \sum_i t_i\Re(w\overline{z_i})=|w|^2.
\]
Choose a zero $a$ for which $\Re(w\overline a)\ge|w|^2$. Then
\[
 |w-a/2|^2=|w|^2-\Re(w\overline a)+|a|^2/4
 \le|a|^2/4.
\]
Together with \eqref{in:early-localization}, this gives $|w-\xi|\le1$.

Schmeisser also established geometric and algebraic special cases.
His 1969 paper treats hulls whose vertices lie on the unit circle;
interior zeros are allowed \cite{Schmeisser1969}. His 1971 paper
develops coefficient estimates and proves Sendov's conjecture for
trinomials and for monic polynomials with real nonpositive
nonleading coefficients \cite{Schmeisser1971}. In Section~3 of
\cite{Schmeisser1977}, he formulated the convex-hull conjecture and
observed that the earlier proofs give this stronger conclusion in
degrees at most five, for hulls with all vertices on the unit circle,
and for trinomials. The same paper proves the assertion for triangular
hulls of arbitrary degree. Problem~4.30 of Hayman and Lingham
\cite{HaymanLingham} records the convex-hull question and its special
cases $p(0)=0$ and monic polynomials with real nonpositive
nonleading coefficients.

\subsection{Subsequent distance estimates}

For Sendov's conjecture, Brown and Xiang \cite{BrownXiang1999}
proved the result in degrees at most eight, and in arbitrary degree
when there are at most eight distinct zeros. The large-degree
approach began with D\'egot's theorem \cite{Degot}: for every fixed
$a\in(0,1)$ there is a degree bound depending on $a$ such that
a polynomial with $p(a)=0$ and all zeros in $\overline\D$ has
a critical point within distance one of $a$. Chalebgwa
\cite{Chalebgwa} made this dependence explicit, obtaining uniform
bounds when $a$ ranges over a compact subinterval of $(0,1)$.
Tao \cite{Tao2022} proved a single absolute, unspecified degree
threshold valid for all zeros. His argument treats separately
sequences of distinguished zeros whose moduli tend to zero,
to an interior point of $(0,1)$, or to one. An explicit absolute
bound $n\ge10^{200000}$ was subsequently given by Zhang
\cite{ZhangExplicit}. Mazur's all-degree argument and Tao's
exposition concern the distance from each actual zero to the
critical points \cite{Mazur,Tao2026}; the convex-hull assertion
appears as Conjecture~18 in \cite{Tao2026}.

A related line of work studies the centroid
\[
 G=\frac1n\sum_{z\in Z(p)}z
   =\frac1{n-1}\sum_{\xi\in Z(p')}\xi.
\]
Paw{\l}owski \cite{Pawlowski} proved, for zeros in $\overline\D$,
\[
 \min_{\xi\in Z(p')}|G-\xi|
 \le\frac{2n^{1/(n-1)}}{n^{2/(n-1)}+1}<1.
\]
Zhang improved this bound to $\sqrt{(n-2)/(n-1)}$
\cite{ZhangCentroid}. For $Z(p)\subseteq\overline\D$, his later paper
\cite{ZhangHull} gives
\[
 \min_{\xi\in Z(p')}|a-\xi|\le\sqrt{1-|a|^2}
       \qquad(a\in\conv Z(p')).
\]
The domain in this estimate is the critical-point hull.
For example, its radius follows from a convex representation
$a=\sum_j t_j\xi_j$:
\[
 \sum_j t_j|\xi_j-a|^2
 =\sum_j t_j|\xi_j|^2-|a|^2\le1-|a|^2.
\]
Schmeisser's question concerns the larger set $\conv Z(p)$.

Two recent estimates at an actual zero are useful here.
For zeros in $\overline\D$, Zhang \cite{ZhangQuadratic} proved the
quadratic Tang--Zhang inequality
\[
 \sum_{\xi\in Z(p')}\frac1{|a-\xi|^2}\ge n-1
       \qquad(a\in Z(p)),
\]
with an infinite summand at a common zero of $p$ and $p'$.
For arbitrary complex polynomials, Bouthat and Fallat
\cite{BouthatFallat}, and Zhang
\cite{ZhangVariance}, proved Borcea's quadratic variance bound
\[
 \min_{\xi\in Z(p')}|a-\xi|\le\sigma_2(p)
       \qquad(a\in Z(p)),\qquad
 \sigma_2(p)^2=\frac1n\sum_{z\in Z(p)}|z-G|^2.
\]
Section~\ref{sec:actual-roots} proves these estimates for
$n\ge6251$ and $n\ge2048$, respectively.

\subsection{The explicit degree bound}

An earlier argument established Schmeisser's conjecture for all
sufficiently large degrees \cite{TangQualitative}.
Theorem~\ref{thm:main} gives the explicit threshold $6251$.
Its proof consists of finite polynomial identities, analytic
inequalities uniform for $n\ge6251$, and a finite collection of
rational and interval verifications.

A point missed by all the critical-point unit disks can be moved
along a ray to a supporting edge of the zero hull. At a rightmost
zero $b$ we study the reciprocal critical points $(b-\xi)^{-1}$.
A matrix representation bounds their first two moments. A saddle
estimate and local primitive estimates sharpen these bounds and
produce nearby zeros of the original polynomial. Near the origin,
a finite Cayley--Hermite--Biehler identity supplies a source estimate;
the radial identity couples it to the same polynomial's moments.
In the other regions, equal-value fibers, product estimates and
positive sums of root counts yield a contradiction.
Section~\ref{sec:assembly} combines the parameter ranges.

The paper gives the analytic estimates, the complete remainders,
the parameter domains and the finite verification criteria.
The electronic appendix supplies the rational data and verification
programs. All zeros and critical points in sums and products are
counted with multiplicity. Terminating decimals denote exact rational
numbers. Each certified strict inequality holds on the indicated
closed parameter cell, including its boundary.

\section{Normalization and finite identities}\label{sec:preliminaries}

Write $D(c,r)=\{z\in\C:|z-c|<r\}$ for the open disk, and
$t_+=\max(t,0)$ for real $t$.
For vectors, $\|\cdot\|$ is the Euclidean norm; for matrices it is
the associated operator norm. We use $A^*$ for the conjugate transpose,
$\|A\|_F^2=\sum_{i,j}|A_{ij}|^2$ for the squared Frobenius norm,
and $A\ge0$ for a Hermitian positive semidefinite matrix.
The symbols $I$ and $\mathbf1$ denote the identity matrix and the
all-ones vector of the required size; $\tr$ and $\diag$ denote
the trace and diagonal-matrix operations.

Write $N=n-1$, and list the zeros of $p$ as $z_1,\ldots,z_n$ and
the critical points as $\xi_1,\ldots,\xi_N$. A \emph{strict hole}
is a point $a\in\conv Z(p)$ for which $|a-\xi_j|>1$ for every $j$.
The negation of Theorem~\ref{thm:main} is precisely the existence of
such a point. We may multiply $p$ by a nonzero constant whenever
convenient.

\begin{lemma}[Gauss--Lucas]\label{pr:gauss-lucas}
Every critical point belongs to $\conv Z(p)$.
\end{lemma}
\begin{proof}
If $w\notin\conv Z(p)$, a separating line and a rotation give
$\Re(e^{-i\theta}(w-z_k))>0$ for every $k$. Thus every
$e^{i\theta}/(w-z_k)$ has positive real part. Their sum is
$e^{i\theta}p'(w)/p(w)$ and cannot vanish. A zero shared by $p$
and $p'$ already lies in the zero hull.
\end{proof}

\begin{lemma}\label{pr:edge-reduction}
Suppose $n\ge6251$ and all zeros lie in $\overline\D$. If a strict
hole exists, after a rotation there is a strict hole $a$ with
$0<a<1$ on an actual supporting edge of a two-dimensional zero hull.
For every critical point,
\begin{equation}\label{pr:critical-halfplane}
 \Re\xi_j<a/2.
\end{equation}
There is an actual simple rightmost zero $b=x+iy$ with $x\ge a$.
After conjugation we may take $y\ge0$.
\end{lemma}
\begin{proof}
Gauss--Lucas implies $|\xi_j|\le1$, so zero cannot be a strict
hole. Rotate a strict hole onto the positive real axis. Then
\[
 2a\Re\xi_j<a^2+|\xi_j|^2-1\le a^2.
\]
Let $A$ be the last point of the zero hull on that ray. For
$a\le A\le1$,
\[
 |A-\xi_j|^2-|a-\xi_j|^2
 =(A-a)(A+a-2\Re\xi_j)\ge0.
\]
Thus the hole persists at $A$. If $A$ is an actual zero, it is either
multiple, in which case it is critical, or simple. In the latter
case Proposition~\ref{ar:energy} gives
$\sum_j|A-\xi_j|^{-2}\ge N$, contrary to all the distances being
greater than one.

A hull consisting of one point has no hole. If the hull is a segment,
Gauss--Lucas puts all critical points on it. The distance between any
two consecutive distinct critical points is at most two, since the
whole segment lies in the unit disk. Their closed unit intervals
therefore have connected union. Each endpoint of the hull belongs
to that union by the preceding actual-zero argument, so the whole
segment is covered. In the remaining case $A$ lies on a supporting
edge whose endpoints are actual zeros. Strict convexity of the unit
disk excludes $A=1$. Rename $A$ as $a$ and repeat the displayed
half-plane inequality at this point. A zero maximizing real part
has real part at least $a$; it cannot be multiple by
\eqref{pr:critical-halfplane}. Conjugation makes its imaginary part
nonnegative.
\end{proof}

For an actual simple rightmost zero $b$, put
\begin{equation}\label{pr:field-notation}
 v_j=(b-\xi_j)^{-1},\quad S=\sum_jv_j=\frac{p''(b)}{p'(b)},
 \quad H=|S|,\quad Q=\sum_j|v_j|^2,\quad
 d=\min_j|b-\xi_j|.
\end{equation}
Rightmostness and Gauss--Lucas give $\Re v_j\ge0$ and
$|v_j|\le1/d$. The equality $S=p''(b)/p'(b)$ counts the complete
critical divisor. If the other original zeros are $u_1,\ldots,u_N$,
factoring $p(z)=(z-b)g(z)$ gives the further identity
\begin{equation}\label{pr:actual-root-identity}
 S=2\sum_{k=1}^N\frac1{b-u_k}.
\end{equation}

\begin{lemma}\label{pr:radial}
The data in \eqref{pr:field-notation} satisfy
\begin{equation}\label{pr:radial-identity}
 N=2\Re(bS)+(1-|b|^2)Q
       -\sum_j\frac{1-|\xi_j|^2}{|b-\xi_j|^2}.
\end{equation}
If all zeros lie in $\overline\D$, $a>0$ is a strict hole, and
$b=x+iy$ is rightmost with $x\ge a$,
then $\Re S>0$ and
\begin{equation}\label{pr:hole-reciprocal}
 Q<\frac{2\Re S}{2x-a}.
\end{equation}
\end{lemma}
\begin{proof}
For each $\xi$ expand
\[
 2\Re\frac b{b-\xi}
 =1+\frac{|b|^2-|\xi|^2}{|b-\xi|^2}
\]
and sum. For the second assertion, Gauss--Lucas and the strict hole give
$2a\Re\xi<a^2+|\xi|^2-1\le a^2$. Thus
$\Re(b-\xi)^{-1}=(x-\Re\xi)/|b-\xi|^2
 >(x-a/2)/|b-\xi|^2$. Summing proves the assertion.
\end{proof}

\begin{lemma}\label{pr:primitive}
In the disk $|z|<d$, with the branch of each logarithm equal to zero
at the origin,
\[
 \frac{p'(b+z)}{p'(b)}=\prod_{j=1}^N(1+zv_j),\qquad
 p(b+z)=p'(b)\int_0^z\prod_{j=1}^N(1+tv_j)\,dt.
\]
For every integer $r\ge2$ and $|z|<d$,
\begin{equation}\label{pr:complete-log-tail}
 \left|\sum_j\Log(1+zv_j)
 -\sum_{k=1}^{r-1}\frac{(-1)^{k+1}z^k}{k}\sum_jv_j^k\right|
 \le\frac{Q|z|^r}{r d^{r-2}(1-|z|/d)}.
\end{equation}
\end{lemma}
\begin{proof}
The product is the factorization of $p'$; integration uses $p(b)=0$.
The logarithms are analytic because $|zv_j|<1$. For $k\ge r$,
$\sum_j|v_j|^k\le Qd^{2-k}$ and $1/k\le1/r$.
Summing the entire remaining geometric series proves the bound.
\end{proof}

When $H>0$ it is useful to write
\begin{equation}\label{pr:normalized-moments}
 D=Hd,\qquad \mu=D^{-1},\qquad q=Q/H^2,\qquad
 \sigma=S/H,\qquad w_j=v_j/H.
\end{equation}
Then $|\sigma|=1$, $\sum w_j=\sigma$,
$\sum|w_j|^2=q$, and $\max|w_j|=\mu$. For $r>0$, the
polynomial $p_r(z)=p(rz)$ at the corresponding zero $b/r$ has
field data $d/r,rS,rH,r^2Q$. Thus $D,q,\mu,\sigma$ are invariant
under positive real rescaling. Rescaled coordinates will be specified
when used.

\section{Estimates at an actual zero}\label{sec:actual-roots}

We prove the quadratic Tang--Zhang inequality \cite{ZhangQuadratic}
and the quadratic variance bound \cite{BouthatFallat,ZhangVariance}
in the degree ranges used below. Both estimates concern an actual
zero of the polynomial.

\subsection{The reciprocal-square estimate}

\begin{proposition}\label{ar:energy}
Let $p$ have degree $n\ge6251$, with all its zeros in $\overline\D$.
For each simple zero $a$ of $p$, one has
\[
 E_a:=\sum_{p'(\xi)=0}\frac1{|a-\xi|^2}\ge N.
\]
The inequality is strict if $|a|<1$.  At a multiple zero we interpret
the corresponding energy as $+\infty$.
\end{proposition}

\begin{proof}
Multiply by a constant to make $p$ monic.
Rotation reduces the interior case to $0<a<1$. Write
\[
 p(z)=(z-a)\prod_{j=1}^N(z-z_j),\quad
 q_j=(a-\xi_j)^{-1},\quad
 \mu=\frac1N\sum_jq_j=x+iy,
\]
\[
 s=\frac1N\sum_j|q_j|^2,\qquad v=s-|\mu|^2.
\]
Suppose $s\le1$. Thus $|\mu|\le1$ and $0\le v\le1-|\mu|^2$.
Put
\[
 F(t)=\prod_j(1-atq_j),\qquad J=\prod_j z_jq_j,
 \qquad \delta=1-a^2.
\]
Integration of $p'$ from $a$ to $0$, followed by its factorization,
gives
\begin{equation}\label{ar:communication-origin}
 (N+1)\int_0^1F(t)\,dt=(-1)^NJ,\qquad |J|\le s^{N/2}\le1.
\end{equation}
Integration from $a$ to $1/a$ gives a second identity:
\begin{equation}\label{ar:communication-polar}
 \prod_j\frac{1-az_j}{a-z_j}
 =\int_0^1\prod_j\bigl(a+\delta t q_j\bigr)\,dt.
\end{equation}
The denominators are nonzero because $a$ is simple. Since
\[
 |1-az_j|^2-|a-z_j|^2=\delta(1-|z_j|^2)\ge0,
\]
the arithmetic--geometric mean inequality yields
\begin{equation}\label{ar:polar}
 1\le\int_0^1\bigl(a^2+2a\delta xt+\delta^2st^2\bigr)^{N/2}\,dt.
\end{equation}
Define
\[
 \beta_s(t)=1-2axt+a^2st^2,\quad
 \beta(t)=1-2axt+a^2t^2,\quad
 \alpha=\frac{N\delta}{2}.
\]
For $0<t<1$, the integrand in \eqref{ar:polar} is strictly smaller than
\(
 \exp\{\alpha[-1+(2-\beta(1))t]\}.
\)
It follows first that $\beta(1)<1$, and hence $x>a/2$.  Writing
$g=1-\beta(1)>0$ and evaluating the exponential integral gives
\begin{equation}\label{ar:gap-equation}
 e^{\alpha g}-e^{-\alpha}>\alpha(1+g).
\end{equation}
Consequently
\begin{equation}\label{ar:polar-gaps}
 0<\beta(1)<\frac{\alpha}{1+\alpha},\qquad
 g>\frac{\log\alpha}{\alpha}\quad(\alpha\ge1).
\end{equation}
For the first bound, put $z=\alpha/(1+\alpha)$. Since
$z<\log(1+\alpha)$,
\[
 e^z-1-z<\frac{z^2e^z}{2}<\frac{\alpha^2}{2(1+\alpha)}.
\]
On the other hand, integrating
$1-e^{-t}>t/(1+t)\ge t/(1+\alpha)$ gives
$\alpha-1+e^{-\alpha}>\alpha^2/[2(1+\alpha)]$.
Thus \eqref{ar:gap-equation} excludes $\alpha g\le z$.
The second bound follows directly from $e^{\alpha g}>\alpha$.

There is also a useful estimate retaining $s$. Put
\[
 \rho=\delta s-1+2ax,\qquad \Lambda=\frac{(N+2)\delta}{2}.
\]
The quadratic in \eqref{ar:polar} is bounded above by its chord
$a^2+\delta(1+\rho)t$.  Necessarily $\rho>0$; otherwise its integral
is strictly less than one. Integration of this chord to the power $N/2$
and $\log(1+u)\le u$ give, for $\Lambda>1$,
\begin{equation}\label{ar:rho}
 \rho\ge\frac{\log\Lambda}{\Lambda},\quad
 \beta_s(1)=s-\rho\le1-\rho,\quad
 x^2\ge\rho.
\end{equation}
Indeed $e^{\Lambda\rho}\ge a^{N+2}+\Lambda(1+\rho)>\Lambda$,
and $\rho\le2ax-a^2=x^2-(x-a)^2$.

\smallskip
\noindent\emph{Two integration estimates.}
For nonnegative $b_j$, the arithmetic--geometric mean inequality implies
\[
 \prod_{k\ne j}b_k
 \le c_N\left(\frac1N\sum_kb_k^2\right)^{(N-1)/2},
 \qquad c_N=\left(\frac N{N-1}\right)^{(N-1)/2}<\frac53.
\]
The exact differentiation formula
\[
 F'=-Na\mu F-a^2t\sum_jq_j^2\prod_{k\ne j}(1-atq_k)
\]
therefore yields
\begin{equation}\label{ar:direct-budget}
 1\le |F(1)|+\frac N{N+1}a+|R|,
 \quad
 |R|\le\frac53Na^2\int_0^1t\beta_s(t)^{(N-1)/2}\,dt,
 \quad |F(1)|\le\beta_s(1)^{N/2}.
\end{equation}

For a second estimate, set $w_j=q_j-\mu$, $D(t)=1-at\mu$, and
$G_\theta(t)=\prod_j[D(t)-at\theta w_j]$.  Since $\sum_jw_j=0$,
\[
 D\partial_\theta G_\theta
 =-a^2t^2\theta\sum_jw_j^2\prod_{k\ne j}(D-at\theta w_k).
\]
Moreover $D\ne0$, and the mean squared modulus of the factors is at
most $\beta_s(t)$. Integration in $\theta$ gives
\[
 |F-D^N|\le\frac{5Na^2t^2v}{6|D|}\beta_s(t)^{(N-1)/2}.
\]
Integrate once more, use \eqref{ar:communication-origin}, and use
$|D|\ge1-axt$. With $r=|\mu|$, the result is
\begin{equation}\label{ar:centered-budget}
 1\le ar+C r(1-r^2)+E,
 \quad
 C=\frac56a^3N(N+1)\int_0^1
       \frac{t^2\beta(t)^{(N-1)/2}}{1-axt}\,dt,
 \quad E=\beta(1)^{(N+1)/2}.
\end{equation}
If $2C<a$, then $ar+Cr(1-r^2)\le a$, since
\[
 a-ar-Cr(1-r^2)=(1-r)[a-Cr(1+r)]\ge0.
\]

\smallskip
\noindent\emph{The range $a^2\ge9/10$.}
Let $\nu=(N-3)/2$. Here $\alpha\le N/20$.
For $0<\alpha\le3$, \eqref{ar:polar-gaps} gives
$\nu g>(N-3)/8>4\log N$.
For $3\le\alpha\le N/20$, it gives
\[
 \nu g>10(1-3/N)\log(N/20)>4\log N.
\]
The last inequalities hold already at $N=700$ and persist because the
derivatives of the differences are
\[
 \frac18-\frac4N>0,\qquad
 \frac{6-30/N}{N}+\frac{30\log(N/20)}{N^2}>0.
\]
Also $9N/160>4\log N$. Thus
\[
 \beta(1)^\nu<N^{-4},\qquad (133/160)^\nu<N^{-4}.
\]
On $0\le t\le1/4$,
$\beta(t)\le e^{-3a^2t/4}$ and $1-axt\ge3/4$.
Extending the positive integral to infinity gives
\[
 C_{[0,1/4]}
 \le\frac{10240}{243}\frac{N(N+1)}{(N-1)^3a^3}<\frac{60}{N}.
\]
For $1/4\le t\le1$, convexity bounds $\beta(t)$ by
$\max(133/160,\beta(1))$, and $\beta(t)\le2(1-axt)$.
Consequently
\[
 C_{[1/4,1]}<\frac{10}{3}N(N+1)N^{-4}<\frac7{N^2},
 \qquad C<\frac{61}{N}<\frac a2.
\]
If $\alpha\ge1$, then $E/\delta<1/(2N^3)<1/4$.
If $0<\alpha\le1$, the bound
$\beta(1)<\min(\alpha,1/2)$ gives
\[
 \frac E\delta<\frac N2
 \frac{\min(\alpha,1/2)^{(N+1)/2}}\alpha
 \le N2^{-(N+1)/2}<\frac14.
\]
For the middle inequality, the quotient increases with $\alpha$
on $(0,1/2]$ and decreases on $[1/2,1]$.
The latter inequality holds for every integer $N\ge16$ by its endpoint
and the ratio test. Equation \eqref{ar:centered-budget} would now give
$1\le a+E<a+(1-a^2)/4<1$.

\smallskip
\noindent\emph{The range $0<a\le\sqrt{9/10}$.}
Now $\delta\ge1/10$. Split the integral in
\eqref{ar:direct-budget} at the minimum of $\beta_s$, clipping this
point to $[0,1]$. On its decreasing part
$\beta_s(t)\le e^{-axt}$, while on its increasing part
$\beta_s(t)\le1-\rho$. Using \eqref{ar:rho} gives
\begin{align}
 |R|_{\rm dec}&\le D_N(a):=
 \min\left(\frac{28}{Na^2},\frac{7\Lambda}{N\log\Lambda}\right),
 \label{ar:D}\\
 |R|_{\rm inc}&\le U_N(a):=
 \frac53a^2N\Lambda^{-(N-1)/((N+2)\delta)},\qquad
 |F(1)|\le V_N(a):=\Lambda^{-N/((N+2)\delta)}.
 \label{ar:UV}
\end{align}
The first alternative in $D_N$ is omitted at $a=0$.
Both alternatives in $D_N$ decrease with $N$.  The exponent defining
$V_N$ has positive derivative
\[
 \frac{N+2\log\Lambda}{(N+2)^2\delta}.
\]
For $a>0$, the logarithmic derivative of $U_N$ is
\[
 \frac{(N+2)^2\delta-N(N-1)-3N\log\Lambda}
 {N(N+2)^2\delta}<0,
\]
because $\delta\le1$, $\log\Lambda>2$, and its numerator is at most
$5N+4-3N\log\Lambda<4-N$.
It therefore suffices to bound $D_N+U_N+V_N$ at $N_0=6250$.

Here are explicit bounds on seven closed intervals, each intersected
with $[0,\sqrt{9/10}]$; the right column is a strict upper bound
for the sum on that whole intersection.
\begin{center}
\begin{tabular}{ccr}
\toprule
$l$&$h$&upper bound\\\midrule
$0$&$1/4$&$659/1000$\\
$1/4$&$3/8$&$396/1000$\\
$3/8$&$1/2$&$346/1000$\\
$1/2$&$5/8$&$192/1000$\\
$5/8$&$3/4$&$54/1000$\\
$3/4$&$7/8$&$11/1000$\\
$7/8$&$19/20$&$6/1000$\\\bottomrule
\end{tabular}
\end{center}
For a row $[l,h]$, put
\[
 \delta_- =\max(1/10,1-h^2),\quad \delta_+=1-l^2,
 \quad L_\pm=(N_0+2)\delta_\pm/2.
\]
The three whole-interval upper bounds are
\[
 \min\left(\frac{28}{N_0l^2},\frac{7L_+}{N_0\log L_+}\right),\quad
 \frac53h^2N_0\exp\left(-\frac{N_0-1}{N_0+2}
                          \frac{\log L_-}{\delta_+}\right),\quad
 \exp\left(-\frac{N_0}{N_0+2}\frac{\log L_-}{\delta_+}\right).
\]
At $l=0$ omit the first alternative.  Their sum is below the displayed
rational upper bound in each row. These seven inequalities are checked
by the rational logarithm and exponential formulas in
Section~\ref{sec:certificates}; the complete checker is
\cert{verify_actual_roots.py}.  Every upper bound is smaller than
$1-h$, and $(19/20)^2>9/10$. Thus
$D_N+U_N+V_N<1-a$, contradicting \eqref{ar:direct-budget}.

At $a=0$, factorization gives
$\prod_j|q_j|=n/\prod_j|z_j|\ge n$, hence $s>1$.
At $a=1$, the identity
\[
 \sum_j\frac1{1-\xi_j}=2\sum_j\frac1{1-z_j}
\]
has real part at least $N$, since each summand on the right has real
part at least $1/2$. Cauchy--Schwarz gives $E_a\ge N$.
Rotation treats all boundary zeros and proves the proposition.
\end{proof}

\subsection{A variance bound}

\begin{proposition}\label{ar:variance}
Let $p$ have degree $n\ge2048$ and zeros $z_1,\ldots,z_n$, with
\[
 \bar z=\frac1n\sum_i z_i,\qquad
 \sigma_2(p)^2=\frac1n\sum_i|z_i-\bar z|^2.
\]
Define $d(a)=\min_{p'(\xi)=0}|a-\xi|$.
For every zero $a$, one has $d(a)\le\sigma_2(p)$.
In particular, if all zeros lie in $\overline\D$, then
\begin{equation}\label{ar:variance-unit}
 d(a)^2\le1-\frac{1-|a|^2}{n}.
\end{equation}
\end{proposition}

\begin{proof}
The cases $\sigma_2(p)=0$ and a multiple zero are immediate.
Translation, scaling, rotation, and multiplication by a nonzero
constant reduce a contrary example to a monic polynomial with
\[
 \sum_i z_i=0,\quad \sum_i|z_i|^2=n,\quad z_1=a\ge0,
 \quad |a-\xi_j|>1\quad(1\le j\le N).
\]
Let $D=\diag(z_i)$, $e=n^{-1/2}(1,\ldots,1)^T$, and let $U$ be an
isometry onto $e^\perp$. The characteristic polynomial of $U^*DU$ is
$p'/n$: its determinant is the $e$-cofactor of $zI-D$, namely
$p(z)e^*(zI-D)^{-1}e=p'(z)/n$. The identity extends across the zeros
by polynomial continuation. Schur triangularization and the Frobenius
norm therefore give
\[
 \sum_j|\xi_j|^2\le\|U^*DU\|_F^2
 =\sum_i|z_i|^2-2\|De\|^2+|e^*De|^2=n-2.
\]
The critical centroid is zero, since
$N^{-1}\sum_j\xi_j=n^{-1}\sum_i z_i=0$. Thus, writing
$\kappa=N^{-1}\sum_j|\xi_j|^2$, we have
$\kappa\le(N-1)/N<1$, and $a>0$.
The original centroid also implies $a^2\le N$.

Set
\[
 q_j=(a-\xi_j)^{-1},\quad \mu=\E q,\quad
 s=\E|q|^2<1,\quad v=s-|\mu|^2,\quad
 \Pi=\prod_j|q_j|,\quad h=1-\sqrt{1-s}.
\]
Here $\E$ denotes the average over all $N$ critical labels. Since
$\E(q^{-1})=a$, $\E|q|^{-2}=a^2+\kappa$, and $\E\xi=0$,
\begin{equation}\label{ar:variance-covariance}
 |1-a\mu|^2\le\kappa v.
\end{equation}
Indeed $a\mu-1=\E[\xi(q-\mu)]$.  Applying weighted Cauchy--Schwarz
to $q^{-1}-\bar q=(1-|q|^2)/q$ also gives
\begin{equation}\label{ar:variance-weighted}
 |a-\mu|^2\le(1-s)(a^2+\kappa-1)
 \le(1-s)(a^2-1/N).
\end{equation}
Consequently $\Re\mu\ge ah$ and
$|\mu|\le a(1+\sqrt{1-s})$.

The product factorizations and the arithmetic--geometric mean
inequality give
\[
 \Pi\le s^{N/2},\quad
 |p(0)|\le a[(n-a^2)/N]^{N/2}\le\min(1,\sqrt e\,a),
\]
\[
 |p'(a)|=n/\Pi,\qquad
 \sum_{i\ne1}|a-z_i|^2=n(1+a^2).
\]
In particular,
\[
 \frac n{s^{N/2}}\le\frac n\Pi=|p'(a)|
 \le\left(\frac{n(1+a^2)}N\right)^{N/2}.
\]
It follows that
\begin{equation}\label{ar:variance-moment-range}
 1\le\frac Nn n^{2/N}\le(1+a^2)s
 \le2(1+\sqrt{1-s}),
 \qquad
 a^2>\frac{2\log n-1}{n}.
\end{equation}
For the upper bound use $ah\le|\mu|\le\sqrt s$ and
$s/h=1+\sqrt{1-s}$. Since $s<1$ and $e^u\ge1+u$,
\[
 1+a^2>\frac Nn n^{2/N}
 \ge\frac{N+2\log n}{n}
 =1+\frac{2\log n-1}{n}.
\]
In particular $M=Na^2>14$ for $n\ge2048$.

Put
\[
 F(t)=\prod_j(1-atq_j),\quad
 \beta(t)=1-2a\Re\mu\,t+a^2st^2,\quad
 \gamma(t)=s\bigl(a^2(1-t)^2+\kappa\bigr).
\]
On the entire unit interval,
\begin{equation}\label{ar:variance-beta}
 \beta(t)\le1-a^2h^2t,
 \qquad \beta(t)\le1-a^2st(1-t).
\end{equation}
Integration from $a$ to $0$ and the differentiation identity used above
give
\begin{equation}\label{ar:variance-direct}
 1\le|F(1)|+\frac Nn|\mu|\Pi|p(0)|+|R|,
\end{equation}
where $c_N=(N/(N-1))^{(N-1)/2}<5/3$ and
\begin{align}
 |F(1)|&\le\min(\beta(1),\kappa s)^{N/2},\notag\\
 |R|&\le c_NNa^2s\int_0^1t\min(\beta(t),\gamma(t))^{(N-1)/2}\,dt.
 \label{ar:variance-remainder}
\end{align}
The $\gamma$ bound follows by rewriting the $j$th differentiated term
as
\[
 q_j^2\prod_{k\ne j}(1-atq_k)
 =q_j\Bigl(\prod_kq_k\Bigr)
       \prod_{k\ne j}(a(1-t)-\xi_k)
\]
and using the critical centroid and second moment before applying the
arithmetic--geometric mean inequality.

We also need the centered identity. Write $D(t)=1-at\mu$ and
$\omega_j=a(\mu-q_j)$. Then $\sum\omega_j=0$ and
$\sum|\omega_j|^2=Na^2v$. Comparing $F$ to $D^N$ gives
\begin{equation}\label{ar:variance-centered}
 1\le |\mu|\Pi|p(0)|+|D(1)|^n+E_c,
 \qquad E_c=na|\mu|\int_0^1|F(t)-D(t)^N|\,dt.
\end{equation}
The interpolation $\prod_j(D+t\theta\omega_j)$ gives both bounds
\begin{align}
 |F-D^N|&\le\frac{c_NNa^2t^2v}{2|D|}\beta^{(N-1)/2}
                  &&(D\ne0),\label{ar:center-first}\\
 |F-D^N|&\le\frac{c_{N,2}N(N-1)a^2t^2v}{2}\beta^{(N-2)/2},
 \quad c_{N,2}=\left(\frac N{N-2}\right)^{(N-2)/2}<3.
 \label{ar:center-second}
\end{align}
For the second estimate differentiate twice in $\theta$, use
$\sum_{j\ne k}|\omega_j\omega_k|\le N(N-1)a^2v$, and integrate
the Taylor remainder with weight $1-\theta$. Thus it remains valid
when $D=0$.

We divide the parameter range into three cases. The numerical endpoint
bounds and their degree monotonicities are checked at the end of the proof.

\smallskip
\noindent\emph{Case 1: $s\le3/4$.}
Equation \eqref{ar:variance-moment-range} gives $a^2s\ge1/4$.
On $[0,3/4]$, \eqref{ar:variance-beta} gives
$\beta(t)\le e^{-a^2st/4}$, so
\[
 |R|_{[0,3/4]}\le\frac{64c_NN}{(N-1)^2a^2s}
 <\frac{432}{N}<\frac14.
\]
On $[3/4,1]$,
$\gamma(t)\le15s/16+2(1+\sqrt{1-s})/16\le57/64$.
Since $a^2s\le4$,
\[
 |R|_{[3/4,1]}\le\frac53N(57/64)^{(N-1)/2}
 <\frac1{16}.
\]
Each of the other two terms in \eqref{ar:variance-direct} is less than
$(3/4)^{N/2}<1/16$. Their sum is less than $7/16$, a contradiction.

\smallskip
\noindent\emph{Case 2: $a\le1/4$.}
Now $s\ge16/17$, $h>3/4$, and $|\mu|\le5a/4$.
The product term in \eqref{ar:variance-direct} is at most
$25a^2/12\le25/192$. On $[0,1/2]$ use
$\beta(t)\le1-a^2t$; on $[1/2,1]$ use convexity and
$\beta\le1-a^2/2$. These give
\[
 |R|_{[0,1/2]}<\frac{27}{4M}<\frac{27}{56},
 \qquad
 |R|_{[1/2,1]}\le\frac58M e^{-(N-1)M/(4N)}<\frac{35}{128}.
\]
Also $|F(1)|\le e^{-9M/32}<1/48$. Thus the right side of
\eqref{ar:variance-direct} is at most
\[
 \frac{25}{192}+\frac{27}{56}+\frac{35}{128}+\frac1{48}
 =\frac{2437}{2688}<1.
\]

\smallskip
\noindent\emph{Case 3: $s>3/4$ and $a>1/4$.}
We have $h>1/2$, $a<2$, and $|\mu|<3a/2$.
Suppose first that $v>0$. From
\eqref{ar:variance-covariance}, $\beta(1)=|D(1)|^2+a^2v$, and
\eqref{ar:variance-beta},
\[
 |D(1)|^n\le v\beta(1)^{(n-2)/2}<v/32.
\]
For $1/4<a\le1/2$, we have $h\ge1-a$ and
$a^2h^3\ge27/1024$. On $[0,1/4]$ use
$|D|\ge7/8$ and $\beta\le e^{-3a^2ht/2}$ in
\eqref{ar:center-first}, obtaining
\[
 E_{c,[0,1/4]}/v
 \le\frac{512c_NnN(|\mu|/a)}{189(N-1)^3a^2h^3}
 <\frac{300}{N}<\frac14.
\]
At the two endpoints of $[1/4,1]$,
\[
 \beta(1/4)\le1-a^2(7/16-a/2)\le1-5/256,
 \qquad
 \beta(1)\le1-a^2(1-a)^2\le1-9/256.
\]
The two losses increase on $[1/4,1/2]$. Convexity gives
$\beta\le1-5/256$ on $[1/4,1]$, where $|D|\ge1/2$.
In \eqref{ar:center-first} the coefficient after multiplication by
$na|\mu|$ is at most
$c_NnNa^3|\mu|\le5nN/32<3n^2/16$. Hence
\[
 E_{c,[1/4,1]}/v\le\frac3{16}n^2e^{-5(n-2)/512}<\frac1{32}.
\]
For $a>1/2$, instead use $a^2h^3>1/32$, $|D|>1/2$ on
$[0,1/4]$, and $|\mu|/a<3/2$. This gives
\[
 E_{c,[0,1/4]}/v
 \le\frac{128c_NnN(|\mu|/a)}{27(N-1)^3a^2h^3}
 <\frac{400}{N}<\frac14.
\]
For $a>1/2$, the endpoint bounds are
\[
 \beta(1/4)\le1-3a^2h/8<1-3/64,\qquad
 \beta(1)\le1-a^2h^2<1-1/16.
\]
Convexity gives $\beta<1-3/64$ on $[1/4,1]$.
Since $a<2$, $|\mu|<1$, and $c_{N,2}<3$, the coefficient in
\eqref{ar:center-second} after multiplication by $na|\mu|$ is
less than $12nN(N-1)<12n^3$. Thus
\[
 E_{c,[1/4,1]}/v\le12n^3e^{-3(n-3)/128}<\frac1{32}.
\]
In either subcase $|D(1)|^n+E_c<5v/16$, whereas
\[
 |\mu|\Pi|p(0)|\le\sqrt{s-v}\le\sqrt{1-v}\le1-v/2.
\]
This contradicts \eqref{ar:variance-centered}.
If $v=0$, all $q_j=\mu$, and $a\mu=1$, $D(1)=0$, $F=D^N$.
The same identity becomes $1\le|\mu|^{N+1}<1$.

Here are elementary checks for the constants used in the three cases.
The inequality $\log2>56/81$ follows from the first two terms of
the positive atanh series. The expression
$(1-1/n)(2\log n-1)$ increases and exceeds $14$ at $n=2048$.
The bounds $e<25/9$ and $e>5/2$ follow from the exponential series.
For $N\ge2047$,
\[
 \frac{N^2}{(N-1)^2}<\frac{101}{100},\qquad
 \frac{(N+1)N^2}{(N-1)^3}<\frac{101}{100},
\]
and both ratios decrease. The two centered coefficients are bounded by
\[
 \frac{512}{189}\frac53\frac32\frac{101}{100}\frac{1024}{27}
 =\frac{6619136}{25515}<300,
 \qquad
 \frac{128}{27}\frac53\frac32\frac{101}{100}\,32
 =\frac{51712}{135}<400.
\]
Positive exponential sums through degrees $7$ and $12$ give
$e^{7161/2047}>32$ and $e^{63/16}>48$.
The remaining endpoint estimates follow from
$2^{17}>32\cdot2048$,
$(5/2)^{19}>6\cdot2048^2$, and
$2^{47}>384\cdot2048^3$.
Each function $n^ke^{-c(n-j)}$ used above decreases on $n\ge2048$,
as its logarithmic derivative is $k/n-c<0$ there.
Finally, in the unit-disk case,
\[
 \sigma_2(p)^2=\frac1n\sum_i|z_i|^2-|\bar z|^2
 \le\frac{n-1+|a|^2}{n},
\]
which proves \eqref{ar:variance-unit}.
\end{proof}

\section{Reciprocal matrices and local fields}\label{sec:matrix-fields}

Throughout this section $b$ is an actual simple rightmost zero.
The notation $d,S,H,Q$ is that of \eqref{pr:field-notation}.
The estimates are independent of the degree.

\subsection{A matrix envelope}

\begin{proposition}\label{mf:envelope}
Put $\eta=d\Re S$, $\tau=d\Im S$, and
\[
 \Delta=3\pi^2-8,\quad K_0=\frac{9\pi^2}{\Delta},\quad
 A_0=2+\frac{24\pi}{\Delta},\quad
 \gamma=\frac{40-3\pi^2}{2\Delta}.
\]
Then
\begin{equation}\label{mf:exact-envelope}
 Qd^2\le K_0+A_0\eta+\gamma\eta^2-\tau^2/2.
\end{equation}
In particular, in the normalization \eqref{pr:normalized-moments},
writing $\sigma=c+is$,
\begin{equation}\label{mf:coarse-envelope}
 q<\frac{21}{5}\mu^2+\frac{11}{2}\mu c+\frac{c^2}{4}-\frac{s^2}{2}.
\end{equation}
\end{proposition}

\begin{proof}
Translate and scale to $b=0,d=1$. Write $v_j=-1/\xi_j$ and let
$w_k=-1/u_k$, where the $u_k$ are all the other original zeros.
Set $M=I+\mathbf1\mathbf1^*$ and
$A=M^{1/2}\diag(w_k)M^{1/2}=K+iJ$, with $K,J$ Hermitian.
Since $\Re w_k\ge0$, we have $K\ge0$. The rank-one determinant
identity and differentiation of $p(z)=zp'(0)\prod_k(1+zw_k)$ give
\[
 \det(I+z\diag(w_k)M)=\frac{p'(z)}{p'(0)}
                    =\prod_j(1+zv_j).
\]
Thus the eigenvalues of $A$, including multiplicities, are the $v_j$.

For a Hermitian matrix $B$ define
$\mathcal T(B)_{ij}=i\operatorname{sgn}(i-j)B_{ij}$, with zero
diagonal. If $x$ is a unit vector and $b_k=|x_k|^2$, diagonal
unitary conjugation removes the coordinate phases and gives
\begin{equation}\label{mf:rankone-char}
 \det(tI-\mathcal T(xx^*))
 =\frac12\left\{\prod_k(t+ib_k)+\prod_k(t-ib_k)\right\}.
\end{equation}
Indeed $tI-\mathcal T(xx^*)$ is the midpoint of an upper and a
lower triangular matrix with diagonals $t+ib_k$ and $t-ib_k$;
their difference has rank one. The determinant is affine along
a rank-one line. Zero coordinates follow by continuity.
After omitting those coordinates, the positive roots in
\eqref{mf:rankone-char} are the solutions of
$\sum_k\arctan(b_k/t)=(2j-1)\pi/2$. The left side strictly
decreases from $p\pi/2$ to zero, where $p$ is the number of
nonzero coordinates. Since $\arctan u\le u$ and $\sum b_k=1$,
the $j$th positive eigenvalue is at most
$L_j=2/((2j-1)\pi)$. The spectrum is symmetric about zero.

Let $\eta=\tr K$ and $k=\tr K^2$. The trace pairing
$\tr(K\mathcal T(xx^*))=-x^*\mathcal T(K)x$, followed by
diagonalization of either sign of $\mathcal T(xx^*)$, shows that
for every $\lambda>0$,
\begin{equation}\label{mf:schur-pair}
 \|\mathcal T(K)\|
 \le\lambda\eta+\sqrt{k}
       \left(\sum_{j\ge1}(L_j-\lambda)_+^2\right)^{1/2}.
\end{equation}
Here the diagonal entries of the conjugated $K$ are nonnegative,
sum to $\eta$, and have sum of squares at most $k$. This proves
the inequality by Cauchy--Schwarz and includes the whole spectrum.

In an upper triangular Schur basis for $A$,
$J=\diag(\Im v_j)+\mathcal T(K)$, whence
$\|J\|\le1+\|\mathcal T(K)\|$. Write $w_j=a_j+ib_j$ and
$\beta=\sum b_j=(\Im S)/2$. Direct multiplication gives
\[
 k=3\sum a_j^2+\Bigl(\sum a_j\Bigr)^2,\quad
 \tr J^2=3\sum b_j^2+\beta^2,\quad
 \|Je\|^2=\frac{N+1}{N}\left(\sum b_j^2+\beta^2\right),
 \qquad e=N^{-1/2}\mathbf1.
\]
In particular,
\[
 \tr J^2=\frac{3N}{N+1}\|Je\|^2-2\beta^2
 \le3\|J\|^2-2\beta^2.
\]
In a Schur basis the off-diagonal squared contributions to
$\tr K^2$ and $\tr J^2$ agree, so
\[
 \tr K^2-\tr J^2
 =\sum_j\bigl((\Re v_j)^2-(\Im v_j)^2\bigr).
\]
Also $0\le\Re v_j\le1$, whence $\sum_j(\Re v_j)^2\le\eta$.
Therefore
\[
 Q=2\sum_j(\Re v_j)^2-k+\tr J^2
 \le2\eta+3(1+\|\mathcal T(K)\|)^2-k-2\beta^2.
\]
Take
\[
 \lambda=\frac{24\pi+(40-3\pi^2)\eta}
                   {6\pi(3\pi+4\eta)},\qquad
 a_\lambda^2=(2/\pi-\lambda)^2+(2/(3\pi)-\lambda)^2.
\]
For $\eta\ge0$, one has $L_3<\lambda<L_1$ and
$3a_\lambda^2<1$. These follow by multiplying the positive
denominators; the endpoint values of the fractional linear
function $\lambda(\eta)$ suffice. Thus the sum in
\eqref{mf:schur-pair} is at most $a_\lambda^2$, including when
$\lambda>L_2$. Maximizing the resulting quadratic in $\sqrt k$
gives
\[
 3(1+\|\mathcal T(K)\|)^2-k
 \le\frac{3(1+\lambda\eta)^2}{1-3a_\lambda^2}.
\]
Put
\[
 R_\eta=18\pi^2+48\pi\eta+(40-3\pi^2)\eta^2.
\]
Direct substitution gives
\[
 1+\lambda\eta=\frac{R_\eta}{6\pi(3\pi+4\eta)},\qquad
 1-3a_\lambda^2
 =\frac{\Delta R_\eta}{6\pi^2(3\pi+4\eta)^2}.
\]
Consequently
\[
 2\eta+\frac{3(1+\lambda\eta)^2}{1-3a_\lambda^2}
 =2\eta+\frac{R_\eta}{2\Delta}
 =K_0+A_0\eta+\gamma\eta^2.
\]
This proves \eqref{mf:exact-envelope}. Finally $\pi>157/50$
implies $K_0<21/5$, $A_0<11/2$, and $\gamma<1/4$.
Divide by $(Hd)^2$ to obtain \eqref{mf:coarse-envelope}.
\end{proof}

\subsection{Logarithmic bounds and saddle moments}

The next inequality controls the logarithm on a disk.

\begin{lemma}\label{mf:log-metric}
For $0<r<1$ put
\[
 h(r)=\frac{r-\log(1+r)}{r^2},\qquad
 j(r)=\frac{-r-\log(1-r)}{r^2}.
\]
For every $z=x+iy$ with $|z|\le r$,
\[
 \log|1+z|\le x-h(r)x^2+j(r)y^2.
\]
\end{lemma}
\begin{proof}
The analytic identity
$(z-\Log(1+z))/z^2=\int_0^1t/(1+tz)\,dt$ holds in $|z|<1$.
For fixed $t$, the image of $|z|\le r$ under $t/(1+tz)$ lies in
the disk with center $t/(1-t^2r^2)$ and radius
$t^2r/(1-t^2r^2)$. The difference of squared radius and squared
distance from the center is
\[
 \frac{t^4(r^2-x^2-y^2)}
 {(1-t^2r^2)((1+tx)^2+t^2y^2)}\ge0.
\]
Integration therefore puts the preceding integral in the disk
with center $c_r=-\log(1-r^2)/(2r^2)$ and radius
$e_r=(\operatorname{atanh}r-r)/r^2$. Multiplication by $z^2$
and taking real parts gives
$\Re\Log(1+z)\le x-(c_r-e_r)x^2+(c_r+e_r)y^2$.
The two coefficients are $h(r)$ and $j(r)$.
\end{proof}

\begin{lemma}\label{mf:saddle-dual}
Let $\Re v_j\ge0$, $|v_j|\le1$, $\eta=\sum\Re v_j$, and
$Q=\sum|v_j|^2>0$. Suppose $z_0=-X+iY$, $|z_0|=s<1$,
$X\ge0$, and $\sum_jv_j/(1+z_0v_j)=0$. Then
\[
 \frac{\eta}{Q}\ge
 \frac{X(1-s)^2}{(1+|Y|)^2-X^2}.
\]
\end{lemma}
\begin{proof}
Conjugate to make $Y\ge0$, and put
$B=X/(1+2Y+s^2)$, $A=(1-2BX)/(1-s)^2$.
For $v=a+ib$, $r=|v|\le1$, and $d_v=|1+z_0v|^2$, expansion gives
\[
 (Aa-Br^2)d_v-(a-Xr^2)
 =a(Ad_v-1+2BXr^2)
   +Br^2[2Y(1+b)+s^2(1-r^2)].
\]
The second summand is nonnegative. Also $d_v\ge(1-sr)^2$ and
$2BX\le2s^2/(1+s^2)\le s<1$. The function
$A(1-sr)^2+2BXr^2$ decreases on $[0,1]$ to its value one
at $r=1$, since $As(1-s)-2BX=(s-2BX)/(1-s)\ge0$.
Thus $\Re(v/(1+z_0v))\le Aa-B|v|^2$. Sum and use the saddle
equation to obtain $\eta/Q\ge B/A$, the stated bound.
\end{proof}

\begin{lemma}\label{mf:saddle-variance}
Let $\Re v_j\ge0$, $|v_j|\le1$, $Q=\sum|v_j|^2>0$, and
$\sum_jv_j/(1+z_0v_j)=0$, with $\Re z_0\le0$.
Write $u_j=v_j/(1+z_0v_j)$ and $Q_u=\sum|u_j|^2$.
Suppose $|z_0|<s<1$, $-\Re z_0<X$, $|\Im z_0|<Y$,
and $X<1-Y$. Here $s$ is an upper bound for the modulus.
Put $U=(1-s)^{-1}$.
Then
\[
 \frac1{1+2Y+s^2}\le\frac{Q_u}{Q}
 \le q_m:=\min\left(U^2,\frac1{1-2X^2-2Y}\right)
\]
when the last denominator is positive, and
\[
 \frac{\sum(\Re u_j)^2}{Q_u}\le\beta
 :=\min\left\{
 X\left(\frac{1+X}{1-s^2}+XU^2\right),
 \frac{X(1-Y+X)}{(1-Y)^2+X^2}\right\}.
\]
\end{lemma}
\begin{proof}
Write $z_0=x_0+iy_0$, $v_j=a_j+ib_j$, and
$d_j=|1+z_0v_j|^2$. Then
\[
 \Re u_j=\frac{a_j+x_0|v_j|^2}{d_j},\quad
 \Im u_j=\frac{b_j-y_0|v_j|^2}{d_j},\quad
 |u_j|^2=\frac{|v_j|^2}{d_j}.
\]
Since $a_j\ge0$, $x_0\le0$, and $|v_j|\le1$,
$d_j\le1+2Y+s^2$. The zero real sum gives
$\sum a_j/d_j=-x_0Q_u$. In the exact identity
\[
 Q=Q_u+2x_0\sum a_j|u_j|^2
       -2y_0\sum b_j|u_j|^2
       +|z_0|^2\sum|v_j|^2|u_j|^2
\]
we have $0\le\sum a_j|u_j|^2\le-x_0Q_u$ and
$|\sum b_j|u_j|^2|\le Q_u$. This proves the mass bounds.

The inverse disk inequality
$(1-|z_0|^2)|u_j|^2+2\Re(z_0u_j)\le1$ gives
$\Re u_j\le(1+X)/(1-s^2)$, while
$\Re u_j\ge-X|u_j|^2\ge-XU^2$.
Split the zero real sum into its positive and negative terms.
The total absolute negative part is at most $XQ_u$. Multiplying
each part by its respective amplitude bound proves the first
variance bound.

For the second, write $z_0=-X_0+iY_0$. Give
$r_j=1/\overline{v_j}$ weights $|u_j|^2/Q_u$, omitting a
zero $v_j$ if necessary. Since
$1/\overline{u_j}=r_j+\overline{z_0}$ and $\sum u_j=0$,
their weighted mean is $X_0+iY_0$. Each $r_j$ has nonnegative
real part and modulus at least one. Put $L=1-|Y_0|$ and
$c=X_0/L<1$. The numbers
$r'_j=(r_j-iY_0)/L=a_j'+ib_j'$ have $a_j'\ge0$,
$a_j'^2+b_j'^2\ge1$, and weighted mean real part $c$.
On this unbounded domain,
\[
 \frac{(a-c)^2}{(a-c)^2+b^2}\le\frac{a+c^2}{1+c^2}.
\]
For $a\ge1$ the right side is at least one. For $0\le a\le1$
replace $b^2$ by $1-a^2$ and use
\[
 (a+c^2)(1-2ac+c^2)-(1+c^2)(a-c)^2
 =a(1-a)(1+c)^2\ge0.
\]
The denominator cannot vanish. Taking weighted means proves
$\sum(\Re u_j)^2/Q_u\le c(1+c)/(1+c^2)$.
This function increases on $[0,1]$, with derivative
$(1+2c-c^2)/(1+c^2)^2>0$. Substituting $c\le X/(1-Y)$
proves the second bound.
\end{proof}

\subsection{An affine reciprocal-field estimate}

\begin{proposition}\label{mf:affine}
For every actual simple rightmost zero,
\begin{equation}\label{mf:affine-bound}
 Qd^2\le20+11Hd.
\end{equation}
Consequently $Qd^2\le13(1+Hd)$ whenever $Hd\ge7/2$.
\end{proposition}

\begin{proof}
Normalize to $b=0,d=1$, and suppose $Q>20+11H$.
Set $B=-\sum v_j^2$, $\eta=\Re S$.
The matrix envelope implies $\eta>0$. Conjugate so that
$\Im S\ge0$, and write
\[
 t=\frac{\Im S}{\eta},\quad
 x=\frac{\Re B}{Q},\quad y=\frac{|\Im B|}{Q}.
\]
We have $B-Q=-2\sum(\Re v_j)v_j$, hence
$|B-Q|\le2\eta$ and $|B|\le Q$. Also
$H/Q<1/11$. The envelope gives
$t^2<2\gamma<(7/10)^2$: if its quadratic coefficient were
nonpositive, its constant and linear coefficients would already
be smaller than $20+11H$. Thus
\begin{equation}\label{mf:affine-domain}
 (t,y)\in[0,7/10]\times[0,1/4],\quad
 (x-1)^2+y^2<\frac4{121(1+t^2)},\quad x^2+y^2\le1.
\end{equation}
Write $B/Q=x+i\epsilon y$, with $\epsilon\in\{-1,1\}$.
Both signs are used when $y=0$.

Here and below directed rational rounding is denoted by
$\lfloor u\rfloor_m=\lfloor mu\rfloor/m$ and
$\lceil u\rceil_m=\lceil mu\rceil/m$. Write
$\operatorname{sf}_m(u)=\lfloor\sqrt u\rfloor_m$ and
$\operatorname{sc}_m(u)=\lceil\sqrt u\rceil_m$; these are computed
by integer square roots. Unless indicated, square-root meshes
have $m=10^6$. Use $p_0=6283/2000<\pi<22/7$ and put
$k_U=9p_0^2/(3p_0^2-8)$,
$a_U=2+24p_0/(3p_0^2-8)$,
$g_U=(40-3p_0^2)/(2(3p_0^2-8))$.
These are upper bounds for $K_0,A_0,\gamma$, since each decreases
on the indicated interval for $\pi$.

Each row of the certificate consists of six rationals
$(t_l,t_h,y_l,y_h,\rho,k)$, with $\rho>0$, $k>0$, and the
phase sign $\epsilon$.
They specify a closed rectangle and a contour. The following
construction gives every bound used to check that row.
Put
\[
 d_l=\operatorname{sf}(1+t_l^2),\quad
 d_h=\operatorname{sf}(1+t_h^2),\quad
 e_0=\frac2{11d_l},\quad
 g_0=g_U-t_l^2/2.
\]
If $y_l>e_0$ the cell is impossible. Otherwise define
$r_0=1-\operatorname{sc}(e_0^2-y_l^2)$.
If $r_0^2+y_l^2\ge1$ or $g_0\le0$ the cell is also impossible
under \eqref{mf:affine-domain}. In a retained cell
cap $y_h$ by
$\widehat y_h=\min(y_h,e_0,\operatorname{sc}(1-r_0^2))$.
The retained rows have $r_0>\widehat y_h\ge y_l$ and $g_0>0$.
Let $h_0$ be the largest $1/100$ mesh point strictly below the
positive root of
\[
 P(h)=k_U-20+(a_U-11d_l)h+g_0h^2.
\]
The checks $h_0>0$, $P(h_0)<0$, and $P(h_0+1/100)\ge0$
imply $\eta>h_0$. Set
\[
 Q_0=20+11d_lh_0,\quad r_Q=\operatorname{sf}_{10^4}(Q_0),\quad
 b_m=\operatorname{sf}(r_0^2+y_l^2).
\]
Choose the first $10^{-6}$ mesh point $a\in[1,2]$ with
\begin{equation}\label{mf:saddle-rouche}
 (ab_m-1)(11-a)>a^2.
\end{equation}
Put $s=a/11$, $U=(1-s)^{-1}$, $e=a^2/[b_m(11-a)]$.
For $L=p''/p'$ the geometric series gives
\[
 L(z)=S+Bz+E(z),\qquad |E(z)|\le\frac{Q|z|^2}{1-|z|}.
\]
On $|z|=aH/Q<s$, \eqref{mf:saddle-rouche} gives
$|Bz|>|S+E(z)|$. Rouch\'e's theorem gives an actual zero $z_0$
of $p''$ with $|z_0|<aH/Q$ and
\[
 |z_0+S/B|<(H/Q)e,\qquad \Re z_0\le0.
\]
The latter assertion is Gauss--Lucas twice; $p'$ is nonzero
in the whole disk $|z|<1$.

For phase bounds set
\[
 q_U=1/d_l,\ c_U=t_h/d_h,\quad
 q_L=1/\operatorname{sc}(1+t_h^2),\
 c_L=t_l/\operatorname{sc}(1+t_l^2),
\]
\[
 R_f=\frac{r_0}{r_0^2+y_l^2},\quad
 Y_f=\frac{\widehat y_h}{r_0^2+\widehat y_h^2},\quad
 R_L=\operatorname{sf}(1-\widehat y_h^2).
\]
On the unit-disk fiber, differentiation in $x>y$ gives
$R_L\le x/(x^2+y^2)\le R_f$ and
$y_l\le y/(x^2+y^2)\le Y_f$. For $\epsilon=1$ take
\[
 D_c=q_UR_f+c_UY_f+e,\qquad
 I_c=\max(0,c_UR_f-q_Ly_l,q_UY_f-c_LR_L)+e;
\]
for $\epsilon=-1$ take
\[
 D_c=\max(0,q_UR_f-c_Ly_l)+e,\qquad
 I_c=c_UR_f+q_UY_f+e.
\]
These follow by multiplying $S/B$ explicitly, and give
$-\Re z_0<(H/Q)D_c$, $|\Im z_0|<(H/Q)I_c$.
Initially put $X=\min(s,D_c/11)$, $Y=\min(s,I_c/11)$.
Lemma~\ref{mf:saddle-dual} permits replacing $X$ by its minimum
with
\[
 \frac{\operatorname{sc}((1-s)^4+4h_*^2(1+Y)^2)-(1-s)^2}
      {2h_*},\qquad h_*=\frac1{11d_l}.
\]
Apply Lemma~\ref{mf:saddle-variance} and denote its upper
mass and variance bounds by $q_m,\beta$. Each row checks
$1-2X^2-2Y>0$, $X<1-Y$, and $\beta<1/2$.

For $u_j=v_j/(1+z_0v_j)$ let $B_0=-\sum u_j^2$.
The geometric series gives
$|B_0-B|/Q<U^2-1=:p_*$. The mass and variance bounds give
\[
 \Re(B_0/Q)\ge\frac{1-2\beta}{1+2Y+s^2},\quad
 |\Im(B_0/Q)|\le2q_m\sqrt{\beta(1-\beta)},\quad
 |B_0/Q|\le q_m.
\]
Use the rational enclosures
\[
 R=\left\lfloor\max\left(r_0-p_*,
             \frac{1-2\beta}{1+2Y+s^2}\right)\right\rfloor_{10^4},
\]
\[
 I=\left\lceil\min(\widehat y_h+p_*,
             2q_m\operatorname{sc}(\beta(1-\beta)))\right\rceil_{10^4},
 \quad M=\lceil q_m\rceil_{10^4},\quad
 G=\operatorname{sf}_{10^4}(2p_0/M).
\]
Thus $R>0$, and
\[
 \left|\int_{-\infty}^{\infty}e^{-B_0t^2/(2Q)}\,dt\right|
 =\sqrt{\frac{2\pi}{|B_0/Q|}}\ge G.
\]

\smallskip
\noindent\emph{The contour.}
Define
\[
 \lambda=\left\lfloor\frac{\rho}{k/r_Q+D_c/11}\right\rfloor_{10^4},
 \quad T_0=\lfloor\rho r_Q\rfloor_{10^4},\quad
 \delta=\operatorname{sc}_{10^5}(\rho^2(1+\lambda^{-2})).
\]
Check $\lambda>1$, $T_0>0$, $U\delta<1$. Put
\[
 T=\max\{T_0,\lambda[k+(H/\sqrt Q)D_c]\},\qquad
 D_c'=k-\sqrt Q\Re z_0.
\]
Then $0<D_c'<T/\lambda$ and $T/\sqrt Q<\rho$.
Indeed $Q>r_Q^2$ and $H/Q<1/11$, so
\begin{align*}
 \frac{T}{\sqrt Q}
 &=\max\left\{\frac{T_0}{\sqrt Q},
       \lambda\left(\frac{k}{\sqrt Q}+\frac{H}{Q}D_c\right)\right\}\\
 &<\max\left\{\frac{T_0}{r_Q},
       \lambda\left(\frac{k}{r_Q}+\frac{D_c}{11}\right)\right\}
 \le\rho.
\end{align*}
Take the vertical segment $z_0+it/\sqrt Q$, $-T\le t\le T$,
and both horizontal segments joining its ends to $\Re z=k/\sqrt Q$.
Relative to $z_0$ every point has modulus less than $\delta$.

Use the directed logarithm coefficients
\begin{equation}\label{mf:directed-log}
 \underline h(r)=\left\lfloor\sum_{\ell=0}^{15}
                      \frac{(-r)^\ell}{\ell+2}\right\rfloor_{10^6},
 \qquad
 \overline j(r)=\left\lceil\sum_{\ell=0}^{15}
          \frac{r^\ell}{\ell+2}+\frac{r^{16}}{18(1-r)}
                       \right\rceil_{10^6}.
\end{equation}
The alternating and geometric remainders prove
$\underline h\le h$ and $\overline j\ge j$.
Replace the two coefficients in Lemma~\ref{mf:log-metric}
before expanding its quadratic form. At $r=U\rho$ and
$r=U\delta$ write the respective pairs as $(h_v,j_v)$ and
$(h_c,j_c)$. Set
\[
 q_l=\lfloor(1+2Y+s^2)^{-1}\rfloor_{10^6},\quad
 A_v=\lfloor h_v-(h_v+j_v)\beta\rfloor_{10^6},
\]
\[
 A_c=\lfloor h_c-(h_c+j_c)\beta\rfloor_{10^6},\quad
 B_c=\lceil j_c-(h_c+j_c)\beta\rceil_{10^6}.
\]
Every row has $q_l,A_v>0$ and $B_c\ge0$.

\smallskip
\noindent\emph{The paired vertical integral.}
For $|t|\le T$ put
$f(t)=\sum_j\Log(1+iu_jt/\sqrt Q)$ and
$g(t)=-B_0t^2/(2Q)$. For a symmetric shell
$a_j\le|t|\le b_j$, put
$\rho_j=\min(\rho,b_j/r_Q)$, taking $\rho_j=\rho$ if $b_j=\infty$,
and $r_j=U\rho_j$. With $h_j=\underline h(r_j)$ and
$j_j=\overline j(r_j)$, set
\[
 \alpha_j=\left\lfloor\min\{R/2,
                q_l[h_j-(h_j+j_j)\beta]\}\right\rfloor_{10^5}>0.
\]
Both $\Re f(\pm t)$ and $\Re g(t)$ are at most $-\alpha_jt^2$.
Writing the even and odd parts of $f-g$ as $E(t),O(t)$,
the logarithmic series gives
\[
 |O(t)|\le\frac{q_mU|t|^3}{\sqrt Q}
       \left(\frac13+\frac{r_j^2}{5(1-r_j^2)}\right),\quad
 |E(t)|\le\frac{q_mU^2t^4}{Q}
       \left(\frac14+\frac{r_j^2}{6(1-r_j^2)}\right).
\]
These bounds retain all odd orders at least five and all even
orders at least six. For arbitrary complex $A,O$,
\[
 e^{A+O}+e^{A-O}-2e^A
 =O^2\int_{-1}^1(1-|s|)e^{A+sO}\,ds.
\]
Apply this with $A=g+E$, and apply the first-difference integral
to $e^{g+E}-e^g$. The real-part bounds hold along both line
segments, so
\[
 |e^{f(t)}+e^{f(-t)}-2e^{g(t)}|
 \le(|O(t)|^2+2|E(t)|)e^{-\alpha_jt^2}.
\]
Use the shell endpoints $0,1,2,3,4,5,6,8,10,12,\infty$.
If $\Delta M_{k,j}$ is an upper bound for
$\int_{a_j}^{b_j}t^ke^{-\alpha_jt^2}\,dt$, the
paired error is at most
\[
 \mathcal E=\left\lceil\sum_j
 \frac{q_m^2U^2c_{o,j}^2\Delta M_{6,j}
          +2q_mU^2c_{e,j}\Delta M_{4,j}}{r_Q^2}
                      \right\rceil_{10^4},
\]
where $c_{o,j}$ and $c_{e,j}$ are the parenthesized coefficients.
Extending a shell past the actual endpoint $T$ only enlarges its
positive error integral.

Here the moments have a purely rational verification. For
$M_k(a)=\int_a^\infty t^ke^{-\alpha t^2}\,dt$,
\[
 M_k(a)=\frac{a^{k-1}e^{-\alpha a^2}}{2\alpha}
          +\frac{k-1}{2\alpha}M_{k-2}(a).
\]
For $a>0$,
$ae^{-\alpha a^2}/(1+2\alpha a^2)\le M_0(a)
 \le e^{-\alpha a^2}/(2\alpha a)$.
Differentiating these proposed bounds and integrating to infinity
proves the two signs. At zero,
$M_4(0)=3\sqrt\pi/(8\alpha^{5/2})$ and
$M_6(0)=15\sqrt\pi/(16\alpha^{7/2})$.
Use directed square roots and $\pi<22/7$.
For each finite endpoint $x=\alpha a^2<82$,
\[
 E_{40}(x)\le e^x\le E_{80}(x)
               +\frac{x^{81}}{81!(1-x/82)},\qquad
 E_m(x)=\sum_{\ell=0}^m\frac{x^\ell}{\ell!}.
\]
Taking reciprocals in the correct directions gives rational
moment bounds; subtract a lower right tail from an upper left
tail. The final infinite shell is included. The certificate
checks every domain and positivity condition in this construction.

\smallskip
\noindent\emph{The two horizontal segments.}
There are two valid bounds, either of which may be used on a row.
For the translated bound, expansion of Lemma~\ref{mf:log-metric}
at $w=(d+i\epsilon_c T)/\sqrt Q$, with
$\epsilon_c\in\{-1,1\}$ and $0\le d\le T/\lambda$, gives
\[
 \Re\sum\Log(1+u_jw)
 \le-q_l(A_c-B_c/\lambda^2)T^2
           +(h_c+j_c)IT^2/\lambda.
\]
Indeed, if $\nu=\sum(\Re u_j)^2/Q_u\le\beta$, the quadratic
part before substitution is
\[
 \frac{Q_u}{Q}\{-[h_c-(h_c+j_c)\nu]T^2
                  +[j_c-(h_c+j_c)\nu]d^2\},
\]
which increases with $\nu$ since $T^2\ge d^2$.
The mixed term has modulus at most $(h_c+j_c)IdT$.
When $A_c-B_c/\lambda^2>0$, the downward-rounded damping is
\[
 \kappa_{\rm tr}=
 \left\lfloor q_l(A_c-B_c/\lambda^2)
                  -(h_c+j_c)I/\lambda\right\rfloor_{10^5}.
\]

For the original-coordinate bound put
$l=1/\lambda$, $J=1+I_c/(\lambda D_c)$, and
$r_o=\operatorname{sc}_{10^5}(\rho^2(l^2+J^2))$.
If $J<2$ and $r_o<1$, take $h_o=\underline h(r_o)$,
$j_o=\overline j(r_o)$, and
$A_o=h_o-(h_o+j_o)/(11d_l)$.
Writing $z=x_c+iy_c$, $V=\sum(\Re v_j)^2\le\eta$ and
$W=\sum\Re v_j\Im v_j$, the exact expanded bound is
\[
 \log|p'(z)/p'(0)|
 \le\eta x_c-(\Im S)y_c
 +[j_oQ-(h_o+j_o)V]x_c^2
 +[(h_o+j_o)V-h_oQ]y_c^2+2(h_o+j_o)Wx_cy_c.
\]
On both segments,
$|x_c|\sqrt Q/T\le l$,
$|y_c|\sqrt Q/T\in[2-J,J]$,
$x_c\le k/\sqrt Q$, and $H/(\sqrt Q T)\le1/(\lambda D_c)$.
As $2|W|\le\widehat y_hQ$, the normalized upper bound is
\[
 j_ol^2+\frac{q_Uk}{\lambda D_cT_0}
 +\max_{2-J\le v\le J}(-A_ov^2+Lv),\qquad
 L=(h_o+j_o)\widehat y_hl+\frac{c_U}{\lambda D_c}.
\]
When $A_o>0$, the maximum is attained at an endpoint or at
$L/(2A_o)$ if that point belongs to the interval.
For normalization at the saddle, its exact equation gives
\[
 \ell_0=\sum_j\left[\Log(1+z_0v_j)
                      -\frac{z_0v_j}{1+z_0v_j}\right].
\]
The series has coefficient $(-1)^m(m-1)/m$ in degree
$m\ge2$. Hence $|\ell_0|\le Q|z_0|^2f(s)$, where
$f(s)=1/(1-s)-\sum_{j\ge0}s^j/(j+2)$.
Use
$f_s=\lceil1/(1-s)-\sum_{j=0}^{15}s^j/(j+2)\rceil_{10^6}$.
Since $Q|z_0|^2/T^2\le a^2/(\lambda^2D_c^2)$, a second
valid damping is
\[
 \kappa_{\rm or}=\left\lfloor
 -\max_{[2-J,J]}(-A_ov^2+Lv)-j_ol^2
 -\frac{q_Uk}{\lambda D_cT_0}
 -\frac{f_sa^2}{\lambda^2D_c^2}\right\rfloor_{10^5}.
\]
Choose the larger valid positive damping $\kappa$.
Every retained row has at least one.

\smallskip
\noindent\emph{The final finite predicates.}
Put
$F_0=\lfloor(G-\mathcal E-1/1000-1/100)/2\rfloor_{10^4}$
and $J_0=\lceil1+I_c/(\lambda D_c)\rceil_{10^4}$.
In addition to the conditions above, each row verifies
\begin{gather}
 F_0>0,\qquad
 E_2(RT_0^2/2)>\frac{2000}{RT_0},\qquad
 \kappa T_0^2>\frac12,\label{mf:affine-final-one}\\
 E_{24}(\kappa T_0^2)>\frac{200T_0}{\lambda},\qquad
 \kappa(k^2+J_0^2T_0^2)>J_0^2,\label{mf:affine-final-two}\\
 E_{24}(\kappa T_0^2)>
       \frac{11}{10}\frac{k^2+J_0^2T_0^2}{kF_0}.
 \label{mf:affine-final-three}
\end{gather}
The two Gaussian tails are bounded by
\[
 2\int_T^\infty e^{-Rt^2/2}\,dt
 \le\frac{2e^{-RT^2/2}}{RT}<\frac1{1000}.
\]
The sum of the moduli of the two horizontal integrals, divided by
$|p'(z_0)|/\sqrt Q$, is bounded by
\[
 \frac{2T}{\lambda}e^{-\kappa T^2}<\frac1{100}.
\]
The first bound decreases for $T>0$; the second decreases for
$\kappa T^2>1/2$. Thus the finite comparisons at $T_0$ prove both
bounds on the whole ray $T\ge T_0$.
The endpoint primitive values therefore differ by more than
$2F_0|p'(z_0)|/\sqrt Q$. At one endpoint
$z_{\epsilon_c}=k/\sqrt Q+i(\Im z_0+\epsilon_c T/\sqrt Q)$,
\[
 |p(z_{\epsilon_c})|>F_0|p'(z_0)|/\sqrt Q,\quad
 |p'(z_{\epsilon_c})|\le e^{-\kappa T^2}|p'(z_0)|,\quad
 \sqrt Q|\Im z_{\epsilon_c}|\le J_0T.
\]
But every original zero lies in $\Re z\le0$, and zero itself
is a simple zero. Therefore
$\Re(p'/p)(z)\ge\Re(1/z)$ in $\Re z>0$, which would imply
\[
 \frac{k}{k^2+J_0^2T^2}<\frac{e^{-\kappa T^2}}{F_0}.
\]
The last two predicates make
$(k^2+J_0^2T^2)e^{-\kappa T^2}$ decreasing for $T\ge T_0$
and already smaller than $(10/11)kF_0$ at $T_0$.
This is a contradiction.

The finite input consists of the four tables
\cert{C11_A20_PHASE_1_DIRECT_QUADRATIC_METRIC_CELLS.csv},
\cert{C11_A20_PHASE_-1_DIRECT_QUADRATIC_METRIC_CELLS.csv},
and their corresponding \cert{SKIPS.csv} tables.
They contain respectively $77,198$ retained rectangles and
$14,14$ excluded rectangles. On each sign they cover exactly
$[0,7/10]\times[0,1/4]$. This is checked by sorting every
rational vertical edge, merging the horizontal intervals on
each intervening open slab, and repeating the check on each
closed edge. All excluded rectangles satisfy the explicit
exclusion above. The program
\cert{affine/replay_affine_c11_a20.py} reconstructs every
quantity from the six rational entries and verifies
\eqref{mf:saddle-rouche}--\eqref{mf:affine-final-three}
with rational arithmetic, including all $2750$ shells.
This proves the required finite assertion and the proposition.
The final consequence follows from $20+11D\le13(1+D)$
for $D\ge7/2$.
\end{proof}

\subsection{Closing the unbounded field range}

The two signed-curvature bounds from Proposition~\ref{lt:quartic}
give the following estimate.

\begin{proposition}\label{mf:D270}
If $D=Hd\ge270$, then $Qd/H<9/4$.
\end{proposition}
\begin{proof}
Put $y=Qd/H=qD$. Proposition~\ref{mf:affine} implies
\[
 y\le13(1+1/D)\le Y_0:=3523/270,\qquad
 q\le13(D+1)/D^2\le3523/72900<97/2000.
\]
We also have $\mu=1/D\le1/270$.
For $0\le r\le97/2000$ define
\[
 v(r)=r+17r^2+137r^3+313r^4,\quad m(r)=1-v(r),
\]
\[
 A(r)=\frac{1+17r+137r^2+313r^3}{m(r)},\qquad
 B(r)=\frac{18+154r+450r^2+313r^3}{m(r)},
\]
\[
 C(r)=1+\left(2+\frac{100}{199}\right)rA(r)^2
 +3rB(r)+\frac{600}{199}r^2A(r)^2
 +4rA(r)+29r+436r^2.
\]
Let $K=p''/p'$. The coarse quartic row gives
\[
 \frac{\Re K'(b)}{H^2}
 <C(q)q^2+\frac{50000}{3}q^5+\frac{650}{9}q\mu.
\]
On the other hand, rightmostness gives
$\Re K'(b)=Q-2\sum(\Re v_j)^2\ge Q-2H/d$.
Consequently
\begin{equation}\label{mf:coarse-bootstrap}
 0<P_D(y):=y^2C(y/D)+\frac{50000}{3}\frac{y^5}{D^3}
                       +\frac{650}{9}y-D(y-2).
\end{equation}

On the whole interval $0\le r\le97/2000$, $v(r)<1$.
The convergent series $1/m=\sum_{k\ge0}v^k$ and all numerators
above have nonnegative coefficients. Thus $C,C',C''\ge0$ there.
For $3\le y\le Y_0$ and $D\ge270$,
\[
 \partial_y^2P_D
 =2C+\frac{4y}{D}C'+\frac{y^2}{D^2}C''
                      +\frac{1000000}{3}\frac{y^3}{D^3}>0,
\]
\[
 \partial_DP_D=-\frac{y^3}{D^2}C'
                    -50000\frac{y^5}{D^4}-(y-2)<0.
\]
All occurrences of $C$ and its derivatives have argument $y/D$.
Direct rational substitution gives
\[
 P_{270}(3)=
 -\frac{343077017913365252920829}{10130931628072339534650}<0,
 \qquad P_{270}(Y_0)<-197.
\]
Convexity bounds $P_{270}$ above by its endpoint chord.
Monotonicity in $D$ then contradicts
\eqref{mf:coarse-bootstrap} if $y\ge3$.

Now $y<3$, hence $q<1/90<7/625$.
The fine quartic row gives
$\Re K'(b)/H^2<(9/4)q^2+13q\mu$.
It follows that
$0<(9/4)y^2+13y-D(y-2)$.
At $D=270$ this convex quadratic has endpoint values
$-1719/64$ and $-843/4$ at $y=9/4$ and $y=3$.
It decreases in $D$ for every $y>2$. Thus $y\ge9/4$ is
impossible. All the rational endpoint inequalities, including
the two companion rows, are checked in
\cert{quartic/verify_d270.py}.
\end{proof}

\subsection{A direct finite-product obstruction}

The following criterion applies to a closed four-dimensional parameter
cell. It uses the original-root factorization in the right half-plane.

\begin{lemma}\label{mf:cubic-moment}
Suppose $\Re w_j\ge0$, $|w_j|\le\mu$, $\sum w_j=\sigma=e^{i\phi}$,
and $\sum|w_j|^2=q$, with $0\le\phi\le\pi/2$.
Put $T=-\sum w_j^2=q(1-\alpha+i\beta)$.
Then
\[
 \alpha^2+\beta^2\le(2\mu\cos\phi/q)^2,\qquad
 (1-\alpha)^2+\beta^2\le1.
\]
If $c=\cos\phi>0$, put $A=q\alpha/2$, $C=-q\beta/2$.
For any upper bound $\mu_h\ge\mu$, the third moment
$P_3=\sum w_j^3$ satisfies
\begin{gather}
 \frac{4A^2}{c}-3\mu_h^2c\le\Re P_3
 \le\min\left(\mu_hA,\mu_h^2c-\frac{4C^2}{c}\right),
 \label{mf:cubic-real}\\
 \left|\Im P_3-\frac{3AC}{c}\right|
 \le\frac32\left(\mu_h^2c-\frac{A^2+C^2}{c}\right)
                         +\mu_h(q-A).
 \label{mf:cubic-imag}
\end{gather}
\end{lemma}
\begin{proof}
Write $w_j=x_j+iy_j$. The identity
$T-q=-2\sum x_jw_j$, weighted Cauchy--Schwarz, and
$\sum x_j=c$ give $|T-q|\le2\mu c$; also $|T|\le q$.
Here $A=\sum x_j^2$, $C=\sum x_jy_j$.
Give the points $(x_j,y_j)$ probability weights $x_j/c$.
Their mean is $(A,C)/c$, and the trace of their covariance
matrix is at most $\mu_h^2-(A^2+C^2)/c^2$.
Now
$\Re P_3=4\sum x_j^3-3\sum x_j|w_j|^2$
gives the lower bound, while
$\Re P_3\le\sum x_j^3\le\mu_h A$ and
$\Re P_3=\sum x_j|w_j|^2-4\sum x_jy_j^2$
give the two upper bounds.
Finally $\Im P_3=3\sum x_j^2y_j-\sum y_j^3$.
The covariance bound controls the first sum about $AC/c$,
and $|\sum y_j^3|\le\mu_h\sum y_j^2=\mu_h(q-A)$
controls the second. Thus points with $x_j=0$ are included.
\end{proof}

\begin{lemma}\label{mf:point-obstruction}
Let $b$ be an actual simple rightmost zero, assume $H=|S|>0$,
and use \eqref{pr:normalized-moments}. Conjugate so that
$\sigma=e^{i\phi}$ with $0\le\phi\le\pi/2$.
Define the actual normalized primitive
\[
 J(w)=\frac{S\,p(b+w/S)}{p'(b)},\quad
 \theta=\frac{-\sum_jw_j^2}{\sigma^2},\quad
 \kappa=\frac{\sum_jw_j^3}{\sigma^3}.
\]
Suppose a closed parameter cell has bounds $q\le q_h$,
$\mu\le\mu_h$, $\phi_l\le\phi\le\phi_h$.
At a nonzero test point $w$, put $W=|w|$ and assume $\mu_hW<1$.
For recorded exact coefficients $\theta_0,\kappa_0$, let
$|\theta-\theta_0|\le h_2$, $|\kappa-\kappa_0|\le h_3$
hold on the whole cell. Put
\[
 B(t)=\sum_{k=4}^8\frac{q_h\mu_h^{k-2}W^k}{k}t^k
       +\frac{q_h\mu_h^7W^9}{9(1-\mu_hW)}t^9,
\]
\[
 J_0=w\int_0^1e^{tw+\theta_0w^2t^2/2+\kappa_0w^3t^3/3}\,dt,
 \qquad g_0=e^{w+\theta_0w^2/2+\kappa_0w^3/3},
\]
\[
 E_J=W\int_0^1 e^{A_0t+B_0t^2+C_0t^3}
     \left(e^{h_2W^2t^2/2+h_3W^3t^3/3+B(t)}-1\right)\,dt,
\]
\[
 E_g=|g_0|\left(e^{h_2W^2/2+h_3W^3/3+B(1)}-1\right),
 \quad
 (A_0,B_0,C_0)=
 \bigl(\Re w,\Re(\theta_0w^2/2),\Re(\kappa_0w^3/3)\bigr).
\]
Write $\phi_m=(\phi_l+\phi_h)/2$,
$\eta_\phi=2\sin((\phi_h-\phi_l)/4)$, $r=|J_0|$, $g=|g_0|$,
and
\[
 P_-=\Re(we^{-i\phi_m})-\eta_\phi W,\quad
 C_+=\Re(e^{i\phi_m}g_0\overline{J_0})+\eta_\phi gr.
\]
If
\begin{equation}\label{mf:point-predicates}
 P_->0,\qquad r>E_J,\qquad
 \frac{C_++gE_J}{r^2-E_J^2}
       +\frac{E_g}{r-E_J}-\frac{P_-}{W^2}<0,
\end{equation}
no actual configuration belongs to that parameter cell.
For a quadratic model, set $\kappa_0=h_3=0$ and begin the
remainder sum at $k=3$; alternatively use
$B(t)=q_h\mu_hW^3t^3/[3(1-\mu_hW)]$.
\end{lemma}

\begin{proof}
For $\Re(w/\sigma)>0$, the original-root factorization
and rightmostness imply
\begin{equation}\label{mf:actual-positivity}
 \Re\{\sigma(J'(w)/J(w)-1/w)\}\ge0.
\end{equation}
Indeed, after multiplication by $H$, the expression is
$p'(z)/p(z)-1/(z-b)=\sum_{u\ne b}(z-u)^{-1}$ at
$z=b+w/S$. Each summand has positive real part.
The critical-point product gives
$J'(w)=\exp(w+\theta w^2/2+\kappa w^3/3+R_4(w))$.
Lemma~\ref{pr:primitive}, with each order through eight kept
separately, gives $|R_4(tw)|\le B(t)$, including every order
at least nine. Integrating the exponential
perturbation proves $|J-J_0|\le E_J$, $|J'-g_0|\le E_g$.

If $|Z-v|\le E<|v|$, inversion maps that disk exactly to
\[
 \left|\frac1Z-\frac{\overline v}{|v|^2-E^2}\right|
 \le\frac{E}{|v|^2-E^2}.
\]
This follows by multiplying the inequality by
$|Z|(|v|^2-E^2)$ and expanding squares.
Apply it to $J$, use the derivative error and
$|e^{i\phi}-e^{i\phi_m}|\le\eta_\phi$, and then subtract
$\Re(\sigma/w)\ge P_-/W^2$.
The resulting upper bound for \eqref{mf:actual-positivity}
is exactly the last expression in \eqref{mf:point-predicates}.
The denominators are positive because $r>E_J$.
\end{proof}

We specify the coefficient enclosures used in this criterion.
For a cell in $(q,\phi,\alpha,\beta)$ take exact midpoints
$q_m,\phi_m,\alpha_m,\beta_m$, set
$C_m=1-\alpha_m+i\beta_m$,
$d_q=(q_h-q_l)/2$, and
$d_C=\sqrt{((\alpha_h-\alpha_l)/2)^2+
                ((\beta_h-\beta_l)/2)^2}$.
Then a valid quadratic radius is
\[
 h_2=d_q(|C_m|+d_C)+q_md_C
       +2q_m|C_m|\sin((\phi_h-\phi_l)/2)
       +|q_mC_me^{-2i\phi_m}-\theta_0|.
\]
Indeed,
\begin{align*}
 &q(1-\alpha+i\beta)e^{-2i\phi}-q_mC_me^{-2i\phi_m}\\
 &\quad=(q-q_m)(1-\alpha+i\beta)e^{-2i\phi}
       +q_m\bigl(\alpha_m-\alpha+i(\beta-\beta_m)\bigr)e^{-2i\phi}\\
 &\qquad\quad+q_mC_m(e^{-2i\phi}-e^{-2i\phi_m}).
\end{align*}
Apply the triangle inequality and
$|e^{-2i\phi}-e^{-2i\phi_m}|
 \le2\sin((\phi_h-\phi_l)/2)$.
For the cubic radius put $c_l=\cos\phi_h$, $c_h=\cos\phi_l$,
$A_l=q_l\alpha_l/2$, $A_h=q_h\alpha_h/2$, and let $C_l,C_h$
be the extrema of $-q\beta/2$ over its four endpoint pairs.
Division is used only when $c_l>0$. Put
$C_*^2=\operatorname{dist}(0,[C_l,C_h])^2$ and
\[
 R_l=4A_l^2/c_h-3\mu_h^2c_h,\qquad
 R_h=\min(\mu_hA_h,\mu_h^2c_h-4C_*^2/c_h).
\]
Let $I_l^0,I_h^0$ be the minimum and maximum of $3AC/c$
over the eight endpoint triples. Set
\[
 V_*=\max(0,\mu_h^2c_h-(A_l^2+C_*^2)/c_h),\qquad
 R_I=\tfrac32V_*+\mu_h(q_h-A_l).
\]
Equations \eqref{mf:cubic-real}--\eqref{mf:cubic-imag} enclose
$P_3$ in $[R_l,R_h]+i[I_l^0-R_I,I_h^0+R_I]$.
Subdivide its entire imaginary interval into $m$ equal closed
strips. For a strip with midpoint $P_m$ and midpoint-to-corner
radius $R_m$, use
\[
 h_3=R_m+2|P_m|\sin(3(\phi_h-\phi_l)/4)
                 +|P_me^{-3i\phi_m}-\kappa_0|.
\]
Every quantity is enclosed in the indicated outward direction.
Every strip is retained, with its endpoints. The values
$m=4,8$ suffice for the two cubic certificates below.

There are three further exclusions for a closed parameter cell.
Let $y_*$ be the threshold $3$ or $9/4$ in the contrary assumption
$q/\mu\ge y_*$, and let $\mu_0$ be the entrance cap, so that
$\mu\le\min(\mu_0,q/y_*)$.
Assume $q_l>0$, $0\le\phi_l\le\phi_h\le\pi/2$ and
$0\le\alpha_l\le\alpha_h<1$, and put
$b_*^2=\operatorname{dist}(0,[\beta_l,\beta_h])^2$.
On the branch $q\le y_*\mu_0$, use
$R_{\mathrm{mom}}=2\cos\phi_l/y_*$ and $\mu_h=q_h/y_*$;
on the branch $q\ge y_*\mu_0$, use
$R_{\mathrm{mom}}=2\mu_0\cos\phi_l/q_l$ and $\mu_h=\mu_0$.
The cell is excluded if any of
\begin{equation}\label{mf:cell-exclusions}
 \alpha_l^2+b_*^2>R_{\mathrm{mom}}^2,\qquad
 (1-\alpha_h)^2+b_*^2>1,\qquad
 q_l>\mathcal E(\mu_h,\phi_l)
\end{equation}
holds, where
\[
 \mathcal E(\mu,\phi)
 =K_0\mu^2+A_0\mu\cos\phi
       +\gamma\cos^2\phi-\tfrac12\sin^2\phi.
\]
Here $K_0,A_0,\gamma$ are the constants in Proposition~\ref{mf:envelope}.
The first two tests use lower bounds for the squared distances in
the two moment disks. For the third, \eqref{mf:exact-envelope} gives
$q\le\mathcal E(\mu,\phi)$, while
\[
 \partial_\mu\mathcal E=2K_0\mu+A_0\cos\phi\ge0,\qquad
 \partial_\phi\mathcal E
 =-\sin\phi\{A_0\mu+(2\gamma+1)\cos\phi\}\le0.
\]
All three tests therefore exclude the whole cell; their strict
inequalities are checked with outward enclosures.

\begin{proposition}\label{mf:D30}
If $D=Hd\ge30$, then $Qd/H<3$.
\end{proposition}
\begin{proof}
Proposition~\ref{mf:D270} treats $D\ge270$. Suppose
$30\le D\le270$ and $q/\mu\ge3$.
Proposition~\ref{mf:affine}, in its consequence for $D\ge7/2$,
and Lemma~\ref{mf:cubic-moment} put the contrary domain in
\[
 q\in[1/90,403/900],\quad
 \mu\le\min(1/30,q/3),\quad
 0\le\phi\le\pi/2,\quad
 0\le\alpha\le2/3,\quad|\beta|\le2/3,
\]
with the two moment-disk inequalities and the exact matrix
envelope. Divide the $q$ interval at $1/10$.
The fixed dyadic trees in
\cert{fields/tang_cubic_family_D30_y3.json}
have $9613$ and $16784$ closed leaves. A leaf is excluded by
\eqref{mf:cell-exclusions} or Lemma~\ref{mf:point-obstruction}.
The certificate has $3237,239,390$ leaves of the first three
types, $3423$ quadratic-model leaves, $8$ single cubic-model
leaves, and $19100$ complete cubic-strip families.
The latter have $79392$ members. Thus there are $82823$
point obstructions in all.

The verifier reconstructs each leaf from the exact dyadic
path; both halves of every subdivision are closed. Bottom-up
sibling coalescence reduces each tree to its root,
rejecting duplicate leaves, missing siblings, ancestor conflicts,
and inconsistent splitting coordinates. All inequalities and
all the point predicates are checked on the entire leaves by
\cert{fields/tang_verify_cubic_family_cover_D30_frozen.py}.
The source file records the exact test points and coefficient
centers. Stored finite decimals defining those centers are
interpreted as their exact binary rational values; their
mismatch from the mathematical centers is included in the
radii above. The integrands are entire, with integration remainders
bounded as in Section~\ref{sec:certificates}.
The certificate excludes the stated closed domain.
Together with $D\ge270$ this proves the proposition.
\end{proof}

\begin{proposition}\label{mf:D25}
If $D=Hd\ge25$, then $Qd/H<3$.
\end{proposition}
\begin{proof}
Only $25\le D\le30$ remains after Proposition~\ref{mf:D30}.
Under the contrary ratio $q/\mu\ge3$, the affine estimate gives
\[
 1/10\le q\le59/125,\quad \mu\le\min(1/25,q/3),\quad
 0\le\alpha\le2/3,\quad|\beta|\le2/3.
\]
In this range $\phi<\pi/4$. Indeed if $\phi\ge\pi/4$,
\eqref{mf:coarse-envelope}, $\cos\phi<71/100$, and
$q\ge3\mu$ give
\[
 0<\frac{21}{5}\mu^2+\frac{181}{200}\mu-\frac18
 \le-\frac{513}{6250}<0.
\]
Use the closed $\phi$ interval $[0,\pi/4]$ and divide
the $q$ interval at $3/25$.
The two dyadic trees in
\cert{fields/tang_D25_cubic_family_cover.json}
have $15303$ and $35311$ leaves. The exclusions consist of
$4752$ moment-disk, $393$ quadratic-disk, and $219$ matrix
leaves, using \eqref{mf:cell-exclusions}.
There are $2246$ quadratic-model leaves, $15$ single
cubic-model leaves, and $42989$ complete strip families;
the families contain $184216$ members.
Every one of the $186477$ point obstructions satisfies
\eqref{mf:point-predicates}.
The checker \cert{fields/tang_verify_D25_cubic_family_frozen.py}
uses the preceding finite-product and covariance formulas and
reconstructs the two closed trees.
Thus it excludes this domain. Its common boundary
with $D\ge30$ is included.
\end{proof}

\begin{proposition}\label{mf:D100}
If $D=Hd\ge100$, then $Qd/H<9/4$.
\end{proposition}
\begin{proof}
By Proposition~\ref{mf:D270}, it remains to consider
$100\le D\le270$ and $q/\mu\ge9/4$.
The contrary domain is contained in
\[
 q\in[1/120,1313/10000],\quad
 \mu\le\min(1/100,4q/9),\quad
 0\le\phi\le\pi/2,\quad
 0\le\alpha\le8/9,\quad|\beta|\le8/9,
\]
with the same two moment disks and matrix envelope.
Split the $q$ interval at $9/400$. The two trees in
\cert{fields/point_D100_y225.json} have $5200$ and $4545$ leaves.
There are $1507$ moment-disk exclusions, $57$ quadratic-disk
exclusions, and $21$ matrix exclusions from
\eqref{mf:cell-exclusions}. On each of the
remaining $8160$ leaves the quadratic version of
Lemma~\ref{mf:point-obstruction} applies, with
$B(t)=q_h\mu_hW^3t^3/[3(1-\mu_hW)]$.
The coefficient radius is the full-cell $h_2$ displayed above.
The checker \cert{fields/tang_independent_D100.py} reconstructs
both trees, all coefficient radii, and every strict
point obstruction. It evaluates the equivalent entire integral
$J_0/g_0=w\int_0^1e^{(t-1)w+(t^2-1)\theta_0w^2/2}\,dt$;
this follows by dividing the defining integral by $g_0$.
Thus every point of the contrary domain is excluded, which
proves the proposition.
\end{proof}

\section{Companion zeros and local fields}\label{sec:local-tools}

Throughout this section $b$ is a simple actual zero, $S=p''(b)/p'(b)\ne0$,
and $H,Q,d,D,q,\mu,\sigma$ have the meanings in
\eqref{pr:field-notation} and \eqref{pr:normalized-moments}. Put
\[
 \alpha_j=\frac1{S(b-\xi_j)},\qquad
 J(w)=\frac{S p(b+w/S)}{p'(b)},\qquad K=p''/p'.
\]
Thus $\sum\alpha_j=1$, $\sum|\alpha_j|^2=q$ and
$\max|\alpha_j|=\mu$. In particular $q>0$. The exact identities are
\begin{equation}\label{lt:finite-product}
 J'(w)=\prod_j(1+\alpha_jw),\qquad
 J(w)=\int_0^w\prod_j(1+\alpha_jv)\,dv.
\end{equation}
When a phase or curvature conclusion is asserted, $b$ is assumed
rightmost among all original zeros. The constructed companions are
zeros of $p$.

\subsection{Logarithmic bounds and circles}

For $m\ge2$ and $\mu|w|<1$, the logarithm vanishing at zero satisfies
\begin{equation}\label{lt:log-tail}
 \log J'(w)=\sum_{k=1}^{m-1}\frac{(-1)^{k+1}w^k}{k}
                    \sum_j\alpha_j^k+E_m(w),\qquad
 |E_m(w)|\le\frac{q\mu^{m-2}|w|^m}{m(1-\mu|w|)}.
\end{equation}
Indeed $\sum_j|\alpha_j|^k\le q\mu^{k-2}$ for $k\ge2$, so the
entire remaining logarithmic series is bounded by a geometric series.
Also, throughout $|w|\le Z<1/\mu$,
\begin{equation}\label{lt:log-derivative}
 \left|\frac{J''}{J'}(w)\right|
 \le1+qZ+\frac{q\mu Z^2}{1-\mu Z}.
\end{equation}
This follows from
$J''/J'=1-w\sum\alpha_j^2+w^2\sum\alpha_j^3/(1+\alpha_jw)$.

\begin{lemma}\label{lt:circle}
Suppose $J'$ is nonzero on $|w|\le Z_{\mathrm{out}}$,
$|J''/J'|\le L$ there, and $|\zeta|+R<Z_{\mathrm{out}}$. If
\[
 \frac{|J(\zeta)|}{|J'(\zeta)|}
       +\frac{e^{LR}-1-LR}{L}<R,
\]
then $|w-\zeta|<R$ contains exactly one zero of $J$, which is simple.
\end{lemma}
\begin{proof}
Since $J''(0)/J'(0)=\sum_j\alpha_j=1$, the hypothesis gives $L\ge1$.
On every segment $\zeta+tv$, $|v|\le R$, integration of the
logarithmic derivative gives
$|J'(\zeta+tv)/J'(\zeta)-1|\le e^{Lt|v|}-1$.
A further integration gives
\[
 |J(\zeta+v)-J(\zeta)-vJ'(\zeta)|
 \le |J'(\zeta)|\frac{e^{L|v|}-1-L|v|}{L}.
\]
On $|v|=R$, Rouch\'e's theorem compares
$J(\zeta+v)$ with $vJ'(\zeta)$. The critical-free containing disk
makes its unique zero simple.
\end{proof}

\subsection{Gaussian quartic centers}

Put $\theta=-\sum\alpha_j^2=K'(b)/S^2$, so $|\theta|\le q$.
For $s=\pm2\pi i$ define
\begin{align}
 \zeta_s={}&s+(s-s^2/2)\theta+(3s-2s^2+s^3/2)\theta^2\notag\\
 &+(15s-21s^2/2+23s^3/6-5s^4/8)\theta^3\notag\\
 &+(105s-72s^2+88s^3/3-22s^4/3+7s^5/8)\theta^4.
 \label{lt:gaussian-center}
\end{align}
If $\zeta_s=\sum_{j=0}^4c_j(s)\theta^j$, write
\[
 t\zeta_s+\theta t^2\zeta_s^2/2
       =st+\sum_{j=1}^9 A_j(s,t)\theta^j.
\]
Thus $A_j=tc_j+(t^2/2)\sum_{k+l=j-1}c_kc_l$, where a coefficient
outside $0\le j\le4$ is zero. Define
\[
 E_0=1,\qquad E_j=\frac1j\sum_{k=1}^{\min(j,9)}kA_kE_{j-k},
 \qquad I_0=0,\quad I_m=\frac{1-mI_{m-1}}s\quad(m\ge1).
\]
For a polynomial $f(t)=\sum f_mt^m$ set
$\mathcal I_s f=\sum f_m I_m$. Since $e^s=1$,
$\mathcal I_s f=\int_0^1e^{st}f(t)\,dt$ at the two indicated values
of $s$. The coefficients of the Gaussian primitive at the center are
\[
 a_j(s)=\mathcal I_s\left(\sum_{k=0}^{\min(4,j)}c_kE_{j-k}\right).
\]
The recurrence gives $a_0=\cdots=a_4=0$ exactly. It can alternatively
be used to determine $c_1,\ldots,c_4$ successively: the new center
coefficient enters the corresponding residual with coefficient one.

The following strict bounds hold for
$6\le p\le44/7$, $0\le t\le1$, with $s=\pm ip$:
\begin{center}\small
\begin{tabular}{cl}
\toprule
polynomials & bounds in increasing order\\\midrule
$c_1,\ldots,c_4$ & $21,133,1030,9000$\\
$A_1,\ldots,A_9$ & $63/10,45,400,3400,86000,320000,$\\
 & $1700000,9100000,39000000$\\
$a_5,\ldots,a_{11}$ & $82000,350000,1900000,13000000,$\\
 & $120000000,3100000000,23000000000$\\\bottomrule
\end{tabular}
\end{center}
Here and below a bound for a complex polynomial means a bound for its
modulus. To verify these inequalities, substitute $s=ip$ in the
recurrences and subtract
the squared modulus from the square of the proposed bound. Convert
the resulting rational polynomial to tensor Bernstein form on
$[6,44/7]\times[0,1]$, omitting the second variable when absent.
If every control coefficient is positive, accept the box; otherwise
bisect it and use the exact de Casteljau averages. At depth $d$ the
path-variable axis is chosen unless $d\equiv2\pmod3$, when the
$p$-axis is chosen. For one-variable polynomials bisect the $p$-axis.
All twenty inequalities terminate with positive coefficients on a total
of twenty-four closed boxes, at depth at most twenty. The recurrences
and bounds specify the certificate;
\cert{check_gaussian_domains_new.py} implements it with rational arithmetic.
The coefficient identities $a_0=\cdots=a_4=0$ are checked by exact
polynomial arithmetic. They apply to the primitive at $s=\pm2\pi i$,
where $e^s=1$. The resulting residual polynomials are bounded on
$6\le p\le44/7$.

Set
\begin{align*}
 Z(q)&=63/10+21q+133q^2+1030q^3+9000q^4,\\
 M(q)&=(63/10)q+45q^2+400q^3+3400q^4+86000q^5\\
     &\quad+320000q^6+1700000q^7+9100000q^8+39000000q^9.
\end{align*}
Let $P_{11}$ be the degree-eleven Taylor polynomial of $Z(q)e^{M(q)}$,
and let $B_5,\ldots,B_{11}$ denote the seven residual bounds in the
table. For $J_\theta(w)=\int_0^we^{z+\theta z^2/2}\,dz$ we have
\begin{equation}\label{lt:gaussian-remainder}
 |J_\theta(\zeta_s)|
 <\sum_{j=5}^{11}B_jq^j+Z(q)e^{M(q)}-P_{11}(q).
\end{equation}
Every coefficient in the last difference is nonnegative. Expand the
exponential along the path, bound the retained coefficients by the table
and every higher coefficient by $Z e^M$.
For exact numerical bounds we use
\begin{equation}\label{lt:exp-upper}
 \mathcal E_{24}(x)=\sum_{k=0}^{24}\frac{x^k}{k!}
       +\frac{x^{25}}{25!(1-x/26)}>e^x\qquad(0<x<26).
\end{equation}
The quotient of the right side of \eqref{lt:gaussian-remainder} by
$q^5$ increases with $q$. Substituting \eqref{lt:exp-upper} proves
the respective coefficients $132200,126000,86200,89000$ at
$q=1/20,97/2000,7/625,9/556$.

\begin{proposition}\label{lt:quartic}\label{lt:fine200}\label{lt:wideD100}
Each row of the following table gives two distinct simple actual zeros,
other than $b$, in the disks $|S(z-b)-\zeta_\pm|<R$.
The disks lie in the indicated critical-free containing disk.
\begin{center}\small
\begin{tabular}{ccclc}
\toprule
row&$q_0$&$\mu_0$&$R(q,\mu)$&$Z_{\mathrm{out}}$\\\midrule
coarse&$97/2000$&$1/270$&$300000q^5+1300q\mu$&$41/5$\\
fine&$7/625$&$1/200$&$120000q^5+220q\mu$&$659/100$\\
wide&$9/400$&$1/100$&$250000q^5+400q\mu$&$174/25$\\
origin&$33/1000$&$8/625$&$3/10$&$15/2$\\
medium&$9/556$&$1/139$&$120000q^5+225q\mu$&$168/25$\\\bottomrule
\end{tabular}
\end{center}
Every row holds on $0<q\le q_0$, $0<\mu\le\mu_0$.
At a rightmost zero the fine and medium rows give
$|\Im S|<4Q/H$, and the wide row gives $|\Im S|<(23/5)Q/H$.
\end{proposition}
\begin{proof}
For the first four rows use $J'=e^{w+\theta w^2/2+E_3(w)}$ and
\eqref{lt:log-tail}. The following constants are exact:
\begin{center}\small
\begin{tabular}{c|rrrrrrr}
\toprule
row&$Z_0$&$M_0$&$E_0$&$B$&$P_0$&$L$&$\beta$\\\midrule
coarse&$39/5$&$101/200$&$43/25$&$126000$&$550$&$143/100$&$27/100$\\
fine&$657/100$&$2/25$&$109/100$&$86200$&$175$&$11/10$&$1/100$\\
wide&$343/50$&$43/250$&$5/4$&$132200$&$250$&$59/50$&$3/50$\\
origin&$719/100$&$7/25$&$71/50$&$132200$&$350$&$32/25$&$11/50$\\\bottomrule
\end{tabular}
\end{center}
At each upper corner, set
$r_3=q_0\mu_0 Z_0^3/[3(1-\mu_0Z_0)]$. The finite predicates are
\begin{gather}
 Z(q_0)<Z_0,\quad M(q_0)<M_0,\quad
 \mathcal E_{24}(M_0+r_3)<E_0,\quad
 \frac{Z_0^4E_0}{12(1-\mu_0Z_0)}<P_0,\label{lt:gaussian-gates1}\\
 Z_0+R_0<Z_{\mathrm{out}}<1/\mu_0,\quad
 1+q_0Z_{\mathrm{out}}+
 \frac{q_0\mu_0Z_{\mathrm{out}}^2}{1-\mu_0Z_{\mathrm{out}}}<L,\notag\\
 \frac{\mathcal E_{24}(LR_0)-1-LR_0}{LR_0}<\beta,
 \label{lt:gaussian-gates2}
\end{gather}
where $R_0=R(q_0,\mu_0)$. The center residual obeys
\[
 |J(\zeta_s)/J'(\zeta_s)|<E_0Bq^5+E_0P_0q\mu.
\]
Indeed the primitive discrepancy is bounded by integrating
$q\mu Z_0^3t^3/[3(1-\mu_0Z_0)]$ for $0\le t\le1$; its integral
contains the factor $1/12$, and $|J'(\zeta_s)|>1/E_0$.
For the first three rows the two coefficient inequalities
$E_0B<(1-\beta)[q^5]R$ and
$E_0P_0<(1-\beta)[q\mu]R$ prove the circle criterion; the brackets
denote coefficients of the indicated monomials in $R$.
In the origin
row their sum at the corner is less than $0.217280<(1-\beta)R_0$.
Every majorant is increasing, so these predicates hold on the
closed rectangle. Source exclusion and separation follow respectively
from
\[
 6-21q_0-133q_0^2-1030q_0^3-9000q_0^4>R_0,
\]
\[
 12(1-q_0-17q_0^2-137q_0^3-313q_0^4)>2R_0.
\]
All displayed predicates are rational after \eqref{lt:exp-upper} is
substituted, and direct evaluation gives positive differences.

For the medium row estimate the discrepancy relative to
the endpoint derivative. Write
$g(t)=t\zeta_s+\theta t^2\zeta_s^2/2$. The coefficient of $\theta$
in $g(t)-g(1)-(t-1)s$ is
$(t-1)s+(t^2-t)s^2/2$ and has modulus less than $63/10$.
To see this, put $u=1-t$ and use $p^2=(2\pi)^2<40$ and
\[
 1-u^2[1+10(1-u)^2]=(1-u)(10u^3-10u^2+u+1).
\]
The cubic on the right has Bernstein coefficients
$(1,7/6,1/2,1/4)$ on $[0,1/2]$,
$(1/4,1/8,5/48,11/32)$ on $[1/2,3/4]$, and
$(11/32,7/12,13/12,2)$ on $[3/4,1]$.
For the remaining coefficients use
$(t-1)c_k+(t^2-1)\sum_{i+j=k-1}c_ic_j/2$.
This gives the relative-path majorant
\begin{align*}
 M_r(q)={}&(63/10)q+(2653/10)q^2+(10442/5)q^3+18282q^4\\
 &+(174349/2)q^5+325990q^6+1727450q^7
       +9270000q^8+40500000q^9.
\end{align*}
At $q_0=9/556$, $\mu_0=1/139$, take
\begin{gather*}
 Z_0=6681/1000,\quad M_0=29/250,\quad M_{r,0}=183/1000,\\
 E_0=57/50,\quad L=28/25,\quad\beta=3/200.
\end{gather*}
The bounds $Z(q_0)<Z_0$, $M(q_0)<M_0$, $M_r(q_0)<M_{r,0}$ and
$e^{M_0+r_3}<E_0$ follow from the same rational predicates.
Dividing the integral discrepancy by the actual $J'(\zeta_s)$ gives
\begin{equation}\label{lt:relative-gaussian-error}
 \left|\frac{J(\zeta_s)-J_\theta(\zeta_s)}{J'(\zeta_s)}\right|
 <\frac{q\mu Z_0^4e^{M_{r,0}+2r_3}}{12(1-\mu_0Z_0)}
 <215q\mu.
\end{equation}
The path remainder and endpoint denominator each contribute a factor
$e^{r_3}$. The model contribution is less than
$101460q^5$. Now $101460<(1-\beta)120000$ and
$215<(1-\beta)225$. The other circle, separation and containment
predicates in \eqref{lt:gaussian-gates2} hold with
$Z_{\mathrm{out}}=168/25$.
These estimates prove the medium row.

At a rightmost $b$, both companions satisfy $\Re(w/\sigma)\le0$.
Choose the sign of $s$ for which
$\Re(s/\sigma)=2\pi|\Im\sigma|$. The triangle inequality gives
\[
 2\pi|\Im\sigma|<q\left(21+133q+1030q^2+9000q^3+R/q\right).
\]
At the fine corner the bracket is less than $24$; at the wide corner
it is less than $28.681<(157/25)(23/5)$; at the medium corner it is
less than $4(157/25)$. Since $2\pi>157/25>6$, the phase bounds follow.
\end{proof}

\subsection{Signed curvature from the two companions}

Write $p_0=2\pi$ for this subsection only, and define
\begin{align*}
 B(\theta)&=1+\theta+(3-p_0^2/2)\theta^2
 +(15-23p_0^2/6)\theta^3
 +(105-88p_0^2/3+7p_0^4/8)\theta^4,\\
 T(\theta)&=\theta+4\theta^2+(21-5p_0^2/4)\theta^3
 +(144-44p_0^2/3)\theta^4.
\end{align*}
Then $\zeta_\pm=(p_0^2/2)T\pm ip_0B$. If the actual roots lie
in disks of radius $R$ about these centers, rightmostness gives
\begin{equation}\label{lt:signed-projection}
 \frac{p_0^2}2\Re(T/\sigma)+p_0|\Im(B/\sigma)|<R.
\end{equation}
For other centers, add to $R$ an upper bound for their displacements
from \eqref{lt:gaussian-center}.

\begin{lemma}\label{lt:signed-moment}
For $r=|\theta|\le1/20$, put
\[
 v=r+17r^2+137r^3+313r^4,\quad m=1-v.
\]
\[
 A_r=\frac{1+17r+137r^2+313r^3}{m},\quad
 B_r=\frac{18+154r+450r^2+313r^3}{m}.
\]
\[
 \mathcal C(r)=1+(2+100/199)rA_r^2+3rB_r
        +(600/199)r^2A_r^2+4rA_r+29r+436r^2.
\]
The two companion disks imply
\begin{equation}\label{lt:curvature-Cq}
 \frac{\Re K'(b)}{H^2}<\mathcal C(q)q^2+\frac{2R}{p_0^2}
 \qquad(q\le1/20).
\end{equation}
It also implies the sharper bound with $\mathcal C(q)$ replaced by
$6/5$. For $q\le3/50$ the replacement $13/10$ is valid provided
an independent circle estimate gives
$|\Im\sigma|<8q$.
\end{lemma}
\begin{proof}
Set $\tau=B/|B|=c+is$ and $\delta=\Im(B/\sigma)$.
The coefficient bounds give $|B|\ge m$, $c>99/100$,
$1-c\le(100/199)s^2$, $|s|\le rA_r$, and
$|s-\Im\theta|\le r^2B_r$. For
$f(\sigma)=\Re(\theta\sigma^2-T/\sigma)$,
\[
 |f(\sigma)-f(\tau)|
 \le r(3+4r+29r^2+436r^3)|\sigma-\tau|.
\]
The exact chord identity is
\[
 \frac{|\sigma-\tau|^2}{\delta^2}
   =\frac{2}{|B|^2(1+\Re(\sigma\overline\tau))}.
\]
On the right semicircle, $\Re(\sigma\overline\tau)\ge-|s|$.
The endpoint comparison at $r=1/20$ gives
$|\sigma-\tau|\le(9/5)|\delta|$ and
\[
 (9/5)r(3+4r+29r^2+436r^3)<7/22<2/p_0.
\]
The case $\delta=0$ follows by continuity: the antipodal point
$-\tau$ has negative real part.

Writing $\theta=x+iy$, expansion of $f(\tau)$ gives
\begin{align*}
 y^2-4x^2&+x(1-c)-2xs^2-y[(s-y)(2c+1)+2y(c-1)]\\
 &+4(x^2-y^2)(1-c)-8xys
 -(21-5p_0^2/4)\Re(\theta^3\overline\tau)
 -(144-44p_0^2/3)\Re(\theta^4\overline\tau).
\end{align*}
Use the bounds above, $|xy|\le r^2/2$, and the coefficient bounds
$|21-5p_0^2/4|<29$, $|144-44p_0^2/3|<436$.
The result is $f(\tau)\le\mathcal C(r)r^2$.
Every term in $\mathcal C$ increases, and
$\mathcal C(97/2000)<9$, $\mathcal C(7/625)<11/5$.
Combining this estimate with \eqref{lt:signed-projection} cancels
$\delta$ and proves \eqref{lt:curvature-Cq}.

For the two sharper constants there is a finite polynomial proof.
Write $\theta=re^{i\phi}$, $u=\cos\phi$, and let $T_j(u)$ be the
Chebyshev polynomials. If $B=\sum b_j\theta^j$ and
$T=\sum d_j\theta^j$, define
\begin{align*}
 M&=\sum_{i,j}b_ib_jr^{i+j}T_{|i-j|}(u),\\
 U&=\sum_{i,j}b_ib_jr^{i+j+1}T_{i+j+1}(u),\qquad
 V=\sum_{i,j}d_ib_jr^{i+j}T_{|i-j|}(u),\\
 P&=(U-V)/r^2,\qquad W=V(M-1)/r^2.
\end{align*}
The last two expressions are polynomials, including at $r=0$.
For $(r_0,k)=(1/20,6/5)$ or $(3/50,13/10)$, let
$m_0=1-r_0-17r_0^2-137r_0^3-313r_0^4$. The polynomials
\begin{equation}\label{lt:signed-polynomials}
 A=kM-P,\qquad (1+m_0)A+W
\end{equation}
are positive on $0\le r\le r_0$, $-1\le u\le1$ and
\[
 3947841760435743/10^{14}\le p_0^2
           \le3947841760435744/10^{14}.
\]
Indeed, tensor Bernstein conversion on the three closed $u$ intervals
$[-1,0]$, $[0,1/2]$, $[1/2,1]$ gives 4950 positive rational
coefficients for each $r_0$. The conversion is exactly
$x^k=\sum_{j=k}^n\binom jk\binom nk^{-1}
\binom njx^j(1-x)^{n-j}$ after affine substitution in each variable.
These polynomials and the three boxes specify the certificate.
Machin's formula and alternating arctangent sums of
lengths twenty and twenty-one give the stated enclosure of $p_0^2$.
Now
\[
 f(\tau)/r^2=\frac1M\left(P-\frac{W}{1+\sqrt M}\right).
\]
If $W\ge0$ use $A>0$; if $W<0$ use $1+\sqrt M\ge1+m_0$ and
the second inequality in \eqref{lt:signed-polynomials}. Thus
$f(\tau)\le kr^2$ on the whole disk, with $r=0$ supplied by continuity.

At the larger cap $r\le q\le3/50$, the phase bound gives
$\Re\sigma>7/8$. The coefficient bounds give
$|B|>21/25$, $|\Im\tau|<1/5$, $\Re\tau>97/100$, hence
$\Re(\sigma\overline\tau)>3011/4000$.
The rational comparison
$(21/25)^2(1+3011/4000)(4/3)^2>2$ gives
$|\sigma-\tau|<(4/3)|\delta|$. At $r=3/50$,
\[
 (4/3)r(3+4r+29r^2+436r^3)=214911/781250<7/22<2/p_0.
\]
The same cancellation proves the second assertion on this larger domain.
\end{proof}

Two consequences needed by the field argument in Section~\ref{sec:matrix-fields}
are
\begin{align}
 \Re K'(b)/H^2&<\mathcal C(q)q^2+(50000/3)q^5+(650/9)q\mu
 &&\text{on the coarse row},\label{lt:coarse-curvature}\\
 \Re K'(b)/H^2&<(9/4)q^2+13q\mu
 &&\text{on the fine row}.\label{lt:fine-curvature}
\end{align}
These follow from the companion disks, since $p_0>6$ and
$11/5+(120000/18)(7/625)^3<9/4$, $220/18<13$.
On the medium row the sharper signed constant instead gives
\begin{equation}\label{lt:medium-curvature}
 \Re K'(b)/H^2<(5/4)q^2+(23/2)q\mu,
\end{equation}
because $2/p_0^2<1250/24649$,
$6/5+(1250/24649)120000(9/556)^3<5/4$, and
$(1250/24649)225<23/2$.

\subsection{Centers retaining four moments}

In addition to $\theta$ put $\kappa=\sum\alpha_j^3$ and
$\lambda=-\sum\alpha_j^4$. Their bounds are
$|\kappa|\le q\mu$, $|\lambda|\le q\mu^2$. Assign weights
$1,2,3$ to $\theta,\kappa,\lambda$. For $s=\pm2\pi i$, the
weight-four center is
\begin{equation}\label{lt:four-center}
 Z_4=\zeta_s+u\kappa+v\theta\kappa+w\theta^2\kappa
                    +r\kappa^2+l\lambda+m\theta\lambda,
\end{equation}
where the four pure $\theta$ terms are those in
\eqref{lt:gaussian-center}, and
\begin{align*}
 u&=-s^3/3+s^2-2s,&
 v&=5s^4/6-13s^3/3+12s^2-20s,\\
 w&=-7s^5/4+157s^4/12-51s^3+131s^2-210s,&
 r&=s^5/3-13s^4/6+26s^3/3-22s^2+40s,\\
 l&=-s^4/4+s^3-3s^2+6s,&
 m&=3s^5/4-19s^4/4+18s^3-51s^2+90s.
\end{align*}
For higher weights define
\[
 J_4(z)=\int_0^ze^{w+\theta w^2/2+\kappa w^3/3+\lambda w^4/4}\,dw.
\]
Let $T_j(z)=\sum_{k=0}^j(-1)^{j-k}j!z^k/k!$. Its formal expansion is
\[
 J_4(z)=\sum_{a,b,c\ge0}
 \frac{\theta^a\kappa^b\lambda^c}{2^a3^b4^c a!b!c!}
 \{e^zT_{2a+3b+4c}(z)-T_{2a+3b+4c}(0)\}.
\]
Starting from $Z_0=s$, define recursively
\begin{equation}\label{lt:center-recursion}
 Z_j=Z_{j-1}-[\text{weight }j]J_4(Z_{j-1}),
        \qquad 1\le j\le8,\qquad e^s=1.
\end{equation}
Here $[\text{weight }j]$ selects the homogeneous component of weight $j$.
The new correction has derivative coefficient one, so this cancels
all coefficients up to the indicated weight. The first four levels
give exactly \eqref{lt:four-center}. The numbers of new terms in
levels $1,\ldots,8$ are $1,2,3,4,5,7,8,10$. Thus
\eqref{lt:center-recursion} uniquely specifies all forty nonconstant
terms of $Z_8$ as rational polynomials in $s$.

For $Z=Z_m(s)$, where $m=4$ or $8$ and $s=\pm2\pi i$, put
\begin{align*}
 A(t)&=t(Z-s)+\theta t^2 Z^2/2+\kappa t^3Z^3/3+\lambda t^4Z^4/4,\\
 A_r(t)&=(t-1)(Z-s)+\theta(t^2-1)Z^2/2
                  +\kappa(t^3-1)Z^3/3+\lambda(t^4-1)Z^4/4.
\end{align*}
These exponents give
\[
 J_4(Z)=Z\int_0^1e^{st+A(t)}\,dt,\qquad
 \frac{J_4(Z)}{J_4'(Z)}
 =Z\int_0^1e^{s(t-1)+A_r(t)}\,dt.
\]
For caps $q_0,\mu_0>0$ with $q\le q_0$, $\mu\le\mu_0$, use
the scaled variables
\[
 (u_1,u_2,u_3)=
 \left(\frac{\theta}{q_0},\frac{\kappa}{q_0\mu_0},
                         \frac{\lambda}{q_0\mu_0^2}\right),
 \qquad |u_1|,|u_2|,|u_3|\le1.
\]
For $\nu\in\mathbb Z_{\ge0}^3$, write
$\mathbf u^\nu=u_1^{\nu_1}u_2^{\nu_2}u_3^{\nu_3}$ and
$\mathrm{wt}(\nu)=\nu_1+2\nu_2+3\nu_3$; $\gamma\le\nu$ means
coordinatewise comparison. For either exponent $A$ or $A_r$, let
$A_\gamma(t)$ be its coefficient of $\mathbf u^\gamma$ after this
substitution, and let $E_\nu(t)$ be the coefficient of $\mathbf u^\nu$
in its exponential. Each $A_\gamma$ has degree at most four in $t$.
The exponential coefficients satisfy
\begin{equation}\label{lt:euler-recursion}
 E_0=1,\qquad
 E_\nu=\frac1{\mathrm{wt}(\nu)}
       \sum_{0<\gamma\le\nu}\mathrm{wt}(\gamma)A_\gamma E_{\nu-\gamma},
 \qquad \nu\ne0.
\end{equation}
The sum is over nonzero multi-indices $\gamma\le\nu$.
This identity is obtained by applying the weighted Euler derivative to
the chosen exponential. Since $e^s=1$, both integrals use the recurrence
for $I_m$ above. The integrated coefficients through weight $m$ vanish
exactly.
For the absolute residual $J_4(Z)$ or the relative residual
$J_4(Z)/J_4'(Z)$, let $F_*$ be the sum of the moduli of its
coefficients in $\mathbf u$ with $m<\mathrm{wt}(\nu)\le M$,
where $M=14$ for $m=4$ and $M=28$ for $m=8$.

On every interval
$[j/128,(j+1)/128]$ express each quartic coefficient in Bernstein
form. Its values lie in the convex hull of its five complex control
coefficients, so the largest control modulus bounds the entire
interval. The same applies to its difference from the endpoint value.
Alternatively the midpoint Taylor bound
$\sum_{k=0}^4|f^{(k)}(m_t)|h^k/k!$, where
$m_t=(2j+1)/256$ and $h=1/256$ are the midpoint and half-length,
bounds that interval. The four-moment rows use the former method and the
weight-eight rows use the latter. There are 253 path polynomials for
$Z_4$ and 1409 for $Z_8$.

For each exponent coefficient take the largest bound over the 128
intervals. Let $a_j$ be the sum of these bounds in weight $j$,
and $z_j$ the sum of the moduli of the coefficients of $Z$ in
$\mathbf u$ in that weight, including $z_0=|s|$. Put
$A_* (v)=\sum a_jv^j$ and $Z_*(v)=\sum z_jv^j$. The infinite tail
after weight $M$ is bounded by
\begin{equation}\label{lt:positive-weight-tail}
 T_*=Z_*(1)e^{A_*(1)}-
          \sum_{j=0}^M[v^j]\{Z_*(v)e^{A_*(v)}\}.
\end{equation}
The positive exponential coefficients are generated by
$e_0=1$, $e_j=j^{-1}\sum_{k=1}^jka_ke_{j-k}$.
All coefficients in the discarded tail are nonnegative. Thus
\eqref{lt:positive-weight-tail} bounds every omitted weight, including
every mixed parameter term. Its exponential is bounded upward by
\eqref{lt:exp-upper}, or by an outward complex ball evaluation with
an explicit remainder bound.

To pass from the model to the actual finite product, put
\[
 r_5=\frac{q_0\mu_0^3Z_*(1)^5}{5(1-\mu_0Z_*(1))},
       \qquad |Z|\le Z_*(1)<1/\mu_0.
\]
An absolute model residual $F_*+T_*$ has actual derivative-relative
bound $(F_*+T_*)e^{M_c+r_5}$, where $M_c$ bounds the center exponent
after $s$. A relative model residual instead has bound
$(F_*+T_*)e^{r_5}$. In both cases the actual product discrepancy is
bounded by
\begin{equation}\label{lt:four-actual-tail}
 \frac{q\mu^3Z_*(1)^6e^{M_r+2r_5}}{30(1-\mu_0Z_*(1))}.
\end{equation}
Here $M_r$ bounds the relative exponent along the path, and $1/30$
includes the integral of $t^5$. The path remainder and endpoint
denominator give the two $r_5$ terms.

There are two scaling rules. If $q=tq_0$, $0<t\le1$, and
$\mu\le\mu_0$, a monomial $\theta^a\kappa^b\lambda^c$ is bounded by
$t^{a+b+c}$ times its corner majorant. The uncancelled $Z_4$ terms
have total degree at least two; the tail above weight fourteen has
total degree at least five.
For $Z_8$ the corresponding degrees are at least three and ten.
These statements follow from $\mathrm{wt}(a,b,c)\le3(a+b+c)$
and the exact low-weight cancellation. They hold term by term in the
multivariate majorant and hence for its infinite tail. When instead
$q\le k/D$, $\mu=1/D$ and $D$ increases, the three parameters
scale with weights $1,2,3$.

\begin{proposition}\label{lt:q04mu04}\label{lt:q05mu02}\label{lt:W8}
For each row and $0<q\le q_0$, $0<\mu\le\mu_0$, the disks
$|S(z-b)-Z_m(s)|<R$, $s=\pm2\pi i$, are disjoint from each other
and from $b$, and contain exactly one simple original zero each,
denoted $z_\pm$ respectively.
Put $w_\pm=S(z_\pm-b)$. Both disks lie in the critical-free
disk $|S(z-b)|<Z_{\mathrm{out}}$. The final column gives
$|\Im S|<kQ/H$ at a rightmost zero.
\begin{center}\small
\begin{tabular}{ccclcc}
\toprule
$m$&$q_0$&$\mu_0$&$R(q,\mu)$&$Z_{\mathrm{out}}$&$k$\\\midrule
4&$1/20$&$1/50$&$9q$&$17/2$&$7$\\
4&$1/25$&$1/25$&$(45/4)q$&$17/2$&$15/2$\\
4&$3/70$&$1/70$&$34q^2+13400q\mu^3$&$17/2$&$21/4$\\
8&$3/50$&$1/50$&$64q^2+95000q\mu^3$&$19/2$&$8$\\
8&$273/5000$&$91/5000$&$22q^2+45000q\mu^3$&$9$&$33/5$\\\bottomrule
\end{tabular}
\end{center}
The first row also gives $|w_\pm|<17/2$ and
$|(w_++w_-)/2|<19/10$ for the same pair.
On the subrectangle $q\le57/1000$, $\mu\le1/51$, the fourth row gives
\begin{equation}\label{lt:Q057}
 |w_\pm|<8917/1000,\qquad
 |(w_++w_-)/2|<2109/1000,\qquad |\Im S|<(37/5)Q/H.
\end{equation}
\end{proposition}
\begin{proof}
Use \eqref{lt:center-recursion}--\eqref{lt:four-actual-tail}. For the
first two rows the exact caps of the center modulus are respectively
$8.030851273687$ and $7.706829806638$. Upper bounds at
the corners are as follows:
\begin{center}\small
\begin{tabular}{c|rrrrr}
\toprule
$(q_0,\mu_0)$&$F_*$&$T_*$&model/actual derivative&product error&circle error\\\midrule
$(.05,.02)$&$.091297$&$.017192$&$.216053$&$.008556$&$.200791$\\
$(.04,.04)$&$.079034$&$.012556$&$.169285$&$.049027$&$.200791$\\\bottomrule
\end{tabular}
\end{center}
The sums in the last three columns are strictly less than $9/20$.
The containing-disk bound is $L=31/20$ in both cases.
The finite sum retains all weights five through fourteen; the
infinite sum is exactly bounded by \eqref{lt:positive-weight-tail}.
The path, derivative and circle inequalities follow from
the deterministic recurrences and all 128 closed path intervals
specified above. At smaller $q$, both actual residuals are at most
$t=q/q_0$ times their corner values. Since
$(e^{LR}-1-LR)/L$ is convex and vanishes at zero, the circle error
is at most $t$ times its corner value. Thus the radii are $t(9/20)$
throughout the parameter rectangles.

For the third row the same construction at its own corner yields
\begin{gather*}
 Z_*(1)<7.629836,\quad |Z_4-s|<1.346650,\\
 M_c<.478584,\quad M_r<.481296,\quad
 F_*<.033415,\quad T_*<.000911.
\end{gather*}
Using the unrounded rational witnesses produced by the path conversion,
the model contribution is less than $30.181q^2$ and the actual
product contribution is less than $11961q\mu^3$.
Both are below nine tenths of the corresponding radius coefficient.
At the corner $R_0<.064124$, $Z_*(1)+R_0<17/2$,
\eqref{lt:log-derivative} is less than $31/20$, and the nonlinear
fraction is less than $.052<1/10$. The model terms and positive tail
are bounded by $(q/q_0)^2$ times their corner values;
\eqref{lt:four-actual-tail} retains its explicit $q\mu^3$.
Thus the third row holds throughout its parameter rectangle.

For the last two rows use the relative exponent $A_r$.
The upper bounds, including all weights at least twenty-nine,
are
\begin{center}\small
\begin{tabular}{c|rrrrrr}
\toprule
row&$F_*$&$T_*$&model error&product error&$L$&nonlinear fraction\\\midrule
$q_0=.06$&$.136221$&$.012814$&$.150033$&$.028511$&$7/4$&$1/3$\\
$q_0=.0546$&$.045190$&$.000674$&$.046017$&$.010525$&$17/10$&$1/5$\\\bottomrule
\end{tabular}
\end{center}
In each row the model error divided by $q_0^2$ is less than
$(1-\beta)$ times the $q^2$ radius coefficient, and the product
error divided by $q_0\mu_0^3$ is less than $(1-\beta)$ times the
other coefficient. The nonlinear fraction is less than
$\beta$. Containment, source exclusion and separation hold at both
corners and improve on the smaller rectangles. Lemma~\ref{lt:circle}
therefore applies.

The coefficient polynomials are generated by
\eqref{lt:center-recursion}; their exponential coefficients are given
by \eqref{lt:euler-recursion}; every path coefficient is bounded on
the stated 128 intervals; and the infinite positive tail is
\eqref{lt:positive-weight-tail}. The electronic appendix contains the
exact rational coefficient lists and programs performing these
operations at 352 or 448 bits with outward bounds. A proposed rational
modulus bound is accepted only if it exceeds its whole interval bound.
The vanishing identities are checked separately over $\mathbb Q[s]$.

Every center correction has positive total parameter degree, and
$R(q,\mu)\le(q/q_0)R_0$. If $\Gamma$ bounds $|Z_m-s|$ at
the corner, rightmostness gives
$2\pi|\Im\sigma|<(q/q_0)(\Gamma+R_0)$.
The rational or outward strict gates
$\Gamma+R_0<2\pi q_0k$ give exactly the last column of the table.
Averaging the two centers cancels all odd powers of $s$. The sum of
the absolute even coefficients at $(.05,.02)$, plus $9/20$, is
less than $1.890068<19/10$. At $(.057,1/51)$ the corresponding
evaluations give respectively
$8.916593$, $2.108213$ and $7.352979$, proving
\eqref{lt:Q057} for the pair $(z_+,z_-)$.
\end{proof}

\subsection{Curvature and the hole field}

\begin{lemma}\label{lt:hole-field}
Assume the unit-disk strict-hole geometry of
Lemma~\ref{pr:edge-reduction}. Suppose $C_0\ge0$ and
\[
 Qd/H<k,\qquad |\Im S|<K_0Q/H,\qquad
 \Re K'(b)<C_0/d^2.
\]
Then
\begin{equation}\label{lt:general-hole-field}
 H>\frac{N-(2C_0+2K_0k)/d}{2(1+a)}.
\end{equation}
\end{lemma}
\begin{proof}
Put $p_b=1-|b|^2\ge0$, $\lambda_b=1-|b-a|^2>0$.
For $v_j=(b-\xi_j)^{-1}$, the strict hole gives
\[
 \lambda_b|v_j|^2+2\Re((b-a)v_j)<1.
\]
This is an open disk whose greatest real coordinate is
$m_R=(1-x+a)/\lambda_b$. Rightmostness gives $\Re v_j\ge0$,
so
\[
 Q-2\sum(\Re v_j)^2=\Re K'(b)<C_0/d^2,\qquad
 Q<2m_RH+C_0/d^2.
\]
Since $\lambda_b\ge p_b$, one has $p_bm_R\le1-x+a$.
Also $p_b<2d$, because $d>1-|b|$ in the strict crescent.
Apply \eqref{pr:radial-identity} and $0\le y\le1$ to obtain
\[
 N<2(x+p_bm_R)H+2K_0Q/H+p_bC_0/d^2
   \le2(1+a)H+(2K_0k+2C_0)/d.
\]
This also covers $p_b=0$ without division by it.
\end{proof}

\begin{proposition}\label{lt:D70}
If $D\ge70$, then
\[
 |\Im S|<(21/4)Q/H,\qquad Qd/H<Y_{70}(D)<8/3,
\]
\[
 Y_{70}(D)=\frac{4D}{D-91/10+\sqrt{(D-91/10)^2-24D}}.
\]
In the strict-hole unit-disk geometry,
\begin{equation}\label{lt:field70}
 H>\frac{N-(596/5)/d}{2(1+a)}.
\end{equation}
If $D\ge100$, the stronger consequence is
\begin{equation}\label{lt:field100}
 H>\frac{N-(1899/20)/d}{2(1+a)}.
\end{equation}
\end{proposition}
\begin{proof}
First apply Proposition~\ref{mf:D25}, giving $Qd/H<3$.
This enters the third row of Proposition~\ref{lt:W8}.
The six additional moment terms give
$|Z_4-\zeta_s|<175q\mu$ on its parameter rectangle.
This follows by taking the modulus of each displayed coefficient in
\eqref{lt:four-center} at $s=2\pi i$; their total at the corner
is less than $.107054<175(3/70)(1/70)$.
The effective Gaussian radius is therefore
\[
 R_{\rm eff}=34q^2+(175+13400/4900)q\mu.
\]
Lemma~\ref{lt:signed-moment} with $6/5$ and the exact comparisons
\[
 6/5+(1250/24649)34<3,\qquad
 (1250/24649)(175+13400/4900)<91/10
\]
give $\Re K'(b)/H^2<3q^2+(91/10)q\mu$.
On the other hand $\sum(\Re v_j)^2\le H/d$, so
$\Re K'(b)/H^2\ge q-2\mu$.
With $y_0=Qd/H$ this yields
\[
 3y_0^2-(D-91/10)y_0+2D>0.
\]
At $D=70$ the polynomial is positive at $2$, negative at $3$,
decreasing on $[2,3]$, and equals $-16/15$ at $8/3$.
These signs persist for all $D\ge70$; the bound $y_0<3$ removes
the upper quadratic branch. The smaller root is $Y_{70}(D)$.
For $Y_{70}=2+v$ its inverse relation is
$D=151/(5v)+211/10+3v$, decreasing on $0<v\le2/3$.
This proves the assertion for all $D\ge70$.
For the hole field use $k=8/3$, $K_0=21/4$ and
$C_0=3k^2+(91/10)k=228/5$ in Lemma~\ref{lt:hole-field}.
The resulting constant is $596/5$. For $D\ge100$,
Proposition~\ref{mf:D100} gives $k=9/4$ instead.
Then $C_0=2853/80$ and the constant is $1899/20$.
\end{proof}

\begin{proposition}\label{lt:medium139}
If $D\ge139$, then
\[
 |\Im S|<4Q/H,\qquad
 Qd/H<Y(D):=\frac{4D}{D-23/2+\sqrt{(D-23/2)^2-10D}}<9/4.
\]
In the unit-disk strict-hole geometry,
\begin{equation}\label{lt:medium-field}
 H>\frac{N-(2637/32)/d}{2(1+a)}.
\end{equation}
\end{proposition}
\begin{proof}
Proposition~\ref{mf:D100} first gives $Qd/H<9/4$, which
enters the medium row of Proposition~\ref{lt:quartic}.
Equation~\eqref{lt:medium-curvature} gives
\[
 (5/4)y_0^2-(D-23/2)y_0+2D>0,\qquad y_0=Qd/H.
\]
At $D=139$ this is positive at $2$, negative at $9/4$, and
decreases on $[2,9/4]$; the signs persist on the degree-independent
ray. The bound $y_0<9/4$ excludes the upper branch. The inverse relation
at $Y=2+v$ is $D=28/v+33/2+(5/4)v$, decreasing for
$0<v\le1/4$. Finally apply Lemma~\ref{lt:hole-field} with
$k=9/4$, $K_0=4$ and $C_0=(5/4)k^2+(23/2)k$.
Its constant is $2637/32$.
\end{proof}

\begin{proposition}\label{lt:W8-field}
Suppose $Qd/H<3$ and $D\ge5000/91$. Then the second weight-eight
row gives $|\Im S|<(33/5)Q/H$ and
$\Re K'(b)<85/d^2$. In the strict-hole unit-disk geometry it follows
that
\begin{equation}\label{lt:field210}
 H>\frac{N-210/d}{2(1+a)}.
\end{equation}
\end{proposition}
\begin{proof}
The assumed ratio and field enter $q\le273/5000$, $\mu\le91/5000$.
The difference $Z_8-\zeta_s$ has coefficient-norm sum less
than $.397022$ at this corner. Adding the radius gives a
total effective Gaussian radius less than $.477420$.
Every term of this center difference has weight at least two.
As $D$ increases, with $t=(5000/91)/D$, its parameters lie in the
corner polydisk scaled by $t,t^2,t^3$. Consequently the effective
radius is at most $t^2$ times its corner bound.
Since $33/5<8$, the phase hypothesis for the $13/10$ bound in
Lemma~\ref{lt:signed-moment} holds. The coefficient sum in
\eqref{lt:center-recursion} gives
\[
 D^2\Re K'(b)/H^2<84.718<85.
\]
This comparison uses all forty terms and the whole radius; the
coefficient bound is a finite directed evaluation of those explicitly
defined polynomials. Lemma~\ref{lt:hole-field} with $k=3$,
$K_0=33/5$ and the sharper displayed coefficient gives a field
constant less than $2(84.718)+6(33/5)<210$.
\end{proof}

\section{Inverse logarithmic counting and transport}\label{sec:inverse-tools}

Use \eqref{pr:field-notation} and \eqref{pr:normalized-moments},
with $b$ an actual simple zero and $S\ne0$.

\subsection{Inverse rectangles}

\begin{lemma}\label{lt:inverse-log}\label{lt:odd-pi}
Suppose $Q/H<V$, $d\ge d_0>R>0$ and $H\ge H_0>0$. Put
\begin{gather*}
 \Lambda=\frac{VR}{1-R/d_0},\qquad
 e_R=V\left(\frac{R^2}{2}+\frac{R^3}{3(d_0-R)}\right),\qquad
 \widehat R=R-e_R,\\
 B_R=\frac{V}{(1-R/d_0)^2(1-\Lambda)^3},\qquad
 p_*=(1+\Lambda)^{-1},\qquad A_*=(1-\Lambda)^{-1}.
\end{gather*}
Assume $0<\Lambda<1$ and $\widehat R>0$. Let $0<\nu<1$, $W>0$,
and choose positive $e,\epsilon,g$ such that
\begin{gather}
 e>e^{-W},\qquad B_R/H_0<\epsilon,\qquad
 \nu B_R\widehat R<g,\qquad
 W^2<(1-\nu^2)(H_0\widehat R)^2,\label{lt:inverse-corner}\\
 \epsilon(1+e)+ge<1-A_*e,\quad
 \epsilon(1+e)+ge^2<p_*-e,\quad
 \epsilon+ge<p_*.\label{lt:inverse-sides}
\end{gather}
For any positive odd integer $\ell$ with
$\ell\pi\le\nu H\widehat R$, the disk $|z-b|<R$ contains at
least $\ell$ distinct simple original zeros. If
$\nu H\widehat R\ge\pi$, its number of original zeros is
strictly greater than
\begin{equation}\label{lt:inverse-density}
 \nu H\widehat R/\pi-2>\nu(7/22)H\widehat R-2.
\end{equation}
The source $b$ is included once in this count.
\end{lemma}
\begin{proof}
Let $L(z)=\log(p'(z)/p'(b))$ on $|z-b|<d$, with $L(b)=0$.
The finite logarithms give on $|z-b|\le R$
\[
 |L'(z)/S-1|<\Lambda,\qquad
 |L(z)-S(z-b)|<He_R.
\]
For the second bound retain the quadratic term and bound the entire
tail from order three by the displayed geometric series. The first
bound integrated along each joining segment proves injectivity: the
quotient $[L(z)-L(w)]/[S(z-w)]$ lies in $|u-1|<\Lambda$.
Rouch\'e's theorem on $|z-b|=R$ shows that the image contains
$|v|<H\widehat R$. Its analytic inverse $\phi$ satisfies
$\phi(0)=b$. Put $A(v)=S\phi'(v)$. Differentiation gives
\[
 A(0)=1,\qquad \Re A>p_*,\qquad |A|<A_*,\qquad |A'|<B_R/H.
\]
The last assertion uses
$|L''|\le Q/(1-R/d_0)^2$ and
$A'=-S L''(\phi)/(L'(\phi))^3$.

The whole rectangle $|\Re v|\le W$, $|\Im v|\le Y=\ell\pi$
lies strictly inside the inverse disk by \eqref{lt:inverse-corner}.
The actual primitive is
\[
 \frac{Sp(\phi(v))}{p'(b)}
  =\int_0^ve^uA(u)\,du=e^vA(v)-1-E(v),\qquad
 E(v)=\int_0^ve^uA'(u)\,du.
\]
Use the broken path $0,-W,-W+i\Im v,v$. For $v=x+iy$,
\[
 |E(v)|<\epsilon(1+e^x)+ge^{-W}.
\]
For $G=e^vA-1$, the left side has $|G|>1-A_*e^{-W}$,
the right side has $e^{-W}|G|>p_*-e^{-W}$, and each horizontal
side has $|G|>1+p_*e^x$, since $\ell$ is odd. Equations
\eqref{lt:inverse-sides} give $|E|<|G|$ on all four sides,
including every corner. Hence Rouch\'e compares the actual
primitive with $G$.

For the count, use the homotopy $A_t=1-t+tA$,
$0\le t\le1$. Then $\Re A_t\ge p_*$ and $|A_t|\le A_*$.
The same left and right estimates, and the strict negative real part
on the odd-$\pi$ sides, keep $e^vA_t-1$ nonzero throughout the
boundary. It has the same number of zeros as $e^v-1$, namely
$\ell$. Injectivity of $\phi$ and the nonvanishing of $p'$ in
the physical disk give distinct simple original zeros.
Finally, the largest permitted odd integer is
$1+2\lfloor(\nu H\widehat R/\pi-1)/2\rfloor$, which is
strictly larger than $\nu H\widehat R/\pi-2$.
\end{proof}

\begin{lemma}\label{lt:quarter-pi}
Use the notation and hypotheses of Lemma~\ref{lt:inverse-log},
replacing the third inequality in \eqref{lt:inverse-sides} by
\[
 \epsilon+ge<\frac{p_*}{1+p_*+A_*}.
\]
For an integer $k\ge0$ with $Y=(2k+1/2)\pi\le\nu H\widehat R$,
put $\mathcal R_Y=\{v:|\Re v|<W,\ |\Im v|<Y\}$.
Then $\phi(\mathcal R_Y)\subset\{z:|z-b|<R\}$ contains $2k+1$
simple original zeros. In particular, when
$\nu H\widehat R/\pi\ge1/2$, the number of original zeros in
$|z-b|<R$ is strictly greater than
$\nu H\widehat R/\pi-3/2$.
\end{lemma}
\begin{proof}
On these horizontal sides $e^{iy}=\pm i$. For $u=e^x>0$,
\[
 |e^vA(v)-1|\ge\max(1-A_*u,p_*u),\qquad
 \frac{\max(1-A_*u,p_*u)}{1+u}\ge\frac{p_*}{1+p_*+A_*}.
\]
The two quotients meet at $u=1/(p_*+A_*)$. This gives the
horizontal Rouch\'e comparison. On $y=\pm Y$, the amplitude homotopy
satisfies
\[
 \Im(e^vA_t-1)=e^x\sin y\,\Re A_t,
\]
which is positive on the upper side and negative on the lower side.
Here $\Re A_t>0$.
The model $e^v-1$ has exactly
$2k+1$ zeros. Choose the largest permitted $k$ to obtain the
last bound.
\end{proof}

The same inverse lemmas allow a sharper logarithmic radius. For any
integer $M\ge2$ one may replace $e_R$ by
\begin{equation}\label{lt:higher-log-loss}
 V\left(\sum_{k=2}^M\frac{R^k}{k d_0^{k-2}}
       +\frac{R^{M+1}}{(M+1)d_0^{M-1}(1-R/d_0)}\right).
\end{equation}
This retains every logarithmic order beyond $M$; all the other
constants and the full-boundary proof are unchanged. To use exact
integer counts, it is sufficient to verify
$\ell\pi_{\rm up}\le\nu H_0\widehat R$ for a rational upper
bound $\pi_{\rm up}>\pi$. If the source root is inserted separately
in a later weighted sum, use $\ell-1$ other roots, since $b$ was
already counted. The one-sided and narrow-strip refinements needed
near the origin are proved in Section~\ref{sec:origin}.

\subsection{Schur coupling and positive cumulative sums}

\begin{lemma}\label{lt:schur}
There is a nonnegative matrix $(c_{ij})$, indexed by all original
zeros and all critical points, with
\[
 \sum_jc_{ij}=1/n,\qquad \sum_ic_{ij}=1/N,\qquad
 N\sum_ic_{ij}z_i=\xi_j.
\]
If $|z_i|\le1$, then
$N\sum_ic_{ij}|z_i-\xi_j|^2\le1-|\xi_j|^2$.
\end{lemma}
\begin{proof}
Let $T=\diag(z_i)$, $e=n^{-1/2}(1,\ldots,1)^T$ and compress
$T$ to $e^\perp$. Its characteristic polynomial is $p'/n$ after
making $p$ monic. Indeed its cofactor in the $e$ direction is
$p(z)e^*(zI-T)^{-1}e=p'(z)/n$, first off the zeros and then
identically. Lift a unitary Schur basis of the compression to an
$n$ by $N$ matrix $U$. Its columns are orthonormal,
$UU^*=I-ee^*$, and the diagonal entries of $U^*TU$ are
$\xi_j$. Set $c_{ij}=|U_{ij}|^2/N$. All three identities follow.
Expanding around the barycenter gives
\[
 N\sum_ic_{ij}|z_i-\xi_j|^2
  =\sum_i|U_{ij}|^2|z_i|^2-|\xi_j|^2\le1-|\xi_j|^2.
\]
\end{proof}

\begin{lemma}\label{lt:transport}
Assume the unit-disk strict-hole geometry of
Lemma~\ref{pr:edge-reduction}, with the chosen simple rightmost zero
$b=x+iy$, $x\ge a$. For $0<a\le A<1$, define
\[
 \Gamma=\frac2{2-A^2},\quad u=\Gamma A,\quad
 q_A=\frac{u}{1-u},\quad E_A=\frac{2\Gamma}{(1-u)^2},
\]
and assume $u<1/2$. For $\xi_j=r_j\omega_j$, $|\omega_j|=1$, put
\[
 g=\frac1N\sum_j\frac1{b-\omega_j},\quad k=\Re g,\quad h=H/N,
 \quad \eta=\sum_{i,j}c_{ij}\frac{|z_i-\xi_j|^2}{|b-\xi_j|^2}.
\]
Then $k>0$, $h>(1-q_A)k$, and $\eta<E_Aak$.
\end{lemma}
\begin{proof}
The hole gives $1-a<r_j\le1$, and with
$t_j=a/2-\Re\omega_j$ it gives
\[
 (1-r_j)(1+r_j-a^2)<2ar_jt_j.
\]
In particular $t_j>0$. Since $2r/(1+r-a^2)$ increases on
$0<r\le1$, $1-r_j<\Gamma at_j$. Also
$t_j<x-\Re\omega_j\le|b-\omega_j|$. Hence
\[
 |b-r_j\omega_j|>(1-u)|b-\omega_j|,\qquad
 \left|\frac1{b-r_j\omega_j}-\frac1{b-\omega_j}\right|
 <q_A\Re\frac1{b-\omega_j}.
\]
Averaging gives $|S/N-g|<q_Ak$ and the field bound. The variance
part of Lemma~\ref{lt:schur}, together with
$1-r_j^2<2\Gamma at_j$, gives
\[
 \eta\le\frac1N\sum_j\frac{1-r_j^2}{|b-r_j\omega_j|^2}
 <\frac{2\Gamma a}{(1-u)^2N}
       \sum_j\frac{t_j}{|b-\omega_j|^2}\le E_Aak.
\]
\end{proof}

\begin{lemma}\label{lt:abel}
Suppose the number $L_j$ of original roots in nested disks
$|z-b|<t_jd$, where $0<t_1<\cdots<t_m<1$, satisfies
$L_j>s_jD-2$ with $0=s_0<s_1<\cdots<s_m$.
If $D\ge D_0>0$, $N\ge N_0$, and
\[
 c=\frac{N_0}{N_0+1}
 \left(\sum_{j=1}^m(1-t_j)^2(s_j-s_{j-1})-\frac2{D_0}\right)>0,
\]
then $\eta>chd$. Under Lemma~\ref{lt:transport},
$d<E_Aa/[(1-q_A)c]$.
\end{lemma}
\begin{proof}
Every root in the $j$th disk has
$|z_i-\xi_k|>(1-t_j)|b-\xi_k|$ for every $k$.
Its entire row in the coupling contributes more than
$w_j/n$, where $w_j=(1-t_j)^2$. Positive summation of cumulative
counts gives
\[
 n\eta>\sum_{j<m}(w_j-w_{j+1})L_j+w_mL_m
 >D\sum_jw_j(s_j-s_{j-1})-2w_1.
\]
Use $w_1<1$, $N/(N+1)\ge N_0/(N_0+1)$, and $D\ge D_0$.
Combining the result with $\eta<E_Aak<E_Aah/(1-q_A)$ proves
the last assertion. The cumulative counts enter the Abel sum with
nonnegative coefficients.
\end{proof}

For root-product weights, let $B_p=\prod_{i=1}^n f_i$, where
$0\le f_i\le1$ for every original zero, and set $-\log0=+\infty$.
If every root in the $j$th nested region has $f_i\le e^{-w_j}$,
where $w_1\ge\cdots\ge w_m\ge0$, and its cumulative count is
at least $K_j$, then
\begin{equation}\label{lt:abel-product}
 -\log B_p\ge w_mK_m+\sum_{j<m}(w_j-w_{j+1})K_j.
\end{equation}
This is the pointwise step-function lower estimate followed by
summation over every root. If $K_j=s_jH_*-2$, the budget is
$H_*\sum_jw_j(s_j-s_{j-1})-2w_1$.

\subsection{Crescent geometry and half-disk factors}

\begin{lemma}\label{lt:crescent}
Under the strict-hole geometry with $0<a\le x$, $b=x+iy$,
$y\ge0$, $|b|\le1$, put $t=\sqrt{1-a^2/4}$.
Every critical point satisfies
\begin{equation}\label{lt:crescent-gap}
 |b-\xi|^2>(x-a/2)^2+(t-y)^2.
\end{equation}
Consequently $d>x-a/2$, $d>1-y$, $d>1-|b|$, and
\begin{equation}\label{lt:quadrature-gap}
 d^2>(1-|b|)^2+a^2/4.
\end{equation}
\end{lemma}
\begin{proof}
Let $e=a/2+it$, write $\xi=e-u+iv$ and put $\lambda=y/t$.
Then $u>0$ by \eqref{pr:critical-halfplane}, and $0\le\lambda<1$
because $y\le\sqrt{1-x^2}\le\sqrt{1-a^2}<t$.
The exact identity
\[
 |b-\xi|^2-|b-e|^2
 =(1-\lambda)(|\xi-a|^2-1)
       +[2(x-a)+a\lambda]u+\lambda(u^2+v^2)>0
\]
proves the crescent comparison. Its difference from
$(1-y)^2$ is $x(x-a)+2y(1-t)\ge0$.
For fixed $r=|b|$ and $a$, the expression
$1+r^2-ax-2t\sqrt{r^2-x^2}$ increases for $a\le x<r$:
its derivative is $-a+2tx/\sqrt{r^2-x^2}>0$.
Continuity treats $x=r$. Thus its minimum occurs at $x=a$.
Since $2r-5a^2/4>0$ and
\[
 (2r-5a^2/4)^2-4(1-a^2/4)(r^2-a^2)
       =a^2(r-1)(r-4)+9a^4/16\ge0,
\]
squaring in the positive direction gives
\eqref{lt:quadrature-gap}. This includes $r=1$.
\end{proof}

\begin{lemma}\label{lt:direct-D70}
Under the preceding hypotheses put $c=x/2$. Then
\[
 D>\frac{Nc}{2\sqrt{2-2c+c^2}-2\sqrt2c(1-c)}.
\]
In particular, if $N\ge6250$ and $x\ge30501/500000$, then $D>70$.
\end{lemma}
\begin{proof}
For fixed $b$, the squared right side of \eqref{lt:crescent-gap}
decreases as $a$ increases to $x$, since its derivative is
$-x+ay/(2\sqrt{1-a^2/4})<0$ for $0<a\le x\le1$.
Put $t=\sqrt{1-c^2}$, $v=t-y\ge0$ and $G=\sqrt{c^2+v^2}$.
Thus $d>G$ and $p_b=1-|b|^2=-3c^2+2tv-v^2\le2v-3c^2-v^2$.
Since $\Re S>cQ$, the radial identity implies
\[
 D>\frac{NcG}{2|b|c+p_b}.
\]
The identities
$c+(1-c)v\le\sqrt{1+(1-c)^2}G$ and
$(3c^2+v^2)^2-8c^2G^2=(v^2-c^2)^2\ge0$,
together with $|b|\le1-p_b/2$, give
\begin{align*}
 2|b|c+p_b&\le2[c+(1-c)v]-(1-c)(3c^2+v^2)\\
 &\le
 [2\sqrt{2-2c+c^2}-2\sqrt2c(1-c)]G.
\end{align*}
The bracket is positive for $0<c\le1/2$.
For $c\ge c_0=30501/1000000$, one has
$2-2c+c^2<49/25<2$. Therefore the denominator is less than
$(14/5)(1-c+c^2)$ and
\[
 D>\frac{5N}{14}\frac{c}{1-c+c^2}
 \ge\frac{476578125000000}{6793005177007}>70.
\]
The rational function increases with $c$ because its derivative is
$(1-c^2)/(1-c+c^2)^2>0$.
\end{proof}

\begin{lemma}\label{lt:same-gap-field}
Assume all original zeros lie in $\overline\D$ and $a>0$ satisfies
$|a-\xi_j|>1$ for every critical point. Let $b=x+iy$ be a simple
rightmost zero with $x\ge a$, $y\ge0$, and $|b|\le1$, as in
Lemma~\ref{lt:crescent}. Put $c=x-a/2$ and $p_b=1-|b|^2$.
If $|\Im S|<K_0Q/H$ and $Nc>2K_0$, then
\[
 D>\frac{Nc^2-2K_0c}{2xc+p_b}.
\]
If $N\ge N_0$, $a\le a_U$, $x\ge x_0>a_U/2$, $p_b\le p_U$, and
$N_0(x_0-a_U/2)>2K_0$, then the uniform bound is obtained by
putting $N=N_0$, $c=x_0-a_U/2$, $x=x_0$, $p_b=p_U$.
\end{lemma}
\begin{proof}
The half-plane inequality gives $d>c$, $Q/H<1/c$. Hence
$N<(2x+p_b/c)H+2K_0/c$ and the first assertion follows.
With $x,p_b$ fixed, the derivative numerator in $c$ is
$2[Nxc^2+p_b(Nc-K_0)]>0$.
After writing $c=x-a_U/2$, the derivative numerator in $x$ is
$2[(Na_U/2+2K_0)c^2+p_b(Nc-K_0)]>0$.
The quotient decreases with $p_b$ and increases with $N$.
These signs also hold at $p_b=0$.
\end{proof}

\begin{lemma}\label{lt:half-disk-factor}
Let $b=x+iy$ be an actual simple rightmost zero, with
$x>0$, $y\ge0$, $0<B_L\le|b|\le B_U<1$,
$x\le X$, and $\gamma=b-a/2$, $0<a\le a_U$.
For a disk radius $R$ put
\[
 P=1-B_U^2,\qquad T=\frac{a_U}{2P},\qquad U=\frac RP,\qquad
 L=1+2TX.
\]
Suppose $R^2<P$ and $0<f<1$ satisfy
\begin{equation}\label{lt:half-disk-gate}
 f^2-(1-f^2B_L^2)(T^2+U^2)>2Uf^2LB_U.
\end{equation}
Then every original zero $z$ with $|z-b|\le R$ obeys
$|(\gamma-z)/(1-\overline z\gamma)|<f$.
\end{lemma}
\begin{proof}
Write $\tau=a/2$, $p_b=1-|b|^2$, $z=b+v$,
$L_\gamma=1-|\gamma|^2=p_b+2\tau x-\tau^2$ and
$A_0=p_b+\tau\overline b$. Rightmostness gives $\Re v\le0$.
The gate implies
\begin{gather*}
 T^2<\frac{f^2}{1-f^2B_L^2}<\frac1{1-B_L^2},\\
 \tau^2\le P^2T^2<\frac{P^2}{1-B_L^2}\le P\le p_b.
\end{gather*}
Thus $L_\gamma>0$, $|\gamma|<1$, and $1-f^2|\gamma|^2>0$.
Direct expansion gives
\[
 f^2|1-\overline z\gamma|^2-|\gamma-z|^2
 =f^2|A_0|^2-\tau^2-(1-f^2|\gamma|^2)|v|^2-2\Re(Vv),
\]
where $V=f^2L_\gamma\overline b+\tau(1-f^2)$.
Thus $\Re V>0$, $\Im V=-f^2L_\gamma y\le0$.
On the half disk $|v|\le R$, $\Re v\le0$,
$\Re(Vv)\le f^2L_\gamma yR$. The resulting lower bound is
\begin{align*}
 f^2p_b^2&-(1-f^2|b|^2)(\tau^2+R^2)
 +2f^2\tau x(p_b-R^2)+f^2\tau^2R^2-2Rf^2L_\gamma y.
\end{align*}
The two terms containing $2f^2\tau x$ and $f^2\tau^2$ are nonnegative.
After division by $p_b^2$, use $p_b\ge P>R^2$ and
\[
 \frac{L_\gamma}{p_b}
 =1+\frac{2\tau x-\tau^2}{p_b}\le1+2TX=L.
\]
The resulting bound is at least
$f^2-(1-f^2B_L^2)(T^2+U^2)-2Uf^2LB_U>0$.
The positive difference also excludes a zero denominator.
\end{proof}

For $N\ge N_0$, if a field estimate has the form $H>(N-c)/F$,
with $c\ge0$, $F>0$ and $(N_0-c)/F>H_*>0$, then
\begin{equation}\label{lt:energy-ray}
 \begin{gathered}
 \frac HN>\frac{1-c/N}{F}\ge\frac{1-c/N_0}{F}
             >\frac{H_*}{N_0},\\
 Q\ge\frac{H^2}{N}>\frac{N}{N_0^2}H_*^2\ge\frac{H_*^2}{N_0}.
 \end{gathered}
\end{equation}
The lower bound $(1-c/N)/F$ is nondecreasing in $N$, so these
estimates apply on the whole ray $N\ge N_0$.

\section{The boundary and the compact range}\label{bc:section}

We prove in this section the assertion of the main theorem whenever the
modulus of the proposed exceptional point is at least $1/10$.  Some of
the moment estimates obtained along the way will also be useful in the
remaining range.  Throughout the section, $P$ is monic of degree $n\ge6251$,
$N=n-1$, and all its zeros $z_1,\ldots,z_n$ belong to $\overline\D$.
Its critical points, with multiplicity, are denoted by
$\xi_1,\ldots,\xi_N$.  We put
\begin{equation}\label{bc:moments-definition}
 \mu=\frac1n\sum_i z_i=\frac1N\sum_j\xi_j,\qquad
 Q=\frac1N\sum_j|\xi_j|^2,\qquad
 L(z)=\frac1N\sum_j\log|z-\xi_j|.
\end{equation}
The letter $Q$ in this section denotes a second moment; reciprocal
energy will always be written $E_b$.

Suppose, towards a contradiction, that a real point $0<a\leq1$ in the
convex hull of the zeros satisfies
\begin{equation}\label{bc:hole}
 |a-\xi_j|>1\quad(1\leq j\leq N).
\end{equation}
Choose an actual zero $b=x+iy$ of largest real part, so $x\geq a$;
reflection permits us to assume $y\geq0$.  Write
\begin{equation}\label{bc:local-notation}
 r=|b-a|,\quad d=\min_j|b-\xi_j|,\quad
 S=\sum_j\frac1{b-\xi_j},\quad H=|S|,\quad
 E_b=\sum_j\frac1{|b-\xi_j|^2}.
\end{equation}
Gauss--Lucas and \eqref{bc:hole} imply
\begin{equation}\label{bc:entrance}
 |\xi_j|\leq1,\quad \Re\xi_j<a/2,\quad
 r^2\leq1-a^2<1,\quad d>1-r,\quad
 \Re S>\frac a2E_b.
\end{equation}
In particular $b$ is simple.  We use the actual-zero energy inequality
$E_b\geq N$ of Proposition~\ref{ar:energy}.  It excludes $r=0$ and also excludes a
strict witness which is itself an original zero.  Consequently all
applications below have $0<a<1$ and $r>0$.

\subsection{A product inequality retaining the original zeros}

The fiber arguments use the following fact for arbitrary equal-value
fibers.

\begin{lemma}[Equal-value perpendicular bisector]\label{bc:equal-fiber}
If $P(v)=P(w)$ and $v\ne w$, the critical points of $P$ cannot all
belong to either one of the two open half-planes bounded by the
perpendicular bisector of $[v,w]$.
\end{lemma}

\begin{proof}
We prove the needed half-plane polarization fact first.  Let $f$ be a
nonzero polynomial of degree at most $m$ with no zero in
$\{\Re z>0\}$.  For $\Re w>0$ put
\[
 T_w f(z)=f(z)+\frac{w-z}{m}f'(z).
\]
This polynomial has degree at most $m-1$ and is nonzero throughout
$\{\Re z>0\}$.  To see the latter assertion, write the zeros of
$f$ as $\zeta_1,\ldots,\zeta_d$, where $d\leq m$, and fix
$z$ with $x=\Re z>0$.  Each $v_j=(z-\zeta_j)^{-1}$ satisfies
$\Re v_j\geq x|v_j|^2$.  The disk
$\{v:\Re v\geq x|v|^2\}$ is convex and contains zero; hence
$v=m^{-1}\sum_{j=1}^d v_j$ is in that disk as well.  If $v=0$,
then $T_w f(z)=f(z)\ne0$.  Otherwise $T_w f(z)=0$ would imply
$w=z-1/v$, whose real part is at most zero, a contradiction.

Let $F(\zeta_1,\ldots,\zeta_m)$ be symmetric and affine in each
variable separately.  If $f(z)=F(z,\ldots,z)$ has no zero in the
right half-plane, neither does $F$ when all its variables are in that
half-plane.  Indeed, replacing one of the remaining equal variables
$z$ by $w$ applies exactly the operation $T_w$, with $m$ replaced by
the number of remaining variables.  Iteration proves the assertion.
Reflection proves its left-half-plane version.

An affine change of variable reduces the asserted bisector to the
imaginary axis and the two equal-value points to $-1,1$.  Consider
\[
 F(\zeta_1,\ldots,\zeta_N)
   =\int_{-1}^{1}\prod_{j=1}^N(t-\zeta_j)\,dt.
\]
Its diagonal is
\[
 F(z,\ldots,z)=\frac{(1-z)^{N+1}-(-1-z)^{N+1}}{N+1}.
\]
Every zero of this polynomial has $|1-z|=|-1-z|$, and hence is on
the imaginary axis.  Half-plane polarization therefore makes $F$
nonzero if all the critical points are strictly on the same side of
that axis.  On the other hand, $F(\xi_1,\ldots,\xi_N)$ is a nonzero
constant multiple of $P(1)-P(-1)=0$.  This contradiction proves the
lemma.
\end{proof}

For $|\zeta|<1$ define
\[
 \mathcal B_P(\zeta)=
 \prod_{i=1}^n\left|\frac{\zeta-z_i}{1-\overline{z_i}\zeta}\right|.
\]
Every factor is at most one.  The following lower bound for the product
retains the reciprocal energy $E_b$ at the original zero $b$.

\begin{lemma}[Adaptive product inequality]\label{bc:blaschke}
Under \eqref{bc:hole}--\eqref{bc:local-notation}, put
\begin{equation}\label{bc:adaptive-def}
 u=\min(r^2,1/2),\qquad t=1-u,\qquad
 \gamma=b-ta,\qquad C=ta\bigl(2x-(1+t)a\bigr).
\end{equation}
Then $|\gamma|<1$, $C\geq a^2u(1-u)>0$, and
\begin{equation}\label{bc:blaschke-bound}
 C\exp(CE_b/2)
 <2(1+|b|)(1+|\gamma|)\mathcal B_P(\gamma)
 \leq8\mathcal B_P(\gamma).
\end{equation}
In particular, if $n\geq6251$, then
\begin{equation}\label{bc:Ccap}
 C<C_*:=\frac1{388}.
\end{equation}
\end{lemma}

\begin{proof}
Set $k=1-|\gamma|^2$.  Since
$|\gamma|^2=r^2+2ua(x-a)+u^2a^2>u$ and
$|b|^2-|\gamma|^2=ta(2x-ta)>0$, we have $0<k<t<1$.
Direct expansion gives, for every critical point,
\begin{align}
 |b-\xi|^2-|1-\overline\gamma\xi|^2
 &=t\bigl(|\xi-a|^2-1\bigr)
   +C+(t-k)(1-|\xi|^2)>C.             \label{bc:pointwise-product}
\end{align}
Also $|1-\overline\gamma\xi|^2-|\gamma-\xi|^2
=k(1-|\xi|^2)\geq0$.

Let $M=1+|b|$ and $\eta=C/(2M)$.  For every $z\in D(b,\eta)$,
\[
 |z-\xi|^2>(|b-\xi|-\eta)^2
 \geq |b-\xi|^2-2M\eta+\eta^2>|\gamma-\xi|^2.
\]
Here $\eta<|b-\xi|/2$, by \eqref{bc:pointwise-product}, and
$\eta<ta=|b-\gamma|$.  Lemma~\ref{bc:equal-fiber}
therefore forbids two equal-value points, one in $D(b,\eta)$ and the
other equal to $\gamma$.  In particular $P(\gamma)\ne0$.
If $|P(\gamma)|<\eta|P'(b)|/n$, the logarithmic derivative of
$P(z)-P(\gamma)$ at $b$ would give
\[
 \left|\sum_{P(\zeta)=P(\gamma)}\frac1{b-\zeta}\right|
 =\frac{|P'(b)|}{|P(\gamma)|}>\frac n\eta.
\]
At least one of these $n$ zeros, counted with multiplicity, would be in
$D(b,\eta)$, again a contradiction.  Hence
\begin{equation}\label{bc:fiber-lower}
 |P(\gamma)|\geq\frac{\eta}{n}|P'(b)|.
\end{equation}

Let $q=P'/n$, and introduce the reciprocal polynomials
\[
 P^*(z)=\prod_i(1-\overline{z_i}z),\qquad
 q^*(z)=\prod_j(1-\overline{\xi_j}z).
\]
From \eqref{bc:pointwise-product} and $-\log(1-v)>v$ for $0<v<1$,
\begin{equation}\label{bc:critical-product}
 \frac{|P'(b)|}{n|q^*(\gamma)|}>
 \exp\left(\frac C2\sum_j|b-\xi_j|^{-2}\right).
\end{equation}
The exact reversed derivative identity is
\[
 \frac{q^*(\gamma)}{P^*(\gamma)}
 =\frac1n\sum_i\frac1{1-\overline{z_i}\gamma}.
\]
For $|w|\leq g<1$ one has $\Re(1-w)^{-1}\geq(1+g)^{-1}$,
as follows by multiplying by $|1-w|^2$.  Thus the modulus of the last
display is at least $(1+|\gamma|)^{-1}$.  Combining it with
\eqref{bc:fiber-lower} and \eqref{bc:critical-product} proves
\eqref{bc:blaschke-bound}.  Finally
\[
 \frac{6250}{776}>\log(8\cdot388)
\]
implies \eqref{bc:Ccap}, since $C\exp(NC/2)$ increases strictly in
$C>0$ and in $N$.  This numerical comparison is certified by the
positive exponential series, as described below.
\end{proof}

The normalization at the origin uses the following movement estimate.

\begin{lemma}[Movement of an equal-value fiber]\label{bc:movement}
Let $\omega=b/|b|$, let $a_0\leq a$, and suppose
$2\Re(\overline\omega\xi_j)<d_0$ for all $j$.  If positive numbers
$s,\delta$ satisfy
\begin{equation}\label{bc:movement-gates}
 a_0-s>\delta,\qquad
 (a_0-s)(a_0+s-d_0)>\delta(4+\delta),
\end{equation}
then $|P(s\omega)|\geq\delta|P'(b)|/n$.
\end{lemma}

\begin{proof}
Writing $v=|b|\geq a_0$, the squared-distance difference at $b$ is
\[
 |b-\xi|^2-|s\omega-\xi|^2
 =(v-s)\bigl(v+s-2\Re(\overline\omega\xi)\bigr)
 >(a_0-s)(a_0+s-d_0).
\]
Moving $b$ by at most $\delta$ decreases this difference by at most
$\delta(4+\delta)$.  Thus every point of $D(b,\delta)$ is farther
from every critical point than $s\omega$ is.  The first condition in
\eqref{bc:movement-gates} puts $s\omega$ outside that disk.  The same
fiber and logarithmic-derivative argument as in
\eqref{bc:fiber-lower} proves the assertion.
\end{proof}

\subsection{Normalization, annuli, and moment estimates}

\begin{proposition}[Three moment estimates]\label{bc:moment-proposition}
Assume $n\geq6251$.  Under \eqref{bc:hole}, each row of
Table~\ref{bc:moment-conclusions} implies $P(0)\ne0$.  With
\begin{equation}\label{bc:B-and-error}
 B=-\frac1n\log|P(0)|,\qquad
 \mathcal E(r;a,N)=\frac{8r}{aN}
       +\frac{r^2}{2(1-r)(1-2r)},
\end{equation}
its conclusions are
\begin{align}
 B&<\frac\ell n+\mathcal E(r;a,N),                 \label{bc:moment-B}\\
 |\mu|&\leq\kappa B+n^{-2},                       \label{bc:moment-mean}\\
 Q&<1-a^2+w\left(\frac\ell N+\mathcal E(r;a,N)\right)
                   +2n^{-2}.                    \label{bc:moment-Q}
\end{align}
\begin{table}[ht]
\centering
\caption{The moment estimates.  The hypotheses use the actual $a,r$.}
\label{bc:moment-conclusions}
\begin{tabular}{@{}cccccc@{}}\toprule
Row&$a\geq$&$r\leq$&$\ell$&$\kappa$&$w$\\\midrule
M0&$.9$&$.057$&$4.42$&$2.5482$&$5.098$\\
M1&$.93$&$.055$&$4.22$&$2.3817$&$4.765$\\
M2&$.95$&$.054$&$4.11$&$2.3108$&$4.623$\\\bottomrule
\end{tabular}
\end{table}
The same formulas hold with the parameters in the four boundary rows
of the tables below, whenever their respective hypotheses hold.
\end{proposition}

We give the normalization argument with parameters; this also supplies
the estimates needed to prove the boundary assertions.  A row specifies
lower $a_0$, upper displacement $e<1/2$, a squared critical-distance
bound $R_c^2<(1-e)^2$, and the constants listed in
Tables~\ref{bc:preliminary-table}--\ref{bc:annulus-table}.
The names B, A, G1, G2 stand, respectively, for the boundary range
$a\geq.975$, and the small displacement ranges
$(a,r)\in[.85,1)\times[0,.02]$,
$[.96,1)\times[0,.04]$, and
$[.965,1)\times[0,.045]$.
For row B the displacement bound is a conclusion of
Lemma~\ref{bc:blaschke}: indeed
\[
 (.975)^2(.053)^2\bigl(1-(.053)^2\bigr)>1/388,
\]
and $u(1-u)$ increases on $[0,1/2]$; hence $r<.053$.

\begin{proof}[Proof of Proposition~\ref{bc:moment-proposition}]
First,
\[
 |\omega-1|^2=\frac{2y^2}{|b|(|b|+x)}\leq r^2/a^2,
 \qquad
 2\Re(\overline\omega\xi_j)<a+2r/a.
\]
All rows satisfy $2e<a_0^2$.  The function $a+2e/a$ therefore
increases on $[a_0,1]$, and is at most $1+2e<d_0$ there.
The two pairs $(s_1,\delta_1),(s_2,\delta_2)$ in the tables satisfy
\eqref{bc:movement-gates}.  Lemma~\ref{bc:movement} applies to both.

Put $k=N/n$, $A_n=\delta_1^{2/n}$, and
$A_-=1-\alpha_0/n$.  The inequality
$\alpha_0/2>\log(1/\delta_1)$ gives $A_n>A_->0$.
Original-zero AM--GM at $s_1\omega$ gives
\begin{equation}\label{bc:directional-mean}
 2s_1\Re(\overline\omega\mu)
 \leq1+s_1^2-A_n e^{2kL(b)}.
\end{equation}
On $R_c\leq s\leq2$ the logarithmic chord
\begin{equation}\label{bc:coarse-chord}
 s^2\leq R_c^2+\frac{22}5\log(s/R_c)
\end{equation}
follows by convexity of $e^{2v}$ in $v=\log s$ and its endpoint
inequality.  The latter follows already from
$\log(1/R_c)\geq(1-R_c^2)/2$ and the displayed certificate below.
Average \eqref{bc:coarse-chord} at $s=|b-\xi_j|$ and use
\eqref{bc:directional-mean}.  For $v=|b|$ this yields
\begin{equation}\label{bc:F-definition}
 Q\leq F(L(b)):=
 \frac v{s_1}(1+s_1^2-A_ne^{2kL(b)})+R_c^2-v^2
     +\frac{22}5(L(b)-\log R_c).
\end{equation}
Let
\[
 T=\frac{1+s_1^2-A_-R_c^2}{s_1},\qquad
 F_0=a_0T+R_c^2-a_0^2.
\]
The row comparisons $T<2a_0$ and
$22/5-2a_0kA_-R_c^2/s_1<0$ show that $F$ decreases on the
entire half-line $L(b)\geq\log R_c$.  Indeed $F''<0$, while
$R_c^{2k}\geq R_c^2$; the quadratic $vT-v^2$ decreases for
$v\geq a_0$.  The resulting chain is
\begin{equation}\label{bc:coarse-correct-chain}
 Q\leq F(L(b))\leq F(\log R_c)\leq F_0<q_c,
 \qquad |\mu|<m_c,
\end{equation}
where the last inequality uses $q_c<m_c^2$.

The same calculation at every $z\omega$, $0\leq z\leq s_2$, bounds
the average squared critical distance by
\[
 z^2+\frac{v-z}{s_1}(1+s_1^2-A_ne^{2kL(b)})+R_c^2-v^2
       +\frac{22}5(L(b)-\log R_c).
\]
Subtract $.9e^{2L(b)}$.  Its derivative in $L(b)$ is negative
throughout the same half-line, because
\[
 2R_c^2\left(kA_-\frac{a_0-s_2}{s_1}+.9\right)>22/5.
\]
At $L(b)=\log R_c$ it is at most
$z^2-zT+F_0-.9R_c^2$.  This is a convex quadratic in $z$;
the two endpoint values, at $0,s_2$, are negative.
Critical-zero AM--GM consequently gives
\begin{equation}\label{bc:primitive-decay}
 |P'(z\omega)|<(.95)^N|P'(b)|\qquad(0\leq z\leq s_2).
\end{equation}
Since $e^{N/20}>200n$ and $s_2<100\delta_2$, integration and the
second movement inequality give
\begin{equation}\label{bc:normalization}
 |P(0)|\geq\frac{\delta_2}{2n}|P'(b)|>0,
 \qquad
 L(b)\leq\frac\ell N-\frac nN B,
\end{equation}
where $e^\ell>2/\delta_2$.  The inequalities
$R_c>1.001p$ and $(1.001)^N>2p/\delta_2$ also give
$|P(0)|>p^n$.  In particular
$0\leq B< -\log R_c+\ell/n<B_c$.

For the annulus iteration, suppose $|\mu|\leq m$,
$Q<q_c$, and
\begin{equation}\label{bc:annulus-gates}
 \rho^2+2\rho m+q_c<s_*^2,\qquad
 \rho<s_*<p,\qquad n(s_*/p)^n<1.
\end{equation}
For every $|z|\leq\rho$, AM--GM gives
$|P'(z)|<ns_*^{N}$.  Integrating on the segment from $0$ to $z$,
\[
 \frac{|P(z)-P(0)|}{|P(0)|}
 <\frac\rho{s_*}n(s_*/p)^n<1.
\]
Thus every original zero has modulus greater than $\rho$.
For $0<\rho\leq s\leq1$, the function
$2\sinh t/t$, $t=-\log s$, increases with $t\geq0$.
It follows that
\[
 s^{-1}-s\leq
 \frac{\rho^{-1}-\rho}{-\log\rho}(-\log s).
\]
Since $-P'(0)/P(0)=\sum_i z_i^{-1}$, any
$\kappa>(\rho^{-1}-\rho)/(-\log\rho)$ therefore gives
\begin{equation}\label{bc:reciprocal-mean}
 |\mu|\leq\kappa B+
 \frac{|P'(0)|}{n|P(0)|}.
\end{equation}
The residual in this formula is bounded by
\[
 \frac{m_c^N}{p^n}=\frac{(m_c/p)^n}{m_c}<2d_s^n<n^{-2},
\]
using $m_c/p<d_s$ and $1/m_c<2$.
Beginning with $m=m_c$,
each successive triple $(\rho,\kappa,m')$ in
Table~\ref{bc:annulus-table} satisfies \eqref{bc:annulus-gates}
and $\kappa B_c+n^{-2}<m'$.  The last triple supplies
\eqref{bc:moment-mean}.

The geometric bounds \eqref{bc:entrance} and $E_b\geq N$ give
\[
 H>Na/2,\qquad E_b/H<2/a,\qquad
 E_b/H^2<4/(Na_0^2),\qquad (Hd)^{-1}<2/[Na_0(1-e)].
\]
The last two bounds are respectively below $7/625$ and $1/270$.
The Fine companion estimate of Proposition~\ref{lt:fine200} therefore gives
\begin{equation}\label{bc:fine-phase}
 |\Im S|<4E_b/H<8/a.
\end{equation}
For $w_0=b-a$, the complete logarithmic series at $b$ converges
absolutely, because $r/d<r/(1-r)<1$.  Its linear term satisfies
\[
 -\Re(w_0S)=-(x-a)\Re S+y\Im S\leq r|\Im S|.
\]
The terms of order at least two have total absolute value, after
averaging, at most
$E_br^2/[2N(1-r/d)]$.  Using $E_b/N<(1-r)^{-2}$ proves
\begin{equation}\label{bc:reverse-log}
 L(a)\leq L(b)+\mathcal E(r;a,N).
\end{equation}
Since $L(a)>0$, \eqref{bc:normalization} proves
\eqref{bc:moment-B}.

Finally $s^2\leq1+w\log s$ on $1\leq s\leq2$, because
$w\log2>3$ and the exponential is convex.  Apply this chord at
$s=|a-\xi_j|$ and use
\eqref{bc:normalization}, \eqref{bc:reciprocal-mean}, and
\eqref{bc:reverse-log}.  We obtain
\[
 Q<1-a^2+w\left(\frac\ell N+\mathcal E(r;a,N)\right)
      +2n^{-2}+(2a\kappa-wn/N)B.
\]
The coefficient of $B$ is negative, since $w>2\kappa$ and $a\leq1$.
This proves \eqref{bc:moment-Q}.
\end{proof}


\begin{table}[ht]
\centering\small
\caption{Preliminary parameters.  In every row $s_1=.25$ and $s_2=.45$.}
\label{bc:preliminary-table}
\begin{tabular}{@{}crrrrrrrr@{}}\toprule
Row&$a_0$&$e$&$R_c^2$&$d_0$&$\delta_1$&$\delta_2$&$\alpha_0$&$\ell$\\\midrule
B & 0.975 & 0.053 & 0.896708 & 1.10611 & 0.01976 & 0.03838 & 7.86 & 3.97\\
A & 0.85 & 0.02 & 0.960299 & 1.04011 & 0.00823 & 0.02383 & 9.61 & 4.44\\
G1 & 0.96 & 0.04 & 0.921499 & 1.08011 & 0.02114 & 0.03857 & 7.73 & 3.96\\
G2 & 0.965 & 0.045 & 0.911924 & 1.09011 & 0.02047 & 0.03836 & 7.79 & 3.97\\
M0 & 0.9 & 0.057 & 0.889148 & 1.11411 & 0.00534 & 0.02433 & 10.48 & 4.42\\
M1 & 0.93 & 0.055 & 0.892924 & 1.11011 & 0.01089 & 0.0297 & 9.05 & 4.22\\
M2 & 0.95 & 0.054 & 0.894815 & 1.10811 & 0.01474 & 0.03346 & 8.45 & 4.11\\
\bottomrule\end{tabular}
\end{table}
\begin{table}[ht]
\centering\small
\caption{Coarse normalization and annulus parameters.}
\label{bc:coarse-table}
\begin{tabular}{@{}crrrrrrr@{}}\toprule
Row&$q_c$&$m_c$&$p$&$s_*$&$B_c$&$d_s$&$w$\\\midrule
B & 0.59717 & 0.77287 & 0.94595 & 0.94414 & 0.055158 & 0.81705 & 4.494\\
A & 0.59041 & 0.76849 & 0.97891 & 0.9771 & 0.020976 & 0.78506 & 4.33\\
G1 & 0.54582 & 0.7389 & 0.95893 & 0.95712 & 0.041521 & 0.77056 & 4.33\\
G2 & 0.56641 & 0.75271 & 0.95394 & 0.95213 & 0.046745 & 0.78907 & 4.354\\
M0 & 0.70869 & 0.84194 & 0.94195 & 0.94014 & 0.059463 & 0.89384 & 5.098\\
M1 & 0.66376 & 0.81482 & 0.94395 & 0.94214 & 0.057312 & 0.86322 & 4.765\\
M2 & 0.63422 & 0.79648 & 0.94495 & 0.94314 & 0.056237 & 0.8429 & 4.623\\
\bottomrule\end{tabular}
\end{table}
\begin{longtable}{@{}clll@{}}
\caption{Successive annulus triples $(\rho,\kappa,m')$.
A dash means that only the resulting inequality $|\mu|\leq\kappa B+n^{-2}$ is retained.}
\label{bc:annulus-table}\\
\toprule Row&$\rho$&$\kappa$&$m'$\\\midrule\endfirsthead
\toprule Row&$\rho$&$\kappa$&$m'$\\\midrule\endhead
B & 0.169 & 3.2333 & 0.178344\\
 & 0.39 & 2.3091 & 0.127367\\
 & 0.427 & 2.2504 & 0.124129\\
 & 0.43 & 2.2462 & 0.123897\\
 & 0.43 & 2.2462 & --\\
A & 0.206 & 2.9424 & 0.061721\\
 & 0.543 & 2.1268 & 0.044613\\
 & 0.558 & 2.1155 & 0.044376\\
 & 0.558 & 2.1155 & --\\
G1 & 0.216 & 2.8802 & 0.11959\\
 & 0.498 & 2.1661 & 0.08994\\
 & 0.523 & 2.1432 & 0.088989\\
 & 0.523 & 2.1432 & --\\
G2 & 0.197 & 3.0035 & 0.1404\\
 & 0.457 & 2.2109 & 0.10335\\
 & 0.486 & 2.1783 & 0.101826\\
 & 0.488 & 2.1762 & 0.101728\\
 & 0.488 & 2.1762 & --\\
M0 & 0.096 & 4.4043 & 0.261894\\
 & 0.229 & 2.8073 & 0.166932\\
 & 0.281 & 2.5822 & 0.153547\\
 & 0.29 & 2.5515 & 0.151721\\
 & 0.291 & 2.5482 & --\\
M1 & 0.125 & 3.7872 & 0.217054\\
 & 0.301 & 2.5165 & 0.144227\\
 & 0.348 & 2.3928 & 0.137138\\
 & 0.353 & 2.3817 & 0.136502\\
 & 0.353 & 2.3817 & --\\
M2 & 0.144 & 3.5093 & 0.197354\\
 & 0.343 & 2.4043 & 0.135212\\
 & 0.385 & 2.318 & 0.130359\\
 & 0.389 & 2.3108 & 0.129954\\
 & 0.389 & 2.3108 & --\\
\bottomrule\end{longtable}
\begin{table}[ht]
\centering\small
\caption{Uniform bounds $Q<q_0$, $B<B_0$, $|\mu|<m_0$, and
$|P(0)|>p_f^n$ furnished by the moment formulas at $n=6251$.}
\label{bc:refined-table}
\begin{tabular}{@{}crrrr@{}}\toprule
Row&$q_0$&$B_0$&$m_0$&$p_f$\\\midrule
B & 0.06001 & 0.0023642 & 0.005311 & 0.997625\\
A & 0.28164 & 0.0009535 & 0.0020177 & 0.999036\\
G1 & 0.08531 & 0.0015932 & 0.0034151 & 0.998396\\
G2 & 0.07689 & 0.0018604 & 0.0040492 & 0.998129\\
M0 & 0.20395 & 0.0027331 & 0.0069651 & 0.997256\\
M1 & 0.14726 & 0.0025497 & 0.0060732 & 0.99744\\
M2 & 0.10888 & 0.0024586 & 0.0056819 & 0.997531\\
\bottomrule\end{tabular}
\end{table}

The numerical statements used in the preceding proof are finite rational
comparisons, with logarithms bounded by convergent positive series.
The remaining comparisons are as follows, with $n_0=6251$,
$N_0=6250$, $k_0=N_0/n_0$, and $A_-=1-\alpha_0/n_0$:
\begin{align}
 &2e<a_0^2,\quad 1+2e<d_0,\quad R_c^2<(1-e)^2,
       \quad\eqref{bc:movement-gates}\text{ for }j=1,2;
                                                        \label{bc:pre-gates-1}\\
 &R_c^2+\tfrac{22}5\bigl(\log2+(1-R_c^2)/2\bigr)>4,
       \quad \alpha_0/2>\log(1/\delta_1),
       \quad e^\ell>2/\delta_2;                         \label{bc:pre-gates-2}\\
 &T<2a_0,\quad F_0<q_c<m_c^2,\quad
       2a_0k_0A_-R_c^2/s_1>22/5;                       \label{bc:pre-gates-3}\\
 &2R_c^2[k_0A_-(a_0-s_2)/s_1+.9]>22/5,
       \quad z^2-zT+F_0<.9R_c^2\quad(z=0,s_2);         \label{bc:pre-gates-4}\\
 &N_0/20>\log(200n_0),\quad s_2<100\delta_2,
       \quad R_c>1.001p,
       \quad N_0\log(1.001)>\log(2p/\delta_2);         \label{bc:pre-gates-5}\\
 &-\log R_c+\ell/n_0<B_c,\quad 1/m_c<2,
       \quad m_c/p<d_s,
       \quad (1-d_s)n_0>\log(2n_0^2);                 \label{bc:pre-gates-6}\\
 &-n_0\log(s_*/p)>\log n_0,\quad
       -n_0\log(s_*/p)>1,\quad(1-d_s)n_0>2;           \label{bc:pre-gates-7}\\
 &\rho^2+2\rho m+q_c<s_*^2,\quad \rho<s_*,\quad
       \rho^{-1}-\rho<-\kappa\log\rho,
       \quad \kappa B_c+n_0^{-2}<m';                 \label{bc:pre-gates-8}\\
 &\frac4{N_0a_0^2}<\frac7{625},\quad
       \frac2{N_0a_0(1-e)}<\frac1{270},\quad
       w>2\kappa,\quad w\log2>3.                     \label{bc:pre-gates-9}
\end{align}
In \eqref{bc:pre-gates-8}, $m=m_c$ at the first stage and is the
preceding $m'$ thereafter; the last inequality is omitted for a dash.
The caps in Table~\ref{bc:refined-table} are checked by substituting
$a=a_0,r=e,n=n_0$ in \eqref{bc:moment-B}--\eqref{bc:moment-Q}, and by
$\kappa B_0+n_0^{-2}<m_0$, $1-B_0>p_f$.
These conditions prove all the displayed estimates for every $n\geq n_0$:
$k,A_-$ increase, the moment errors decrease, and
$n(s_*/p)^n$ and $2n^2d_s^n$ decrease after their stated derivative
bounds.

\subsection{Support and the normal of a boundary edge}

The following transfer is a finite identity between the original and
critical empirical measures.  It retains the normalization difference
between $n$ and $N$.

\begin{lemma}[Moment support bound]\label{bc:poisson-support}
Suppose $P(0)\ne0$, $B\leq B_0$, $Q\leq q<1$, and $|\mu|\leq m$.
Let $0<h<1$ and put
\[
 \Lambda=\frac{(1-q)^2}{1+q+2m}.
\]
If
\begin{equation}\label{bc:support-test}
 (N\Lambda+1)(1-h)^2>2(1-h)+2nB_0,
\end{equation}
then the convex hull of the original zeros contains $\overline{D(0,h)}$
in its interior.
\end{lemma}

\begin{proof}
Rotate a proposed contrary support direction to the real axis, so that
$\Re z_i\leq h$ for all $i$.  Define
\[
 C_o(z)=\prod_i(1-z_i z),\quad C_c(z)=\prod_j(1-\xi_jz),
 \quad f(z)=\frac1n\sum_i\frac1{1-z_i z},\quad
 g(z)=\frac1N\sum_j\frac1{1-\xi_jz}.
\]
The exact identities $C_c=C_o-zC_o'/n$ and $C_c/C_o=f$ give
\begin{equation}\label{bc:poisson-identity}
 N(g-f)=f-1-zf'/f.
\end{equation}
The support assumption, and Gauss--Lucas for the critical points,
make every term well defined at $z=1$.  Also
\[
 |f'(1)|\leq\frac1n\sum_i|1-z_i|^{-2},\qquad
 \Re f(1)\geq(1-h)\frac1n\sum_i|1-z_i|^{-2}.
\]
Consequently $f(1)\ne0$ and $|f'(1)/f(1)|\leq(1-h)^{-1}$.
Set $p_o=2\Re f(1)-1$, $p_c=2\Re g(1)-1$.  Then
\[
 p_o=\frac1n\sum_i\frac{1-|z_i|^2}{|1-z_i|^2}
 \leq\frac{2B_0}{(1-h)^2},\qquad
 Np_c=np_o-1-2\Re(f'(1)/f(1)).
\]
Here we used $1-|z_i|^2\leq-2\log|z_i|$.
For the other measure, Cauchy--Schwarz gives
\[
 p_c\left(\frac1N\sum_j|1-\xi_j|^2\right)
 \geq\left(\frac1N\sum_j\sqrt{1-|\xi_j|^2}\right)^2
 \geq(1-Q)^2.
\]
Its denominator equals $1+Q-2\Re\mu$, whence $p_c\geq\Lambda$.
Combining these estimates contradicts \eqref{bc:support-test}.
Since every direction was allowed, all support values exceed $h$.
\end{proof}

We use the inverse-log quantities of Lemma~\ref{lt:inverse-log}.
Suppose, at the actual rightmost zero $b$,
\begin{equation}\label{bc:field-hypotheses}
 H>H_0>0,\quad E_b/H<V,\quad d^{-1}<M.
\end{equation}
For $R>0$ define
\begin{align}
 \lambda&=\frac{VR}{1-MR},&
 e_R&=V\left(\frac{R^2}2+\frac{MR^3}{3(1-MR)}\right),
 &\widehat r&=R-e_R,                                  \label{bc:inverse-def-1}\\
 B_2&=\frac{V}{(1-MR)^2(1-\lambda)^3},&
 p_-&=(1+\lambda)^{-1},&p_+&=(1-\lambda)^{-1}.
                                                               \label{bc:inverse-def-2}
\end{align}
Choose $0<\nu<1$, $W>0$, and $e_W>e^{-W}$, and put
\[
 \epsilon=1.001B_2/H_0,\qquad g=1.001\nu B_2\widehat r.
\]
The sufficient comparisons are
\begin{align}
 &MR<1,\quad 0<\lambda<1,\quad \widehat r>0,
       \quad W^2<(1-\nu^2)(H_0\widehat r)^2,
                                                    \label{bc:inverse-gates-1}\\
 &\epsilon(1+e_W)+ge_W<1-p_+e_W,
                                                    \label{bc:inverse-gates-2}\\
 &\epsilon(1+e_W)+ge_W^2<p_--e_W,
       \qquad \epsilon+ge_W<p_-.                   \label{bc:inverse-gates-3}
\end{align}
They ensure that the analytic logarithm
$\mathcal L(z)=\log(P'(z)/P'(b))$, fixed at $b$, is injective on
$D(b,R)$ and has an inverse $\phi$ on $D(0,H\widehat r)$.
Moreover
\begin{equation}\label{bc:inverse-location}
 |\phi(w)-b-w/S|<e_R,\qquad
 \frac{SP(\phi(w))}{P'(b)}=\int_0^we^v A(v)\,dv,\quad A=S\phi',
\end{equation}
with $\Re A>p_-$, $|A|<p_+$ and $|A'|<B_2/H$.
On every rectangle with vertical sides $\Re w=\pm W$ and odd
multiples of $\pi$ as horizontal sides, contained in
$|\Im w|\leq\nu H\widehat r$, the complete side estimates and
the homotopy $(1-s)+sA$ give the same zero count as $e^w-1$.
In particular the number of original zeros in $D(b,R)$ is greater than
\begin{equation}\label{bc:inverse-count}
 \nu H\widehat r/\pi-2>
 \nu(7/22)H\widehat r-2,
\end{equation}
provided the last expression is positive.  The loss two includes the
choice of the outermost odd heights.  When
$\nu H_0\widehat r>4\pi$, the top adjacent strip and the bottom
adjacent strip each contain one distinct simple original zero, with
\begin{equation}\label{bc:high-strips}
 |\Re w|<W,\qquad
 \Im w_+>\nu H\widehat r-4\pi,\qquad
 \Im w_-<-\nu H\widehat r+4\pi.
\end{equation}
These statements concern zeros of $P$ itself, through
\eqref{bc:inverse-location}.

\begin{lemma}[Localization of an edge normal]\label{bc:edge-normal}
Suppose $a$ is on a supporting edge of the original hull, the hull
contains a positive centered disk, and $|b-a|\leq e$.  Suppose
\eqref{bc:field-hypotheses}--\eqref{bc:inverse-gates-3} hold,
$\nu H_0\widehat r>4\pi$, and $|\Im S|<J_0$.
Set $s=J_0/H_0<1$ and
\begin{align}
 E_0&=W/H_0+s\nu\widehat r+e_R,\nonumber\\
 L_0&=(1-s^2)(\nu\widehat r-4\pi/H_0)-sW/H_0-e_R.
                                                        \label{bc:normal-errors}
\end{align}
If $L_0>e$, the outward normal $(\cos\theta,\sin\theta)$ of this
edge satisfies
\begin{equation}\label{bc:normal-bound}
 \cos\theta>0,\qquad |\tan\theta|<K:=\frac{E_0}{L_0-e}.
\end{equation}
In particular its distance from the origin is greater than
$a/\sqrt{1+K^2}$.
\end{lemma}

\begin{proof}
Write $\sigma=S/H$.  Since $\Re S>0$, one has
$|\Im\sigma|<s$ and $\Re\sigma>1-s^2$.  The two zeros supplied
by \eqref{bc:high-strips}, and $w/S=\overline\sigma w/H$, satisfy
\[
 \Re z_\pm>x-E_0,\qquad
 \Im z_+>y+L_0,\qquad \Im z_-<y-L_0.
\]
These bounds follow by applying \eqref{bc:inverse-location} on the
strips.
The positive support disk gives $a\cos\theta>0$.  Every original
zero obeys
\[
 \cos\theta(\Re z-a)+\sin\theta\Im z\leq0.
\]
Use $z_+$ when $\sin\theta\geq0$ and $z_-$ otherwise, and use
$x-a\geq0$, $|y|\leq e$.  This gives
$|\sin\theta|(L_0-e)<E_0\cos\theta$, proving the lemma.
\end{proof}

\begin{lemma}[Geometry of the whole edge and radial paths]
\label{bc:edge-geometry}
Suppose the edge through $a$ has distance at least $h>0$ from the
origin and its normal satisfies \eqref{bc:normal-bound}.  Set
\[
 c=\frac{h-K\sqrt{1-h^2}}{\sqrt{1+K^2}}.
\]
Every point $z$ of the edge has $\cos(\arg z)\geq c$ and
\begin{equation}\label{bc:whole-edge-gap}
 |z-\xi_j|>1-K-\sqrt{1-h^2}.
\end{equation}
Furthermore, suppose $a\geq a_1$, $0<\gamma<1$, and $0<r_0<h$.
If
\begin{equation}\label{bc:four-corners}
 1-v^2+2t(vc-\gamma)-t^2(1-\gamma^2)>0
 \quad(v=a_1,1;\ t=r_0,1),
\end{equation}
then the whole radial segment from $r_0u/|u|$ to each endpoint $u$
of the edge satisfies $|z-\xi_j|>\gamma|z|$.
\end{lemma}

\begin{proof}
If $h_e$ is the distance of the line and $\theta$ is its normal
angle, every point of its unit-disk chord satisfies
$|\arg z-\theta|\leq\arccos h_e\leq\arccos h$.
Combining this with $|\tan\theta|<K$ gives the bound $c$.
In normal and tangent coordinates the chord has tangent coordinate
in $[-\sqrt{1-h_e^2},\sqrt{1-h_e^2}]$, while the tangent
coordinate of $a$ has modulus $h_e|\tan\theta|$.  Therefore
$|z-a|\leq\sqrt{1-h^2}+K$, which proves
\eqref{bc:whole-edge-gap} by \eqref{bc:hole}.
For $z=t e^{i\arg u}$, the difference
\[
 (1-\gamma t)^2-|a-z|^2
 \geq1-a^2+2t(ac-\gamma)-t^2(1-\gamma^2)
\]
is bounded below by a polynomial concave separately in $a$ and $t$.
Its minimum on $[a_1,1]\times[r_0,1]$ is therefore at a corner.
Under \eqref{bc:four-corners}, $|a-z|<1-\gamma t$, and
\eqref{bc:hole} gives the claimed radial gap.
\end{proof}

\subsection{A potential estimate with an error proportional to the moment}

The following estimate retains a factor $Q$ in the interpolation error.

\begin{lemma}[Endpoint potential estimate]\label{bc:endpoint-potential}
Let $0<r_0<h$, and let $C_0,v_0,\beta,J,s_0,g,e_0$
be positive parameters.
Suppose an edge through a strict witness $a$ has distance at least $h$
from the origin, and all its points have critical distance greater than
$\gamma>0$.  Suppose, at its two endpoints $u$, the radial paths
from $r_0u/|u|$ to $u$ have critical distance greater than
$\gamma|z|$.  Assume
\[
 Q\leq q<1,\quad |\mu|\leq m<\min(h,r_0),\quad
 |P(0)|>p_f^n,\qquad p_f>.99.
\]
If $n\geq1000$, the comparisons
\eqref{bc:outer-gates}, \eqref{bc:endpoint-log-gates},
\eqref{bc:inner-gates}, and \eqref{bc:hessian-test} hold, and
$12+J/(1-Jq)+1/99<24$, then
\begin{equation}\label{bc:potential-24}
 -\log|P(0)|<24Q.
\end{equation}
\end{lemma}

\begin{proof}
Put $v=Q-|\mu|^2=N^{-1}\sum_j|\xi_j-\mu|^2\geq0$.
The exact centered resolvent identity is
\begin{equation}\label{bc:centered-resolvent}
 \frac1N\sum_j\frac1{z-\xi_j}-\frac1{z-\mu}
 =\frac1N\sum_j
 \frac{(\xi_j-\mu)^2}{(z-\mu)^2(z-\xi_j)}.
\end{equation}
On each outer radial path its modulus is at most $C_0v$ if
\begin{equation}\label{bc:outer-gates}
\begin{gathered}
 \frac1{(r_0-m)^2\gamma r_0}<C_0,\qquad
 1-\frac{m}{r_0(r_0-m)}>v_0,\qquad
 v_0\frac Nn-C_0q>\beta,\\
 \frac{C_0}{\beta^2(h-m)}<J,\qquad Jq<1.
\end{gathered}
\end{equation}
The second comparison bounds from below the inward logarithmic speed
of the centered radius $|z-\mu|$.  Indeed, for $z=t\omega$,
\[
 \Re\frac\omega{t\omega-\mu}
 \geq\frac1t-\frac{m}{t(t-m)}
 \geq1-\frac{m}{r_0(r_0-m)}.
\]
Let $q_P=P'/n$, $C_u=q_P(u)/(u-\mu)^N$, and
$w_u=r_0u/|u|$.  Integrate the logarithmic derivative in
\eqref{bc:centered-resolvent} inward from $u$ a distance $s$.
The centered monomial decays at least as $e^{-Nv_0s}$, and the
logarithm of its ratio to $q_P$ has modulus at most
$NC_0vs$.  Since $e^u-1\leq ue^u$ for $u\geq0$, the
absolute difference integral is at most
\[
 |P'(u)|\int_0^\infty nC_0vs e^{-\beta ns}\,ds
 =\frac{C_0v|P'(u)|}{\beta^2n}.
\]
Consequently
\begin{equation}\label{bc:outer-primitive}
\begin{gathered}
 \int_{w_u}^uP'(z)\,dz
 =C_u\bigl((u-\mu)^n-(w_u-\mu)^n\bigr)+E_{\rm out},\\
 |E_{\rm out}|\leq Jv|C_u(u-\mu)^n|.
\end{gathered}
\end{equation}

For the endpoint logarithm set $\zeta_j=(\xi_j-\mu)/(u-\mu)$.
Their mean is zero.  The parameter comparisons we use are
\begin{equation}\label{bc:endpoint-log-gates}
\begin{gathered}
 \frac{1+m}{h-m}<1.1,\quad
 \log\frac{1+m}{\gamma}<1.2,\quad
 \log\frac2{h-m}<1.2,\quad
 \frac{4(1.2+1.1)}{(h-m)^2}<12.
\end{gathered}
\end{equation}
They imply $|\zeta_j|<1.1$ and
$\gamma/(1+m)<|1-\zeta_j|<2/(h-m)$.
For $|\zeta|\leq1/2$, the complete logarithmic series gives
$|\log|1-\zeta|+\Re\zeta|\leq|\zeta|^2$.
For $|\zeta|>1/2$, the corresponding bound is
$4(1.2+1.1)|\zeta|^2$.  Averaging, and using the zero mean,
gives
\begin{equation}\label{bc:centered-log}
 |L(u)-\log|u-\mu||\leq12v.
\end{equation}

For the inner segment, write $\eta_j=\xi_j-\mu$ and, for fixed $z$,
consider
$f_z(t)=\prod_j(z-\mu-t\eta_j)$.  Its first derivative at $t=0$
vanishes.  For every product missing two labels, AM--GM gives
\[
 \prod_{j\ne i,k}|z-\mu-t\eta_j|
 \leq e\bigl(|z-\mu|^2+t^2v\bigr)^{(N-2)/2},
\]
because $(N/(N-2))^{(N-2)/2}\leq e$.  Also
$\sum_{i\ne k}|\eta_i\eta_k|\leq N^2v$.
Taylor's integral remainder therefore shows that
\[
 |q_P(z)-(z-\mu)^N|\leq2N^2vs_0^{N-2}
 \qquad(|z|\leq r_0),
\]
provided $(r_0+m)^2+q<s_0^2$.  The same argument applied to
$\prod_j(1-t\eta_j/(u-\mu))$ gives
\[
 |C_u|\leq e^{gn},\qquad
 |C_u-1|\leq2N^2ve^{gn}
\]
under $q/[2(h-m)^2]<g$ and $e/[2(h-m)^2]<2$.
If $(r_0+m)e^g<s_0$ and $e^gs_0^2<1$, it follows that
\[
 |q_P(z)-C_u(z-\mu)^N|\leq4N^2vs_0^{N-2}
 \qquad(|z|\leq r_0).
\]
Combining its integral from $0$ to $w_u$ with
\eqref{bc:outer-primitive} and $P(u)=0$ gives
\begin{equation}\label{bc:endpoint-identity}
\begin{gathered}
 0=P(0)+C_u(u-\mu)^n(1+\theta_u)+E_u,\quad
 |\theta_u|\leq Jv,\\
 |E_u|\leq4n^3vs_0^{n-3}+e^{gn}|\mu|^n.
\end{gathered}
\end{equation}
The condition $m<.006$ gives $|\mu|^n\leq Q(.006)^{n-2}$.
The inner conditions are
\begin{align}
 &(r_0+m)^2+q<s_0^2,\quad
   q/[2(h-m)^2]<g,\quad e<4(h-m)^2,\nonumber\\
 &(r_0+m)e^g<s_0,\quad e^gs_0^2<1,\quad
   4/s_0^3<6,\quad s_0/.99<e_0,\quad m<.006,\nonumber\\
 &e^g(.006)/.99<.01,\quad
   6(1000)^3e_0^{1000}+10^6(.01)^{1000}<.01,\quad
   e_0(1001/1000)^3<1.                            \label{bc:inner-gates}
\end{align}
They imply, for every $n\geq1000$,
\begin{equation}\label{bc:Q-error}
 |E_u/P(0)|\leq Q\bigl(6n^3e_0^n+10^6(.01)^n\bigr)\leq Q/100.
\end{equation}
Since $Jq<1$ and $|E_u/P(0)|<1$, we may take logarithms in
\eqref{bc:endpoint-identity}.  Using
$\log|C_u(u-\mu)^n|=NL(u)+\log|u-\mu|$, we obtain
\[
 n\left|L(u)-\frac1n\log|P(0)|\right|
 \leq\left(12+\frac J{1-Jq}+\frac1{99}\right)Q\leq24Q.
\]
The inequalities are strict when $Q>0$.

The assumed inequality \eqref{bc:hessian-test} implies $2h^2-1>0$.
The edge has unit tangent $v_e=ie^{i\theta}$ and support coordinates
$z=e^{i\theta}(h_e+it)$, with $h_e\geq h$ and
$h_e^2+t^2=|z|^2\leq1$.  Hence $h_e^2-t^2\geq2h_e^2-1>0$, and
\[
 -\Re\frac{v_e^2}{z^2}
 =\frac{h_e^2-t^2}{(h_e^2+t^2)^2}
 \geq h_e^2-t^2\geq2h_e^2-1\geq2h^2-1.
\]
The identity
\[
 (z-\xi)^{-2}=z^{-2}+2\xi z^{-3}
       +\frac{\xi^2(3z-2\xi)}{z^3(z-\xi)^2},
 \qquad 3z-2\xi=z+2(z-\xi),
\]
therefore bounds the second derivative of $L$ along the whole edge
from below by
\begin{equation}\label{bc:hessian-test}
 2h^2-1-\frac{2m}{h^3}
       -q\left(\frac1{h^2\gamma^2}+\frac2{h^3\gamma}\right)>0.
\end{equation}
Thus $L$ is convex on that edge.  Its two endpoint estimates imply
$nL(a)\leq\log|P(0)|+24Q$.  Since $L(a)>0$, this is
\eqref{bc:potential-24}.  The argument also covers $v=0$.
\end{proof}

\begin{lemma}[Annular cubic rigidity]\label{bc:cubic-rigidity}
Suppose all original zeros satisfy $\rho\leq|z_i|\leq1$, where
$0<\rho\leq1$, and all critical points belong to
$\{z:|z|\leq1,\ |z-1|\geq1\}$.  Suppose $C\geq1$, $0<q<1$,
$Q\leq q$, and $|P(0)|\geq e^{-CQ}$.  Then $Q=0$, and
$P(z)=z^n+c$ with $|c|=1$, provided
\begin{equation}\label{bc:rigidity-floor}
 n\geq\left\lceil\max\left\{
 \frac{256}{(1-q)^2},\ \frac{8Cq}{1-q},\
 C\left(\frac8\rho+\frac{16}{\rho^2}+\frac6{\rho^3}\right)+11
 \right\}\right\rceil.
\end{equation}
\end{lemma}

\begin{proof}
For $z=x+iy$ and $s=|z|^2$ the identity
\[
 \Re(-z-z^2-z^3/4)-s/4
 =(1+x)(3s/4-x-x^2)
\]
has nonnegative right side on the stated lens.  Indeed
$-1\leq x\leq1/2$, $s\geq x^2$, and $s\geq2x$; use the first
lower bound when $x\leq0$ and the second when $x\geq0$.
Writing $t_k=N^{-1}\sum_j\xi_j^k$, we deduce
\begin{equation}\label{bc:cubic-positive}
 Q\leq4|t_1|+4|t_2|+|t_3|.
\end{equation}

Assume $Q>0$ and put $a_j=P^{(j)}(0)/P(0)$, $j=1,2,3$.
Differentiate the product for $P'$ and apply AM--GM to every remaining
subproduct.  Since $(N/(N-j+1))^{(N-j+1)/2}\leq e^{(j-1)/2}<3$,
\[
 |a_j|\leq3e^{Cq}n^3Q^{(n-3)/2}.
\]
Set $\delta=1-q$.  The first two conditions in
\eqref{bc:rigidity-floor} imply
\[
 3e^{Cq}n^4q^{(n-5)/2}<1.
\]
$\log n\leq\sqrt n$ and $\log q\leq-\delta$ give
\begin{align*}
 \log3+Cq+4\log n+\tfrac{n-5}2\log q
 &\leq\log3+Cq-\delta n/4+5\delta/2\\
 &\leq\log3-\delta n/8+5\delta/2\\
 &\leq\log3-32/\delta+5\delta/2<0.
\end{align*}
It follows that $|a_j|\leq Q/n$.

Let $R_k=\sum_i z_i^{-k}$ and $s_k=n^{-1}\sum_i z_i^k$.  The exact
identities
\[
 R_1=-a_1,\quad R_2=a_1^2-a_2,\quad
 R_3=-(a_3-3a_1a_2+2a_1^3)/2
\]
give $|R_k|\leq2Q/n$ for $k=1,2,3$.  For an original modulus
$s\in[\rho,1]$, integrate the derivative of $e^{kt}-e^{-kt}$ from
$0$ to $t=-\log s$ to obtain
\[
 |z_i^k-\overline{z_i^{-k}}|
 =s^{-k}-s^k\leq2k\rho^{-k}(-\log s).
\]
Sum over all original zeros.  The radial budget
$\sum_i-\log|z_i|\leq CQ$ implies
\begin{equation}\label{bc:power-moment-small}
 |s_k|\leq A_kQ/n,\qquad A_k=2kC\rho^{-k}+1,\quad k=1,2,3.
\end{equation}
Here $n\geq256$ justifies replacing the term $2/n$ by one.

Newton's identities for $P$ and $P'/n$ give
\[
 t_1=s_1,\qquad
 t_2=\frac{(n-2)s_2+s_1^2}{n-1},\qquad
 t_3=\frac{(n-3)s_3+3s_1s_2-s_1^3}{n-1}.
\]
Using $|s_1|\leq1$, \eqref{bc:cubic-positive}, and
\eqref{bc:power-moment-small}, we obtain
\[
 Q\leq\left(B_1+9+\frac{E_1}{n-1}\right)\frac Qn,
\]
where
\[
 B_1=C(8/\rho+16/\rho^2+6/\rho^3),\qquad
 E_1=10C/\rho+12C/\rho^2+8.
\]
But $B_1+10-E_1=C(-2/\rho+4/\rho^2+6/\rho^3)+2>0$.
Since $n\geq B_1+11$, the coefficient in parentheses is less than
$B_1+10<n$, a contradiction.  Thus $Q=0$.
Now $P'=nz^{n-1}$, and the product hypothesis and the unit-disk roots
force the constant term of $P$ to have modulus one.
\end{proof}

\begin{proposition}[Boundary and small displacement ranges]\label{bc:boundary}
For every $n\geq6251$, a hull point satisfying \eqref{bc:hole} is
impossible in each of the following cases:
\begin{equation}\label{bc:boundary-scopes}
\begin{array}{c|c|c}
 &a&r\\ \hline
 \mathrm B&[.975,1]&\text{arbitrary}\\
 \mathrm A&[.85,1]&[0,.02]\\
 \mathrm {G1}&[.96,1]&[0,.04]\\
 \mathrm {G2}&[.965,1]&[0,.045].
\end{array}
\end{equation}
The latter three conclusions hold for any choice of actual rightmost
zero satisfying the indicated displacement bound.
\end{proposition}

\begin{proof}
We first treat a boundary witness $a<1$ which is not an original
zero.  The normalization and moment proof above gives
Table~\ref{bc:refined-table}.  Apply Lemma~\ref{bc:poisson-support}
using the weak support radii in Table~\ref{bc:geometry-table}.
For row A, repeat the support test and the actual-parameter moment
inequality in the following order:
\begin{equation}\label{bc:feedback}
\begin{array}{c|c|c}
 \text{lower }a&\text{new upper }Q&\text{next support radius}\\ \hline
 .9305&.13831&.9454\\
 .9454&.11036&.9478\\
 .9478&.10582&.9481\\
 .9481&.10525&\text{--}.
\end{array}
\end{equation}
Each new moment bound is obtained from \eqref{bc:moment-Q} by
increasing only the lower bound for its explicit $a^2$ term.
The error $\mathcal E(.02;.85,N)$ is retained throughout, and the
caps for $B,|\mu|$ remain unchanged.
Put $a_1=.9481$, $q=.10525$ in row A;
in the other three rows put $a_1=a_0$, $q=q_0$.

At $b$ the finite resolvent expansion
\[
 \frac SN=\frac1b+\frac{\mu}{b^2}
                  +\frac1N\sum_j\frac{\xi_j^2}{b^2(b-\xi_j)}
\]
gives
\begin{equation}\label{bc:moment-field}
 H>Nc_f,\quad E_b/H<V,\quad d^{-1}<M,\qquad
 c_f=1-\frac m{a_1}-\frac q{a_1(1-e)},\quad
 V=\frac1{c_f(1-e)^2},\quad M=\frac1{1-e}.
\end{equation}
The listed parameters give $c_f>0$.  Since
$1-m/|b|-q/(|b|d)>c_f$ and $|b|\leq1$,
\[
 H\geq\frac N{|b|}\left(1-\frac m{|b|}-\frac q{|b|d}\right)>Nc_f.
\]
Use $H_0=N_0c_f$, $\nu=.99$, $W=4$ and the radii $R$ in
Table~\ref{bc:geometry-table}.  All comparisons
\eqref{bc:inverse-gates-1}--\eqref{bc:inverse-gates-3} hold, as does
$\nu H_0\widehat r>4\pi$.  Equation \eqref{bc:fine-phase} supplies
$J_0=8/a_0$, with the original $a_0$ also in row A.
Lemma~\ref{bc:edge-normal} consequently gives $K$.

For B, G1, and G2 the direct comparisons
$a_0^2>h^2(1+K^2)$ show that this particular edge has distance
greater than the listed $h$.  Row A uses its support radius $.9481$.
All four rows satisfy
$1-K-\sqrt{1-h^2}>\gamma$ and the four corner inequalities
\eqref{bc:four-corners} with $r_0=.9$.
Lemma~\ref{bc:edge-geometry} therefore supplies the gaps on the
whole edge and the endpoint radial paths.

The endpoint parameters in Table~\ref{bc:endpoint-table} are positive,
and $r_0=.9<h$ in all four rows.
All the hypotheses of Lemma~\ref{bc:endpoint-potential} hold.
In particular the Hessian bounds and endpoint coefficients in
Table~\ref{bc:boundary-margins} are strictly positive and strictly
less than $24$, respectively.  We obtain
$-\log|P(0)|<24Q$.
The additional annulus radii $\rho_f$ in the table follow from
\[
 \rho_f^2+2\rho_fm+q<s_f^2,\qquad
 \rho_f<s_f,\qquad
 s_f/p_f<s_*/p,
\]
and the already proved whole-degree bound $n(s_*/p)^n<1$.
Thus every original zero satisfies $|z_i|>\rho_f$.

Every critical point belongs to the lens in
Lemma~\ref{bc:cubic-rigidity}, since
\[
 |\xi-1|^2-|\xi-a|^2=(1-a)(1+a-2\Re\xi)>0.
\]
The three bounds in \eqref{bc:rigidity-floor}, with $C=24$,
$\rho=\rho_f$, and the final $q$, are all less than $900$.
Lemma~\ref{bc:cubic-rigidity} gives $Q=0$.  All critical points
then equal zero, contradicting \eqref{bc:hole} for $a<1$.

The weak support test already excludes a hull of dimension less than
two.  A boundary vertex is an original zero and was excluded by
the energy inequality; every other boundary point lies on a genuine
edge as used above.  If $a=1$, strict convexity of $\overline\D$
makes $1$ an original zero, again impossible.
Finally, for an interior witness move outward on its real ray to the
boundary point $A$.  Since $a\leq A\leq x$,
\[
 |A-\xi|^2-|a-\xi|^2=(A-a)(A+a-2\Re\xi)\geq0,\qquad
 |b-A|\leq|b-a|.
\]
The strict hole and each applicable displacement hypothesis persist.
The boundary result therefore proves all of \eqref{bc:boundary-scopes}.
\end{proof}

\begin{table}[ht]
\centering\small
\caption{Support and geometric parameters for the four exclusion rows.}
\label{bc:geometry-table}
\begin{tabular}{@{}crrrrrrr@{}}\toprule
Row&weak $h$&final $h$&$R$&$\gamma$&$\rho_f$&$s_f$&final $q$\\\midrule
B & 0.9239 & 0.95027 & 0.113 & 0.455 & 0.955 & 0.993 & 0.06001\\
A & 0.9305 & 0.9481 & 0.061 & 0.583 & 0.934 & 0.993 & .10525\\
G1 & 0.9351 & 0.94709 & 0.092 & 0.51 & 0.943 & 0.993 & 0.08531\\
G2 & 0.9308 & 0.94828 & 0.1 & 0.49 & 0.947 & 0.993 & 0.07689\\
\bottomrule\end{tabular}
\end{table}
\begin{table}[ht]
\centering\small
\caption{Endpoint parameters, with $r_0=.9$ in every row.}
\label{bc:endpoint-table}
\begin{tabular}{@{}crrrrrrr@{}}\toprule
Row&$C_0$&$v_0$&$\beta$&$J$&$s_0$&$g$&$e_0$\\\midrule
B & 3.056 & 0.9932 & 0.8086 & 4.952 & 0.93797 & 0.03371 & 0.947455\\
A & 2.369 & 0.9973 & 0.7468 & 4.495 & 0.95869 & 0.0589 & 0.968384\\
G1 & 2.716 & 0.9955 & 0.7626 & 4.954 & 0.94956 & 0.048 & 0.959162\\
G2 & 2.83 & 0.9947 & 0.7759 & 4.984 & 0.94572 & 0.04323 & 0.955283\\
\bottomrule\end{tabular}
\end{table}
\begin{table}[ht]
\centering\small
\caption{Outward numerical bounds for the boundary geometry.
The comparisons themselves use rational intervals.}
\label{bc:boundary-margins}
\begin{tabular}{@{}crrr@{}}\toprule
Row&upper bound for $K$&lower Hessian bound&upper endpoint coefficient\\\midrule
B&.229122900&.165248328&19.055896748\\
A&.092635534&.024898554&20.541111328\\
G1&.164970033&.026449841&20.590324175\\
G2&.187982215&.064807027&20.090774931\\\bottomrule
\end{tabular}
\end{table}

The first support comparison and each feedback comparison
check both \eqref{bc:support-test} at $n_0$ and its slope
\[
 \Lambda(1-h)^2-2B_0>0.
\]
The four inverse applications check
\eqref{bc:inverse-gates-1}--\eqref{bc:inverse-gates-3} for both signs
of the strip, as well as $s<1$, $\nu\widehat r-4\pi/H_0>0$,
and $L_0>e$.  The geometric comparisons are $h>.9$,
$1-K-\sqrt{1-h^2}>\gamma$, $c>0$, and the four instances of
\eqref{bc:four-corners}.  The endpoint checks are exactly
\eqref{bc:outer-gates}, \eqref{bc:endpoint-log-gates},
\eqref{bc:inner-gates}, the stricter lower bound $.01$ in
\eqref{bc:hessian-test}, and
$12+J/(1-Jq)+1/99<24$.
The final annulus and the three inequalities in
\eqref{bc:rigidity-floor} are checked as in the proof.
These tests, together with \eqref{bc:pre-gates-1}--\eqref{bc:pre-gates-9},
form the boundary certificate.

As $n$ increases, the fixed moment caps remain valid, while the field
$Nc_f$ increases.  Thus the inverse corner inequality improves and
$\epsilon$ decreases.  In \eqref{bc:normal-errors}, $E_0$ decreases
and $L_0$ increases, so $K$ decreases.  The geometric inequalities
therefore persist.  The outer rate $v_0N/n-C_0q$ increases, and the
inner errors have the explicit decreasing consecutive-degree ratios
in \eqref{bc:inner-gates}.  This proves the boundary conclusions on
the whole infinite degree range, including all stated endpoints.

\subsection{A direct entrance for the compact range}

\begin{lemma}[Crescent entrance]\label{bc:crescent}
Under \eqref{bc:hole}, if $a\geq.1$ and $n\geq6251$, then
\begin{equation}\label{bc:crescent-D}
 D:=Hd>\frac{5N}{14}\frac{c}{1-c+c^2}>117.17,
 \qquad c=x/2.
\end{equation}
Thus Propositions~\ref{mf:D100} and \ref{lt:D70} apply.
\end{lemma}

\begin{proof}
A closest point to $b=x+iy$, $x\geq a$, $y\geq0$, in the closed
crescent $\{|\xi|\leq1,\ |\xi-a|\geq1\}$ is
$a/2+i\sqrt{1-a^2/4}$.  Indeed a minimizer must be on one of its
two circular boundary arcs.  On $|\xi|=1$, maximizing
$\Re(\overline b\xi)$ over $\Re\xi\leq a/2$ gives the stated
endpoint, because the direction of $b$ lies to its right.
On $|\xi-a|=1$, the same argument maximizes
$\Re((\overline b-a)(\xi-a))$ over
$\Re(\xi-a)\leq-a/2$; $b-a$ has nonnegative real and imaginary
parts, giving the same endpoint.  Since the strict hole excludes
that endpoint,
\[
 d^2>1+x^2+y^2-ax-2y\sqrt{1-a^2/4}.
\]
The right side decreases in $a$ for $0<a\leq x\leq1$, since its
derivative is $-x+ay/[2\sqrt{1-a^2/4}]<0$.
Set $c=x/2$, $v=\sqrt{1-c^2}-y\geq0$, and
$G=\sqrt{c^2+v^2}$.  Comparing at $a=x$ gives $d>G$.
Also, if $p=1-|b|^2$, then
\[
 p=-3c^2+2v\sqrt{1-c^2}-v^2\leq2v-3c^2-v^2.
\]
The exact radial identity, obtained by summing
$2\Re(b/(b-\xi))+(1-|b|^2)/|b-\xi|^2
=1+(1-|\xi|^2)/|b-\xi|^2$, is
\begin{equation}\label{bc:radial-identity}
 N=2\Re(bS)+(1-|b|^2)E_b
       -\sum_j\frac{1-|\xi_j|^2}{|b-\xi_j|^2}.
\end{equation}
Furthermore $\Re S>cE_b$, because $a\leq x$.
Using $|b|\leq1-p/2$ in \eqref{bc:radial-identity}, we obtain
\[
 \frac NH\leq2+p(c^{-1}-1),\qquad
 Hd>\frac{NcG}{2c+(1-c)p}.
\]
The second inequality is strict because $d>G$.
By Cauchy--Schwarz and the identity
$(3c^2+v^2)^2-8c^2(c^2+v^2)=(v^2-c^2)^2$,
\begin{align*}
 \frac{2c+(1-c)p}{G}
 &\leq\frac{2c+2(1-c)v}{G}
       -(1-c)\frac{3c^2+v^2}{G}\\
 &\leq2\sqrt{2-2c+c^2}-2\sqrt2\,c(1-c).
\end{align*}
For $.05\leq c\leq1/2$, $\sqrt{2-2c+c^2}<7/5<\sqrt2$.
This proves the first inequality in \eqref{bc:crescent-D}.
The function $c/(1-c+c^2)$ increases on that interval.
Its value at $c=.05$, with $N=6250$, proves the second.
\end{proof}

\begin{proposition}[The bridge interval]\label{bc:bridge}
For every $n\geq6251$, no hull witness satisfies
\eqref{bc:hole} with $.1\leq a\leq.11$.
\end{proposition}

\begin{proof}
Lemma~\ref{bc:crescent} gives $D>100$.  Throughout the argument,
$C$ and $\gamma$ are the quantities in \eqref{bc:adaptive-def}.
The logarithmic form of Lemma~\ref{bc:blaschke} is
\begin{equation}\label{bc:log-bridge}
 -\log\mathcal B_P(\gamma)+CE_b/2<\log(8/C).
\end{equation}

\emph{The range $|b|\leq.2$.}
Here $d>.8$ and Proposition~\ref{mf:D100} gives
$E_bd/H<9/4$.  In particular
$V=E_b/H<9/(4\cdot.8)$ and
$H>N(.8)/(9/4)$, using $E_b\geq N$.
Apply \eqref{bc:inverse-gates-1}--\eqref{bc:inverse-gates-3}
with $R=.1$, $\nu=.99$, $W=6$, $e_W=1/400$,
$M=1/.8$, and $H_0=N_0(.8)/(9/4)$.
The image radius is $\widehat r=379/4480$, and the count
\eqref{bc:inverse-count} is greater than $1831/32$ at the floor.
Since $|\gamma|\leq|b|$, each counted zero has factor at most
\[
 \frac{R+.11}{1-(.2+R)(.2)}<\frac14.
\]
As $e<4$, its logarithmic weight exceeds one.
The count divided by $n$ increases, and exceeds $1/120$ for all
$n\geq6251$.

The actual-zero variance inequality \eqref{ar:variance-unit}
gives $d^2\leq1-(1-|b|^2)/n$, and therefore $d<1-.47/n$.
Since $d+r>1$, we have $r>.4/n$.
Since $a\geq.1$ and $u=\min(r^2,1/2)>(.4/n)^2$,
the monotonicity of $u(1-u)$ on $[0,1/2]$ gives
\[
 C\geq a^2u(1-u)
 >(.1)^2\left(\frac{.4}{n}\right)^2
       \left(1-\left(\frac{.4}{n}\right)^2\right)
 >\frac1{1000n^2}.
\]
Hence
\[
 \log(8/C)<10+2\log n<n/120
       <-\log\mathcal B_P(\gamma),
\]
a contradiction to \eqref{bc:log-bridge}.
For the middle inequality, at $n=6251$ one has
$\log n<9$ and $n/120>28$; its difference increases since
$1/120>2/n$.

\emph{The range $.2<|b|\leq.8$.}
Now $r>.09$, so
$C>C_m=(.1)^2(.09)^2(1-(.09)^2)$.
The field conclusion of Proposition~\ref{lt:D70}, with $D>100$
from Lemma~\ref{bc:crescent}, gives
\[
 H>\frac{N-(1899/20)/(.2)}{2(1+.11)}.
\]
Use $d>.2$, $V=9/(4\cdot.2)$, $M=1/.2$,
$R=.04$, and the same $\nu,W,e_W$.
The image radius is $.0295$, the count at the floor exceeds
$26254139/1184000$, and every counted factor is at most
\[
 \frac{.04+.11}{1-(.8+.04)(.8)}<.46.
\]
The comparison $e^{.75}<1/.46$ gives a product budget exceeding
$16.63>12$, while $e^{12}>8/C_m$.
This again contradicts \eqref{bc:log-bridge}; its count increases
with $N$.

\emph{The range $.8<|b|\leq.94$.}
Here $r>.69$, so $u\geq.4761$, $1/2\leq t\leq.5239$, and
\[
 C>C_u=(.1)^2(.4761)(.5239)=24942879/10^{10}.
\]
Retain only the factor of the original zero $b$ itself.
Its numerator is $ta\leq.057629$, while the real part of its
denominator is
$1-|b|^2+tax\geq1-.94^2+.5(.1)^2=.1214$.
Thus that factor is at most $57629/121400<.475$, and its logarithmic
weight exceeds $.74$.  Since
\[
 .74+6250C_u/2>8.08,\qquad e^{8.08}>8/C_u,
\]
the product factor and the energy term for the same $C$ contradict
\eqref{bc:log-bridge}.

\emph{The range $|b|>.94$.}
Now $r>.83$, hence $u=t=1/2$.
Put $a_*=.03$, $x_* =.1$, and $s=3a_*/4$.
Concavity of $a(x-3a/4)$ on $[a_*,x]$, and $x>3a_*$, give
\[
 C=a(x-3a/4)\geq a_*(x-s).
\]
In fact the value at $a=x$ exceeds the value at $a=a_*$ because
$(x-a_*)(x-3a_*)>0$.
The comparison $e^{8.15}>8/[a_*(x_*-s)]$ and
\eqref{bc:log-bridge} imply
$E_b/N<k/(x-s)$, where $k=2(8.15)/(Na_*)$.
Also $H^2\leq NE_b$ by Cauchy--Schwarz.
The phase conclusion of Proposition~\ref{lt:D70} is
$|\Im S|<(21/4)E_b/H$; the half-plane bound
$E_b/H<2/x$ follows directly from $\Re S>(x/2)E_b$.
Use these estimates in the exact identity \eqref{bc:radial-identity}.
Writing $p_{\max}=1-.94^2$, we find the necessary inequality
\begin{equation}\label{bc:high-necessary}
 1<2x\sqrt{\frac{k}{x-s}}
       +\frac{p_{\max}k}{x-s}
       +\frac{4(21/4)}{Nx}.
\end{equation}
Under $v=\sqrt{x-s}$ the right side is a positive linear combination
of $v+s/v$, $v^{-2}$ and $(v^2+s)^{-1}$.
Each is convex throughout the interval in question: for the last
term this follows from $3v^2>s$.
Its maximum is therefore at $x=x_*$ or $x=1$.
At each of those endpoints, with $N=N_0$, the exact comparisons
\[
 R_x:=1-\frac{p_{\max}k}{x-s}-\frac{21}{Nx}>0,\qquad
 R_x^2>\frac{4x^2k}{x-s}
\]
make the right side of \eqref{bc:high-necessary} less than one.
Every term decreases with $N$, which proves the contradiction in
the whole high-modulus range.  The four cases include their common
endpoints and exhaust every possible $b$.
\end{proof}

\subsection{Nested original-zero counts in the middle range}

\begin{proposition}[The compact range]\label{bc:compact}
For every $n\geq6251$, no hull witness satisfies
\eqref{bc:hole} with $.1\leq a\leq.975$.
\end{proposition}

\begin{proof}
We prove the part $.11\leq a\leq.975$ by the closed cover in
Table~\ref{bc:compact-cover}; Proposition~\ref{bc:bridge} supplies
the remaining interval.  A row specifies $A\leq a\leq B$ and
$r_-(n)\leq r\leq r_+$.
All numbers in the tables are exact rationals.

We first construct a lower bound $H(n)=\alpha(n-1)-\beta$ for $H$,
together with constants $V,M$ as in \eqref{bc:field-hypotheses}.
There are two constructions.

For a trace row put
\[
 M=\frac1{1-r_+},\qquad c=\frac{A}{2(1+r_+)^2}.
\]
Lemma~\ref{bc:crescent} gives $D>100$.
Since $1-|b|^2\leq2d$, \eqref{bc:radial-identity} and
$E_bd/H<9/4$ give $H>2N/13$ and
$D>2N(1-r_+)/13$.
Set
\begin{equation}\label{bc:Dfloor}
 D_f=\frac{2N_0(1-r_+)}{13}-10^{-6}.
\end{equation}
For the six trace rows $D_f\geq201$.
The Fine signed-moment inequality \eqref{lt:fine-curvature}, derived
from Lemma~\ref{lt:signed-moment},
on its domain $E_b/H^2\leq7/625$, $(Hd)^{-1}\leq1/200$,
applies here: $E_b/H^2<(9/4)/201<7/625$.
For its ratio consequence, set
$v_j=[H(b-\xi_j)]^{-1}$, $q_b=E_b/H^2$,
$\mu_b=(Hd)^{-1}$, and $y_b=E_bd/H=q_b/\mu_b$.
Then
\[
 q_b-2\mu_b
 \leq-\Re\sum_jv_j^2
 <(9/4)q_b^2+13q_b\mu_b,
\]
because $0\leq\Re v_j\leq\mu_b$ and
$\sum_j\Re v_j\leq1$.
Consequently
\begin{equation}\label{bc:ratio-quadratic}
 (9/4)y_b^2-(D-13)y_b+2D>0.
\end{equation}
The quadratic is strictly decreasing on $[2,9/4]$ for $D\geq201$.
The values $\kappa$ in Table~\ref{bc:trace-ratios} satisfy
\begin{equation}\label{bc:kappa-gate}
 2<\kappa<9/4,\qquad
 (9/4)\kappa^2-(D_f-13)\kappa+2D_f<0.
\end{equation}
As its value at $\kappa>2$ decreases with $D$, and $y_b<9/4$ is
already known, \eqref{bc:ratio-quadratic} implies $y_b<\kappa$.
We may therefore take
\[
 V=\min(2/A,\kappa/(1-r_+)),\qquad
 H_{\rm init}=\max(N_0A/2,2N_0/13).
\]
The comparisons $V/H_{\rm init}<7/625$ and
$M/H_{\rm init}<1/270$ justify the signed estimate
\[
 \Re\sum_j(b-\xi_j)^{-2}>-(9/4)V^2-13VM.
\]
Each real reciprocal coordinate
$X_j=\Re(b-\xi_j)^{-1}$ belongs to $(c,M)$:
\[
 2(a-\Re\xi)-a|a-\xi|^2
 =a(1-|\xi|^2)+(1-a^2)(a-2\Re\xi)>0,
\]
and $|b-\xi|<(1+r_+)|a-\xi|$.
Using $X_j^2\leq(c+M)X_j-cM$ and $E_b\geq N$ in
$E_b=2\sum_jX_j^2-\Re\sum_j(b-\xi_j)^{-2}$ gives
\begin{equation}\label{bc:trace-field}
 H>\alpha N-\beta,\quad
 \alpha=\frac{1+2cM}{2(c+M)},\qquad
 \beta=\frac{(9/4)V^2+13VM}{2(c+M)}.
\end{equation}

For a moment row take $(\ell,w,\kappa)$ from M0, M1, or M2 as
indicated in Table~\ref{bc:compact-cover}.  Define, at the fixed floor,
\begin{align}
 e_r&=\frac{8r_+}{AN_0}
            +\frac{r_+^2}{2(1-r_+)(1-2r_+)},\nonumber\\
 q_m&=1-A^2+w(\ell/N_0+e_r)+2/n_0^2,\qquad
 m_m=\kappa(\ell/n_0+e_r)+1/n_0^2.                \label{bc:compact-moments}
\end{align}
These are upper bounds for $Q,|\mu|$ throughout the row
and for every larger degree, by Proposition~\ref{bc:moment-proposition}.
Set
\begin{equation}\label{bc:compact-moment-field}
 \alpha=1-\frac{m_m}A-\frac{q_m}{A(1-r_+)},\quad
 \beta=0,\quad V=\frac1{\alpha(1-r_+)^2},\quad
 M=\frac1{1-r_+}.
\end{equation}
The finite comparisons give $\alpha>0$; the resolvent estimate
\eqref{bc:moment-field} gives $H>\alpha N$.

For each row and each indicated integer $j$, use $R_j=j/1000$ and
the parameters $\nu_j,W_j$ specified in
Tables~\ref{bc:nu-table} and \ref{bc:W-table}.
All intermediate integers between the displayed endpoints are included.
Use \eqref{bc:inverse-def-1}--\eqref{bc:inverse-def-2} with this row's
$V,M$, and set $H_0=\alpha N_0-\beta$.
Every one of the comparisons
\eqref{bc:inverse-gates-1}--\eqref{bc:inverse-gates-3} holds, together
with $\nu_jH_0\widehat r_j>22/7$.
Equation \eqref{bc:inverse-count} consequently gives, for the actual
number $\mathcal N_j$ of original zeros in $D(b,R_j)$,
\begin{equation}\label{bc:nested-count}
 \mathcal N_j>k_j(n):=
      \nu_j(7/22)\widehat r_j\bigl(\alpha(n-1)-\beta\bigr)-2>0.
\end{equation}

We estimate their Blaschke factors at $\gamma$ from
\eqref{bc:adaptive-def}.  In every middle row
$r^2\leq r_+^2<1/2$, so $t=1-r^2$.
For a counted zero $z=b+v$, rightmostness gives $\Re v\leq0$.
It follows that
\[
 |z-\gamma|^2=|v+ta|^2\leq R_j^2+B^2.
\]
Substitution of $C=ta[2x-(1+t)a]$ gives the denominator identity
\begin{equation}\label{bc:correlated-denominator}
 \Re(\overline b\gamma)
 =\frac{C(1+r^2)}{2(1-r^2)}
        +r^2+\frac{a^2r^2(1-r^2)}2.
\end{equation}
Also
$\Re(\overline v\gamma)=(x-ta)\Re v+y\Im v\leq r_+R_j$;
here $x-ta\geq0$.  Thus, with
\begin{equation}\label{bc:Delta-def}
 \Delta_j=
 \frac{C_*(1+r_+^2)}{2(1-r_+^2)}
       +(1+B^2/2)r_+^2+R_jr_+,
\end{equation}
we have $\Re(\overline z\gamma)<\Delta_j$.
The row comparisons give $0<\Delta_j<1$ and hence
\begin{equation}\label{bc:factor-bound}
 \left|\frac{\gamma-z}{1-\overline z\gamma}\right|
 \leq\frac{\sqrt{B^2+R_j^2}}{1-\Delta_j}\leq\eta_j<1.
\end{equation}
Here $\eta_j$ is the smallest multiple of $10^{-12}$ whose square is
at least $(B^2+R_j^2)/(1-\Delta_j)^2$.

Put $t_j=(1-\eta_j)/(1+\eta_j)$ and define the rational weights by
\begin{equation}\label{bc:weights}
 w_j=10^{-12}\left\lfloor10^{12}
       \left(2\sum_{k=0}^{23}\frac{t_j^{2k+1}}{2k+1}\right)
       \right\rfloor.
\end{equation}
All omitted terms of the logarithmic series are positive, so
$0<w_j<-\log\eta_j$.
For each row the $w_j$ decrease and the $k_j(n_0)$ increase with $j$.
The physical disks are nested, so finite Abel summation gives
\begin{align}
 -\log\mathcal B_P(\gamma)
 &>w_J\mathcal N_J+
         \sum_{j<J}(w_j-w_{j+1})\mathcal N_j\nonumber\\
 &>w_Jk_J(n)+
         \sum_{j<J}(w_j-w_{j+1})k_j(n)
   =\sum_jw_j\bigl(k_j(n)-k_{j-1}(n)\bigr)
   =:\mathcal A(n),                               \label{bc:Abel}
\end{align}
where $k_0=0$ and the indices are understood in their successive
order within a row.  The first inequality assigns to each counted
zero the largest available weight.  The second uses nonnegative
coefficients and lower bounds on nested disk counts.

Set
\[
 C_-(n)=A^2r_-(n)^2(1-r_-(n)^2).
\]
Because $r^2<1/2$, $C\geq C_-(n)$.
For a fixed radial lower endpoint retain the energy bonus
$(n-1)C_-/2$ in \eqref{bc:log-bridge}; for a shrinking lower
endpoint $r_-(n)=c/n$ discard that positive term.
The targets $T_0$ in Table~\ref{bc:compact-margins} satisfy
\begin{equation}\label{bc:compact-final-gates}
 e^{T_0}>\frac8{C_-(n_0)},\qquad
 \mathcal A(n_0)+\text{bonus}>T_0.
\end{equation}
Put $\lambda_j=\nu_j(7/22)\widehat r_j$ and
\[
 W_*=\sum_jw_j(\lambda_j-\lambda_{j-1}),\qquad\lambda_0=0.
\]
The exact telescoping identity is
\[
 \mathcal A(n)=W_*\bigl(\alpha(n-1)-\beta\bigr)-2w_{\rm first}.
\]
The comparisons verify that $\mathcal A(n)+\text{bonus}=A_*n-B_*$
with $A_*,B_*>0$.
For fixed $r_-$ the left side increases, while $\log(8/C_-)$ is
constant, so \eqref{bc:compact-final-gates} contradicts
\eqref{bc:log-bridge} for every $n\geq n_0$.
For $r_-(n)=c/n$,
$C_-(n)\geq C_-(n_0)(n_0/n)^2$, while $\mathcal A(n)/n$
increases.  As $T_0>2$ in these four rows,
\[
 \log\frac8{C_-(n)}
 <T_0+2\log(n/n_0)
 \leq T_0(n/n_0)<\mathcal A(n),
\]
again a contradiction on the entire degree range.
The field entrances and every inverse contour condition also improve
with $n$, since $H(n)$ increases and the other contour constants are
fixed.

It remains to verify that the closed cells just excluded leave no
gaps.  For each interval in $a$, its final far radius $r_f$ is listed
below:
\begin{equation}\label{bc:far-table}
\begin{array}{c|rrrrrrrrrr}
 A&.11&.2&.45&.75&.85&.9&.93&.95&.96&.965\\ \hline
 r_f&.65&.27&.115&.069&.06&.057&.055&.054&.053&.053.
\end{array}
\end{equation}
Each satisfies $r_f^2<1/2$ and
$A^2r_f^2(1-r_f^2)>C_*$.  Lemma~\ref{bc:blaschke} therefore
excludes every $r\geq r_f$, including $r^2\geq1/2$.
For the four shrinking lower endpoints, let
$\overline B=B+.001$.  If $r\leq c/n$, then $|b|<\overline B$,
and \eqref{ar:variance-unit} gives
\[
 d^2<1-(1-\overline B^2)/n,\qquad
 d<1-\frac{1-\overline B^2}{2n}<1-c/n.
\]
Here the four exact comparisons are $B+c/n_0<\overline B$ and
$c<(1-\overline B^2)/2$.
This contradicts $d+r>1$.
For $a\geq.85$ the interval $r\leq.02$ is excluded by row A of
Proposition~\ref{bc:boundary}.  For $a\geq.96,r\leq.04$, and
$a\geq.965,r\leq.045$, use rows G1 and G2, respectively.
All witness and radial endpoints are shared by closed exclusions.
Together with Proposition~\ref{bc:bridge}, this proves
Proposition~\ref{bc:compact}.
\end{proof}

\begin{table}[ht]
\centering\small
\caption{The thirteen closed middle cells.  Every integer in a displayed
$j$ interval is used, with $R_j=j/1000$.}
\label{bc:compact-cover}
\begin{tabular}{@{}crrrrrrl@{}}\toprule
Row&$A$&$B$&$r_-(n)$&$r_+$&$j$ interval&number&field\\\midrule
Ia & 0.11 & 0.2 & $0.45/n$ & 0.1 & 5--266 & 262 & Trace\\
Ib & 0.11 & 0.2 & 0.1 & 0.65 & 6--91 & 86 & Trace\\
II & 0.2 & 0.45 & $0.35/n$ & 0.27 & 5--212 & 208 & Trace\\
III & 0.45 & 0.75 & $0.2/n$ & 0.115 & 5--262 & 258 & Trace\\
IV & 0.75 & 0.85 & $0.12/n$ & 0.069 & 5--278 & 274 & Trace\\
V & 0.85 & 0.9 & 0.02 & 0.06 & 5--281 & 277 & Trace\\
VI & 0.9 & 0.93 & 0.02 & 0.057 & 5--299 & 295 & M0\\
VII & 0.93 & 0.95 & 0.02 & 0.055 & 5--244 & 240 & M1\\
VIIIa1 & 0.95 & 0.96 & 0.02 & 0.025 & 5--248 & 244 & M2\\
VIIIa2 & 0.95 & 0.96 & 0.025 & 0.035 & 5--236 & 232 & M2\\
VIIIb & 0.95 & 0.96 & 0.035 & 0.054 & 5--211 & 207 & M2\\
IXa & 0.96 & 0.965 & 0.04 & 0.053 & 5--193 & 189 & M2\\
IXb & 0.965 & 0.975 & 0.045 & 0.053 & 5--150 & 146 & M2\\
\bottomrule\end{tabular}
\end{table}
\begin{table}[ht]
\centering\small
\caption{Trace-row ratio parameters.  In each row $D_f$ is given by
\eqref{bc:Dfloor}, so that the same lower bound is used in
\eqref{bc:kappa-gate}.}
\label{bc:trace-ratios}
\begin{tabular}{@{}crrrrrr@{}}\toprule
Row&Ia&Ib&II&III&IV&V\\\midrule
$\kappa$&2.04150406 & 2.11136288 & 2.0514852 & 2.04222636 & 2.04008696 & 2.0396935\\
\bottomrule\end{tabular}
\end{table}
\begin{table}[ht]
\centering\small
\caption{Targets and rational lower bounds for the final compact margins.}
\label{bc:compact-margins}
\begin{tabular}{@{}crr@{}}\toprule
Row&$T_0$&$\mathcal A(n_0)+\mathrm{bonus}-T_0>$\\\midrule
Ia & 25.574 & 193.4195\\
Ib & 11.111 & 8.6177\\
II & 24.88 & 44.1758\\
III & 24.378 & 16.8682\\
IV & 24.378 & 3.1044\\
V & 10.23 & 7.9635\\
VI & 10.116 & 5.4864\\
VII & 10.051 & .5545\\
VIIIa1 & 10.008 & .4089\\
VIIIa2 & 9.562 & .7142\\
VIIIb & 8.89 & 1.4735\\
IXa & 8.602 & 1.6317\\
IXb & 8.356 & .4620\\
\bottomrule\end{tabular}
\end{table}

The next two tables specify every contour choice by intervals of
integers.  The notation $[u,v]:k$ assigns the value $k$ to every
$j=u,u+1,\ldots,v$; a single integer denotes a singleton interval.
Assignments in a row are disjoint and cover exactly its $j$ interval
in Table~\ref{bc:compact-cover}.

\begin{longtable}{@{}lp{.84\textwidth}@{}}
\caption{The exact values of $10^4\nu_j$.}\label{bc:nu-table}\\
\toprule Row&Assignments\\\midrule\endfirsthead
\toprule Row&Assignments\\\midrule\endhead
Ia & $5:9975,\quad 6:9982,\quad 7:9987,\quad 8:9989,\quad 9:9991$\\
 & $10:9993,\quad [11,12]:9994,\quad [13,14]:9995,\quad [15,17]:9996,\quad [18,24]:9997$\\
 & $[25,257]:9998,\quad [258,264]:9997,\quad [265,266]:9996$\\
Ib & $6:9891,\quad 7:9919,\quad 8:9937,\quad 9:9950,\quad 10:9959$\\
 & $11:9965,\quad 12:9971,\quad 13:9975,\quad 14:9978,\quad 15:9980$\\
 & $16:9982,\quad 17:9984,\quad 18:9986,\quad 19:9987,\quad 20:9988$\\
 & $21:9989,\quad [22,23]:9990,\quad [24,25]:9991,\quad [26,27]:9992,\quad [28,30]:9993$\\
 & $[31,34]:9994,\quad [35,40]:9995,\quad [41,48]:9996,\quad [49,52]:9990,\quad [53,57]:9991$\\
 & $[58,64]:9992,\quad [65,68]:9993,\quad [69,70]:9986,\quad [71,75]:9987,\quad [76,78]:9988$\\
 & $[79,81]:9980,\quad [82,84]:9981,\quad 85:9971,\quad [86,88]:9972,\quad [89,90]:9961$\\
 & $91:9932$\\
II & $5:9967,\quad 6:9976,\quad 7:9982,\quad 8:9986,\quad 9:9989$\\
 & $10:9990,\quad 11:9992,\quad 12:9993,\quad [13,14]:9994,\quad [15,16]:9995$\\
 & $[17,20]:9996,\quad [21,29]:9997,\quad [30,188]:9998,\quad [189,206]:9997,\quad [207,211]:9996$\\
 & $212:9994$\\
III & $5:9980,\quad 6:9986,\quad 7:9989,\quad 8:9991,\quad 9:9993$\\
 & $10:9994,\quad [11,12]:9995,\quad [13,15]:9996,\quad [16,21]:9997,\quad [22,253]:9998$\\
 & $[254,262]:9997$\\
IV & $5:9984,\quad 6:9989,\quad 7:9991,\quad 8:9993,\quad 9:9994$\\
 & $10:9995,\quad [11,13]:9996,\quad [14,19]:9997,\quad [20,276]:9998,\quad [277,278]:9997$\\
V & $5:9985,\quad 6:9989,\quad 7:9992,\quad 8:9993,\quad 9:9994$\\
 & $10:9995,\quad [11,13]:9996,\quad [14,18]:9997,\quad [19,279]:9998,\quad [280,281]:9997$\\
VI & $5:9989,\quad 6:9992,\quad 7:9994,\quad 8:9995,\quad [9,10]:9996$\\
 & $[11,15]:9997,\quad [16,299]:9998$\\
VII & $5:9991,\quad 6:9993,\quad 7:9995,\quad [8,9]:9996,\quad [10,13]:9997$\\
 & $[14,244]:9998$\\
VIIIa1 & $5:9992,\quad 6:9994,\quad 7:9995,\quad [8,9]:9996,\quad [10,12]:9997$\\
 & $[13,248]:9998$\\
VIIIa2 & $5:9992,\quad 6:9994,\quad 7:9995,\quad [8,9]:9996,\quad [10,12]:9997$\\
 & $[13,236]:9998$\\
VIIIb & $5:9992,\quad 6:9994,\quad 7:9995,\quad [8,9]:9996,\quad [10,13]:9997$\\
 & $[14,211]:9998$\\
IXa & $5:9992,\quad 6:9994,\quad 7:9995,\quad 8:9996,\quad [9,12]:9997$\\
 & $[13,193]:9998$\\
IXb & $5:9992,\quad 6:9994,\quad 7:9995,\quad 8:9996,\quad [9,12]:9997$\\
 & $[13,150]:9998$\\
\bottomrule\end{longtable}

\begin{longtable}{@{}lp{.84\textwidth}@{}}
\caption{The exact values of $W_j$.}\label{bc:W-table}\\
\toprule Row&Assignments\\\midrule\endfirsthead
\toprule Row&Assignments\\\midrule\endhead
Ia & $[5,129]:1,\quad [130,184]:2,\quad [185,212]:3,\quad [213,233]:4,\quad [234,248]:5$\\
 & $[249,257]:6,\quad [258,264]:8,\quad [265,266]:10$\\
Ib & $[6,48]:1,\quad [49,68]:2,\quad [69,78]:3,\quad [79,84]:4,\quad [85,88]:5$\\
 & $[89,90]:6,\quad 91:8$\\
II & $[5,104]:1,\quad [105,149]:2,\quad [150,171]:3,\quad [172,188]:4,\quad [189,199]:5$\\
 & $[200,206]:6,\quad [207,211]:8,\quad 212:10$\\
III & $[5,127]:1,\quad [128,181]:2,\quad [182,209]:3,\quad [210,229]:4,\quad [230,244]:5$\\
 & $[245,253]:6,\quad [254,261]:8,\quad 262:10$\\
IV & $[5,134]:1,\quad [135,191]:2,\quad [192,220]:3,\quad [221,242]:4,\quad [243,257]:5$\\
 & $[258,268]:6,\quad [269,276]:8,\quad [277,278]:10$\\
V & $[5,135]:1,\quad [136,193]:2,\quad [194,222]:3,\quad [223,244]:4,\quad [245,260]:5$\\
 & $[261,270]:6,\quad [271,279]:8,\quad [280,281]:10$\\
VI & $[5,181]:1,\quad [182,253]:2,\quad [254,289]:3,\quad [290,299]:4$\\
VII & $[5,195]:1,\quad [196,244]:2$\\
VIIIa1 & $[5,217]:1,\quad [218,248]:2$\\
VIIIa2 & $[5,213]:1,\quad [214,236]:2$\\
VIIIb & $[5,204]:1,\quad [205,211]:2$\\
IXa & $[5,193]:1$\\
IXb & $[5,150]:1$\\
\bottomrule\end{longtable}

\subsection{Finite arithmetic verification}
\label{bc:arithmetic}

For the boundary calculations, all rational expressions in
Tables~\ref{bc:preliminary-table}--\ref{bc:endpoint-table} and
\eqref{bc:feedback} are evaluated with interval endpoints in $\mathbb Q$.
The following rules make the arithmetic finite and reproducible.
For $0\leq x\leq4$, put $S_{96}(x)=\sum_{j=0}^{96}x^j/j!$.
Then
\[
 S_{96}(x)<e^x<
 S_{96}(x)+\frac{x^{97}}{97!}\frac1{1-x/98}
 \quad(x>0);
\]
negative arguments are handled by reciprocal intervals.
For logarithms, first multiply the argument by a power of two to
place it in $[1,2)$.  If $t=(x-1)/(x+1)$ in that interval, then
\[
 2\sum_{j=0}^{95}\frac{t^{2j+1}}{2j+1}
 \leq\log x\leq
 2\sum_{j=0}^{95}\frac{t^{2j+1}}{2j+1}
 +\frac{2t^{193}}{193(1-t^2)}.
\]
Apply the same rule to $\log2$ when undoing the power-of-two change.
Square roots are enclosed between consecutive multiples of $10^{-56}$,
using integer square roots after multiplying the radicand by $10^{112}$.
For $\pi$, use
\[
 \pi=16\arctan(1/5)-4\arctan(1/239)
\]
with 100 alternating terms for each arctangent and the first omitted
term as its remainder bound.  These rules suffice to verify every
strict comparison in the boundary certificate.  Integer powers, even
those with exponent 1000, are evaluated exactly.

For the middle contours, only rational arithmetic, integer square
roots, and positive exponential and logarithmic sums are needed.
The choices for $e_W$ in \eqref{bc:inverse-gates-2}--\eqref{bc:inverse-gates-3}
are as follows:
\begin{center}\small
\begin{tabular}{@{}rrrr@{}}\toprule
$W$&$10^{12}e_W$&$W$&$10^{12}e_W$\\\midrule
1&367879441172&5&6737947000\\
2&135335283237&6&2478752177\\
3&49787068368&8&335462628\\
4&18315638889&10&45399930\\\bottomrule
\end{tabular}
\end{center}
The positive sums $S_{100}(W)=\sum_{j=0}^{100}W^j/j!$ satisfy
$S_{100}(W)>1/e_W$, and hence $e_W>e^{-W}$.
Likewise $S_{100}(T_0)>8/C_-(n_0)$ certifies the first comparison
in \eqref{bc:compact-final-gates}; all other terms in the final
budget are rational.

An explicit verification consists of the following finite steps.
For each of the thirteen rows, form $V,M,\alpha,\beta$ from
\eqref{bc:trace-field} or \eqref{bc:compact-moment-field} and verify
all its preceding field conditions.  For each integer $j$ in its
listed interval, read $\nu_j,W_j$ from the two assignment tables,
form \eqref{bc:inverse-def-1}--\eqref{bc:inverse-def-2}, and verify
\eqref{bc:inverse-gates-1}--\eqref{bc:inverse-gates-3},
$\nu_jH_0\widehat r_j>22/7$, and $k_j(n_0)>0$.
Next form \eqref{bc:Delta-def} and the outward radical
\eqref{bc:factor-bound}; verify $0<\Delta_j<1$, $0<\eta_j<1$,
and
\[
 (\eta_j-10^{-12})^2<
 \frac{B^2+R_j^2}{(1-\Delta_j)^2}\leq\eta_j^2.
\]
Compute the weight by \eqref{bc:weights}, and check its positivity,
its nonincrease with $j$, and the strict increase of $k_j(n_0)$.
Finally verify the two equivalent forms of the Abel sum in
\eqref{bc:Abel}, the positive slope and intercept loss $A_*,B_*$,
and the strict margins in Table~\ref{bc:compact-margins}.
There are exactly 2918 such contours, all specified by the printed
tables.

The two bridge contours use $(\nu,W,e_W)=(.99,6,1/400)$ with the
parameters given in Proposition~\ref{bc:bridge}.  Their remaining
checks are \eqref{bc:inverse-gates-1}--\eqref{bc:inverse-gates-3},
the displayed factor and count bounds, and the scalar comparisons
in its four cases.  Upper bounds for $e^x$ at
$x=1,.75,.74$ follow from
\[
 e^x<\sum_{j=0}^{24}\frac{x^j}{j!}
       +\frac{x^{25}}{25!}\frac1{1-x/26};
\]
all the other bridge exponential comparisons use $S_{100}$ from below.
The high-modulus radical tests are the two positive-side squared
comparisons following \eqref{bc:high-necessary}.
Thus the entire compact argument reduces to the displayed finite
integer and rational operations, followed by the explicit monotonicity
arguments for $n\geq6251$.

\section{The origin region}\label{sec:origin}

In this section we exclude a rightmost zero close to the upper endpoint of
the unit disk.  The estimates apply down to arbitrarily small positive
scales.

\begin{proposition}[Origin region]\label{or:theorem}
Let $P$ have degree $n\ge6251$, with all its zeros in $\overline\D$.
Suppose that $a>0$ and
\[
             |a-\xi|>1\qquad(P'(\xi)=0).
\]
Let $b=x+iy$ be a rightmost zero of $P$, where $x\ge a$ and $y\ge0$.
Then
\begin{equation}\label{or:scale-conclusion}
       \max(x,1-y)>\rho_*:=\frac{30501}{500000}.
\end{equation}
All zeros and critical points in this statement, and in the sums below,
are counted with their algebraic multiplicities.
\end{proposition}

\subsection{Normalization and radial identities}

Suppose, to the contrary, that
\[
 \rho=\max(x,1-y)\le\rho_*,\qquad
 F(u)=P(i+\rho u),\qquad
 \beta=X-iT=\frac{b-i}{\rho},\qquad \alpha=\frac a\rho.
\]
Then
\begin{equation}\label{or:normalization}
 0<\alpha\le X\le1,\qquad 0<T\le1,\qquad \max(X,T)=1.
\end{equation}
Here $T>0$: a point $x+iy\in\overline\D$ with $x>0$ cannot have
$y=1$.  Every zero $u$ of $F$ satisfies
\begin{equation}\label{or:original-domain}
                  \Re u\le X,\qquad |i+\rho u|\le1.
\end{equation}
Gauss--Lucas and the strict hole give, for each critical point $v=r-it$,
\begin{equation}\label{or:critical-domain}
 \frac\rho2|v|^2\le t<\frac\rho2|v-\alpha|^2,
                  \qquad r<\frac\alpha2.
\end{equation}
The root $\beta$ is simple, since its real part is strictly larger than
that of every critical point.  Put
\begin{equation}\label{or:notation}
 \begin{gathered}
 N=n-1,\quad N_0=6250,\quad n_0=6251,\quad c_\rho=\rho/2,\quad
 I(u)=-\Im u-\frac\rho2|u|^2,\\
 S=\frac{F''(\beta)}{F'(\beta)}=\sum_v\frac1{\beta-v},\quad
 h=|S|,\quad Q=\sum_v\frac1{|\beta-v|^2},\quad
 d=\min_v|\beta-v|,\\
 y_0=\frac Qh,\quad q=\frac Q{h^2},\quad
 \mu=\frac1{dh},\quad D=dh,\quad \tau=\frac{N\rho}{2}.
 \end{gathered}
\end{equation}
We use $y_0$ for the moment ratio, to distinguish it from the physical
coordinate $y$ of $b$.  The original and critical insets $I$ are
nonnegative.  Moreover,
\begin{equation}\label{or:horizontal}
 \Re S>0,\quad h>0,\quad
 Q<\frac{2\Re S}{X},\quad h^2\le NQ,\quad \frac hn<y_0.
\end{equation}
Indeed, $X-r>X/2$ in every summand of $\Re S$; the penultimate
inequality is Cauchy--Schwarz.  Also, by
\eqref{or:critical-domain},
\begin{equation}\label{or:inset-critical}
 I(v)<\rho\alpha(\alpha/2-r)\le\rho X(X/2-r).
\end{equation}
For the last comparison subtract the two quadratic expressions: their
difference is $(X-\alpha)((X+\alpha)/2-r)\ge0$.

For $b=i+\rho\beta$ and any $z\ne\beta$,
\begin{equation}\label{or:radial-one}
 \Re\frac b{\beta-z}
       =\frac\rho2+\frac{I(z)-I(\beta)}{|\beta-z|^2}.
\end{equation}
Summing over the critical points, and writing
$E_c=\sum_v I(v)/|\beta-v|^2$, gives
\begin{equation}\label{or:radial}
 \tau=\rho X\Re S-(1-\rho T)\Im S-E_c+I(\beta)Q.
\end{equation}
Since $E_c\ge0$, $|b|\le1$, and $I(\beta)<T$, it follows that
\begin{equation}\label{or:radial-unphased}
                         \tau<h(1+Ty_0).
\end{equation}
On the other hand \eqref{or:inset-critical} gives
\[
            E_c<\rho(X\Re S-X^2Q/2).
\]
Substitution into \eqref{or:radial}, retaining the term $-E_c$, yields
\begin{equation}\label{or:radial-sharp}
 T(1-\rho T/2)Q<\tau+(1-\rho T)\Im S.
\end{equation}
Thus, if $|\Im S|<J$,
\begin{align}
 \tau&<\rho h+J+Ty_0h,\label{or:radial-phase}\\
 \left(\frac hn\right)^2
 &<\left(\frac hN\right)^2
 <\frac{\rho/2+J/N}{T(1-\rho T/2)}.\label{or:radial-ratio}
\end{align}

At zero inset we have the following bound.
\begin{lemma}[A boundary source]\label{or:zero-inset}
If $X=1$ and $I(\beta)=0$, then $h>N/2$.
\end{lemma}
\begin{proof}
Now $|b|=1$.  At a simple zero the original-root factorization gives
$S=2\sum_{u\ne\beta}(\beta-u)^{-1}$.  If $z=i+\rho u$ is another
physical zero, then
\[
 2\Re\frac b{b-z}-1=\frac{1-|z|^2}{|b-z|^2}\ge0.
\]
Consequently
\[
 \Re(bS)=2\rho\sum_{z\ne b}\Re\frac b{b-z}\ge N\rho.
\]
The critical-root version of \eqref{or:radial-one} says
$\Re(bS)=N\rho/2+E_c$.  Hence
\[
 N\rho/2\le E_c<\rho(\Re S-Q/2),
 \qquad h\ge\Re S>N/2+Q/2>N/2.
\]
\end{proof}

\subsection{A finite Cayley--Hermite--Biehler inequality}

The identities \eqref{or:horizontal}--\eqref{or:radial-ratio} and the
next two source lemmas hold for every finite degree $n\ge2$,
$0<\rho<1$, $0<\alpha\le X\le1$, and $0<T\le1$ satisfying
\eqref{or:original-domain}--\eqref{or:critical-domain} at an actual
simple rightmost zero.  They also apply at the larger scales in the
next section.

\begin{lemma}[Finite polar source estimate]\label{or:HB}
For real $s$ suppose that
\begin{equation}\label{or:Cayley-entrance}
             \Re\frac{s}{1+ic_\rho s}>X.
\end{equation}
Define
\begin{equation}\label{or:source-definitions}
 \begin{gathered}
 W(s)=\sum_{F(u)=0}\frac{I(u)}{|s-(1+ic_\rho s)u|^2},\qquad
 C(s)=\sum_{F'(v)=0}\frac{I(v)}{|s-(1+ic_\rho s)v|^2},\\
 E(s)=\sum_{F(u)=0}\frac{I(u)}{|s-u|^2},\qquad
 \theta=\frac{n\rho}{2(1+c_\rho^2s^2)}.
 \end{gathered}
\end{equation}
Then
\begin{equation}\label{or:HB-inequality}
         W\ge E,\qquad \theta W\le(C+\rho)(W+\theta).
\end{equation}
In particular, if an upper bound $k$ for $C+\rho$ satisfies $k<\theta$,
then $W\le k/(1-k/\theta)$.
\end{lemma}
\begin{proof}
The maps
\[
 w=\frac u{1-ic_\rho u},\qquad
 u=\frac w{1+ic_\rho w},\qquad d_p(w)=1+ic_\rho w
\]
send the disk in \eqref{or:original-domain} to the closed lower
half-plane, because
\[
       -\Im w=\frac{I(u)}{|1-ic_\rho u|^2}.
\]
Consider the finite polynomials
\begin{equation}\label{or:polar-polynomials}
 G(w)=d_p(w)^nF(w/d_p(w)),\qquad
 J_p(w)=d_p(w)G'(w)-nic_\rho G(w)
       =d_p(w)^{n-1}F'(w/d_p(w)).
\end{equation}
A root $u=-i/c_\rho$ makes its
transformed linear factor constant.  It has $I(u)=0$ and contributes
zero to the sums in \eqref{or:source-definitions}.  All other labels
and multiplicities are preserved.  The coefficient $n$ in
\eqref{or:polar-polynomials} is the original degree, even if $G$ has
smaller degree.  Both $G$ and $J_p$ have their finite roots in the
closed lower half-plane.

By \eqref{or:Cayley-entrance}, $G(s)\ne0$.  Put
\[
 g=G'/G,\qquad a(w)=-nic_\rho/d_p(w),\qquad
 H_s(w)=G'(w)+a(s)G(w).
\]
Direct subtraction in the Cayley map proves
\[
 -\Im g(s)=W,\quad -\Im(J_p'/J_p)(s)=C,\quad
 -\Im a(s)=\theta,
 \qquad |g(s)+a(s)|\ge W+\theta>0.
\]
In particular, $H_s(s)$ and $J_p(s)$ are nonzero.  Differentiating
the polar identity before freezing $a$ gives
\begin{equation}\label{or:frozen-polar}
 \frac{H_s'}{H_s}(s)
  =\frac{J_p'}{J_p}(s)-\frac{ic_\rho}{d_p(s)}
                      -\frac{a'(s)}{g(s)+a(s)}.
\end{equation}
Here $|a'(s)|=c_\rho\theta$ and $|ic_\rho/d_p(s)|\le c_\rho$.
The last two terms therefore each have modulus at most $c_\rho$,
and
\begin{equation}\label{or:frozen-upper}
                   -\Im(H_s'/H_s)(s)\le C+\rho.
\end{equation}

The real-rootedness needed below follows from a finite argument.
If a degree $m$ polynomial has all roots strictly below
the real axis, a continuous argument of its value on the real axis
decreases strictly through a total angle $m\pi$: its derivative is
the negative sum of the positive Poisson kernels of the roots.
After any fixed rotation, each crossing of an integer multiple of
$\pi$ is a real zero of its imaginary part.  These crossings account
for the degree of that imaginary part, which is $m$, or $m-1$ when
its leading coefficient vanishes.  Hence that imaginary part is
real-rooted.  Moving roots a small distance downward and taking a
coefficient limit proves the assertion for roots on the boundary as
well, allowing repeated roots and degree drop.  A nonzero coefficient
limit of real-rooted polynomials is real-rooted, with roots allowed to
escape to infinity.

Rotate $G$ so that $H_s(s)>0$ is real, and write $G=U+iV$.  Let
$p=\Re a(s)$ and
\[
                       B=V'+pV-\theta U=\Im H_s.
\]
At the selected point $s$,
\[
 B(s)=0,\qquad H_s(s)V(s)=(W+\theta)|G(s)|^2>0,\qquad
 U'V-UV'=W|G|^2.
\]
In particular $V$ is not identically zero.  By the preceding argument
it is real-rooted, so
\[
 V'^2-VV''\ge0.
\]
This last inequality follows directly by logarithmic differentiation
away from the zeros, and by continuity at them.  The exact Wronskian
identity
\[
 BV'-B'V=\theta(U'V-UV')+V'^2-VV''
\]
now gives
\[
 \theta W\le[-\Im(H_s'/H_s)(s)](W+\theta)
             \le(C+\rho)(W+\theta).
\]
Finally, for $u=r-it$,
\begin{equation}\label{or:denominator-expansion}
 |s-(1+ic_\rho s)u|^2
   =|s-u|^2-\rho s^2t+\frac{\rho^2s^2}{4}|u|^2
   \le |s-u|^2,
\end{equation}
since $t\ge\rho|u|^2/2$.  Termwise comparison gives $W\ge E$.
Rearranging with $k<\theta$ gives the solved form.
\end{proof}

We will also use the unsolved form.  If $W\ge b_0>0$, then
\begin{equation}\label{or:parallel-source}
                   C+\rho\ge\frac{\theta b_0}{\theta+b_0}.
\end{equation}
The function $uv/(u+v)$ is increasing in each positive argument.  Thus
\eqref{or:parallel-source} remains useful without an upper bound on
$h$ or a positive solved denominator.

\subsection{The critical-label comparison}

For $s\ge\alpha$, the critical domain yields the exact identity
\begin{align}
 &|s-(1+ic_\rho s)v|^2
       -(1-\rho^2s^2/2)|s-v|^2\notag\\
 &\quad=\rho s^2\left(\frac\rho2|v-\alpha|^2-t\right)
 +\frac{\rho^2s^2}{2}(s-\alpha)(s+\alpha-2r)
 +\frac{\rho^2s^2}{4}|v|^2>0.\label{or:quadratic-polar}
\end{align}
All three terms have their stated sign on the whole unbounded critical
domain.  In particular, if $\rho\le r_+$ and $s\le s_+$, with
$r_+^2s_+^2/2<1$, set
\begin{equation}\label{or:polar-constants}
 f=\frac1{1-r_+^2s_+^2/2},\qquad
 D_p=1+(r_+s_+/2)^2.
\end{equation}
Then
\[
 C<f C_{\rm real},\qquad
 C_{\rm real}=\sum_v\frac{I(v)}{|s-v|^2},\qquad
 \theta\ge\frac{n\rho}{2D_p}.
\]

\begin{lemma}[A negative square-moment term]\label{or:critical-comparison}
For $s=X+\delta$ with $\delta\ge0$, put
\begin{equation}\label{or:kappa}
 \kappa_X(T,\delta)=\frac{X^2}{2}
 -\frac{X\max(0,T^2-X\delta-\delta^2)^2}
             {4(T^2+\delta^2)(X/2+\delta)}.
\end{equation}
Then
\begin{equation}\label{or:critical-comparison-formula}
 C_{\rm real}<\rho\{X\Re S-\kappa_X(T,\delta)Q\}.
\end{equation}
For $X=1$, the coefficient decreases with $T$, and is nonnegative for
$0<T\le1$, $\delta\ge0$.  For $T=1$ and
$\delta\in\{3/10,1/2,71/100\}$ there is also the bound
\begin{equation}\label{or:small-X-comparison}
 C_{\rm real}<m\rho\Re S\quad(0<X\le m),\qquad
                         m=(1-\delta)^2.
\end{equation}
\end{lemma}
\begin{proof}
Let $A=X-r>X/2$, $B=A+\delta$, $b_0=X/2+\delta$, and
$L=T^2+\delta^2$.  Use \eqref{or:inset-critical} and clear the
positive denominators of one summand in
\eqref{or:critical-comparison-formula}.  The coefficient available
after division by $X$ is
\[
 \frac X2+(A-X/2)
       \frac{2\delta A+\delta^2+2Tt-T^2}{(A+\delta)^2+t^2}.
\]
If the second numerator is nonnegative there is no loss.  Otherwise,
with $z_0=2\delta A+\delta^2-T^2<0$,
\[
 \frac{z_0+2Tt}{B^2+t^2}-\frac{z_0}{B^2}
       =\frac{2TtB^2-z_0t^2}{B^2(B^2+t^2)}\ge0.
\]
It suffices to take $t=0$.  The possible loss is
\[
 \frac{(B-b_0)(L-2\delta B)}{B^2}.
\]
As a function of $w=1/B\in[0,1/b_0]$ this is
$(1-b_0w)(Lw-2\delta)$.  Its maximum is
$\max(0,L-2\delta b_0)^2/(4Lb_0)$, by completing the square; when
$L-2\delta b_0\le0$ its maximum is zero at $w=1/b_0$.
This proves the stated coefficient.  For $X=1$ the subtracted term
increases with $T^2$ on its active interval.  Its numerator is at
most $T^4\le T^2\le(1+2\delta)(T^2+\delta^2)$, which also proves
nonnegativity of $\kappa_1$.

For \eqref{or:small-X-comparison}, put $T=1$.  The critical inset
divided by $\rho$ is less than $XA-X^2/2$.  If $A\ge m$ and $X\le m$,
this is at most $mA-m^2/2$.  The required comparison then reduces to
a quadratic in $A$ whose discriminant is
$[(1-\delta)^2-m][(1+\delta)^2+m]=0$; its remaining $t$ and $t^2$
coefficients are nonnegative.  If $A\le m$, use
$XA-X^2/2\le A^2/2$.  The constant term after clearing denominators is
\[
 P(A)=2m(A+\delta)^2-A(A^2+1).
\]
Its derivative is $-3A^2+4mA+4m\delta-1$.  In each of the three
specified cases,
$16m^2+12(4m\delta-1)<0$, so $P$ decreases on the real line.
Furthermore $P(m)=m\delta^4>0$.  The other terms are nonnegative
($2At$ and $(2m-A)t^2$ after the same division), proving the assertion.
\end{proof}

For a closed parameter cell suppose that
$C_{\rm real}<\rho h A$, where $0<A\le\overline A$, and that
$h\ge H$, $h/n\le U$.  With the constants in
\eqref{or:polar-constants}, if
\begin{equation}\label{or:HB-cell-denominator}
              B=2D_p(f\overline A U+1/n_0)<1,
\end{equation}
then Lemma~\ref{or:HB} gives
\begin{equation}\label{or:HB-cell}
                  \frac Wh<\frac{r_+(f\overline A+1/H)}{1-B}.
\end{equation}
If a ceiling $h<H_+$ has previously been proved, we may instead use
$B=2D_p(f\overline A H_++1)/n_0$.  A source lower bound $W\ge b_0$
gives, without either ceiling,
\begin{equation}\label{or:direct-K}
 \rho h A>K(\rho),\qquad
 K(\rho)=\frac1f\left\{\frac{n_0\rho b_0}
                         {n_0\rho+2D_p b_0}-\rho\right\}.
\end{equation}
The expression in braces is concave in $\rho$, since its second
derivative is $-4n_0^2D_pb_0^2/(n_0\rho+2D_pb_0)^3<0$.
It increases with $n_0$ when the original degree is retained.  Thus
the smaller endpoint value on a closed scale interval is a valid
lower bound for the entire interval.  In each application $K>0$,
so $A>0$ and $h>K/(\rho A)$.

\subsection{Moving sources and a uniform crescent gap}

When $X=1$ and $I=I(\beta)>0$, set
\begin{equation}\label{or:moving-test}
 D_0=1-c_\rho(T+I),\qquad s=\frac{1+I}{D_0}.
\end{equation}
The contribution of the actual zero $\beta$ is exactly
\begin{equation}\label{or:moving-weight}
 \frac{I}{|s-(1+ic_\rho s)\beta|^2}
       =\frac{D_0}{2I}\ge\frac{(1-\rho)^2}{2T}.
\end{equation}
Indeed $D_0=|1-ic_\rho\beta|^2$, and expansion gives
$D_0|s-(1+ic_\rho s)\beta|^2=2I^2$.
Moreover $(1-\rho)^2\le D_0<1$, $1<s<2/(1-\rho)^2<3$, and
\begin{align}
 D_0^2(s-1-c_\rho^2s^2)
 &=I\{1+2c_\rho(1-T)+c_\rho^2(T^2-4T-1)
                      +2c_\rho^3(1+T^2)\}>0.\label{or:moving-domain}
\end{align}
The bracket is at least $1-4c_\rho^2$: after subtraction the
remaining terms are
$2c_\rho(1-T)+c_\rho^2(T-1)(T-3)+2c_\rho^3(1+T^2)$.
Thus the test satisfies \eqref{or:Cayley-entrance}.

On a cell $\rho\le r_+$, $T\le t_+$ put
\begin{equation}\label{or:moving-field}
 b_m=\frac{(1-r_+)^2}{2t_+},\quad
 f_m=(1-9r_+^2/2)^{-1},\quad D_m=1+9r_+^2/4,\quad
 H_m=\frac{n_0b_m/(n_0r_++2D_mb_m)-1}{f_m}.
\end{equation}
Since $C<f_m\rho h$, the unsolved inequality
\eqref{or:parallel-source} proves $h>H_m$.  At zero inset it remains
valid with $h\ge H_m$.
At $X=1$ and $I(\beta)=0$,
\[
             T=\frac\rho2(1+T^2)\le\rho<1.
\]
To justify that endpoint, multiply every physical zero and the physical
witness by $1-\epsilon$.  For a sufficiently small positive $\epsilon$
all the finitely many strict holes persist, and $T'<1$ in the
$X=1$ normalization, where
\begin{equation}\label{or:homothety}
 \rho'=(1-\epsilon)\rho,\quad
 T'=T+\frac{\epsilon}{(1-\epsilon)\rho},\quad
 \beta'-v'=\beta-v,\quad S'=S,\quad h'=h.
\end{equation}
The scaled source has positive inset.  Applying the formula with a
depth endpoint tending down to $t_+$ and taking the rational limit
gives $h\ge H_m$.

We also need a distance bound which keeps the actual depth.  The lower
intersection of the hole circle with $r=\alpha/2$ has depth
\[
             \eta=\frac{1-\sqrt{1-\rho^2\alpha^2/4}}{\rho}.
\]
The inequality
\begin{equation}\label{or:height-gate}
 \frac2{125}-\rho_*\left(\frac14+\frac1{15625}\right)>0
\end{equation}
gives $0\le\eta<\alpha^2/125$ by squaring two positive sides.
To bound the distance on the critical crescent, write
$u=r-\alpha\le-\alpha/2$ on the lower hole arc.  The scalar
product with the vector $(X-\alpha,T-1/\rho)$ is
\[
 (X-\alpha)u+(1/\rho-T)\sqrt{\rho^{-2}-u^2}.
\]
Its derivative is positive for $u\le0$.  Hence the nearest permitted
arc point has the largest real coordinate.  The upper arc and exterior
radial segments cannot be nearer; if $T<\eta$, the horizontal
separation alone suffices.  We obtain
\begin{align}
 X=1&:\quad d^2>\frac14+\max(T-1/125,0)^2,\label{or:gap-X}\\
 T=1&:\quad d^2>G(X):=1+\frac{117}{500}X^2.\label{or:gap-T}
\end{align}
For the second line, the stronger lower bound is
$(X-\alpha/2)^2+(1-\alpha^2/125)^2
 \ge X^2/4+(1-X^2/125)^2
 =G(X)+X^4/15625$.

\subsection{From inverse contours to source moments}

We use the inverse and symmetric contour estimates of
Lemmas~\ref{lt:inverse-log} and~\ref{lt:quarter-pi},
and prove two refinements for one-sided strips.

Suppose $Q/h\le Y$, $1/d\le m$, and $h\ge H$.  For $0<mR<1$ put
\begin{equation}\label{or:inverse-parameters}
 \begin{gathered}
 \lambda=\frac{YR}{1-mR},\qquad
 e_R=Y\left\{\sum_{j=2}^{16}\frac{m^{j-2}R^j}{j}
                  +\frac{m^{15}R^{17}}{17(1-mR)}\right\},\\
 \widehat R=R-e_R,\qquad
 B_R=\frac{Y}{(1-mR)^2(1-\lambda)^3}.
 \end{gathered}
\end{equation}
The simpler upper bound
$e_R=Y\{R^2/2+mR^3/[3(1-mR)]\}$ is also permissible and will be
used for the base range.  These bound the complete logarithmic series
of the actual critical product; the last displayed term includes
every omitted order.  For example
\[
 \frac1S\log\frac{F'(\beta+z)}{F'(\beta)}=z+E_R(z),\qquad
 |E_R(z)|\le e_R\quad(|z|\le R),
\]
and the derivative differs from $1$ by at most $\lambda$.
If $\lambda<1$ and $\widehat R>0$, this logarithm is injective and
its image contains the disk of radius $\widehat R$.
In the unscaled coordinate
$w=\log(F'(\beta+z)/F'(\beta))$, the inverse position
$\phi(w)=\beta+z$ is defined on $|w|<h\widehat R$, with
$\phi(0)=\beta$ and $F'(\phi(w))/F'(\beta)=e^w$.
Set $A(w)=S\phi'(w)$.  The normalized primitive
$SF(\phi(w))/F'(\beta)$ has derivative $e^w A(w)$, and
\[
        \Re A\ge(1+\lambda)^{-1},\quad
        |A|\le(1-\lambda)^{-1},\quad |A'|\le B_R/h.
\]

Choose $0<\nu<1$, $W_0>0$, and rational bounds
\begin{equation}\label{or:contour-gates}
 \begin{gathered}
 B_R/H<\epsilon_0,\quad \nu B_R\widehat R<g_0,\quad
 0<p_0\le(1+\lambda)^{-1}<1<(1-\lambda)^{-1}\le a_0,\\
 e_0>e^{-W_0},\qquad
 W_0^2<(1-\nu^2)(H\widehat R)^2,\\
 \epsilon_0(1+e_0)+g_0e_0<1-a_0e_0,\qquad
 \epsilon_0(1+e_0)+g_0e_0^2<p_0-e_0.
 \end{gathered}
\end{equation}
For an odd-$\pi$ contour require
\begin{equation}\label{or:odd-gate}
      \epsilon_0+g_0e_0<p_0,\qquad \nu H\widehat R>\pi.
\end{equation}
For a quarter-period contour replace these requirements by
\begin{equation}\label{or:quarter-gate}
 \epsilon_0+g_0e_0<\frac{p_0}{1+p_0+a_0},\qquad
                        \nu H\widehat R>\pi/2.
\end{equation}
These are conditions on the containing disk, all corners, both vertical
sides, both horizontal sides, and the amplitude homotopy.  The resulting
number $L_R$ of actual original roots in $|u-\beta|<R$ obeys
\begin{equation}\label{or:counts}
 \frac{L_R}{h}>\frac{\nu\widehat R}{\pi}-\frac{\ell_0}{h},
 \qquad
 \ell_0=2\ \hbox{in \eqref{or:odd-gate}},\quad
 \ell_0=3/2\ \hbox{in \eqref{or:quarter-gate}}.
\end{equation}
The trimmed odd height is the largest odd multiple of $\pi$ not
exceeding the available height, including equality.  The full side
estimates already cover such an endpoint.  A one-sided odd-$\pi$ strip
between the fractions $\gamma$ and $\nu$ of the available height,
$0<\gamma<\nu$, gives
\begin{equation}\label{or:one-sided-count}
                    \frac{L_{\gamma,\nu}}h>
           \frac{(\nu-\gamma)\widehat R}{2\pi}-\frac2h.
\end{equation}

\begin{lemma}[One-sided inverse strips]\label{or:one-sided}
Under \eqref{or:contour-gates} and \eqref{or:odd-gate}, the upper
strip $\gamma h\widehat R<\Im w<\nu h\widehat R$ and its reflected
lower strip each have the count \eqref{or:one-sided-count}.
They exclude the source $\beta$.  Writing
$S/h=c_0+is_0$ with $c_0>0$, each upper-strip root lies above
$\beta$ by more than
\begin{equation}\label{or:strip-displacement}
                c_0\gamma\widehat R-e_R-|s_0|W_0/h;
\end{equation}
each lower-strip root lies below it by more than the same quantity.
\end{lemma}
\begin{proof}
Put $\mathcal J(w)=SF(\phi(w))/F'(\beta)$.
The estimates preceding \eqref{or:contour-gates} give
\[
 A(0)=1,\quad \Re A\ge p_0,\quad |A|\le a_0,\quad
 |A'|<\epsilon_0,\quad
 \mathcal J(w)=\int_0^w e^zA(z)\,dz.
\]
All corners of the containing rectangle lie in the inverse disk.
Integration by parts along the entire path
$0\longrightarrow-W_0\longrightarrow-W_0+i\eta
 \longrightarrow x+i\eta$ gives
\[
 \mathcal J(w)=e^wA(w)-1-\mathcal E(w),\qquad
 |\mathcal E(w)|<\epsilon_0(1+e^x)+g_0e^{-W_0}.
\]
The vertical segment contributes at most
$(B_R/h)|\eta|e^{-W_0}\le\nu B_R\widehat R e^{-W_0}$;
this gives the $g_0e^{-W_0}$ term.
On the left side the modulus of $e^wA(w)-1$ is at least
$1-a_0e^{-W_0}$, and on the right side it is at least
$p_0e^{W_0}-1$.  The two inequalities in the last line of
\eqref{or:contour-gates} control these whole sides, respectively.
On every horizontal line $\eta=(2k+1)\pi$ the modulus is at least
$1+p_0e^x$.  It dominates the error, because
\[
 \epsilon_0(1+e^x)+g_0e^{-W_0}
 \le(\epsilon_0+g_0e^{-W_0})(1+e^x)
 <p_0(1+e^x)\le1+p_0e^x.
\]
Thus Rouch\'e applies on every subrectangle bounded by two such
horizontal lines.  The homotopy
$A_t=1+t(A-1)$ preserves $\Re A_t\ge p_0$ and $|A_t|\le a_0$,
so $e^wA_t(w)-1$ stays nonzero on this boundary.
The zero count is therefore exactly that of $e^w-1$.

For the upper strip take the smallest odd multiple $L$ of $\pi$
at least $\gamma h\widehat R$, and the largest odd multiple $U$
at most $\nu h\widehat R$.  If $L<U$, the subrectangle contains
exactly $(U-L)/(2\pi)$ roots, and
\[
 \frac{U-L}{2\pi}>
            \frac{(\nu-\gamma)h\widehat R}{2\pi}-2.
\]
If the right side is negative, the claimed lower count is automatic;
if it is nonnegative, the strip width is at least $4\pi$ and
$L<U$ necessarily.  The same argument applies to $-U,-L$.
The boundaries contain no zeros, including when an endpoint is exactly
an odd multiple of $\pi$.  The inverse is injective and its image
is critical-free, so all counted roots are distinct simple actual
roots.  Since $\gamma>0$, neither strip contains $w=0$.
Finally
\[
             |\phi(w)-\beta-w/S|\le e_R,
 \qquad \Im(w/S)=\frac{c_0\Im w-s_0\Re w}{h},
\]
which proves \eqref{or:strip-displacement} and its lower-strip
counterpart.
\end{proof}

\begin{lemma}[Real widths and an upper root]\label{or:narrow-strip}
If $|\Im S|<J$, the roots counted in \eqref{or:counts} satisfy
\begin{equation}\label{or:width}
 |\Re(u-\beta)|\le\ell_R:=\min\left(R,
         \frac{W_0}{H}+\frac{J\nu\widehat R}{H}+e_R\right).
\end{equation}
For a row satisfying \eqref{or:odd-gate}, a positive upper strip
also supplies an actual root above $\beta$
by more than
\begin{equation}\label{or:depth}
 \left(1-\frac{J^2}{H^2}\right)\nu\widehat R-e_R
       -\frac{4\pi}{H}-\frac{JW_0}{H^2},
\end{equation}
whenever this expression is positive.  If it exceeds the depth $T$,
the original disk condition is contradicted.
\end{lemma}
\begin{proof}
In the counting rectangle $|\Re w|\le W_0$,
$|\Im w|\le\nu h\widehat R$.  Since
\[
 \Re(w/S)=\frac{c_0\Re w+s_0\Im w}{h},\qquad
 |s_0|<J/h,
\]
the inverse error bound gives \eqref{or:width}.
For the upper root let $U$ be the largest odd multiple of $\pi$
not exceeding $\nu h\widehat R$.
The adjacent rectangle $U-2\pi<\Im w<U$ contains exactly one
actual root by the complete boundary argument of
Lemma~\ref{or:one-sided}.  When \eqref{or:depth} is positive,
$\nu H\widehat R>4\pi$ and $J/H<1$, so that this rectangle has
positive height above zero.  Its root has
$\Im w>\nu h\widehat R-4\pi$.
Using $c_0=\sqrt{1-s_0^2}\ge1-(J/H)^2$, $c_0\le1$, and $h\ge H$
in the imaginary inverse displacement gives exactly
\eqref{or:depth}.  Every original zero has nonpositive imaginary
part in the normalized coordinates, which proves the last assertion.
\end{proof}

We now convert these counts into lower bounds for $E(s)$.  Suppose
$T\in[t_-,t_+]$, $\rho\le r_+$, and $s=X+\delta$.  A counted region
with radius $R$ and real width $\ell$ has $|\Re u|\le1$ in every
application below.  For $X=1$ this follows from rightmostness and
$\ell<2$; for $T=1$ it follows from $R<1$ and $0<X\le1$.
Its whole-region lower weight is
\begin{equation}\label{or:weight}
 \omega(R,\ell;\delta,t_-,t_+,r_+)
 =\max\left(0,\min_{t\in\{\max(0,t_--R),\,t_++R\}}
 \frac{t-r_+(1+t^2)/2}{(\delta+\ell)^2+t^2}\right).
\end{equation}
The derivative numerator of the quotient, writing
$A=(\delta+\ell)^2>0$, is
$A-t^2+r_+(1-A)t$.  Its two roots have opposite signs, and its
positive root is a maximum.  Thus the minimum on a nonnegative
interval is at an endpoint.  If that minimum is negative we instead
use $I(u)\ge0$.

Take increasing radii $R_j$; replace their widths by running maxima
so the regions are nested.  Let $s_j$ be any increasing rational lower
bounds for $\nu_j\widehat R_j/\pi$, and replace the weights by the
maximum of all later weights.  This gives valid
$w_1\ge\cdots\ge w_k\ge0$ on the nested regions.  The source root
has a separate lower weight $b_0$.  Subtracting that root once from
each odd-$\pi$ count, and applying finite Abel summation, gives
\begin{equation}\label{or:Abel-source}
 \frac Eh>
     \sum_{j=1}^k w_j(s_j-s_{j-1})+\frac{b_0-3w_1}{h},
                         \qquad s_0=0.
\end{equation}
Indeed, if $M_j$ are the actual cumulative counts excluding $\beta$,
then $M_j>hs_j-3$ and
$\sum w\ge w_kM_k+\sum_{j<k}(w_j-w_{j+1})M_j$.
The coefficients of the cumulative counts are nonnegative, so the bound
remains valid when an initial $hs_j-3$ is negative.
If $b_0-3w_1<0$, replace its denominator
in \eqref{or:Abel-source} by $H$.  If it is nonnegative and $h<600$,
replace it by $600$; if no upper field is known, discard the
nonnegative term.  This accounts for both signs of the source
compensation.
For a one-sided strip the corresponding calculation is simpler:
if its roots have weight at least $a>0$ and $\beta$ has weight at
least $b_0$, then, since they are distinct from $\beta$,
\begin{equation}\label{or:strip-source-compensation}
 \frac Eh>
 a\frac{(\nu-\gamma)\widehat R}{2\pi}+\frac{b_0-2a}{h}.
\end{equation}
Thus $b_0\ge2a$ makes the last term nonnegative.

For quarter-period counts, one may instead use the cumulative lower
bounds $k_j=\nu_j\widehat R_j/\pi-3/(2H)$ directly, when they and
their successive increments are positive.  Then
\begin{equation}\label{or:Abel-quarter}
                    \frac Eh>\sum_jw_j(k_j-k_{j-1}),\qquad k_0=0.
\end{equation}

\subsection{The base range $0<\rho\le1/100$}

All constants in this subsection use $r_b=1/100$.  Set
\begin{equation}\label{or:base-polar}
 f_m=\frac{1001}{1000},\qquad f_X=f_T=\frac{10001}{10000},
                    D_b=1+9r_b^2/4.
\end{equation}
These dominate the quadratic polar factors for, respectively, the
moving tests $s<3$, the fixed test $s=101/100$, and the fixed test
$s=X+3/10\le13/10$.  The fixed inverse-Cayley conditions are
$(r_b(101/100)/2)^2<1/100$ and
$(r_b(13/10)/2)^2<3/10$.  Lemma~\ref{or:critical-comparison} gives
$C_{\rm real}<\rho h$ for the $X$ tests, and
$C_{\rm real}<\rho\max(X,49/100)h$ for the $T$ test.

There are $55$ inverse rows in this range.  The columns in
Table~\ref{or:base-contours} are
\[
 (Y,m,R,H,\nu,W_0,\epsilon_0,g_0,p_0,a_0,e_0).
\]
The simpler $e_R$ following \eqref{or:inverse-parameters} is used.
\begin{table}[ht]
\centering\scriptsize
\caption{All inverse rows for the base range.  Each $j$-row occurs for
every integer $6\le j\le30$.}\label{or:base-contours}
\begin{tabular}{@{}lrrrrrrrrrrr@{}}
\toprule
row&$Y$&$m$&$R$&$H$&$\nu$&$W_0$&$\epsilon_0$&$g_0$&$p_0$&$a_0$&$e_0$\\
\midrule
$X_1$&2&2&$3/20$&600&$99/100$&4&$1/25$&3&$7/10$&$7/4$&$1/50$\\
$X_D$&2&2&$13/100$&600&$99/100$&4&$1/25$&3&$7/10$&$7/4$&$1/50$\\
$X_{450,j}$&2&2&$j/200$&450&$23/25$&$W_j$&$11/100$&$g_j$&$7/10$&$7/4$&$e_j$\\
$X_{200,j}$&2&2&$j/200$&200&$9/10$&$W_j$&$11/100$&$g_j$&$7/10$&$7/4$&$e_j$\\
$X_{78}$&2&2&$3/20$&78&$97/100$&$11/5$&$3/10$&3&$7/10$&$7/4$&$1/9$\\
$T_{150}$&$25/11$&$100/99$&$3/20$&150&$9/10$&4&$1/6$&3&$7/10$&$7/4$&$1/50$\\
$T_{103}$&$25/11$&$100/99$&$3/20$&103&$9/10$&4&$1/6$&3&$7/10$&$7/4$&$1/50$\\
\bottomrule
\end{tabular}
\end{table}
Here $(W_j,g_j,e_j)=(2,1/5,1/7)$ for $6\le j\le9$, and
$(4,3,1/50)$ for $10\le j\le30$.  Only the two $j$-families use
quarter-period contours; all other rows use odd-$\pi$ contours.
Every row satisfies \eqref{or:contour-gates} and its appropriate
horizontal and height condition.  These are rational inequalities
apart from $e^{-W_0}$; the latter is bounded by the reciprocal of the
positive Taylor sum through degree $48$.  Thus Table~\ref{or:base-contours}
and the displayed formulae specify the finite verification.

\subsubsection{The base case $X=1$: the first ceiling}

Suppose $h\ge600$.  Then $d>1/2$ and $y_0<2$, so
$q<2/600<7/625$ and $\mu<2/600<1/200$.
Proposition~\ref{lt:fine200} therefore gives
\begin{equation}\label{or:base-high-phase}
 |\Im S|<4y_0<8,\qquad
 |\Im(S/h)|<1/70,\qquad \Re(S/h)>999/1000.
\end{equation}
The upper one-sided strip of row $X_1$, with $\gamma=4/5$, has
positive count, since
\[
 \frac{(99/100-4/5)\widehat R_1}{2(22/7)}-\frac2{600}>0.
\]
Its actual root lies above $\beta$ by more than
$(4/5)(999/1000)\widehat R_1-e_{R_1}-4/(70\cdot600)>67/1000$.
Thus $T>67/1000$.

First let $67/1000\le T\le3/10$.  Equation~\eqref{or:radial-ratio},
with $J=9$, gives $u=h/n<311/1000=:U$ for
every $N\ge6250$.  At $s=101/100$,
\[
       \kappa_1(T,1/100)\ge\kappa=\frac{29591}{63600}>0.
\]
As $y_0>u$, put $a(u)=f_X(1-\kappa u)$ and
$B(u)=2D_b(a(u)u+1/n_0)$.  On $0\le u\le U$ the factors
$1-\kappa u$, $1-2\kappa u$, and $1-B(u)$ are positive.
The derivative numerator of
$(a(u)+1/600)/(1-B(u))$ is
\[
 -f_X\kappa(1-2D_b/n_0)
 +2D_b\{f_X^2(1-\kappa u)^2+(f_X/600)(1-2\kappa u)\}.
\]
It has a positive lower bound obtained by replacing the first term
by $-f_X\kappa$ and $u$ by $U$ in the remaining terms.
Therefore the solved source estimate is valid throughout this
interval and
\begin{equation}\label{or:base-low-depth-upper}
                  W/h<r_b\frac{a(U)+1/600}{1-B(U)}<1/50.
\end{equation}

Use the lower one-sided strip of $X_D$, with $\gamma=7/25$, and set
\[
 \begin{split}
 \Delta&=\frac7{25}\frac{999}{1000}\widehat R_D-e_{R_D}
          -\frac4{70\cdot600}-\frac{3481}{5000}r_b>0,\\
 a_D(T)&=\frac{T+\Delta}{(T+157/1200)^2},\qquad
 b_D(T)=\frac{T-r_b(1+T^2)/2}{T^2+1/10000}.
 \end{split}
\]
These are full strip and source weights.  The strip depth estimate
comes from \eqref{or:base-high-phase}; its roots have
$|u|<\sqrt{1.09}+.13<59/50$ and
$|s-u|<\sqrt{T^2+.0001}+.13\le T+157/1200$.
The polynomial
\[
 [T-r_b(1+T^2)/2](T+157/1200)^2
             -2(T+\Delta)(T^2+1/10000)
\]
has all its Bernstein coefficients positive on
$[67/1000,3/10]$.  Hence $b_D(T)>2a_D(T)$ there.
Adding $\beta$ to the strip count in \eqref{or:one-sided-count}
compensates the loss $2/h$ and yields
\[
 \frac Eh>\frac{(99/100-7/25)\widehat R_D}{2(22/7)}a_D(T)>1/50.
\]
For the final comparison it suffices to check the two endpoints:
the only positive stationary point of $a_D$ is a maximum.
This contradicts \eqref{or:base-low-depth-upper}.

If $3/10\le T\le1$, equation~\eqref{or:radial-ratio} gives
$h/n<147/1000$.  Its source denominator is less than $3/10$, and
\[
 \frac Wh<
 \frac{r_b(f_X+1/600)}{1-2D_b(f_X(147/1000)+1/n_0)}<2/125.
\]
The symmetric $X_1$ count has lower weight
\[
       \frac{T-3/20-5r_b/4}{(T+901/6000)^2}.
\]
Indeed $|u|<\sqrt2+3/20<\sqrt{5/2}$ and the displayed denominator
bounds the complete disk.  Multiplication by
$(99/100)(7/22)\widehat R_1-2/600$ gives a lower bound exceeding
$1/50$ at both endpoints, and therefore on the whole interval, again
by its single-maximum derivative.  Thus $h\ge600$ is impossible.

\subsubsection{The base case $X=1$: the remaining fields}

We now have $h<600$.  Lemma~\ref{or:zero-inset} excludes $I(\beta)=0$,
since $N/2\ge3125$.
Apply the moving-source estimate at this positive inset.
The following table gives entire-domain denominator and depth bounds:
\begin{equation}\label{or:base-moving-table}
\begin{array}{c|c|c}
 \text{known ceiling }h<H_+& (C+\rho)/\theta<B_+&T>t_+\\ \hline
 600&193/1000&13/200\\
 450&29/200&9/100\\
 200&13/200&21/100
\end{array}
\end{equation}
For each row the two exact checks are
\[
 \frac{2D_b(H_+f_m+1)}{6251}<B_+,\qquad
 \frac{(1-r_b)^2}{2t_+}>
                   \frac{r_b(H_+f_m+1)}{1-B_+}.
\]
Thus they hold for every positive $\rho\le r_b$.

Suppose first $h\ge450$ and use all $25$ rows $X_{450,j}$.
For each radius, apply \eqref{or:weight} with $\ell=R$ and
$\delta=1/100$ on each of the two closed depth intervals
$[13/200,1/5]$ and $[1/5,1]$.  Use one of these depth intervals
for the entire nested family.  The cumulative bounds
\[
 k_j=\nu_j(7/22)\widehat R_j-\frac{3}{2H}
\]
and their successive increments are positive, and the resulting
weights decrease.  Formula~\eqref{or:Abel-quarter} applies.
After excluding this field, repeat with $H=200$ and depth intervals
$[9/100,1/5]$, $[1/5,1]$.  The exact inequalities, rounded outwards
here to shorter rational numbers, are
\begin{equation}\label{or:base-middle-table}
\begin{array}{c|c|c|c}
 H&[t_-,t_+]&E/h\text{ exceeds}&W/h\text{ is less than}\\ \hline
 450&[.065,.2]&.042&.0125\\
 450&[.2,1]&.029&.0125\\
 200&[.09,.2]&.041&.0118\\
 200&[.2,1]&.0245&.0118
\end{array}
\end{equation}
The source upper in each row is exactly
$r_b(f_X+1/H)/(1-B_+)$ with the previously established $B_+$.
These comparisons prove successively $h<450$, $T>.09$, $h<200$,
and $T>.21$.

On $[.21,1]$ the source weight is at least
\[
                   b_X=\frac{9900}{10001}.
\]
To check the whole interval, the numerator difference after clearing
the positive denominators factors as $(1-T)(AT-B)$, where
$A=1-(r_b/2)(1-1/10000)$,
$B=(r_b/2)(1-1/10000)+1/10000$, and $A(.21)>B>0$.
Let
$a_{78}=\omega(3/20,3/20;1/100,.21,1,r_b)$.
If $h\ge78$, the row $X_{78}$ and source compensation give
\[
 \frac Eh>
 a_{78}\frac{97}{100}\frac7{22}\widehat R_{78}
            +\frac{\min(0,b_X-3a_{78})}{78}>.01176,
 \qquad
 \frac Wh<\frac{r_b(f_X+1/78)}{1-.065}<.01084.
\]
Thus $h<78$.  Finally,
\[
              W<\frac{r_b(78f_X+1)}{1-.065}<.846<b_X\le E,
\]
which excludes the entire base case $X=1$.

\subsubsection{The base case $T=1$}

A direct physical comparison gives $d>99/100$.  In fact, with
$d_0=99/100$,
\[
 (1-d_0\rho)^2-|b-a|^2
 \ge\rho\{2(1-d_0)-(2-d_0^2)\rho\}>0,
\]
and $|a-\xi|>1$ then implies $|b-\xi|>d_0\rho$.
If $h\ge150$, this first gives $dh>100$, so
Proposition~\ref{mf:D100} yields $y_0<25/11$.  The bounds
$q<25/(11\cdot150)<9/400$ and
$\mu<1/(.99\cdot150)<1/100$ permit Proposition~\ref{lt:wideD100},
giving $|\Im S|<(23/5)y_0<11$.
Equation~\eqref{or:radial-ratio} implies $h/n<21/250$.
The row $T_{150}$ then gives
\[
 \frac Eh>.01599>.01211>
 \frac{r_b(f_T+1/150)}{1-2D_b(f_T(21/250)+1/n_0)}>\frac Wh.
\]
Thus $h<150$, whence $(C+\rho)/\theta<1/20$.

If $h\ge103$, the bound $103(.99)>100$ permits
Proposition~\ref{mf:D100}.
The row $T_{103}$ gives
\[
 \frac Eh>.01145>.01063>
              \frac{r_b(f_T+1/103)}{.95}>\frac Wh.
\]
Hence $h<103$.  The source $\beta$ itself always has weight at least
$b_T=99/109$.  For $X\le.77$, the uniform small-$X$ comparison
gives
\[
            W<\frac{r_b(103f_T(.77)+1)}{.95}<.846<b_T,
\]
a contradiction.

It remains to consider $X>.77$.  If $h\ge79$,
\eqref{or:horizontal} and $d>.99$ give
\begin{equation}\label{or:base-quartic-entry}
             q<\frac2{.77\cdot79}<\frac{33}{1000},
             \qquad \mu<\frac1{.99\cdot79}<\frac8{625}.
\end{equation}
The quartic construction of Proposition~\ref{lt:quartic} supplies two
distinct simple actual roots at distance less than $15/(2h)$ from
$\beta$.  For this construction, at $q_b=33/1000$, $\mu_b=8/625$,
$\nu_b=q_b\mu_b$,
use
\[
 (Z_0,M_0,E_0,P_0,Z_{\mathrm{out}},L_0,R_0,\eta_0)
   =(719/100,7/25,71/50,350,15/2,32/25,3/10,11/50).
\]
The center and path majorants of that construction are bounded by
\[
 \begin{split}
 Z(q)&=63/10+21q+133q^2+1030q^3+9000q^4,\\
 M(q)&=(63/10)q+45q^2+400q^3+3400q^4+86000q^5\\
     &\quad+320000q^6+1700000q^7+9100000q^8+39000000q^9.
 \end{split}
\]
They satisfy $Z(q_b)<Z_0$, $M(q_b)<M_0$,
$\mu_bZ_{\mathrm{out}}<1$, $Z_0+R_0<Z_{\mathrm{out}}$, and
\begin{align}
 \exp\left(M_0+\frac{\nu_bZ_0^3}{3(1-\mu_bZ_0)}\right)&<E_0,
 &\frac{Z_0^4E_0}{12(1-\mu_bZ_0)}&<P_0,\notag\\
 E_0(132200q_b^5+P_0\nu_b)&<(1-\eta_0)R_0,
 &1+q_bZ_{\mathrm{out}}+
 \frac{\nu_bZ_{\mathrm{out}}^2}{1-\mu_bZ_{\mathrm{out}}}&<L_0,
 \label{or:quartic-gates}\\
 \frac{e^{L_0R_0}-1-L_0R_0}{L_0}&<\eta_0R_0.&&\notag
\end{align}
The complete Gaussian residual is $132200q_b^5$, on its valid
domain $q_b<1/20$; the integrated actual cubic logarithmic tail
retains the factors $1/3$ and $1/4$ in
\eqref{or:quartic-gates}.  The last two residual bounds give strict
Rouch\'e inequalities on both entire circles.  Those circles avoid
the source and are disjoint because
\[
 6-21q_b-133q_b^2-1030q_b^3-9000q_b^4>R_0,
 \quad
 12(1-q_b-17q_b^2-137q_b^3-313q_b^4)>2R_0.
\]
Rouch\'e's theorem therefore gives the two distinct simple zeros of $F$.

Since $15/(2\cdot79)<19/200$, each companion has weight at least
\[
 \omega(19/200,19/200;3/10,1,1,r_b)
                         =\frac{8672039}{10840400}.
\]
Adding the source once gives
$E\ge b_T+2(8672039/10840400)>r_b(103f_T+1)/.95>W$.
Thus $h<79$, and the remaining source comparison
$W<r_b(79f_T+1)/.95<.843<b_T\le E$ is again impossible.
This proves Proposition~\ref{or:theorem} for all $0<\rho\le1/100$.

\subsection{The higher strip with $T=1$: entering the local estimates}

Assume throughout this subsection that
\begin{equation}\label{or:T-domain}
           1/100\le\rho\le\rho_*,\quad T=1,\quad 0<X\le1.
\end{equation}
The closed lower endpoint overlaps the base range.  By
\eqref{or:gap-T}, $d>\sqrt{G(X)}>1$.
For a test $s=X+\delta$ write
\[
 \kappa_\delta(X)=\kappa_X(1,\delta),\qquad
 A_\delta=X-\kappa_\delta(X)y_0,
 \quad b_\delta(X)=\frac{1-\rho_*(1+X^2)/2}{1+\delta^2}.
\]
On the active interval $1-X\delta-\delta^2\ge0$ the coefficient factors as
\begin{equation}\label{or:T-kappa-factor}
 \kappa_\delta(X)=
 \frac{X[X-(1-\delta)^2][X+(1+\delta)^2]}
                 {2(1+\delta^2)(X+2\delta)};
\end{equation}
outside it $\kappa_\delta(X)=X^2/2$.  With
$D_\delta(X)=1+\rho_*^2(X+\delta)^2/4$ and
$f_\delta(X)=(1-\rho_*^2(X+\delta)^2/2)^{-1}$, define
\begin{equation}\label{or:T-source-function}
 K_\delta(r,X)=\frac1{f_\delta(X)}
 \left\{\frac{n_0r b_\delta(X)}{n_0r+2D_\delta(X)b_\delta(X)}-r\right\}.
\end{equation}
Equation~\eqref{or:direct-K} gives
$\rho hA_\delta>\min_{r\in\{1/100,\rho_*\}}K_\delta(r,X)$.
All tests used here have $1/200\le\delta\le1$; the inverse-Cayley
condition follows from $\rho_*^2<1/200$.

For $\delta=1/2$, replacing the source and polar factors by their
$X=1$ bounds gives
\begin{equation}\label{or:Kh}
 K_h=\min_{r\in\{1/100,\rho_*\}}K_{1/2}(r,1)>.6843,
 \qquad \rho h A>K_h,
\end{equation}
where $A=1/4$ if $X\le1/4$, by
\eqref{or:small-X-comparison}, and
$A=X-\kappa_{1/2}(X)y_0$ otherwise.  Here
\[
              \kappa_{1/2}(X)=\frac{X(4X-1)(4X+9)}{40(X+1)}.
\]
If $X\le2/5$, \eqref{or:Kh} gives
$D>h>K_h/((2/5)\rho_*)>25$.
For $2/5\le X\le1$, the one-variable inequality
\begin{equation}\label{or:T-first-square}
            4(227/200)\kappa_{1/2}-(X+\kappa_{1/2})^2>0
\end{equation}
implies $A(1+y_0)<227/200$ by completion of the square in $y_0$.
Multiplication of \eqref{or:Kh} and
\eqref{or:radial-unphased} then gives
$h^2>K_hN/[2(227/200)]>25^2$.
Thus Proposition~\ref{mf:D25} is applicable in all cases, and
\begin{equation}\label{or:T-ratio3}
                  y_0<3/d<3/\sqrt{G(X)}\le3.
\end{equation}

The same moment ratio sharpens the joint bound to
\begin{equation}\label{or:T-joint-gap}
                  \frac{A(1+y_0)}{G(X)}\le\frac{21}{20}.
\end{equation}
Here are whole-interval certificates for this assertion.  For
$X\le1/4$ use $(1/4)(1+3)=1$.  Put
$Y(X)=3(1-11X^2/100)$.  Squaring positive sides gives
$Y(X)>3/\sqrt{G(X)}$ on $0<X\le.49$; after removal of $X^2$ the
required polynomial is strictly positive.  The ratio
$\kappa_{1/2}(X)/X$ increases, and seven times its value at $.41$
is less than one.  Hence on $[1/4,.41]$ the product
$(X-\kappa_{1/2}y_0)(1+y_0)$ increases up to its allowed endpoint,
and the certificate is
\[
 (21/20)G(X)-(X-\kappa_{1/2}Y(X))(1+Y(X))>0.
\]
On $[.41,1]$ completion of the square gives the certificate
$4(21/20)G(X)\kappa_{1/2}-(X+\kappa_{1/2})^2>0$.
Consequently
\begin{equation}\label{or:T-D45}
                 D^2>\frac{K_hN}{2(21/20)}>45^2.
\end{equation}

Write $S/h=c_0+is_0$, $c_0>0$, $c_0^2+s_0^2=1$.
The universal envelope in Proposition~\ref{mf:envelope} is
\begin{equation}\label{or:T-matrix}
 q<\frac{21}{5}\mu^2+\frac{11}{2}\mu c_0
                         +\frac{c_0^2}{4}-\frac{s_0^2}{2}.
\end{equation}
If $|s_0|\ge.7$, then $c_0<18/25$, and the right side is at most
\[
 \frac{21/5}{45^2}+\frac{(11/2)(18/25)}{45}
                       +\frac{1-.7^2}{4}-\frac{.7^2}{2}<0.
\]
Thus $|\Im S|<.7h$.

Now take $\delta=3/5$.  On $[0,1]$ the coefficient is always active.
The rational functions
$K_{3/5}(1/100,X)-K_{3/5}(\rho_*,X)$ and
$K_{3/5}(\rho_*,X)$ are positive.  Write
\[
 K_{.6}(X)=K_{3/5}(\rho_*,X),\quad
 A_{.6}=X-\kappa_{3/5}(X)y_0,\quad
 R_{.6}(y_0)=.7+.3\rho_*+(199/200)y_0.
\]
Equation~\eqref{or:radial} gives $N\rho/2<hR_{.6}(y_0)$.
Combining this with $\rho hA_{.6}>K_{.6}(X)$ gives
\[
 h^2>\frac{NK_{.6}(X)}{2A_{.6}R_{.6}(y_0)},\qquad
 D^2>\frac{NK_{.6}(X)G(X)}{2A_{.6}R_{.6}(y_0)}.
\]
The two rectangle inequalities are
\begin{align}
 N_0K_{.6}(X)G(X)-2\cdot51^2 A_{.6}R_{.6}(y_0)&>0,
       &&0\le X\le1, 0\le y_0\le3,\label{or:T-D51-poly}\\
 N_0K_{.6}(X)-2\cdot45^2 A_{.6}R_{.6}(y_0)&>0,
       &&0\le X\le1, 0\le y_0\le3.\label{or:T-h45-poly}
\end{align}
They prove $D>51$ and $h>45$.  An exact finite positivity procedure
for these two rational quadratics and all the preceding one-variable
certificates is given below.  If $|s_0|\ge2/3$, then $c_0<3/4$, and
\[
 \frac{21/5}{51^2}+\frac{(11/2)(3/4)}{51}
                                  +\frac5{36}-\frac29<0.
\]
Therefore
\begin{equation}\label{or:T-prephase}
                      h>45,\quad D>51,\quad |\Im S|<(2/3)h.
\end{equation}

\subsubsection{Entry to the larger companion rectangle}

Set $\delta=.71$ and $K(X)=K_\delta(\rho_*,X)$.
The same endpoint comparison proves that $K(X)$ is the smaller
scale endpoint; also $K(X)\ge K_*:=K(1)>0$.
Using \eqref{or:T-prephase} in the exact radial identity gives
\[
 \frac{N\rho}{2}<hR,\qquad
 R=\rho X+\frac23(1-\rho)
                   +\left(1-\frac\rho2(1+X^2)\right)y_0.
\]
Consequently the same source gives
\begin{equation}\label{or:T-q-square}
 q^2<\frac{2y_0^2[X-\kappa_\delta(X)y_0]R}{NK(X)}.
\end{equation}
For $X\le.49$, $R$ decreases with $\rho$.  At $\rho=.01$ put
$\beta_0=33/50+X/100$, $I_0=199/200-X^2/200$.
If $r(X)=\kappa_\delta(X)/X$, the derivative in $y_0$ of
$y_0^2(X-\kappa_\delta y_0)(\beta_0+I_0y_0)$, divided by $Xy_0$,
is
\[
       2\beta_0+3(I_0-r\beta_0)y_0-4rI_0y_0^2.
\]
The ratio $r$ increases: after positive denominators are removed
its derivative numerator is
$\delta^4+6\delta^2+4\delta X+X^2+1>0$.
Let $r_+$ be its value at $.49$, and take the endpoint extrema of
$\beta_0,I_0$ on $[0,.49]$.  The last display is bounded below by
\[
 2\beta_-+3(I_--r_+\beta_+)y_0-4r_+I_+y_0^2.
\]
This concave quadratic is positive at $0$ and $3$.  It follows that
the permitted $y_0$ may be replaced by $Y(X)$ from above.  The entire
interval is then certified by
\begin{equation}\label{or:T-q-low-poly}
 N_0K(X)-\frac{2}{(.057)^2}Y(X)^2
       [X-\kappa_\delta(X)Y(X)](\beta_0+I_0Y(X))>0
                   \quad(0\le X\le.49).
\end{equation}
For $X\ge.49$ the active interval has
$\kappa_\delta(X)\ge(12/25)X^2$; its inactive value $X^2/2$ is
stronger.  Completion of the square gives
$y_0[X-\kappa_\delta(X)y_0]\le25/48$.
Also \eqref{or:T-ratio3} gives $y_0<73/25$.
Thus the right side of \eqref{or:T-q-square} is at most
\[
 \frac{2(25/48)(73/25)
       \{2/3+\rho_*/3+(199/200)(73/25)\}}{N_0K_*}<(.057)^2.
\]
We have proved
\begin{equation}\label{or:T-Q057-entry}
                            q<57/1000,\qquad \mu<1/51.
\end{equation}
Proposition~\ref{lt:W8} applies in the form \eqref{lt:Q057}.
Its two distinct actual roots
$u_\pm$ have normalized positions $w_\pm=S(u_\pm-\beta)$ satisfying
\begin{equation}\label{or:T-Q057-pair}
 |w_\pm|<8917/1000,\quad
 \left|\frac{w_++w_-}{2}\right|<2109/1000,\quad
                         |\Im S|<(37/5)y_0.
\end{equation}
The mean estimate uses the even part of the two centers at the same
moment parameters, as specified in that proposition.

\subsection{The higher $T=1$ strip: two field ceilings}

If $h\ge105$, then $D>100$ by \eqref{or:gap-T}.
Propositions~\ref{mf:D100} and \ref{lt:wideD100} give
\begin{equation}\label{or:T-first-phase}
         y_0<9/4,\quad q<9/400,\quad \mu<1/100,\quad
                         |\Im S|<(23/5)y_0.
\end{equation}
For a closed cell $\rho\in[p,r]$, $y_0\in[l,u]$, take
$J=(23/5)u$ and a rational $H\ge105$ with either $H=105$ or
\begin{equation}\label{or:T-field-H}
                  H\le\frac{N_0p/2-J}{r+u}.
\end{equation}
Then $h\ge H$.  Put
\begin{equation}\label{or:T-radial-cell}
             B_r=\frac{r/2+J/N_0}{1-r/2}.
\end{equation}
A positive rational $U$ is a valid bound for $h/n$ if at least one
of $U\ge u$, $U^2\ge B_r$, or $Ul\ge B_r$ holds; the last is
used only for $l>0$.  These follow respectively from
\eqref{or:horizontal}, \eqref{or:radial-ratio}, and
\eqref{or:radial-sharp}.

Retain the same-$X$ distance information.  A cap $0<X_+\le1$ is
valid if $X_+=1$ or
\begin{equation}\label{or:T-Xcap}
                         l^2G(X_+)\ge(9/4)^2.
\end{equation}
For the chosen $\delta$, certify
\begin{equation}\label{or:T-Abar}
 \overline A-X+\kappa_\delta(X)v>0
       \quad(0\le X\le X_+,\quad v\in\{l,u\}).
\end{equation}
Split the $X$ interval at $(1-\delta^2)/\delta$ to use the correct
active or inactive formula.  Since the expression is affine in
$v$, this verifies the whole $(X,y_0)$ cell.  The constants in
\eqref{or:polar-constants} use $s_+=1+\delta$ and $r_+=r$.
After checking \eqref{or:HB-cell-denominator},
\eqref{or:HB-cell} is the source upper bound.

For each contour use \eqref{or:inverse-parameters} with $Y=u$, $m=1$,
and set
\begin{equation}\label{or:automatic-contour}
 \epsilon_0=\frac{1001B_R}{1000H},\quad
 g_0=\frac{1001\nu B_R\widehat R}{1000},\quad
 p_0=(1+\lambda)^{-1},\quad a_0=(1-\lambda)^{-1},\quad
 e_0=\left(\sum_{j=0}^{80}\frac{W_0^j}{j!}\right)^{-1}.
\end{equation}
Verify \eqref{or:contour-gates} and \eqref{or:odd-gate}, using
$\pi<22/7$.  Use \eqref{or:width}, rounded upwards, and the weights
\eqref{or:weight} with $t_-=t_+=1$.  A valid source weight is
$b_0=(1-r)/(1+\delta^2)$.  The source lower bound is
\begin{equation}\label{or:T-cloud-lower}
  \sum_jw_j(s_j-s_{j-1})+\frac{\min(0,b_0-3w_1)}H.
\end{equation}
The source upper is \eqref{or:HB-cell}.  Table~\ref{or:T-cover}
specifies the two full covers; the rational choices of cells and
contours are in the electronic appendix.  On every cell of the first
cover, \eqref{or:T-cloud-lower} strictly exceeds that upper.  The
coverage and arithmetic procedures are specified in
\S\ref{or:finite-verification}.  This proves $h<105$.

Suppose next that $90\le h<105$.  Since $D>90$,
Proposition~\ref{lt:D70} gives, with $z=Qd/h$,
\[
 z<8/3,\quad |\Im S|<(21/4)y_0,\quad
                   3z^2-(D-91/10)z+2D>0.
\]
At $D=90$ this quadratic is negative at both $z=8/3$ and $z=49/20$.
These values decrease with $D$, since $2-z<0$.  The prior
$z<8/3$ excludes the larger root branch, so
\begin{equation}\label{or:T-ratio245}
                           z<49/20,\qquad y_0<49/20.
\end{equation}
Use the nested-count procedure with $H\ge90$, $J=(21/4)u$ and
$49/20$ in place of $9/4$ in \eqref{or:T-Xcap}.
It excludes $26$ closed cells of the second cover.  The remaining
nine cells use the already proved $h<105$ and the following pair
estimate.

\begin{lemma}[Two actual sources]\label{or:pair-source}
Suppose $T=1$, $0<\rho\le r<1$, $R>0$, and $0<m<1$.
Suppose two distinct actual zeros other than $\beta$ have
displacements $v_\pm=u_\pm-\beta$ satisfying
$|v_\pm|<R$, $|(v_++v_-)/2|<m$.  Suppose also that their individual
insets belong to $[3/4,6/5]$.  For a test
$s=X+\delta$ with $\delta>0$, put
\[
 I_\Sigma=2(1-r)(1-m)-rR^2,\qquad
 D_\Sigma=2(1+\delta^2)+4\sqrt{1+\delta^2}\,m+2R^2.
\]
Assume $I_\Sigma>0$.  Then
\begin{equation}\label{or:pair-lower}
        E>\frac{1-r}{1+\delta^2}
                    +\frac{(I_\Sigma/2+49/50)^2}{D_\Sigma}.
\end{equation}
\end{lemma}
\begin{proof}
Rightmostness gives $\Re v_\pm\le0$.  Expanding the inset,
\[
 I(\beta+v)=I(\beta)-(1-\rho)\Im v
                             -\rho X\Re v-\frac\rho2|v|^2.
\]
Since $I(\beta)\ge1-\rho$ and $|v_++v_-|<2m$,
\[
 I(u_+)+I(u_-)
 >2(1-\rho)(1-m)-\rho R^2
 \ge2(1-r)(1-m)-rR^2=I_\Sigma>0.
\]
The second inequality uses $2(1-m)+R^2>0$.
The squared distances satisfy
\[
 \begin{split}
 |s-u_+|^2+|s-u_-|^2
 &=2(1+\delta^2)
   -2\Re\bigl((\delta+i)\overline{v_++v_-}\bigr)
   +|v_+|^2+|v_-|^2\\
 &<D_\Sigma.
 \end{split}
\]
On $[3/4,6/5]$ one has $\sqrt t\ge t/2+49/100$.
Indeed $t-(t/2+49/100)^2$ is concave and positive at both endpoints.
Cauchy--Schwarz gives
\[
 \frac{I(u_+)}{|s-u_+|^2}+\frac{I(u_-)}{|s-u_-|^2}
 \ge\frac{(\sqrt{I(u_+)}+\sqrt{I(u_-)})^2}
              {|s-u_+|^2+|s-u_-|^2}
 >\frac{(I_\Sigma/2+49/50)^2}{D_\Sigma}.
\]
Adding the contribution of $\beta$ proves \eqref{or:pair-lower}.
\end{proof}

By \eqref{or:T-ratio245}, the field $h\ge90$ gives
$q<(49/20)/90<.05$, $\mu<1/90<.02$.
Proposition~\ref{lt:q05mu02} therefore supplies
\[
                          R=17/180,\qquad m=19/900.
\]
Here $R>0$, $0<m<1$, and, for each of the nine cells,
\[
 I_\Sigma\ge2(1-\rho_*)(1-m)-\rho_*R^2
             =\frac{367561059}{200000000}>0.
\]
The individual insets lie in $[3/4,6/5]$, since
\begin{equation}\label{or:pair-inset-gates}
          (1-r)(1-R)-rR^2/2>3/4,\qquad 1+R<6/5.
\end{equation}
For each pair cell verify \eqref{or:T-Abar} and
$B_{105}=2D_p(105f\overline A+1)/n_0<1$.  Its source upper is
\begin{equation}\label{or:T-pair-upper}
                  W<\frac{r(105f\overline A+1)}{1-B_{105}}.
\end{equation}
The lower bound \eqref{or:pair-lower} exceeds
\eqref{or:T-pair-upper} in all nine cells.  This proves $h<90$.

\begin{table}[ht]
\centering
\caption{Closed covers for the higher $T=1$ strip.  Each cover uses
the three scale slabs $[.01,.025]$, $[.025,.045]$, and
$[.045,.061002]$.}\label{or:T-cover}
\begin{tabular}{@{}lccccc@{}}
\toprule
field assumption&moment interval&closed cells&count cells&pair cells&contours\\
\midrule
$h\ge105$&$[0,9/4]$&28&28&0&1950\\
$90\le h<105$&$[0,49/20]$&35&26&9&1881\\
\bottomrule
\end{tabular}
\end{table}

\subsection{The higher $T=1$ strip: the terminal pairs}

For $\delta=.71$ use $A=(1-\delta)^2$ when
$X\le(1-\delta)^2$, and $A=X-\kappa_\delta(X)y_0$ otherwise.
The source coefficient satisfies
\begin{equation}\label{or:Ay-terminal}
                             y_0 A\le11/20.
\end{equation}
Here is the whole-domain verification.  Put $m_0=(1-\delta)^2$.
For $X\le m_0$, use $3m_0<11/20$.  On $[m_0,.35]$ six times
the endpoint value of $\kappa_\delta/X$ is less than one, so
$y_0(X-\kappa_\delta y_0)$ increases up to $y_0=3$; certify
$11/20-3(X-3\kappa_\delta)>0$ on this full interval.
Above $.35$, the active coefficient satisfies
$\kappa_\delta\ge(5/11)X^2$, and the inactive coefficient is stronger.
Completion of the square then gives \eqref{or:Ay-terminal}.

If $|\Im S|<J_0y_0$, equation~\eqref{or:radial} implies
\[
                   \rho(N-2h)<2y_0(h+J_0).
\]
For $h_L\le h\le h_U<N_0/2$, we have $N-2h>0$, and hence
\[
 \rho hA<\frac{2y_0A\,h(h+J_0)}{N-2h}
       \le\frac{2(11/20)h_U(h_U+J_0)}{N_0-2h_U}.
\]
Use $f=(1-\rho_*^2(1.71)^2/2)^{-1}$,
$D_p=1+(\rho_*1.71/2)^2$, and
$B_U=2D_p(fh_U+1)/n_0<1$.  The resulting finite source upper is
\begin{equation}\label{or:T-terminal-upper}
 W<\frac{f\,2(11/20)h_U(h_U+J_0)/(N_0-2h_U)+\rho_*}{1-B_U}.
\end{equation}
The following two intervals cover $45<h<90$, including their seam:
\begin{equation}\label{or:T-terminal-table}
\begin{array}{c|c|c|c|c}
 \text{field interval}&\text{companion estimate}&Z&M&J_0\\ \hline
 (45,70]&\eqref{lt:Q057}&8917/1000&2109/1000&37/5\\
 {}[70,90)&\ref{lt:q05mu02}&17/2&19/10&7
\end{array}
\end{equation}
The second row applies because $q<3/70<.05$ and $\mu<1/70<.02$;
the first applies by \eqref{or:T-Q057-entry}.
In each row put $R=Z/h_L$, $m=M/h_L$.
In both rows $R>0$, $0<m<1$, and the inset gates
\eqref{or:pair-inset-gates} hold with $r=\rho_*$.
The corresponding values of $I_\Sigma$ are
\[
 I_\Sigma=
 \begin{cases}
  \dfrac{201103346909179}{112500000000000}>0,&(45,70],\\[3pt]
  \dfrac{715839963}{392000000}>0,&[70,90).
 \end{cases}
\]
Lemma~\ref{or:pair-source}, using
$\sqrt{1+.71^2}<123/100$, gives the lower
\begin{equation}\label{or:T-terminal-lower}
 \frac{1-\rho_*}{1+.71^2}+
 \frac{\{(1-\rho_*)(1-m)-\rho_*R^2/2+49/50\}^2}
 {2(1+.71^2)+4(123/100)m+2R^2}.
\end{equation}
It exceeds \eqref{or:T-terminal-upper} by more than $.6165$ in the
first row and more than $.0527$ in the second.  These rational
comparisons hold throughout the stated field intervals.
They contradict $W\ge E$ and exclude
\eqref{or:T-domain}.

\subsection{The higher strip with $X=1$: a lower field bound}

Assume now
\begin{equation}\label{or:X-domain}
             1/100\le\rho\le\rho_*,\quad X=1,\quad0<T\le1.
\end{equation}
In this case $d>1/2$, $0<y_0<2$, and
$y_0<2\Re(S/h)$ by \eqref{or:horizontal}.  We first prove
\begin{equation}\label{or:X-h50}
                                   h>50.
\end{equation}

If $T\le1/8$, apply \eqref{or:moving-field} with
$r_+=\rho_*$, $t_+=1/8$.  The exact expression $H_m$ exceeds $54$,
so $h>50$ at positive inset.  At zero inset,
$\rho=2T/(1+T^2)$ and
$\rho_*<2(1/8)/(1+(1/8)^2)$ imply $T<1/8$ strictly.
The finite homothety \eqref{or:homothety} can therefore be chosen
with $T'<1/8$.  It gives the same strict bound for the unchanged $h$.

Suppose $1/8\le T\le1$ and, for contradiction, $h\le50$.
The unphased radial identity gives
$N\rho/2<h(1+Ty_0)<50(1+2T)$, hence
\[
                           \rho<300/N\le R_*=6/125.
\]
Use $\delta=3/10$ and the polar constants
\[
 f=(1-R_*^2(1.3)^2/2)^{-1},\quad D_p=1+(R_*1.3/2)^2,
                 B_{50}=2D_p(50f+1)/n_0<1.
\]
The source has the whole-interval lower weight
\[
 b_0=\min_{t\in\{1/8,1\}}
                 \frac{t-R_*(1+t^2)/2}{t^2+.09}>0.
\]
Write $S/h=c_0+is_0$, $c_0>0$, and set
$k=\kappa_1(1,.3)=13719/34880$.  The critical comparison gives
$C<f\rho h(c_0-ky_0)$, and $c_0-ky_0>0$ follows from $C\ge0$
and its strictness.  Under $h\le50$, the denominator $1-B_{50}$
is positive, so the source inequality gives
\begin{equation}\label{or:X-Rayleigh-source}
 \rho h(c_0-ky_0)>K:=\frac{(1-B_{50})b_0-R_*}{f}
                 =\frac{1279084413230726}{1580047607421875}>0.
\end{equation}

Put $v_0=y_0/c_0\in(0,2)$, so $c_0-ky_0=c_0(1-kv_0)$.
Equation~\eqref{or:radial}, using $I(\beta)\le T\le1$, gives
\[
 \frac{N\rho}{2}<h\{c_0(R_*+v_0)+\sqrt{1-c_0^2}\}.
\]
Multiply by \eqref{or:X-Rayleigh-source}.
The product in the denominator is the Rayleigh form, at the unit
vector $(c_0,\sqrt{1-c_0^2})$, of
\[
             (1-kv_0)
             \begin{pmatrix}R_*+v_0&1/2\\1/2&0\end{pmatrix}.
\]
Its largest eigenvalue is less than $M=4/5$ for every
$0\le v_0\le2$.  Indeed the determinant of $M$ times the identity
minus this matrix is $a_2v_0^2+a_1v_0+a_0$, where
\[
 a_2=Mk-k^2/4>0,\quad
 a_1=-M(1-kR_*)+k/2,\quad a_0=M^2-MR_*-1/4,
 \qquad 4a_2a_0-a_1^2>0.
\]
This quadratic is positive on the whole real line; the other principal
minor is $M>0$.  Also $1-2k>0$.  Thus
\[
 h^2>\frac{KN_0}{2M}
      =\frac{639542206615363}{202246093750}>2500,
\]
contradicting $h\le50$.  This proves \eqref{or:X-h50}.
Consequently $q<2/h<.04$ and $\mu<2/h<.04$.
Proposition~\ref{lt:q04mu04} gives
\begin{equation}\label{or:X-low-phase}
                              |\Im S|<(15/2)y_0.
\end{equation}
The homogeneous version includes all positive $q$ in this rectangle:
the radius is $(45/4)q$, and the center bound
$Z=7.706829806638$ satisfies $Z+.45<(13/5)\pi$.

\subsection{The higher $X=1$ strip: field bounds on closed cells}

For a cell
\begin{equation}\label{or:X-cell}
       \rho\in[p,r],\qquad y_0\in[l,u],\qquad T\in[t_-,t_+],
\end{equation}
use a phase bound $J=\gamma u$, with $\gamma=4$ on a
high-field branch and $\gamma=15/2$ on the low-field branch.
For a test $\delta$, put
\[
 k=\kappa_1(t_+,\delta),\quad \overline A=1-kl,
 b_0=\omega(0,0;\delta,t_-,t_+,r),
\]
and use \eqref{or:polar-constants} with $s_+=1+\delta$.
Where $b_0>0$ the direct source gives
\begin{equation}\label{or:X-K}
 \rho h(1-ky_0)>K,\qquad
 K=\frac1f\min_{z\in\{p,r\}}
          \left\{\frac{n_0z b_0}{n_0z+2D_pb_0}-z\right\}>0.
\end{equation}
The cell conditions include positivity of $\overline A$, the polar
denominator, and $K$.  Set
\[
 G_0=\max_{l\le y_0\le u}(1-ky_0)(1+t_+y_0),\qquad
 G_1=\max_{l\le y_0\le u}y_0(1-ky_0).
\]
Each is a concave quadratic maximum, evaluated at both endpoints
and its vertex when the vertex belongs to the closed interval.
Equations~\eqref{or:X-K} and \eqref{or:radial-unphased} imply
\begin{equation}\label{or:X-source-fields}
             h>\frac{K}{r\overline A},\qquad
             h^2>\frac{KN_0}{2G_0}.
\end{equation}
For the low-field phase \eqref{or:X-low-phase}, the exact radial
identity also gives
$\rho(N-2h)<2y_0(Th+15/2)$.
Multiplying by the same positive factor $h(1-ky_0)$ gives
\[
          K(N-2h)<2y_0(1-ky_0)h(Th+15/2),\qquad h<N/2.
\]
Thus a rational $H<N_0/2$ is a valid strict lower field if
\begin{equation}\label{or:X-phase-field}
                   K(N_0-2H)\ge2G_1H(t_+H+15/2).
\end{equation}
Indeed a contrary $h\le H$ makes its left side no smaller, and its
right side no larger, than those in \eqref{or:X-phase-field}.

Besides these source bounds, every cell has the lower field bounds
\begin{equation}\label{or:X-radial-fields}
 \frac{N_0p}{2(1+t_+u)},\qquad
 \frac{N_0p/2-J}{r+t_+u},\qquad H_m
\end{equation}
from \eqref{or:radial-unphased}, \eqref{or:radial-phase}, and
\eqref{or:moving-field}.  At zero inset the last is non-strict,
which suffices for a choice $H\le H_m$.  The moving test supplies
$H_m$; the final moment calculation includes $\beta$ once.

For every inverse row choose a rational $m>0$ satisfying
\begin{equation}\label{or:X-m-cell}
              m^2\{1/4+\max(t_--1/125,0)^2\}\ge1.
\end{equation}
By \eqref{or:gap-X} this bounds $1/d$ on the whole cell.  Use
$Y=u+10^{-9}$ in \eqref{or:inverse-parameters}, and use
\eqref{or:automatic-contour}.  The rational enlargement of $u$
ensures the required strict derivative inequalities also at a closed
moment endpoint.  All rows satisfy \eqref{or:contour-gates} and
\eqref{or:odd-gate}.  Widths, slopes and depths are computed by
\eqref{or:width}, \eqref{or:counts}, and \eqref{or:depth}.
The physical nesting conditions are $R_j<R_{j+1}<2$ and increasing
slopes; running maxima of the widths remain at most the corresponding
radii.  The regions therefore have $|\Re u|\le1$ as required by
\eqref{or:weight}.

If a depth lower bound exceeds $t_+$, the cell is impossible.  This
positive depth condition forces the height budget above $4\pi$,
so Lemma~\ref{or:narrow-strip} supplies such a root.
Otherwise compute the nested source lower in \eqref{or:Abel-source}.
For an unbounded high-field cell use the correction
$\min(0,b_0-3w_1)/H$; for a low-field cell use
\begin{equation}\label{or:X-signed-correction}
 \begin{cases}
  (b_0-3w_1)/H,&b_0-3w_1<0,\\
  (b_0-3w_1)/600,&b_0-3w_1\ge0.
 \end{cases}
\end{equation}
The latter form is used only after the ceiling $h<600$ is proved.

\subsubsection{The high-field cover}

Suppose $h\ge600$.  The bounds $q<2/600<7/625$ and
$\mu<2/600<1/200$ permit
Proposition~\ref{lt:fine200}, giving $|\Im S|<4y_0$.
On each cell take $J=4u$ and $H$ below at least one of $600$ and
the three lower fields in \eqref{or:X-radial-fields}.
For $t_->0$ put
\begin{equation}\label{or:X-high-radial}
       B_r=\frac{r/2+J/N_0}{t_-(1-rt_-/2)}.
\end{equation}
A positive $U$ bounds $h/n$ if $U\ge u$, $U^2\ge B_r$, or
$Ul\ge B_r$ for $l>0$.  Apply \eqref{or:HB-cell} with
$\overline A=1-\kappa_1(t_+,\delta)l$.  The high cover has
$61$ closed cells covering $[0,2]\times[0,1]$ in $(y_0,T)$ for
the entire scale interval $[.01,.061002]$.  Each cell has either
an actual-root depth contradiction, or a strict source lower bound
greater than \eqref{or:HB-cell}.  Cells with $t_-=0$ use the
depth alternative.  It follows that $h<600$.

\subsubsection{The low-field cover}

Use \eqref{or:X-h50} and \eqref{or:X-low-phase}.
The six closed scale slabs are
\begin{equation}\label{or:X-slabs}
 [.01,.013],\ [.013,.017],\ [.017,.023],\ [.023,.033],\
 [.033,.045],\ [.045,.061002].
\end{equation}
On each cell take $J=(15/2)u$ and choose $H>J$ below at least one
of: $50$; a bound in \eqref{or:X-radial-fields}; a bound in
\eqref{or:X-source-fields}; or a value certified by
\eqref{or:X-phase-field}.  All source tests used for these choices
satisfy $1/200\le\delta\le3/4$ and
$(r(1+\delta)/2)^2<\delta$.
If $H\ge600$, the cell contradicts the preceding ceiling.
Otherwise it has either a depth contradiction or a nested
moment contradiction.  For the latter the upper denominator is
\begin{equation}\label{or:X-low-denominator}
                   B=\frac{2D_p(600f\overline A+1)}{n_0}<1.
\end{equation}
Thus every low-field source upper uses the bound $h<600$.

The six slabs have respectively $36,38,41,38,41,46$ closed cells.
Together with the $61$ high cells this makes $301$ cells.  Their
three contradiction types are $196$ moment comparisons, $57$ actual
root depth comparisons, and $48$ field comparisons.  There are
$14891$ complete inverse contours.  With the rational rounding
conventions below, the smallest moment reserve is
\begin{equation}\label{or:X-reserve}
       \frac{84739835370681964317289}
            {240500000000000000000000000}>.0003523485.
\end{equation}
Every external endpoint and every internal depth, moment, and scale
seam belongs to the closed cover.  This excludes \eqref{or:X-domain}.

\subsection{Finite verification and the entire degree ray}
\label{or:finite-verification}

We finish by specifying the finite verification behind the interval
comparisons above.  The electronic appendix contains the exact
rational parameter files and executable verifiers.  The files specify
cells, tests, radii and trims; the following procedure recomputes
their bounds.

\begin{enumerate}
\item Every integer, fraction and finite decimal is read as an exact
  rational number.  The base rows are precisely
  Table~\ref{or:base-contours}.  The higher $T$ rows are the two files
  \cert{T6251_H105_PROPOSAL.json} and
  \cert{T6251_H90_PROPOSAL.json}; each record specifies
  $(p,r,l,u,H,\delta,\overline A,U,X_+,J)$ and either the list
  $(R,\nu,W_0)$ or the pair alternative.  The higher $X$ rows are
  \cert{X6251_high_0.01_0.061002.json} and the six
  \cert{X6251_MOVING_low_} files indexed by the consecutive endpoints
  of \eqref{or:X-slabs}.  Each record specifies $(l,u,t_-,t_+,H)$,
  its contradiction type, its source tests, and its contour choices
  $(R,\nu,W_0,m)$.
\item First verify the declared parameter domains, the inverse-Cayley
  condition, every local entrance, and the asserted lower
  field using the formulae preceding the relevant cover.  For $T$
  cells verify \eqref{or:T-Xcap}, \eqref{or:T-Abar}, and the
  radial upper.  For $X$ cells verify \eqref{or:X-m-cell} and every
  source positivity and quadratic maximum used for its field.
\item For every contour recompute $\lambda,e_R,\widehat R,B_R$,
  then all of \eqref{or:contour-gates} and the appropriate horizontal
  and height inequalities.  In the higher strips use
  \eqref{or:automatic-contour}.  The displayed Taylor reciprocal is
  strictly greater than $e^{-W_0}$.  Exponential upper bounds in the
  quartic gates use, for $0\le x<32$,
  \[
    e^x\le\sum_{j=0}^{30}\frac{x^j}{j!}
           +\frac{x^{31}}{31!(1-x/32)}.
  \]
  The geometric tail proves the inequality for every omitted order.
\item In the higher covers, round slopes and weights down, and real
  widths up, to the rational lattice $10^{-12}\mathbb Z$.
  Use $7/22<1/\pi$ for slopes and $\pi<22/7$ for depth losses.
  For every $X$ row also verify $Hs_j>1$ after rounding its slope.
  Square roots needed for radial uppers or pair denominators are
  rounded up to $10^{-16}\mathbb Z$, checking the squared enclosure
  by integer arithmetic.  Take running maxima of the widths and
  future maxima of the weights, verify their nesting, and compute
  the Abel sums exactly with the signed correction
  \eqref{or:X-signed-correction} when appropriate.  The $X$ source
  upper is rounded up to $10^{-12}\mathbb Z$ before subtraction.
\item Each cell using nested counts must have $0<B<1$ and a positive
  lower-minus-upper reserve.  Each depth cell must have its rounded lower depth strictly
  greater than $t_+$.  Each low-field exclusion must have $H\ge600$.
  Each pair cell must satisfy $0<r<1$, $R>0$, $0<m<1$, $\delta>0$, $I_\Sigma>0$,
  both inset gates, both endpoint chord inequalities, its squared
  radical enclosure, and the strict
  comparison \eqref{or:pair-lower}--\eqref{or:T-pair-upper}.
\item Covering is checked as an assertion about closed rectangles.
  For each $X$ file, sort all depth endpoints.  At every endpoint
  and one point in each intervening open interval, collect the
  $[l,u]$ intervals of all cells whose closed depth interval contains
  the point.  Their sorted union must cover $[0,2]$ without a gap.
  Membership is constant between adjacent depth endpoints, so this
  proves coverage of the whole rectangle.  Check exact adjacency
  and the exterior endpoints of the six scale slabs.  For each
  $T$ cover, sort all moment endpoints and check that the active
  scale intervals cover $[.01,.061002]$ on every intervening open
  fiber; the closed cells include their endpoint fibers as well.
\end{enumerate}

Here is a finite positivity method for all rational functions used in
the $T$ entrances, source coefficients and terminal bound.  Reduce the
function to a polynomial numerator and denominator over $\mathbb Q$.
Verify nonzero endpoint values, use the Euclidean signed-remainder
sequence (Sturm's rule) to show that neither has a root in the closed
interval, and check that their signs agree at one endpoint.  At an
explicit vanishing factor such as $X^2$ or $1-X$, remove that factor
symbolically and check its endpoint separately.  This proves every
one-variable positivity statement in
\eqref{or:T-first-square}--\eqref{or:T-q-low-poly},
\eqref{or:T-Abar}, and \eqref{or:Ay-terminal} on its specified
interval.  The base source-compensation polynomial is instead
expanded in the Bernstein basis of its specified interval; all its
coefficients are positive.  To make this test explicit, if
$P(a+(b-a)t)=\sum_{j=0}^d c_jt^j$, its Bernstein coefficients are
\[
                   b_k=\sum_{j=0}^k c_j
                                   \frac{\binom{k}{j}}{\binom{d}{j}}.
\]
Nonnegative basis functions whose sum is one then prove positivity
throughout $[a,b]$.

For the two rectangle inequalities
\eqref{or:T-D51-poly}--\eqref{or:T-h45-poly}, collect the exact
coefficients as $a(X)y_0^2+b(X)y_0+c(X)$.  Subdivide the closed
$X$ interval dyadically.  On each leaf the same one-variable procedure
verifies the positive values at $y_0=0,3$ and one of the following
sufficient conditions: $b>0$; $6a+b<0$; $a<0$; or $a>0$ and
$4ac-b^2>0$.  The first condition suffices because if $a\ge0$ the
quadratic is increasing, whereas if $a<0$ its two positive endpoint
values suffice by concavity.  The second has the analogous decreasing
or concave justification.  The third and fourth follow from concavity
and completion of the square, respectively.
The coefficients $a,b$ have a factor $X$; checking their quotients
and $c(0)>0$ handles the entire $X=0$ edge.  Both children of each
subdivision are retained.  The resulting sorted leaves have exact
adjacent endpoints $0$ and $1$, proving the entire rectangle.
This proves the two coupled source and matrix inequalities over their
full rectangles.

These procedures verify the $55$ base contours and all base
inequalities; the higher $T$ procedure verifies $63$ closed cells,
$3831$ contours and $9$ pair cells; the higher $X$ procedure verifies
the $301$ closed cells and reserve \eqref{or:X-reserve}.
The least higher-$T$ reserve among cells using nested counts exceeds
$.0001411$.

Finally every argument extends from $N_0=6250$ to the entire ray
$N\ge N_0$.  All local companion, inverse-log and crescent estimates
are independent of the degree.  The direct source in
\eqref{or:direct-K} increases with $n$; radial lower fields increase
with $N$; the terms $J/N$ and $1/n$ in upper bounds decrease.
The coupled inequality \eqref{or:X-phase-field} strengthens with
$N$, and the two terminal uppers decrease with $N$.
The fixed ceilings and closed parameter covers therefore
work for every $n\ge6251$.  The base range and the two higher strips
exhaust \eqref{or:normalization}, proving
Proposition~\ref{or:theorem}.

\section{Small positive witnesses}\label{sec:small}

Throughout this section $N\ge N_0=6250$, and
\[
 \rho_0=\frac{30501}{500000},\qquad A_0=\frac{\rho_0}{100}.
\]
We work in the strict-hole configuration of
Lemma~\ref{pr:edge-reduction}, with $0<a\le1/10$ and an actual
simple rightmost zero $b=x+iy$, $x\ge a$, $y\ge0$.
Put $r=|b|$, $p_b=1-r^2$, $\tau=1-y$.
Use $d,S,H,Q$ as in \eqref{pr:field-notation}, and
$D,q,\mu$ as in \eqref{pr:normalized-moments}.

\subsection{The first field bound}

\begin{lemma}\label{sm:normalized-D25}
Let $F(u)=p(i+\rho u)$ be a normalized strict-hole configuration
with an actual simple rightmost zero $\beta=X-i$, where
$\alpha=a/\rho$, $0<\alpha\le X\le1$,
$1/100\le\rho\le9/100$, and $n\ge6251$.
Let $v_j$ list all zeros of $F'$, counted with multiplicity.
Write $h=|F''(\beta)/F'(\beta)|$ and
$d_u=\min_j|\beta-v_j|$. Then $d_u>1$ and $D=d_uh>25$.
\end{lemma}
\begin{proof}
Put
\[
 R=9/100,\qquad Q_u=\sum_j|\beta-v_j|^{-2},\qquad Y=Q_u/h.
\]
The critical domain is
$r_j<\alpha/2$ and
$\rho|v_j|^2/2\le t_j<\rho|v_j-\alpha|^2/2$ for
$v_j=r_j-it_j$. Lemma~\ref{lt:crescent}, scaled back to these
coordinates, gives
\[
 d_u^2>(X-\alpha/2)^2+(1-\eta)^2,\qquad
 \eta=\frac{\rho\alpha^2}{4(1+\sqrt{1-\rho^2\alpha^2/4})}
       \le R\alpha^2/4.
\]
Since $0\le\eta\le R\alpha^2/4<1$, it follows that
\[
 d_u^2>G(\alpha):=(X-\alpha/2)^2+(1-R\alpha^2/4)^2.
\]
This polynomial decreases with $\alpha$ on $0\le\alpha\le X$:
its derivative is
$-X+(1/2-R)\alpha+(R^2/4)\alpha^3<0$.
Hence $d_u^2>1+(1/4-R/2)X^2+(R^2/16)X^4>1$.

Apply the finite polar estimate in Lemma~\ref{or:HB} at
$s=X+1/2$. The inverse Cayley test is strictly to the right of
the original zeros, since $(R(1+1/2)/2)^2<1/2$.
The source root contributes at least
$b_0=4(1-R)/5$. By the full critical comparison
\eqref{or:critical-comparison-formula} and
\eqref{or:small-X-comparison},
\[
 C_{\rm real}<\rho h\mathcal A(X,Y),\qquad
 \mathcal A(X,Y)=
 \begin{cases}1/4,&0<X\le1/4,\\
 X-\kappa(X)Y,&1/4\le X\le1,
 \end{cases}
\]
\[
 \kappa(X)=\frac{X(4X-1)(4X+9)}{40(X+1)}.
\]
Use $f=[1-(9/8)R^2]^{-1}$ and $D_p=1+(3R/4)^2$ in
\eqref{or:parallel-source}. The function $uv/(u+v)$ is increasing
in both variables, so
\begin{equation}\label{sm:normalized-source}
 \rho h\mathcal A(X,Y)>K_0,\qquad
 K_0=\min_{\rho\in\{1/100,R\}}
 \frac1f\left(\frac{6251\rho b_0}{6251\rho+2D_pb_0}-\rho\right)
 =\frac{239012578698741}{379195048000000}>0.
\end{equation}
The expression minimized here is concave in $\rho$
and increases with the degree, as shown by differentiating it in
\eqref{or:direct-K}. Thus the bound holds throughout the scale interval
for every $n\ge6251$. The unphased radial inequality is
$N\rho/2<h(1+Y)$.

If $X\le7/25$, then $\mathcal A\le7/25$, whence
$h>K_0/[R(7/25)]>25$.
If $X\ge7/25$, completion of the square gives
$\mathcal A(X,Y)(1+Y)<31/10$ for every real $Y$.
For an exact proof, put $N_k=X(4X-1)(4X+9)$ and $D_k=40(X+1)$.
The degree-six polynomial
\[
 (62/5)N_kD_k-(XD_k+N_k)^2
\]
has the following positive Bernstein coefficients on $[7/25,1]$:
\[
 \frac{120508031}{244140625},\quad
 \frac{1411226329}{1953125},\quad
 \frac{185744659}{78125},\quad
 \frac{416205433}{78125},\quad
 \frac{30945547}{3125},\quad
 \frac{407837}{25},\quad24527.
\]
Multiplying the two inequalities for the same $X,Y$ therefore gives
$h^2>K_0N/(2(31/10))>625$. Since $d_u>1$, both branches give
$D>25$.
\end{proof}

\begin{lemma}\label{sm:first-D25}
The implication $D>25$ holds on each of the two regions
\[
 \max(x,1-y)>\rho_0,\qquad |b|\le47/50.
\]
Both regions also satisfy $D>70$.
\end{lemma}
\begin{proof}
If $x\ge\rho_0$, Lemma~\ref{lt:direct-D70} gives $D>70$.
We treat $x\le\rho_0$.

Return to physical coordinates. Lemma~\ref{lt:crescent} gives
$d>\tau$. If $1/100\le\tau\le9/100$ and $x\le\tau$, the
preceding lemma applies with $\rho=\tau$, $X=x/\tau$,
$\alpha=a/\tau$; its $D$ is exactly the physical $D$.
For $\tau\ge9/100$, use the matrix envelope in
Proposition~\ref{mf:envelope}. Writing $c=\Re S/H$,
$s=|\Im S|/H$, the rational bound $\pi>333/106$ gives
$\gamma<241/1000$, and its constants satisfy
\[
 Qd^2<21/5+(11/2)Dc+(241/1000)D^2c^2-D^2s^2/2.
\]
Retaining the corresponding negative coefficient $741/1000$ and
completing the square in the radial identity yields
\[
 N<\frac{2Dx}{d}+\frac{p_b}{d^2}
       \left(21/5+(11/2)D+(241/1000)D^2\right)
       +\frac{y^2}{(741/1000)p_b}.
\]
On the whole $x\le\rho_0$, $\tau\ge R=9/100$ range,
\[
 x/d<\rho_0/R,\qquad p_b/d^2<(2-R)/R,\qquad
 y^2/p_b\le(1-R)/(2R).
\]
For the last bound, $x<\tau$ gives $p_b>2\tau(1-\tau)$;
the case $y=0$ is immediate. Under $D\le25$ the preceding upper
bound for $N$ is $640685569/102600<6250$, a contradiction.

Outside the origin box, the remaining small-$x$ range has
$\tau>\rho_0>x$. For $r\le47/50$ one has $\tau\ge.06$.
If $.06\le\tau\le\rho_0$, then
\[
 (.94)^2<\rho_0^2+(1-\rho_0)^2,\qquad
 x^2\le(.94)^2-(1-\tau)^2<\tau^2,
\]
because $\tau^2+(1-\tau)^2$ decreases on this interval.
Thus every case enters either the normalized source or the
large-$\tau$ matrix argument. This proves $D>25$ on both regions.

Proposition~\ref{mf:D25} therefore gives $Qd/H<3$.
Outside the origin box and with $x\le\rho_0$, the radial identity
first gives $D>6250/(2/\rho_0+96)>48$. The matrix envelope excludes
$|s|\ge2/3$, since $c<3/4$, $c^2\le5/9$ would give
\[
 \frac{21/5}{48^2}+\frac{(11/2)(3/4)}{48}
       +\frac{241}{1000}\frac59-\frac29<0.
\]
The phase-sensitive radial bound then gives
\[
 D>\frac{6250\rho_0}{2\rho_0+(4/3)(1-\rho_0)+3(2-\rho_0)}>53.
\]
On $r\le.94$ the initial bound is $D>6250(.06)/8=46.875$.
The matrix envelope excludes $|s|\ge.7$, using $c<.715$ and
$c^2\le.51$. Consequently
\[
 D>\frac{6250(.06)}{2\rho_0+1.4(.94)+3(1.94)}>51.
\]
Both regions satisfy $q<.06$ and $\mu<.02$, so the first
weight-eight row of Proposition~\ref{lt:W8} applies.
Its phase bound $|\Im S|<8Q/H$ and the radial identity give respectively
\[
 D>\frac{6250-48/\rho_0}{98}
   =\frac{27771875}{498183},\qquad
 D>\frac{6250-48/.06}{2\rho_0/.06+97}
   =\frac{27250000}{495167}.
\]
Both exceed $5000/91$, so Proposition~\ref{lt:W8-field} applies.
Its field bound with constant $210$ gives
\[
 D>\frac{6250\rho_0-210}{11/5}=\frac{13701}{176}>70
\]
outside the origin box, and $D>[6250(.06)-210]/(11/5)=75>70$
on $r\le.94$. The geometric estimate covers the large-$x$ branch.
Every numerical comparison increases with $N$, so the full degree
ray is included.
\end{proof}

\subsection{Witnesses below $A_0$}

\begin{proposition}\label{sm:fine}
There is no strict hole with $0<a\le A_0$ and
$\max(x,1-y)>\rho_0$.
\end{proposition}
\begin{proof}
The crescent gives
\[
 d>\rho_0-a/2\ge d_0:=\rho_0-A_0/2.
\]
Lemma~\ref{sm:first-D25} first gives $D>70$.
Then \eqref{lt:field70}, followed in order by
Proposition~\ref{mf:D100}, \eqref{lt:field100}, and
Proposition~\ref{lt:medium139}, gives the strict successive bounds
\[
 D>\frac{6250d_0-596/5}{2(1+A_0)}>100,\qquad
 D>\frac{6250d_0-1899/20}{2(1+A_0)}>140,
\]
\[
 D>\frac{6250d_0-2637/32}{2(1+A_0)}>148.
\]
The last step uses $140>139$.
With $\kappa=11067/5000$, the exact inequality
\[
 (5/4)\kappa^2-(148-23/2)\kappa+296<0,\qquad 2<\kappa<9/4,
\]
and the smaller quadratic branch in Proposition~\ref{lt:medium139}
give $Qd/H<\kappa$.

For $R=td$ define
\[
 \Lambda=\frac{\kappa t}{1-t},\quad
 f=t-\kappa\left(t^2/2+t^3/[3(1-t)]\right),\quad
 B=\frac{\kappa(1-t)}{[1-(1+\kappa)t]^3}.
\]
Thus $H\widehat R=Df$, and the inverse lemma applies with
$\epsilon=(1001/1000)B/148$,
$g=(1001/1000)\nu Bf$, $p_*=(1+\Lambda)^{-1}$ and
$A_*=(1-\Lambda)^{-1}$.
The exact data \cert{FINE6251_INTEGER_DATA.json} give forty-one
triples $(j,m,W)$, with $t=j/200$ and $\nu=m/10000$.
For every row the finite predicates are
\begin{gather}
 0<t,\nu,\Lambda<1,\quad f>0,\quad
 W^2<(1-\nu^2)(148f)^2,\quad\nu(148)f>22/7,\notag\\
 \epsilon(1+e)+ge<1-A_*e,\quad
 \epsilon(1+e)+ge^2<p_*-e,\quad\epsilon+ge<p_*.
 \label{sm:profile-gates}
\end{gather}
Here, for $W=1,\ldots,10$, the rational bounds $e>e^{-W}$ are
\[
 1/2, 1/7, 1/20, 1/50, 1/140, 1/400,
 1/1000, 1/2900, 1/8000, 1/22000.
\]
Their validity follows from the positive exponential sum through
degree one hundred. The data also have increasing $t$ and increasing
$\nu f$. Equations~\eqref{sm:profile-gates} are exactly the full
inverse, corner, four-side and homotopy conditions of
Lemma~\ref{lt:inverse-log}; the resulting cumulative count is
$L_j>\nu_j(7/22)Df_j-2$.

Put $\nu_0f_0=0$ and
\[
 c=\frac{6250}{6251}\left\{
 \frac7{22}\sum_j(1-t_j)^2(\nu_jf_j-\nu_{j-1}f_{j-1})-\frac2{148}
 \right\}>0.
\]
Lemmas~\ref{lt:transport} and \ref{lt:abel}, with $A=A_0$, give
$d<Ja$, where $J=E_{A_0}/[(1-q_{A_0})c]$.
Exact substitution of the forty-one rows gives
$J+1/2<77<100$.
Thus $d<Ja<100a-a/2\le\rho_0-a/2$, contrary to the crescent
lower bound. The finite loss $2/148$ was charged once through
positive cumulative summation. The factor $N/(N+1)$ was bounded
below by $6250/6251$, so the proof holds for every $N\ge6250$.
\end{proof}

\subsection{The low and middle modulus ranges}

\begin{proposition}\label{sm:low-middle}
There is no strict hole with $A_0\le a\le1/10$ and $|b|\le47/50$.
\end{proposition}
\begin{proof}
Lemma~\ref{sm:first-D25} gives $D>70$. Since $d>.06$,
\eqref{lt:field70} gives $D>[6250(.06)-596/5]/(11/5)>100$.
Thus Proposition~\ref{mf:D100} and \eqref{lt:field100} apply.
The simultaneous crescent bound
$d^2>(.06)^2+a^2/4$ gives $D>139$: after positive squaring the
required inequality
\[
 6250\sqrt{(.06)^2+a^2/4}>1899/20+278(1+a)
\]
is exactly
\[
 9688341a^2-(1036801/5)a+613319/400>0.
\]
Its minimum over the entire real line is
$65688734375/155013456>0$.
Consequently Proposition~\ref{lt:medium139} supplies
\begin{equation}\label{sm:middle-input}
 Qd/H<9/4,\qquad |\Im S|<4Q/H,\qquad
 H>\frac{N-C_F/d}{2(1+a)},\qquad C_F=2637/32.
\end{equation}

Put $r_a=|b-a|$. If $r_a=0$, Proposition~\ref{ar:energy}
contradicts the strict hole. Otherwise use the adaptive
shift from Lemma~\ref{bc:blaschke}:
\[
 u=\min(r_a^2,1/2),\quad t=1-u,\quad
 \gamma=b-ta,\quad C=ta[2x-(1+t)a]\ge a^2u(1-u)>0.
\]
It gives
\begin{equation}\label{sm:product}
 -\log B_p(\gamma)+CQ/2<\log(8/C),\qquad
 B_p(\gamma)=\prod_{p(z_i)=0}
              \left|\frac{\gamma-z_i}{1-\overline{z_i}\gamma}\right|.
\end{equation}
A zero product gives a contradiction. Every unselected factor
is at most one.

\smallskip\noindent\emph{The range $r\le.11$.}
Take $d_*=.89$, $V=(9/4)/.89$ and $H_*=6250(.89)/(9/4)$,
using Proposition~\ref{ar:energy} and the prior ratio.
The complete inverse row $R=.1$, $W=8$, $\nu=.99$ satisfies
\eqref{lt:inverse-corner}--\eqref{lt:inverse-sides}. Its count
$K=\nu(7/22)H_*\widehat R-2$ obeys $K/6251>1/100$.
The count lower bound has the form $cN-2$, so its quotient by $n=N+1$
increases. Every counted root has adaptive factor at most
\[
 \frac{R+.1}{1-(.11+R)(.11)}<1/4<e^{-1}.
\]
Hence $-\log B_p>n/100$. Proposition~\ref{ar:variance} gives
\[
 d\le\sqrt{1-(1-.11^2)/n}<1-.49/n.
\]
If $r_a\le.4/n$, the strict hole instead gives
$d>1-r_a\ge1-.4/n$. Thus $r_a>.4/n$, and monotonicity of
$u(1-u)$ on $[0,1/2]$ gives
\[
 C>A_0^2(.4/n)^2[1-(.4/n)^2]>1/(20000000n^2).
\]
The last scalar inequality holds at $n=6251$ and strengthens as $n$
increases. Positive exponential sums give $e^{20}>8(20000000)$
and $e^{10}>6251$. Equation~\eqref{sm:product} would imply
$n/100<20+2\log n$, which fails at $6251$ and thereafter, since
the derivative of the difference is $1/100-2/n>0$.

\smallskip\noindent\emph{The range $.11<r\le.82$.}
Use
\[
 d_*=.18,\quad V=(9/4)/.18,\quad
 H_*=[6250-C_F/.18]/2.2,\quad R=11/250,\quad\nu=.99,\quad W=6.
\]
The full inverse predicates hold. If $z=b+v$ is counted,
$\Re v\le0$ by rightmostness. Therefore
$|\gamma-z|^2=|v+ta|^2\le R^2+.01$ and the denominator is
at least $1-(.82+R)(.82)>0$. The exact squared comparison bounds
the factor by $3/8$. Since $e^{.98}<8/3$, the count gives
$-\log B_p>.98K>21.5$.
Also $r_a>.11-a$ and
\[
 C>a^2(.11-a)^2[1-(.11-a)^2]>4\cdot10^{-9}.
\]
Indeed the concave function $a(.11-a)$ exceeds $1/15000$ at
both endpoints of $[A_0,.1]$, and the remaining factor exceeds
$.9$. The exact positive sum $e^{21.5}>8/(4\cdot10^{-9})$
contradicts \eqref{sm:product}.

\smallskip\noindent\emph{The range $.82\le r\le.94$.}
Here $r_a^2\ge(.82-.1)^2>1/2$, so the shift is
$\gamma=b-a/2$ and $C=ax-3a^2/4\ge a^2/4$.
We use twenty-three product rows, specified by rational inputs in
\cert{LOW_MIDDLE6251_PRODUCT_PROPOSAL.json}.
For each row $B_L\le r\le B_U$, $a_L\le a\le a_U$, $x\le X$,
choose the displayed $d_*,V,H_*$ from that data, and nested
physical radii $R_j$. Each full contour uses
\[
 \Lambda_j=\frac{VR_j}{1-R_j/d_*},\quad
 \widehat R_j=R_j-V\left(\frac{R_j^2}{2}
                         +\frac{R_j^3}{3(d_*-R_j)}\right),\quad
 B_j=\frac{V}{(1-R_j/d_*)^2(1-\Lambda_j)^3},
\]
with $\epsilon_j=(1001/1000)B_j/H_*$ and
$g_j=(1001/1000)\nu_jB_j\widehat R_j$.
All predicates \eqref{lt:inverse-corner}--\eqref{lt:inverse-sides}
are checked at $H_*$, together with
$\nu_jH_*\widehat R_j>22/7$.
Lemma~\ref{lt:half-disk-factor} supplies a whole-disk factor $f_j<1$.
The weights $w_j<-\log f_j$ are rational lower bounds from the
positive atanh series after dyadic scaling into $[1/2,1)$.
They decrease with radius. All count slopes increase, and positive
Abel summation \eqref{lt:abel-product} gives the budget
\[
 \mathcal B=H_*\sum_jw_j(s_j-s_{j-1})-2w_1,\qquad
 s_j=\nu_j(7/22)\widehat R_j,\quad s_0=0.
\]
For a positive whole-row lower bound $C_{\min}$ and a rational
target $T_*$, the final predicates are
\begin{equation}\label{sm:middle-budget}
 e^{T_*}>8/C_{\min},\qquad
 \mathcal B+6250C_{\min}/2>T_*.
\end{equation}
The energy term follows from Proposition~\ref{ar:energy} and
$C\ge C_{\min}$. Exponential bounds are finite positive sums or
\eqref{lt:exp-upper}; radical enclosures are accepted by exact squares.

For $.82\le r\le.925$ take $d_*=.075$, $V=(9/4)/d_*$,
$H_*<[6250-C_F/d_*]/[2(1+a_U)]$, $X=.925$,
$C_{\min}=a_L^2/4$. The seven closed witness intervals have endpoints
\[
 A_0, .003, .01, .02, .04, .06, .08, .1.
\]
The respective strict reserves in \eqref{sm:middle-budget} exceed
$3.03,4.54,4.41,1.91,1.45,1.79,2.99$.

For $.925\le r\le.94$, $x\ge\rho_0$, take $X=.94$ and choose
$d_*$ below at least one of
\[
 .06,\quad \rho_0-a_U/2,\quad \sqrt{(.06)^2+a_L^2/4}.
\]
Each displayed quantity is a crescent lower bound for $d$. Choose
$V\ge\min((9/4)/d_*,1/(\rho_0-a_U/2))$ and
$H_*<[6250-C_F/d_*]/[2(1+a_U)]$.
The lower bound $C_{\min}$ is the exact minimum of
$a\max(a,\rho_0)-3a^2/4$ on the full row. Below $\rho_0$
the function is concave; above it the function is $a^2/4$ and
increases. The eleven intervals have endpoints
\[
 A_0, .003, .01, .02, .03, .04, .05, .06, .07, .08, .09, .1.
\]
All eleven strict reserves exceed $.067$.

It remains to consider $.925\le r\le.94$, $x\le\rho_0$.
Put $Y=\tau Q/H$. If $Y\le5/4$, the phase and radial identity give
\begin{equation}\label{sm:low-moment-field}
 H>\frac{N-8(5/4)/.06}{2\rho_0+(2-.06)(5/4)}.
\end{equation}
This uses $\tau\ge.06$, $p_b/\tau\le2-\tau\le1.94$, and
$|\Im S|<4Q/H$.
Take $d_*=.06$, $V=(5/4)/.06$, an integer $H_*$ below
\eqref{sm:low-moment-field}, $X=\rho_0$ and $C_{\min}=a_L^2/4$.
The five intervals with endpoints
$A_0,.003,.01,.02,.04,\rho_0$ have reserves exceeding
$3.68,5.43,5.44,3.01,2.31$.
These twenty-three rows comprise 1807 nested complete contours;
together with the two earlier single contours, there are 1809.

Finally, in this small-$x$ region $\tau<.078$ since
$y^2\ge.925^2-\rho_0^2>(1-.078)^2$.
If $\tau\ge\rho_0$ then $x\le\tau$; if $.06\le\tau\le\rho_0$
the seam used in Lemma~\ref{sm:first-D25} again gives $x<\tau$.
The normalization $u=(z-i)/\tau$ therefore gives $T=1$,
$0<\alpha=a/\tau\le X=x/\tau\le1$, $d_u>1$ and
$Y=Q_u/h<9/4$. The only branch not covered by the preceding
product rows lies in
\[
 .06\le\tau\le.078,\qquad 0<\alpha\le X\le1,\qquad
 5/4\le Y\le9/4.
\]
The bounds $D>139$ and \eqref{sm:middle-input} supply the hypotheses
of Lemma~\ref{sm:T-moment-extension}. In normalized coordinates the
field bound reads $h>[N\tau-C_F]/[2(1+\tau)]$, since $d_u>1$
and $a\le\tau$. The lemma excludes this entire normalized range.

All modulus and witness seams are closed. Each field lower bound
is affine increasing in $N$, while every fixed count and energy
lower bound strengthens along the degree ray. The only logarithmic
degree budget, in the low range, was checked with its derivative.
Thus every configuration in the proposition is covered.
\end{proof}

\subsection{The high modulus range}

\begin{proposition}\label{sm:high}
There is no strict hole with $A_0\le a\le1/10$, $|b|>.94$,
outside the origin box $\max(x,1-y)\le\rho_0$.
\end{proposition}
\begin{proof}
If $x<\rho_0$, the exact comparisons
\[
 .94^2>(.0433)^2+(1-\rho_0)^2,\qquad
 .94^2-\rho_0^2>(1-.062)^2,\qquad .0433/.062>.69
\]
give $\rho_0<\tau<.062$ and $.69<x/\tau<1$.
The normalization $u=(z-i)/\tau$ therefore lies entirely in
Lemma~\ref{sm:T-highX-extension}, with $T=1$,
$\rho_0\le\tau\le.062$, $.69\le X\le1$.
We may henceforth assume $x\ge\rho_0$.

Lemma~\ref{lt:direct-D70} gives $D>70$ geometrically.
Proposition~\ref{lt:D70} then gives phase $K_0=21/4$ and the
field \eqref{lt:field70}. Since $|b-a|>.84$, the adaptive shift
is $\gamma=b-a/2$, and \eqref{sm:product} holds with
$C=ax-3a^2/4>0$.
Put $q_N=Q/N$. Cauchy--Schwarz and the radial identity give
\begin{equation}\label{sm:phase-product}
 1<2x\sqrt{q_N}+p_bq_N+2K_0Q/(NH).
\end{equation}

\smallskip\noindent\emph{The first witness cap.}
Suppose $a\ge a_0$ and $x\ge x_L>3a_0$. Concavity of
$a\mapsto ax-3a^2/4$ on $[a_0,x]$ gives
$C\ge a_0(x-3a_0/4)$: the other endpoint is $x^2/4$, which
is larger when $x>3a_0$.
Choose $L$ with $e^L>8/[a_0(x_L-3a_0/4)]$, and put
$s=3a_0/4$, $k=2L/(N_0a_0)$.
The product inequality gives $q_N<k/(x-s)$.
The reciprocal half-plane gives $Q/H<2/x$.
Thus \eqref{sm:phase-product} requires
\begin{equation}\label{sm:initial-cap-function}
 1<2x\sqrt{\frac{k}{x-s}}+
       (1-B_L^2)\frac{k}{x-s}+\frac{4K_0}{N_0x}
       \qquad(r\ge B_L).
\end{equation}
This expression is convex as a function of $v=\sqrt{x-s}$:
its terms are $2\sqrt k(v+s/v)$,
$(1-B_L^2)k/v^2$, and $4K_0/[N_0(v^2+s)]$.
The final term has nonnegative second derivative when $3v^2\ge s$.
It suffices to check the two endpoints of the full interval in $x$.
For
\[
 B_L=.94,\quad x_L=\rho_0,\quad x_U=1,\quad
 a_0=.013,\quad L=1879/200,\quad K_0=21/4,
\]
both endpoint values are strictly below one. These are rational
checks: the remainder after subtracting the last two terms is
positive, and its square exceeds $4x^2k/(x-s)$.
The logarithmic comparison uses a positive exponential sum.
We conclude $a<.013$.

The cap $a<.013$ gives $d>\rho_0-.0065=.054502$, so
\eqref{lt:field70} gives $D>100$.
Thus Proposition~\ref{mf:D100} and \eqref{lt:field100} apply.
For $r\ge.946$,
Lemma~\ref{lt:same-gap-field}, with $a_U=.013$, $x_0=\rho_0$,
$K_0=21/4$ and $p_U=1-.946^2$, gives
\[
 D>321306321875/1995240393>160.
\]
Proposition~\ref{lt:medium139} therefore supplies phase four,
field constant $C_F=2637/32$, and the smaller ratio branch $Y(D)$
throughout that modulus range.

\smallskip\noindent\emph{Eight successive Schur estimates.}
For a known witness cap $a\le a_U$ and modulus floor $B_L$,
Lemma~\ref{lt:same-gap-field} with phase four gives a whole-region
lower bound for $D$. Choose $D_0$ strictly below that value, then
choose $\kappa$ with
\[
 2<\kappa<9/4,\qquad
 (5/4)\kappa^2-(D_0-23/2)\kappa+2D_0<0.
\]
The smaller branch implies $Qd/H<\kappa$.
Use the complete Schur profiles as in
\eqref{sm:profile-gates}, replacing $148$ and $11067/5000$ by
$D_0,\kappa$. Put $\nu_0f_0=0$. If
\[
 c=\frac{6250}{6251}\left\{
 \frac7{22}\sum_j(1-t_j)^2(\nu_jf_j-\nu_{j-1}f_{j-1})
 -\frac2{D_0}\right\}>0,
\]
then, with $J=E_{a_U}/[(1-q_{a_U})c]$, Lemma~\ref{lt:abel}
gives $d<Ja$ and hence
$x<(J+1/2)a$. The complete sequence is
\begin{center}\small
\begin{tabular}{rrrrrrr}
\toprule
$B_L$&prior $a_U$&$D_0$&$\kappa$&$M$&upper bound for $J+1/2$&rows\\\midrule
$.946$&$.013$&$162$&$2.1928$&$75$&$74.260$&$42$\\
$.955$&$.01108$&$198$&$2.1545$&$67$&$66.009$&$44$\\
$.97$&$.00779$&$301$&$2.0985$&$57$&$56.441$&$49$\\
$.98$&$.00572$&$442$&$2.0659$&$52$&$51.715$&$52$\\
$.99$&$.00430$&$781$&$2.0367$&$49$&$48.244$&$54$\\
$.995$&$.00296$&$1256$&$2.0226$&$47$&$46.493$&$56$\\
$.997$&$.00216$&$1651$&$2.0172$&$46$&$45.798$&$57$\\
$.998$&$.00173$&$1951$&$2.0145$&$46$&$45.515$&$57$\\\bottomrule
\end{tabular}
\end{center}
All 411 triples $(j,m,W)$ are rational integer data in
\cert{HIGH6251_INTEGER_DATA.json}, with $t=j/200$,
$\nu=m/10000$. For each, the complete inverse, every side and
corner, the amplitude homotopy, positive odd height, and positive
Abel increment are verified by \eqref{sm:profile-gates}.
Thus the conclusion of each row is $x<Ma$ on its whole domain.

After each row we improve the witness cap, in the same order.
Suppose $a\ge a_0$ and $x<Ma$, and put $x_0=Ma_0$,
$\gamma_M=1/M-3/(4M^2)$. Concavity gives
\begin{equation}\label{sm:joint-C}
 C\ge
 \begin{cases}
 a_0(x-3a_0/4),&\rho_0\le x\le x_0,\\
 \gamma_Mx^2,&x\ge x_0.
 \end{cases}
\end{equation}
The two expressions agree at $x_0$. For the first branch use
$k=2L/(N_0a_0)$, $s=3a_0/4$, and for the second use
$k_q=2L/(N_0\gamma_M)$. The cap $a<a_U$ gives
the phase error $8/[N_0(x-a_U/2)]$ in
\eqref{sm:phase-product}. The first branch is convex in
$\sqrt{x-s}$ when $3(x-s)>s-a_U/2$; the second branch requires
\[
 1<2\sqrt{k_q}+(1-B_L^2)k_q/x^2+\frac8{N_0(x-a_U/2)},
\]
whose right side decreases with $x$.
Hence positive squared endpoint checks at $x=\rho_0,x_0$, together
with $e^L>8/[a_0(\rho_0-3a_0/4)]$, exclude $a\ge a_0$.
The rows are
\begin{center}\small
\begin{tabular}{rrrr}
\toprule
$B_L$&$M$&new cap $a_0$&$L$\\\midrule
$.946$&$75$&$.01108$&$9.528$\\
$.955$&$67$&$.00779$&$9.834$\\
$.97$&$57$&$.00572$&$10.115$\\
$.98$&$52$&$.00430$&$10.382$\\
$.99$&$49$&$.00296$&$10.738$\\
$.995$&$47$&$.00216$&$11.043$\\
$.997$&$46$&$.00173$&$11.26$\\
$.998$&$46$&$.00146$&$11.426$\\\bottomrule
\end{tabular}
\end{center}
The prior cap of each Schur row is exactly the conclusion of the
preceding cap row, and its modulus domain is contained in the
preceding one. Thus these estimates form a finite ordered chain.

\smallskip\noindent\emph{Root counts in physical disks.}
The closed initial domains are as follows. Each contains
all $x$ from $\rho_0$ to the stated upper endpoint, and all
$a\in[A_0,a_U]$ compatible with the configuration.
\begin{center}\small
\begin{tabular}{crrrr}
\toprule
modulus interval&$a_U$&$M$&upper $x$&closed leaves\\\midrule
$[.94,.946]$&$.013$&--&$1$&1\\
$[.946,.955]$&$.01108$&75&$.831$&1\\
$[.955,.97]$&$.00779$&67&$.52193$&1\\
$[.97,.98]$&$.00572$&57&$.32604$&1\\
$[.98,.99]$&$.00430$&52&$.2236$&2\\
$[.99,.995]$&$.00296$&49&$.14504$&3\\
$[.995,.997]$&$.00216$&47&$.10152$&4\\
$[.997,.998]$&$.00173$&46&$.07958$&4\\
$[.998,.9985]$&$.00146$&46&$.06716$&3\\\bottomrule
\end{tabular}
\end{center}
The data give eleven midpoint splits, each along the modulus or
$x$ coordinate, with both closed children retained. This gives
twenty closed leaves. The binary rule regenerates the whole cover,
including the boundaries of all initial domains and their common seams.

On a leaf $B_L\le r\le B_U$, $x_L\le x\le x_U$, put
$d_*=x_L-a_U/2$, $V=1/d_*$ and $p_{\max}=1-B_L^2$.
The available field bounds at the degree floor are
\begin{align}
 H_{\mathrm{hole}}&=\frac{6250-C_{\mathrm{field}}/d_*}{2(1+a_U)},\notag\\
 H_{\mathrm{phase}}&=\frac{6250-2K_0/d_*}
 {\max\{2x_L+p_{\max}/(x_L-a_U/2),
          2x_U+p_{\max}/(x_U-a_U/2)\}}.
 \label{sm:high-leaf-field}
\end{align}
The denominator in the second bound is the maximum of a convex
function on the full interval. In the first initial domain use
$K_0=21/4$, $C_{\mathrm{field}}=1899/20$; in all others use
$K_0=4$, $C_{\mathrm{field}}=2637/32$.
All numerators are positive. Each leaf declares a field mode and an
integer $H_*$ strictly below the corresponding expression in
\eqref{sm:high-leaf-field}.
Its full physical radius $R$, $\nu=m/10000$ and integer $W$
satisfy the inverse predicates with $d_*,V,H_*$ as above.
Thus the actual cumulative count is greater than
$\nu(7/22)H_*\widehat R-2$.

For its factor $f<1$ the leaf verifies
\eqref{lt:half-disk-gate} with $P=1-B_U^2$,
$T=a_U/(2P)$, $U=R/P$ and $L=1+2Tx_U$.
In the first initial domain take $C_{\min}=A_0(x_L-3A_0/4)$.
In every other initial domain use both lower constraints on
$a$:
\[
 C_{\min}=\max\{A_0(x_L-3A_0/4),\gamma_Mx_L^2\}.
\]
Every field in \eqref{sm:high-leaf-field} has the form $(N-c)/F$
with $c\ge0$. Equation~\eqref{lt:energy-ray} therefore gives
$Q>H_*^2/6250$ throughout the full degree ray. For a rational
$\ell<-\log f$ the complete product lower budget is
\begin{equation}\label{sm:high-budget}
 \ell\{\nu(7/22)H_*\widehat R-2\}
       +\frac{C_{\min}H_*^2}{12500}.
\end{equation}
Each leaf has a rational target $T_*$ with
$e^{T_*}>8/C_{\min}$ and with
\eqref{sm:high-budget} strictly greater than $T_*$.
The logarithmic lower weight is checked by
$\mathcal E_{24}(\ell)<1/f$, and the target by the positive
exponential sum through degree one hundred. All twenty whole-leaf
reserves exceed $.2558$. The factor estimate holds for every
original zero in each physical disk, and all remaining product
factors are at most one.

\smallskip\noindent\emph{The unit-circle endpoint.}
On $r\ge.9985$, the preceding chain gives $x<46a$ and
$a<.00146$. Hence $C\ge\gamma_{46}x^2$.
Take $L=11521/1000$; the exact positive exponential gate gives
$e^L>8/(\gamma_{46}\rho_0^2)$.
The product, phase and radial identities would require
\[
 1<2\sqrt{k_q}+(1-.9985^2)k_q/x^2
                 +\frac8{6250(x-.00146/2)},\qquad
 k_q=\frac{2L}{6250\gamma_{46}}.
\]
The right side decreases for all $x\ge\rho_0$. At $x=\rho_0$
it is less than $.990541<1$; equivalently the positive remainder
after the two rational terms has square exceeding $4k_q$ by more
than $.0158$. This excludes the whole endpoint range, including
$p_b=0$ and $|b|=1$.

The fixed input file contains all eight Schur profiles, the eight
successive caps, the nine initial domains, all eleven exact splits,
the twenty leaf witnesses and the endpoint parameters.
\cert{check_high6251_exact.py} reconstructs all the displayed
rational predicates and every inherited-domain equality.
The first fields and the subsequent same-gap bounds increase with
$N$; the Schur factor uses the smallest $N/(N+1)$; and the counts
and field-energy budgets strengthen by \eqref{lt:energy-ray}.
Thus the proof covers every $N\ge6250$.
\end{proof}

\section{Two supplementary scale intervals}\label{sec:small-extensions}

We prove the two normalized source estimates used at the joins in
Section~\ref{sec:small}. Throughout this section we use the strict-hole
configuration of Lemma~\ref{pr:edge-reduction}, with an actual simple
rightmost zero $b$. Assume $N=n-1\ge N_0=6250$, put $n_0=6251$,
and write
\[
 F(u)=p(i+\rho u),\qquad \beta=X-i,\qquad
 0<\alpha=a/\rho\le X\le1.
\]
Thus $\rho=1-\Im b$ and $X=\Re b/\rho$. Use
$S=F''(\beta)/F'(\beta)$, $h=|S|$, $Q=\sum_v|\beta-v|^{-2}$,
$d=\min_v|\beta-v|$, $D=dh$, and $y_0=Q/h$, where $v$ runs over
all zeros of $F'$, counted with multiplicity.
The physical quantities $H_{\rm phys}=|p''(b)/p'(b)|$ and
$d_{\rm phys}=\min_{p'(\xi)=0}|b-\xi|$ satisfy
\[
 h=\rho H_{\rm phys},\qquad d_{\rm phys}=\rho d.
\]
Every critical label $v=r-it$ satisfies
\begin{equation}\label{sm:extension-domain}
 r<\alpha/2,\qquad \rho|v|^2/2\le t<\rho|v-\alpha|^2/2.
\end{equation}
The inset is $I(u)=-\Im u-\rho|u|^2/2$. In particular,
$I(\beta)=1-\rho(1+X^2)/2\ge1-\rho$.
The finite polar quantities $W,C,E,\theta$ are those in
\eqref{or:source-definitions}; their proof applies for every
$0<\rho<1$ used below.

\subsection{A closed-cell source criterion}

Suppose
\[
 \rho\in[p,r],\quad y_0\in[l,u],\quad
 X\in[X_-,X_+]\subseteq[0,1],\quad d>1,\quad
 h\ge H,\quad |\Im S|<J.
\]
Here $0<p\le r<1$, $0\le l\le u$, and $H,J>0$.
The radial identities \eqref{or:radial} and
\eqref{or:inset-critical}, with $T=1$, give
\begin{gather}
 \frac{N\rho}{2}<\rho h+J+y_0h,\label{sm:extension-radial-lower}\\
 (1-\rho/2)Q<\frac{N\rho}{2}+J.\label{sm:extension-radial-upper}
\end{gather}
For the second inequality the negative critical-inset term is retained:
\[
 E_c<\rho X\Re S-\rho X^2Q/2.
\]
Its substitution cancels the $X^2$ terms in $I(\beta)Q$.
Consequently, putting
\[
 B_r=\frac{r/2+J/N_0}{1-r/2},
\]
one may take any positive rational $U$ satisfying at least one of
\begin{equation}\label{sm:extension-U}
 U\ge u,\qquad U^2\ge B_r,\qquad Ul\ge B_r\quad(l>0)
\end{equation}
as an upper bound for $h/n$. The first alternative follows from
$h^2\le NQ$; the second uses the same inequality and
\eqref{sm:extension-radial-upper}; the third uses
$(h/n)y_0=Q/n<B_r$.

Choose $\delta>0$ with $r^2(1+\delta)^2/2<1$ and test at
$s=X+\delta$. Define
\begin{equation}\label{sm:extension-kappa}
 \kappa_\delta(X)=\frac{X^2}{2}
 -\frac{X\max(0,1-X\delta-\delta^2)^2}
           {4(1+\delta^2)(X/2+\delta)}.
\end{equation}
Assume the whole-cell inequalities
\begin{equation}\label{sm:extension-A}
 [r(1+\delta)/2]^2<\delta,\qquad
 X-\kappa_\delta(X)y_0<\overline A
\end{equation}
with $\overline A>0$. The first places the inverse Cayley test
strictly to the right of all original zeros, because
$s/(1+(\rho s/2)^2)>X$ follows from
$\delta>X(\rho s/2)^2$. The second and
Lemma~\ref{or:critical-comparison} give
$C_{\rm real}<\rho h\overline A$.
Set
\[
 f=\frac1{1-r^2(1+\delta)^2/2},\quad
 D_p=1+[r(1+\delta)/2]^2,\quad
 B=2D_p(f\overline A U+1/n_0).
\]
If $0<B<1$, the finite polar inequality yields
\begin{equation}\label{sm:extension-HB-upper}
 \frac Eh\le\frac Wh<
       \mathcal U:=\frac{r(f\overline A+1/H)}{1-B}.
\end{equation}
Both finite corrections $1/H$ and $1/n_0$ occur in this expression.

For a lower bound for $E/h$, choose a rational radius
$0<R<1$, a number $0<\nu<1$, and a positive integer $W_0$.
Put
\begin{gather}
 \lambda=\frac{uR}{1-R},\qquad
 e_R=u\left\{\sum_{k=2}^{16}\frac{R^k}{k}
                    +\frac{R^{17}}{17(1-R)}\right\},\qquad
 \widehat R=R-e_R,\label{sm:extension-inverse}\\
 B_R=\frac{u}{(1-R)^2(1-\lambda)^3},\qquad
 \epsilon=\frac{1001B_R}{1000H},\qquad
 g=\frac{1001}{1000}\nu B_R\widehat R,\qquad
 p_0=\frac1{1+\lambda},\quad a_0=\frac1{1-\lambda}.
 \notag
\end{gather}
The complete logarithmic tail appears in $e_R$.
For $W_0\in\{1,2,3,4,5,6,8,10\}$ use $e_0=1/d_{W_0}$, where
\[
\begin{array}{c|rrrrrrrr}
 W_0&1&2&3&4&5&6&8&10\\ \hline
 d_{W_0}&2&7&20&50&148&403&2900&22000 .
\end{array}
\]
In each case $\sum_{k=0}^{69}W_0^k/k!>d_{W_0}$, so
$e_0>e^{-W_0}$. Require
\begin{gather}
 0<R<1,\quad 0<\lambda<1,\quad \widehat R>0,\quad
 W_0^2<(1-\nu^2)(H\widehat R)^2,\label{sm:extension-gates-a}\\
 \epsilon(1+e_0)+ge_0<1-a_0e_0,\qquad
 \epsilon(1+e_0)+ge_0^2<p_0-e_0,\label{sm:extension-gates-b}\\
 \epsilon+ge_0<p_0,\qquad
 0<p_0<1<a_0,\qquad \nu H\widehat R>22/7.
 \label{sm:extension-gates-c}
\end{gather}
These are the entire-contour hypotheses
\eqref{or:contour-gates}--\eqref{or:odd-gate}.
The inverse-log and amplitude comparisons therefore give an
actual original-root count
$L_R>h\nu\widehat R/\pi-2$.
The phase bound confines these roots to
\begin{equation}\label{sm:extension-width}
 |u-\beta|<R,\qquad
 |\Re(u-\beta)|\le
 \ell_R:=\min(R,W_0/H+J\nu\widehat R/H+e_R).
\end{equation}
This is a bound on the entire inverse rectangle.

Use increasing radii $R_j$, round the slopes
$\nu_j(7/22)\widehat R_j$ down and the widths up to the
$10^{-12}$ mesh, and replace the widths by running maxima.
Write the resulting slopes and widths as $s_j,\ell_j$.
They are checked to satisfy
$0<s_1<\cdots<s_k$, $R_1<\cdots<R_k<1$, and
$0\le\ell_1\le\cdots\le\ell_k\le R_k$, with each
$\ell_j\le R_j$.
Rightmostness and $R_j<1$ give $|\Re u|\le1$ for the counted roots.
Their weight in $E(s)$ is at least
\begin{equation}\label{sm:extension-weight}
 \omega_j=\max\left(0,\min_{t\in\{1-R_j,\,1+R_j\}}
       \frac{t-r(1+t^2)/2}{(\delta+\ell_j)^2+t^2}\right).
\end{equation}
This controls the whole depth interval: the derivative numerator
of the displayed quotient is
$A-t^2+r(1-A)t$, $A=(\delta+\ell_j)^2>0$.
It has one positive root, which is a maximum. If an endpoint
lower bound is negative, the true inset is still nonnegative.
Round $\omega_j$ down to the same mesh and put
$w_j=\max_{m\ge j}\omega_m$.
The nested regions make these nonincreasing weights valid.

The actual source $\beta$ has weight at least
$b_0=(1-r)/(1+\delta^2)$. Formula \eqref{or:Abel-source} gives
\begin{equation}\label{sm:extension-Abel}
 \frac Eh>
 \mathcal L:=\sum_{j=1}^kw_j(s_j-s_{j-1})
                   +\frac{\min(0,b_0-3w_1)}{H},
 \qquad s_0=0.
\end{equation}
The source is removed exactly once from each cumulative count.
The finite loss three is consequently retained. If the source
compensation is positive, discarding it remains valid without
an upper bound on $h$. In particular, the first count lower
bound need not itself be positive.

\begin{lemma}\label{sm:extension-cell}
Under the preceding whole-cell hypotheses, a collection satisfying
\eqref{sm:extension-gates-a}--\eqref{sm:extension-gates-c},
the stated nesting and rounding conditions, and
$\mathcal L>\mathcal U$, excludes the entire closed cell.
\end{lemma}
\begin{proof}
Equations \eqref{sm:extension-HB-upper} and
\eqref{sm:extension-Abel} are contradictory.
\end{proof}

Equation \eqref{sm:extension-A} has a polynomial verification.
The dependence on $y_0$ is affine, so it suffices to check
$y_0=l,u$. On the active interval
$X<(1-\delta^2)/\delta$,
\begin{equation}\label{sm:extension-kappa-active}
 \kappa_\delta(X)=
 \frac{X[X-(1-\delta)^2][X+(1+\delta)^2]}
             {2(1+\delta^2)(X+2\delta)}.
\end{equation}
On the other interval it is $X^2/2$.
Clear the positive denominators in
$\overline A-X+\kappa_\delta(X)l$ and
$\overline A-X+\kappa_\delta(X)u$ separately on those two
closed pieces. The rational Bernstein rule in
Section~\ref{sec:certificates}, with exact dyadic subdivisions,
proves positivity. The two closed pieces include their common
endpoint, and both $y_0$ endpoints are checked for either sign
of $\kappa_\delta$.

\subsection{The moment interval up to \texorpdfstring{$\rho=39/500$}{rho=39/500}}

\begin{lemma}\label{sm:T-moment-extension}
There is no strict-hole configuration in the normalized
range
\[
 \frac3{50}\le\rho\le\frac{39}{500},\qquad
 0<\alpha\le X\le1,\qquad \frac54\le y_0\le\frac94
\]
satisfying
\begin{equation}\label{sm:moment-extension-input}
 D\ge139,\quad Qd/h<9/4,\quad |\Im S|<4y_0,\quad
 H_{\rm phys}>
 \frac{N-(2637/32)/d_{\rm phys}}{2(1+a)}.
\end{equation}
\end{lemma}
\begin{proof}
The source criterion requires $d>1$. The lower intersection of the two
critical-boundary circles at $r=\alpha/2$ has depth
\[
 \eta=\frac{1-\sqrt{1-\rho^2\alpha^2/4}}{\rho}
      <\frac{\rho\alpha^2}{4}.
\]
The closest-point comparison gives
$d^2>(X-\alpha/2)^2+(1-\eta)^2$.
To see why it controls the whole crescent, on the lower hole arc put
$u=r-\alpha\le-\alpha/2$. The relevant scalar product is
\[
 (X-\alpha)u+(1/\rho-1)\sqrt{\rho^{-2}-u^2}.
\]
Its derivative is positive for $u\le0$. The largest allowed real
coordinate therefore gives the nearest permitted arc point.
The upper arc and the outward radial parts are farther away.
Consequently
\[
 d^2>1+\frac{X^2}{4}-\frac{\rho\alpha^2}{2}
 \ge1+(1/4-\rho/2)X^2>1,
\]
since $1/4-(39/500)/2>0$.

Using $d>1$, $a=\rho\alpha\le\rho$, and $N\rho>2637/32$,
the last hypothesis in \eqref{sm:moment-extension-input} yields
\begin{equation}\label{sm:moment-extension-field}
 h>\frac{N\rho-2637/32}{2(1+\rho)}.
\end{equation}
The bound $D\ge139$ is an input hypothesis; the table below gives
lower bounds $h\ge H$, starting with $H=137$.

Use the five scale intervals
\[
 [.06,.064],\ [.064,.068],\ [.068,.072],\
 [.072,.075],\ [.075,.078]
\]
and the four moment intervals
\[
 Y_1=[5/4,3/2],\quad Y_2=[3/2,7/4],\quad
 Y_3=[7/4,2],\quad Y_4=[2,9/4].
\]
Every pair is a closed cell. On the $j$th moment interval
put $J=4u$ and use the following exact parameters:
\begin{center}
\begin{tabular}{c|rrrr}
\toprule
$j$&1&2&3&4\\\midrule
$\delta$&$1/2$&$11/20$&$3/5$&$13/20$\\
$10^6\overline A$&428583&358161&304796&262971\\
$J$&6&7&8&9\\
\bottomrule
\end{tabular}
\end{center}
Take $X_-=0$, $X_+=1$ in every cell. The lower fields and upper
ratios $U$ are:
\begin{center}\small
\begin{tabular}{c|r|rrrr}
\toprule
$[\rho_-,\rho_+]$&$H$&
$10^{12}U_1$&$10^{12}U_2$&$10^{12}U_3$&$10^{12}U_4$\\
\midrule
$[.060,.064]$&137&27239669422&22809917356&19645808737&17272727273\\
$[.064,.068]$&148&28952380953&24237405107&20869565218&18343685301\\
$[.068,.072]$&159&30672199171&25670816045&22098399526&19419087137\\
$[.072,.075]$&170&31966753247&26749783550&23023376624&20228571429\\
$[.075,.078]$&179&33265348596&27832119321&23951241267&21040582727\\
\bottomrule
\end{tabular}
\end{center}
For each cell $[p,r]\times[l,u]$, its $H$ is at most one of
\[
 \frac{N_0p-2637/32}{2(1+r)},\qquad
 \frac{N_0p/2-J}{r+u}.
\]
The first comparison uses \eqref{sm:moment-extension-field}
and the second uses \eqref{sm:extension-radial-lower}.
Thus $h\ge H$ in each cell.
The displayed $U$ satisfy \eqref{sm:extension-U}.
Every polynomial comparison in \eqref{sm:extension-A} is
verified by \eqref{sm:extension-kappa-active}.

The fixed contour lists have, in the order of this table,
\[
\begin{array}{rrrr}
55&49&42&37\\
56&50&44&39\\
56&50&45&40\\
58&50&45&40\\
58&51&47&42
\end{array}
\]
members, for a total of $954$.
Their exact triples $(R,\nu,W_0)$ are the rows with family
\textup{M} in
\cert{small_extensions/small_extensions_contours.csv};
the indices begin at zero, with the scale index preceding the
moment index. They are also recorded in
\cert{TMOMENT6251_PROPOSAL.json}.
The exact checker \cert{certify_Tmoment6251_complete.py}
reconstructs \eqref{sm:extension-inverse}--\eqref{sm:extension-Abel},
including every complete contour and both polynomial pieces.
Every cell satisfies
\[
 \mathcal L-\mathcal U>\frac{14}{10000}>0.
\]
Lemma~\ref{sm:extension-cell} excludes the twenty cells.

The scale and moment intervals include all their common endpoints.
For $N\ge N_0$, both lower fields increase, while $J/N$ and
$1/n$ in the upper bounds decrease. All contour, polynomial,
and source-weight bounds are fixed on their closed cells.
The contradiction therefore holds on the entire degree ray.
\end{proof}

\subsection{The high-\texorpdfstring{$X$}{X} interval up to
\texorpdfstring{$\rho=31/500$}{rho=31/500}}

\begin{lemma}\label{sm:T-highX-extension}
There is no strict-hole configuration with
\[
 \rho_L:=\frac{30501}{500000}\le\rho\le
 \rho_U:=\frac{31}{500},\qquad
 X_0:=\frac{69}{100}\le X\le1,\qquad
 0<\alpha\le X .
\]
\end{lemma}
\begin{proof}
The distance calculation preceding \eqref{or:gap-T} remains
valid on this interval. The endpoint inequality
\[
 \frac2{125}-\rho_U\left(\frac14+\frac1{15625}\right)>0
\]
gives $\eta<\alpha^2/125$ on the whole interval.
The nearest-point comparison above gives
\[
 d^2>(X-\alpha/2)^2+(1-\alpha^2/125)^2
 \ge1+\frac{117}{500}X^2=:G(X).
\]
Write $G_0=G(X_0)$.
Lemma~\ref{sm:normalized-D25} applies because
$1/100<\rho_L\le\rho\le\rho_U<9/100$; it gives $D>25$.
Proposition~\ref{mf:D25} therefore gives
\[
 y_0<3/d<3/\sqrt{G_0}<57/20.
\]
The unphased radial identity gives
\[
 h>\frac{N_0\rho_L}{2(1+57/20)}>49,\qquad
 D^2>49^2G_0>51^2.
\]
In the matrix envelope \eqref{mf:coarse-envelope}, put
$S/h=c+is$ with $c\ge0$.
If $|s|\ge7/10$, then $c<18/25$ and its right side is at most
\[
 \frac{21/5}{51^2}
 +\frac{(11/2)(18/25)}{51}
 +\frac{1-(7/10)^2}{4}-\frac{(7/10)^2}{2}<0.
\]
Thus $|\Im S|<(7/10)h$ and $\mu=1/D<1/50$.

\smallskip
\noindent\emph{A joint source and moment estimate.}
Take $\delta=3/5$ and set
\[
 b_0=\frac{1-\rho_U}{1+\delta^2},\quad
 f=[1-\rho_U^2(1+\delta)^2/2]^{-1},\quad
 D_p=1+[\rho_U(1+\delta)/2]^2.
\]
The source $\beta$ and the finite polar inequality give
\begin{equation}\label{sm:highX-K}
 \rho h A_\delta>K,\qquad
 A_\delta=X-\kappa_\delta(X)y_0,\qquad
 K=\min_{\rho\in\{\rho_L,\rho_U\}}
 \frac{n_0\rho b_0/(n_0\rho+2D_pb_0)-\rho}{f}.
\end{equation}
The expression minimized is concave in $\rho$ and increases
with $n$, as in \eqref{or:direct-K}. Exact evaluation gives
\[
 K=\frac{1793507148275409343}{2882630131640625000}
       >\frac{31}{50}.
\]
Hence $A_\delta>0$.
On the entire $X$ interval,
\begin{equation}\label{sm:highX-kappa}
 \kappa_\delta(X)>c_0X^2,\qquad c_0=12/25.
\end{equation}
The active endpoint $(1-\delta^2)/\delta=16/15$ is beyond one,
so this is a single rational-polynomial comparison.

Put $v=Xy_0$. Then $0<v<1/c_0$ and
$A_\delta<X(1-c_0v)$.
The phase bound in the radial identity gives
$N\rho/2<hR$, where
\[
 R:=\beta_0+I_0y_0,\qquad
 \beta_0=7/10+(3/10)\rho_U,\qquad
 I_0=1-\rho_L(1+X_0^2)/2.
\]
Multiplying this estimate and \eqref{sm:highX-K} for the same
configuration gives, for $q=Q/h^2=y_0/h$,
\begin{equation}\label{sm:highX-joint}
 q^2<\frac{2y_0^2A_\delta R}{NK}
 \le\frac{2v^2(1-c_0v)
            [\beta_0/X_0+I_0v/X_0^2]}{N_0K}.
\end{equation}
The second step uses the nonnegative factor $1-c_0v$ and
the decrease of $1/X,1/X^2$ on $X\ge X_0$.
The polynomial
\begin{equation}\label{sm:highX-q-polynomial}
 N_0K-800v^2(1-c_0v)
             [\beta_0/X_0+I_0v/X_0^2]
\end{equation}
is positive throughout $[0,1/c_0]$.
Clearing its positive denominator and using exact Bernstein
subdivision proves this sign, together with
\eqref{sm:highX-kappa}. Thus $q<1/20$.
Proposition~\ref{lt:q05mu02} applies with $q<1/20$ and $\mu<1/50$.
It gives two distinct original zeros, different
from $\beta$, of the form $u_\pm=\beta+w_\pm/S$, with
\begin{equation}\label{sm:highX-pair}
 |w_\pm|<17/2,\qquad
 |(w_++w_-)/2|<19/10,\qquad |\Im S|<7y_0.
\end{equation}

\smallskip
\noindent\emph{The entire range $h\ge95$.}
Suppose $h\ge95$. Then $D^2>95^2G_0>10000$.
Proposition~\ref{mf:D100}, followed by the wide row of
Proposition~\ref{lt:wideD100}, gives
\[
 Qd/h<9/4,\quad q<9/400,\quad \mu<1/100,\quad
 |\Im S|<(23/5)y_0.
\]
In particular $y_0^2G_0<81/16$ and $y_0<427/200$.
We apply Lemma~\ref{sm:extension-cell} on the ten moment
intervals in the following table, taking the entire scale
interval $[\rho_L,\rho_U]$ in each row and $X\in[X_0,1]$.
Here $J=(23/5)u$, where $u$ is the upper moment endpoint.
\begin{center}\small
\begin{tabular}{c|rcrrr}
\toprule
$[l,u]$&$H$&$\delta$&$10^6\overline A$&$10^{12}U$&contours\\
\midrule
$[0,1/4]$&607&$1/200$&1000001&179392392668&149\\
$[1/4,1/2]$&335&$1/5$&924589&129486068112&117\\
$[1/2,3/4]$&230&$3/10$&803342&65122807018&95\\
$[3/4,1]$&175&$7/20$&679683&43668386653&78\\
$[1,5/4]$&140&$9/20$&528001&32941176471&64\\
$[5/4,3/2]$&117&$1/2$&428583&26504850362&52\\
$[3/2,7/4]$&100&$11/20$&355372&22213966289&44\\
$[7/4,15/8]$&95&$3/5$&285411&19094795814&39\\
$[15/8,2]$&95&$3/5$&256512&17872445821&37\\
$[2,427/200]$&95&$3/5$&227612&16806687307&34\\
\bottomrule
\end{tabular}
\end{center}
Each $H$ is either $95$ or at most
$(N_0\rho_L/2-J)/(\rho_U+u)$. Thus $h\ge H$ follows from
the hypothesis $h\ge95$ and the radial identity.
The $U$ satisfy \eqref{sm:extension-U}.
At both moment endpoints, the entire source polynomials
are checked on $[X_0,1]$, using the exact active/inactive
split in \eqref{sm:extension-kappa-active}.

The $709$ exact triples are the family \textup{X} rows of
\cert{small_extensions/small_extensions_contours.csv},
with the cells in the printed order; they are also in
\cert{T6251_HIGHX_PROPOSAL.json}.
The checker \cert{certify_T6251_highX_complete.py}
reconstructs all the predicates in
Lemma~\ref{sm:extension-cell}. In every one of the ten cells,
\[
 \mathcal L-\mathcal U>\frac6{10000}>0.
\]
These ten cells exclude the entire range $h\ge95$.

\smallskip
\noindent\emph{The pair when $49<h<95$.}
It remains to use the pair \eqref{sm:highX-pair}.
For $\delta=3/5$,
\[
 y_0A_\delta<v-c_0v^2\le\frac1{4c_0}=\frac{25}{48}.
\]
The radial identity and $|\Im S|<7y_0$ give
$\rho(N-2h)<2y_0(h+7)$.
Since $A_\delta>0$ and $N-2h>0$,
\[
 \rho h A_\delta
 <\frac{2(25/48)95\cdot102}{6250-190}.
\]
Let $B_{95}=2D_p(95f+1)/6251<1$.
Here $A_\delta<X\le1$ by \eqref{sm:highX-kappa}, so this
is a valid bound for the complete polar denominator.
Lemma~\ref{or:HB} gives
\begin{equation}\label{sm:highX-terminal-upper}
 W<\frac{f\,2(25/48)95\cdot102/(6250-190)+\rho_U}{1-B_{95}}.
\end{equation}

Put $z_\pm=u_\pm-\beta$, $R=17/98$, and $m=19/490$.
The pair satisfies
$|z_\pm|<R$, $|(z_++z_-)/2|<m$, and $\Re z_\pm\le0$.
Its insets lie in $[3/4,6/5]$: for the lower bound use
\[
 I(\beta+z)\ge(1-\rho_U)(1-R)-\rho_UR^2/2>3/4,
\]
and for the upper bound use $I(\beta+z)\le1+R<6/5$.
On this whole interval,
$\sqrt I\ge I/2+49/100$.
Indeed both sides are nonnegative, and the difference
$(I/2+49/100)^2-I$ is a convex quadratic negative at
the two endpoints.
The bound on the pair mean gives
\begin{align*}
 I_++I_-&>I_\Sigma:=2(1-\rho_U)(1-m)-\rho_UR^2,\\
 I_\Sigma&=\frac{43251409}{24010000}>0,\\
 |s-u_+|^2+|s-u_-|^2
 &<D_\Sigma:=2(1+\delta^2)+4(7/6)m+2R^2,
\end{align*}
where $(7/6)^2>1+\delta^2$.
Since $I_\Sigma/2+49/50>0$, Cauchy--Schwarz on these two terms,
together with the separate source $\beta$, yields
\begin{equation}\label{sm:highX-terminal-lower}
 E>\frac{1-\rho_U}{1+\delta^2}
                    +\frac{(I_\Sigma/2+49/50)^2}{D_\Sigma}.
\end{equation}
The exact lower bound in \eqref{sm:highX-terminal-lower}
exceeds the upper bound in \eqref{sm:highX-terminal-upper} by
\[
 \frac{9014845822117732608757183663183337}
      {97049415152918630536485588216000000}>\frac{92}{1000}.
\]
This contradicts $W\ge E$.

The source bound, the polynomial
\eqref{sm:highX-q-polynomial}, the geometric and phase
bounds, and the two terminal bounds are rationally checked
by \cert{certify_T6251_highX_bootstrap.py}; its code is also
called by the complete ten-cell verifier.
For $N\ge6250$, the initial fields and the source strengthen,
while the finite polar and radial corrections decrease.
On the terminal range $h<95$, the denominator satisfies
$N-2h>N-190\ge6060$. All other
domains and contours are fixed, and every scale and moment
endpoint has been included. This proves the lemma on the
entire stated degree ray.
\end{proof}

\section{Completion of the proof}\label{sec:assembly}

Retain $N_0=6250$, $\rho_0=30501/500000$, and
$A_0=\rho_0/100$ from Section~\ref{sec:small}.
\begin{proof}[Proof of Theorem~\ref{thm:main}]
Suppose the theorem fails. By
Lemma~\ref{pr:edge-reduction}, there is a strict supporting-edge hole
$0<a<1$ and an actual simple rightmost zero $b=x+iy$ with
$x\ge a$, $y\ge0$. Put $r=|b|$ and $\rho=\max(x,1-y)$.

For $0<a\le A_0$, Proposition~\ref{sm:fine} excludes
$\rho>\rho_0$, and Proposition~\ref{or:theorem} excludes
$0<\rho\le\rho_0$. For $A_0\le a\le1/10$,
Proposition~\ref{sm:low-middle} applies when $r\le47/50$,
and Proposition~\ref{sm:high} applies when $r>47/50$ outside
the origin box. The remaining origin box is already excluded.
Propositions~\ref{bc:bridge} and~\ref{bc:compact} cover
$1/10\le a\le39/40$, and Proposition~\ref{bc:boundary}
covers $39/40\le a<1$. Thus all values of $a$ have been excluded.

The parameter intervals are closed and the final contradictions strict.
In particular, the high small-witness estimate includes $|b|=1$;
the compact estimate includes the interval $[1/10,11/100]$; and the
two additional small-witness source ranges cover their entire common
boundaries with the origin range. Every estimate is uniform for
$N\ge N_0$. This proves the theorem.
\end{proof}

\Needspace{8\baselineskip}
The affine invariance of the argument gives the following formulation.
\begin{corollary}\label{cor:disk}
Let $p$ have degree at least $6251$ and suppose its zeros lie in a
closed disk of radius $R\ge0$. Every point of $\conv Z(p)$ lies
within distance $R$ of a critical point.
\end{corollary}
\begin{proof}
For $R>0$, translate the center to zero and divide the variable by
$R$, then apply Theorem~\ref{thm:main}. For $R=0$ all zeros and
all critical points coincide.
\end{proof}

\section{Finite certificates and reproducibility}\label{sec:certificates}

The electronic appendix contains the fixed rational inputs, verification
programs, a file index, and a manifest of SHA-256 hashes. The programs
recompute the inequalities from these inputs by the methods below.

\subsection{Elementary enclosures}

All rational calculations use integer numerators and positive integer
denominators. For a positive integer $m$, a square-root enclosure is
obtained by taking the integer square root of $\lfloor m^2a/b\rfloor$
for a rational $a/b\ge0$;
division by $m$ gives the lower bound and addition of $1/m$ gives
an upper bound unless the square is exact. Thus the directions of
the rounded radical bounds are determined by integer inequalities.

The enclosures $6283/2000<333/106<\pi<22/7$ follow from Machin's identity
\[
 \pi=16\arctan(1/5)-4\arctan(1/239).
\]
Indeed $\tan(4\arctan(1/5))=120/119$, and subtraction of
$\arctan(1/239)$ gives tangent one with angle in $(0,\pi/2)$.
If $A_m(x)=\sum_{k=0}^{m-1}(-1)^kx^{2k+1}/(2k+1)$, the
alternating-series bounds give
\[
 16A_6(1/5)-4A_3(1/239)<\pi
 <16A_7(1/5)-4A_2(1/239).
\]
Both comparisons with the displayed rational endpoints are exact.
Finer enclosures are obtained by retaining more terms.

For $x\ge1$, first divide $x$ by a power of two to put it in $[1,2)$.
With $t=(x-1)/(x+1)$ on this interval,
\[
 2\sum_{j=0}^{m-1}\frac{t^{2j+1}}{2j+1}
 \le\log x
 \le2\sum_{j=0}^{m-1}\frac{t^{2j+1}}{2j+1}
       +\frac{2t^{2m+1}}{(2m+1)(1-t^2)}.
\]
The same formula at $x=2$ encloses the range-reduction contribution.
The formulas for $\log(1\pm r)$ in
\eqref{mf:directed-log} retain their complete alternating or
positive geometric remainder. For an integer $m\ge0$, put
$E_m(x)=\sum_{j=0}^m x^j/j!$. If $0\le x<m+2$, then
\begin{equation}\label{ce:exp}
 E_m(x)\le e^x\le E_m(x)
                +\frac{x^{m+1}}{(m+1)!(1-x/(m+2))}.
\end{equation}
Negative exponentials are bounded by reciprocals of positive
enclosures. The seven rows in Proposition~\ref{ar:energy}, for
example, use the lower logarithm bound with $48$ terms and
$e^{-x}\le1/E_{160}(x)$. Their checker uses only rational arithmetic.

The interval computations use directed real and complex balls.
An asserted strict inequality means that the entire resulting
enclosure has the required sign. Input floating-point test points are converted to
their exact binary rationals before use.

\subsection{Entire integrals}

The only numerical integrations in the field certificates are of
entire functions of the form $\exp(P(t))$, with $P$ a polynomial,
over a finite real interval. The preceding analytic lemmas bound the
finite-product errors separately from these model integrals.

There is a direct series verification of each such integral. After
an affine change of variable, write
$P(t)=p_0+\sum_{j=1}^d p_jt^j$ on $[0,1]$ and
$\exp(P(t)-p_0)=\sum_{k\ge0}a_kt^k$.
The coefficients are determined by
\[
 a_0=1,\qquad
 ka_k=\sum_{j=1}^{\min(d,k)}jp_ja_{k-j}.
\]
For any rational $R>1$, Cauchy's estimate gives
$|a_k|\le M_RR^{-k}$, where
$M_R=\exp(\sum_{j=1}^d|p_j|R^j)$.
Consequently
\[
 \left|\int_0^1e^{P(t)}\,dt
       -e^{p_0}\sum_{k=0}^{m}\frac{a_k}{k+1}\right|
 \le |e^{p_0}|
       \frac{M_RR^{-(m+1)}}{(m+2)(1-1/R)}.
\]
Subdividing the interval, or increasing $m$, gives arbitrarily
small error bounds. Formula \eqref{ce:exp} encloses $M_R$.
The field checkers also supply validated analytic quadrature
using Arb/Acb; all integrands satisfy its entire-function
hypothesis. The separate series implementations in the
electronic appendix give the coefficient recurrence and the
Cauchy remainder explicitly.

For the Gaussian shells in Proposition~\ref{mf:affine}, the rational
integration-by-parts and Mills bounds in its proof suffice.
An additional verifier evaluates the Gaussian tails by the
complementary error function with interval bounds.

\subsection{Polynomial bounds on whole boxes}

The companion centers are defined by finite recurrences, so their
coefficient identities are checked by exact polynomial arithmetic.
For the Gaussian center, the moments
at $s=\pm2\pi i$ obey $I_0=0$ and
$I_m=(1-mI_{m-1})/s$. This derives the center and residual
coefficients exactly. The identity $e^s=1$ is used only at
$s=\pm2\pi i$; the resulting polynomial bounds may then be
verified on the larger real interval $6\le |s|\le44/7$.

After rescaling each real variable to $[0,1]$, a polynomial of
degree at most $d$ has the Bernstein representation
\[
 \sum_{j=0}^dc_jt^j
 =\sum_{k=0}^d\left(\sum_{j=0}^k
           \frac{\binom{k}{j}}{\binom{d}{j}}c_j\right)
          \binom{d}{k}t^k(1-t)^{d-k}.
\]
The basis functions are nonnegative and sum to one. Therefore
the extrema of the coefficient list bound the polynomial
on the whole interval. Tensor products give the same result
on a box. Exact dyadic subdivision and the de Casteljau
recurrence preserve these bounds. Bounds on complex moduli
are proved by applying this rule to the difference between
a proposed squared bound and the squared modulus.
The common coefficient-domain verifier checks all twenty
Gaussian polynomial bounds by this method.

Where a one-variable polynomial sign is sharper than its
Bernstein enclosure, the programs use its exact Sturm chain.
Polynomial division is over the rationals; signs at both
endpoints are checked separately. The variation difference
counts roots in the open interval. Combined with a sign at
one point and the endpoint checks, this proves the desired
sign on the whole closed interval.

\subsection{Coverage and degree dependence}

The rectangular covers in this paper use either a rational
edge decomposition or a dyadic tree. In the former, checking
every open slab between consecutive rational edges and every
edge itself proves closed coverage; the horizontal intervals
are merged in sorted order. In the latter, every deepest
pair of leaves must be the two halves of the same coordinate
at a common parent. Replacing such a pair by its parent
must eventually produce the complete root box.

All inequalities involving the degree are extended from
$N_0=6250$ by the stated analytic monotonicities. One
frequent step is worth making explicit. If
$H>(N-c)/F$, with $F>0$ and $0\le c<N_0$, then
$H/N>(1-c/N)/F$ increases as a lower bound with $N$.
Together with $Q\ge H^2/N$ this gives, for every $N\ge N_0$,
\[
 Q>\frac{(N_0-c)^2}{N_0F^2}.
\]
This is the energy lower bound used in the product estimates.

The principal finite covers have the following sizes.
The count of contours includes complete boundaries, with
the inverse-map and amplitude comparisons both checked.
\begin{center}\small
\begin{tabular}{lrr}
\toprule
Component&closed cells or rows&additional finite objects\\
\midrule
Affine field&$275+28$&$2750$ paired shells\\
$D\ge25$ transition&$50614$&$186477$ point obstructions\\
$D\ge30$ transition&$26397$&$82823$ point obstructions\\
$D\ge100$ transition&$9745$&$8160$ point obstructions\\
Origin base&$9$ path ranges&$55$ inverse contours\\
Origin $T=1$ strip&$63$&$3831$ inverse contours\\
Origin $X=1$ strip&$301$&$14891$ inverse contours\\
Fine small-witness range&$1$&$41$ inverse contours\\
Low and middle small-witness range&$23$&$1809$ inverse contours\\
High small-witness range&$20$&$411+20$ contour constructions\\
Moment extension&$20$&$954$ inverse contours\\
High-$X$ extension&$10$&$709$ inverse contours\\
Boundary and auxiliary strips&$8$&whole-strip inequalities\\
Compact range&$13$&$2918+2$ inverse contours\\
\bottomrule
\end{tabular}
\end{center}

The file \cert{verification/INDEX.json} maps each mathematical
component to its fixed data and executable checker. The entry
point \cert{verification/run_all.py} runs the included
verifications and writes new results to a separate results
directory. The input hashes are checked before execution.
The programs use Python's exact rational arithmetic, SymPy
for exact polynomial arithmetic, and, where indicated,
\cert{python-flint} for rigorous balls.
The LaTeX source, all fixed inputs, and the verification code are included.

\begin{thebibliography}{99}

\bibitem{BouthatFallat}
L. Bouthat and S. Fallat, \emph{A proof of Borcea's quadratic variance
conjecture}, preprint (2026), \href{https://arxiv.org/abs/2610.12499v1}{arXiv:2610.12499v1}.

\bibitem{BrownXiang1999}
J. E. Brown and G. Xiang, \emph{Proof of the Sendov conjecture for
polynomials of degree at most eight}, J. Math. Anal. Appl.
\textbf{232} (1999), no.~2, 272--292.
\href{https://doi.org/10.1006/jmaa.1999.6267}{doi:10.1006/jmaa.1999.6267}.

\bibitem{Chalebgwa}
T. P. Chalebgwa, \emph{Sendov's conjecture: A note on a paper of
D\'egot}, Anal. Math. \textbf{46} (2020), no.~3, 447--463.
\href{https://doi.org/10.1007/s10476-020-0050-x}{doi:10.1007/s10476-020-0050-x}.

\bibitem{Degot}
J. D\'egot, \emph{Sendov conjecture for high degree polynomials},
Proc. Amer. Math. Soc. \textbf{142} (2014), 1337--1349.
\href{https://arxiv.org/abs/1106.4126v2}{arXiv:1106.4126v2}.

\bibitem{Gacs1971}
F. Gacs, \emph{On polynomials whose zeros are in the unit disk},
J. Math. Anal. Appl. \textbf{36} (1971), 627--637.
\href{https://doi.org/10.1016/0022-247X(71)90045-X}{doi:10.1016/0022-247X(71)90045-X}.

\bibitem{Hayman1967}
W. K. Hayman, \emph{Research Problems in Function Theory},
Athlone Press, London, 1967.

\bibitem{HaymanLingham}
W. K. Hayman and E. F. Lingham, \emph{Research Problems in Function
Theory}, fiftieth anniversary ed., Springer, Cham, 2019.

\bibitem{Mazur}
L. Mazur, \emph{A Simplified Proof Candidate for Sendov's Conjecture
with Exact Rational Certificates}, public revision~14, August 2, 2026,
\href{https://www.proofatlas.ai/research/sendov-conjecture-proof-candidate/}{ProofAtlas}.

\bibitem{MeirSharma1969}
A. Meir and A. Sharma, \emph{On Ilyeff's conjecture},
Pacific J. Math. \textbf{31} (1969), no.~2, 459--467.
\href{https://msp.org/pjm/1969/31-2/pjm-v31-n2-p20-s.pdf}{Original article}.

\bibitem{Pawlowski}
P. Paw{\l}owski, \emph{On the zeros of a polynomial and its derivatives},
Trans. Amer. Math. Soc. \textbf{350} (1998), 4461--4472.

\bibitem{Rubinstein1968}
Z. Rubinstein, \emph{On a problem of Ilyeff}, Pacific J. Math.
\textbf{26} (1968), no.~1, 159--161.
\href{https://msp.org/pjm/1968/26-1/pjm-v26-n1-p18-s.pdf}{Original article}.

\bibitem{Schmeisser1969}
G. Schmeisser, \emph{Bemerkungen zu einer Vermutung von Ilieff},
Math. Z. \textbf{111} (1969), 121--125.
\href{https://doi.org/10.1007/BF01111192}{doi:10.1007/BF01111192}.

\bibitem{Schmeisser1971}
G. Schmeisser, \emph{Zur Lage der kritischen Punkte eines Polynoms},
Rend. Sem. Mat. Univ. Padova \textbf{46} (1971), 405--415.
\href{https://www.numdam.org/item/RSMUP_1971__46__405_0/}{Original article}.

\bibitem{Schmeisser1977}
G. Schmeisser, \emph{On Ilieff's conjecture},
Math. Z. \textbf{156} (1977), 165--173.
\href{https://doi.org/10.1007/BF01178761}{doi:10.1007/BF01178761}.

\bibitem{Schmeisser2003}
G. Schmeisser, \emph{Majorization of the critical points of a polynomial
by its zeros}, Comput. Methods Funct. Theory \textbf{3} (2003), 95--103.

\bibitem{TangQualitative}
Q. Tang, \emph{Schmeisser's conjecture in large degree}, preprint
(2026), \href{https://hexagonmath.org/2611.00162v1}{Hexagon:2611.00162v1}.

\bibitem{Tao2022}
T. Tao, \emph{Sendov's conjecture for sufficiently-high-degree
polynomials}, Acta Math. \textbf{229} (2022), 347--392.
\href{https://doi.org/10.4310/acta.2022.v229.n2.a3}{doi:10.4310/acta.2022.v229.n2.a3}.

\bibitem{Tao2026}
T. Tao, \emph{A digestion of the proof of Sendov's conjecture},
What's New, August 12, 2026.
\url{https://terrytao.wordpress.com/2026/08/12/a-digestion-of-the-proof-of-sendovs-conjecture/}.

\bibitem{ZhangCentroid}
T. Zhang, \emph{A refinement of Pawlowski's result}, Proc. Amer. Math.
Soc. \textbf{154} (2026), no.~2, 775--782.
\href{https://doi.org/10.1090/proc/17471}{doi:10.1090/proc/17471}.

\bibitem{ZhangExplicit}
T. Zhang, \emph{Sendov's conjecture holds for every degree
$n\ge10^{200000}$}, preprint (2026),
\href{https://arxiv.org/abs/2609.20256v1}{arXiv:2609.20256v1}.

\bibitem{ZhangHull}
T. Zhang, \emph{When do the convex hulls of zeros and critical points
coincide?}, J. Math. Anal. Appl. \textbf{566} (2027), no.~1, 131002,
published online 2026.
\href{https://doi.org/10.1016/j.jmaa.2026.131002}{doi:10.1016/j.jmaa.2026.131002}.

\bibitem{ZhangQuadratic}
T. Zhang, \emph{Beyond Sendov's conjecture: the quadratic Tang--Zhang
inequality}, preprint (2026),
\href{https://arxiv.org/abs/2609.19126v1}{arXiv:2609.19126v1}.

\bibitem{ZhangVariance}
T. Zhang, \emph{Borcea's $2$-variance conjecture}, preprint (2026),
\href{https://arxiv.org/abs/2610.02035v1}{arXiv:2610.02035v1}.

\end{thebibliography}

\end{document}
